\documentclass[11pt]{article}
\newcounter{none}

    \usepackage[breakable]{tcolorbox}
    \usepackage{parskip} % Stop auto-indenting (to mimic markdown behaviour)
    

    % Basic figure setup, for now with no caption control since it's done
    % automatically by Pandoc (which extracts ![](path) syntax from Markdown).
    \usepackage{graphicx}
    % Keep aspect ratio if custom image width or height is specified
    \setkeys{Gin}{keepaspectratio}
    % Maintain compatibility with old templates. Remove in nbconvert 6.0
    \let\Oldincludegraphics\includegraphics
    % Ensure that by default, figures have no caption (until we provide a
    % proper Figure object with a Caption API and a way to capture that
    % in the conversion process - todo).
    \usepackage{caption}
    \DeclareCaptionFormat{nocaption}{}
    \captionsetup{format=nocaption,aboveskip=0pt,belowskip=0pt}

    \usepackage{float}
    \floatplacement{figure}{H} % forces figures to be placed at the correct location
    \usepackage{xcolor} % Allow colors to be defined
    \usepackage{enumerate} % Needed for markdown enumerations to work
    \usepackage{geometry} % Used to adjust the document margins
    \usepackage{amsmath} % Equations
    \usepackage{amssymb} % Equations
    \usepackage{textcomp} % defines textquotesingle
    % Hack from http://tex.stackexchange.com/a/47451/13684:
    \AtBeginDocument{%
        \def\PYZsq{\textquotesingle}% Upright quotes in Pygmentized code
    }
    \usepackage{upquote} % Upright quotes for verbatim code
    \usepackage{eurosym} % defines \euro

    \usepackage{iftex}
    \ifPDFTeX
        \usepackage[T1]{fontenc}
        \IfFileExists{alphabeta.sty}{
              \usepackage{alphabeta}
          }{
              \usepackage[mathletters]{ucs}
              \usepackage[utf8x]{inputenc}
          }
    \else
        \usepackage{fontspec}
        \usepackage{unicode-math}
    \fi

    \usepackage{fancyvrb} % verbatim replacement that allows latex
    \usepackage{grffile} % extends the file name processing of package graphics
                         % to support a larger range
    \makeatletter % fix for old versions of grffile with XeLaTeX
    \@ifpackagelater{grffile}{2019/11/01}
    {
      % Do nothing on new versions
    }
    {
      \def\Gread@@xetex#1{%
        \IfFileExists{"\Gin@base".bb}%
        {\Gread@eps{\Gin@base.bb}}%
        {\Gread@@xetex@aux#1}%
      }
    }
    \makeatother
    \usepackage[Export]{adjustbox} % Used to constrain images to a maximum size
    \adjustboxset{max size={0.9\linewidth}{0.9\paperheight}}

    % The hyperref package gives us a pdf with properly built
    % internal navigation ('pdf bookmarks' for the table of contents,
    % internal cross-reference links, web links for URLs, etc.)
    \usepackage{hyperref}
    % The default LaTeX title has an obnoxious amount of whitespace. By default,
    % titling removes some of it. It also provides customization options.
    \usepackage{titling}
    \usepackage{longtable} % longtable support required by pandoc >1.10
    \usepackage{booktabs}  % table support for pandoc > 1.12.2
    \usepackage{array}     % table support for pandoc >= 2.11.3
    \usepackage{calc}      % table minipage width calculation for pandoc >= 2.11.1
    \usepackage[inline]{enumitem} % IRkernel/repr support (it uses the enumerate* environment)
    \usepackage[normalem]{ulem} % ulem is needed to support strikethroughs (\sout)
                                % normalem makes italics be italics, not underlines
    \usepackage{soul}      % strikethrough (\st) support for pandoc >= 3.0.0
    \usepackage{mathrsfs}
    

    
    % Colors for the hyperref package
    \definecolor{urlcolor}{rgb}{0,.145,.698}
    \definecolor{linkcolor}{rgb}{.71,0.21,0.01}
    \definecolor{citecolor}{rgb}{.12,.54,.11}

    % ANSI colors
    \definecolor{ansi-black}{HTML}{3E424D}
    \definecolor{ansi-black-intense}{HTML}{282C36}
    \definecolor{ansi-red}{HTML}{E75C58}
    \definecolor{ansi-red-intense}{HTML}{B22B31}
    \definecolor{ansi-green}{HTML}{00A250}
    \definecolor{ansi-green-intense}{HTML}{007427}
    \definecolor{ansi-yellow}{HTML}{DDB62B}
    \definecolor{ansi-yellow-intense}{HTML}{B27D12}
    \definecolor{ansi-blue}{HTML}{208FFB}
    \definecolor{ansi-blue-intense}{HTML}{0065CA}
    \definecolor{ansi-magenta}{HTML}{D160C4}
    \definecolor{ansi-magenta-intense}{HTML}{A03196}
    \definecolor{ansi-cyan}{HTML}{60C6C8}
    \definecolor{ansi-cyan-intense}{HTML}{258F8F}
    \definecolor{ansi-white}{HTML}{C5C1B4}
    \definecolor{ansi-white-intense}{HTML}{A1A6B2}
    \definecolor{ansi-default-inverse-fg}{HTML}{FFFFFF}
    \definecolor{ansi-default-inverse-bg}{HTML}{000000}

    % common color for the border for error outputs.
    \definecolor{outerrorbackground}{HTML}{FFDFDF}

    % commands and environments needed by pandoc snippets
    % extracted from the output of `pandoc -s`
    \providecommand{\tightlist}{%
      \setlength{\itemsep}{0pt}\setlength{\parskip}{0pt}}
    \DefineVerbatimEnvironment{Highlighting}{Verbatim}{commandchars=\\\{\}}
    % Add ',fontsize=\small' for more characters per line
    \newenvironment{Shaded}{}{}
    \newcommand{\KeywordTok}[1]{\textcolor[rgb]{0.00,0.44,0.13}{\textbf{{#1}}}}
    \newcommand{\DataTypeTok}[1]{\textcolor[rgb]{0.56,0.13,0.00}{{#1}}}
    \newcommand{\DecValTok}[1]{\textcolor[rgb]{0.25,0.63,0.44}{{#1}}}
    \newcommand{\BaseNTok}[1]{\textcolor[rgb]{0.25,0.63,0.44}{{#1}}}
    \newcommand{\FloatTok}[1]{\textcolor[rgb]{0.25,0.63,0.44}{{#1}}}
    \newcommand{\CharTok}[1]{\textcolor[rgb]{0.25,0.44,0.63}{{#1}}}
    \newcommand{\StringTok}[1]{\textcolor[rgb]{0.25,0.44,0.63}{{#1}}}
    \newcommand{\CommentTok}[1]{\textcolor[rgb]{0.38,0.63,0.69}{\textit{{#1}}}}
    \newcommand{\OtherTok}[1]{\textcolor[rgb]{0.00,0.44,0.13}{{#1}}}
    \newcommand{\AlertTok}[1]{\textcolor[rgb]{1.00,0.00,0.00}{\textbf{{#1}}}}
    \newcommand{\FunctionTok}[1]{\textcolor[rgb]{0.02,0.16,0.49}{{#1}}}
    \newcommand{\RegionMarkerTok}[1]{{#1}}
    \newcommand{\ErrorTok}[1]{\textcolor[rgb]{1.00,0.00,0.00}{\textbf{{#1}}}}
    \newcommand{\NormalTok}[1]{{#1}}

    % Additional commands for more recent versions of Pandoc
    \newcommand{\ConstantTok}[1]{\textcolor[rgb]{0.53,0.00,0.00}{{#1}}}
    \newcommand{\SpecialCharTok}[1]{\textcolor[rgb]{0.25,0.44,0.63}{{#1}}}
    \newcommand{\VerbatimStringTok}[1]{\textcolor[rgb]{0.25,0.44,0.63}{{#1}}}
    \newcommand{\SpecialStringTok}[1]{\textcolor[rgb]{0.73,0.40,0.53}{{#1}}}
    \newcommand{\ImportTok}[1]{{#1}}
    \newcommand{\DocumentationTok}[1]{\textcolor[rgb]{0.73,0.13,0.13}{\textit{{#1}}}}
    \newcommand{\AnnotationTok}[1]{\textcolor[rgb]{0.38,0.63,0.69}{\textbf{\textit{{#1}}}}}
    \newcommand{\CommentVarTok}[1]{\textcolor[rgb]{0.38,0.63,0.69}{\textbf{\textit{{#1}}}}}
    \newcommand{\VariableTok}[1]{\textcolor[rgb]{0.10,0.09,0.49}{{#1}}}
    \newcommand{\ControlFlowTok}[1]{\textcolor[rgb]{0.00,0.44,0.13}{\textbf{{#1}}}}
    \newcommand{\OperatorTok}[1]{\textcolor[rgb]{0.40,0.40,0.40}{{#1}}}
    \newcommand{\BuiltInTok}[1]{{#1}}
    \newcommand{\ExtensionTok}[1]{{#1}}
    \newcommand{\PreprocessorTok}[1]{\textcolor[rgb]{0.74,0.48,0.00}{{#1}}}
    \newcommand{\AttributeTok}[1]{\textcolor[rgb]{0.49,0.56,0.16}{{#1}}}
    \newcommand{\InformationTok}[1]{\textcolor[rgb]{0.38,0.63,0.69}{\textbf{\textit{{#1}}}}}
    \newcommand{\WarningTok}[1]{\textcolor[rgb]{0.38,0.63,0.69}{\textbf{\textit{{#1}}}}}
    \makeatletter
    \newsavebox\pandoc@box
    \newcommand*\pandocbounded[1]{%
      \sbox\pandoc@box{#1}%
      % scaling factors for width and height
      \Gscale@div\@tempa\textheight{\dimexpr\ht\pandoc@box+\dp\pandoc@box\relax}%
      \Gscale@div\@tempb\linewidth{\wd\pandoc@box}%
      % select the smaller of both
      \ifdim\@tempb\p@<\@tempa\p@
        \let\@tempa\@tempb
      \fi
      % scaling accordingly (\@tempa < 1)
      \ifdim\@tempa\p@<\p@
        \scalebox{\@tempa}{\usebox\pandoc@box}%
      % scaling not needed, use as it is
      \else
        \usebox{\pandoc@box}%
      \fi
    }
    \makeatother

    % Define a nice break command that doesn't care if a line doesn't already
    % exist.
    \def\br{\hspace*{\fill} \\* }
    % Math Jax compatibility definitions
    \def\gt{>}
    \def\lt{<}
    \let\Oldtex\TeX
    \let\Oldlatex\LaTeX
    \renewcommand{\TeX}{\textrm{\Oldtex}}
    \renewcommand{\LaTeX}{\textrm{\Oldlatex}}
    % Document parameters
    % Document title
    \title{REPORT}
    
    
    
    
    
    
    
% Pygments definitions
\makeatletter
\def\PY@reset{\let\PY@it=\relax \let\PY@bf=\relax%
    \let\PY@ul=\relax \let\PY@tc=\relax%
    \let\PY@bc=\relax \let\PY@ff=\relax}
\def\PY@tok#1{\csname PY@tok@#1\endcsname}
\def\PY@toks#1+{\ifx\relax#1\empty\else%
    \PY@tok{#1}\expandafter\PY@toks\fi}
\def\PY@do#1{\PY@bc{\PY@tc{\PY@ul{%
    \PY@it{\PY@bf{\PY@ff{#1}}}}}}}
\def\PY#1#2{\PY@reset\PY@toks#1+\relax+\PY@do{#2}}

\@namedef{PY@tok@w}{\def\PY@tc##1{\textcolor[rgb]{0.73,0.73,0.73}{##1}}}
\@namedef{PY@tok@c}{\let\PY@it=\textit\def\PY@tc##1{\textcolor[rgb]{0.24,0.48,0.48}{##1}}}
\@namedef{PY@tok@cp}{\def\PY@tc##1{\textcolor[rgb]{0.61,0.40,0.00}{##1}}}
\@namedef{PY@tok@k}{\let\PY@bf=\textbf\def\PY@tc##1{\textcolor[rgb]{0.00,0.50,0.00}{##1}}}
\@namedef{PY@tok@kp}{\def\PY@tc##1{\textcolor[rgb]{0.00,0.50,0.00}{##1}}}
\@namedef{PY@tok@kt}{\def\PY@tc##1{\textcolor[rgb]{0.69,0.00,0.25}{##1}}}
\@namedef{PY@tok@o}{\def\PY@tc##1{\textcolor[rgb]{0.40,0.40,0.40}{##1}}}
\@namedef{PY@tok@ow}{\let\PY@bf=\textbf\def\PY@tc##1{\textcolor[rgb]{0.67,0.13,1.00}{##1}}}
\@namedef{PY@tok@nb}{\def\PY@tc##1{\textcolor[rgb]{0.00,0.50,0.00}{##1}}}
\@namedef{PY@tok@nf}{\def\PY@tc##1{\textcolor[rgb]{0.00,0.00,1.00}{##1}}}
\@namedef{PY@tok@nc}{\let\PY@bf=\textbf\def\PY@tc##1{\textcolor[rgb]{0.00,0.00,1.00}{##1}}}
\@namedef{PY@tok@nn}{\let\PY@bf=\textbf\def\PY@tc##1{\textcolor[rgb]{0.00,0.00,1.00}{##1}}}
\@namedef{PY@tok@ne}{\let\PY@bf=\textbf\def\PY@tc##1{\textcolor[rgb]{0.80,0.25,0.22}{##1}}}
\@namedef{PY@tok@nv}{\def\PY@tc##1{\textcolor[rgb]{0.10,0.09,0.49}{##1}}}
\@namedef{PY@tok@no}{\def\PY@tc##1{\textcolor[rgb]{0.53,0.00,0.00}{##1}}}
\@namedef{PY@tok@nl}{\def\PY@tc##1{\textcolor[rgb]{0.46,0.46,0.00}{##1}}}
\@namedef{PY@tok@ni}{\let\PY@bf=\textbf\def\PY@tc##1{\textcolor[rgb]{0.44,0.44,0.44}{##1}}}
\@namedef{PY@tok@na}{\def\PY@tc##1{\textcolor[rgb]{0.41,0.47,0.13}{##1}}}
\@namedef{PY@tok@nt}{\let\PY@bf=\textbf\def\PY@tc##1{\textcolor[rgb]{0.00,0.50,0.00}{##1}}}
\@namedef{PY@tok@nd}{\def\PY@tc##1{\textcolor[rgb]{0.67,0.13,1.00}{##1}}}
\@namedef{PY@tok@s}{\def\PY@tc##1{\textcolor[rgb]{0.73,0.13,0.13}{##1}}}
\@namedef{PY@tok@sd}{\let\PY@it=\textit\def\PY@tc##1{\textcolor[rgb]{0.73,0.13,0.13}{##1}}}
\@namedef{PY@tok@si}{\let\PY@bf=\textbf\def\PY@tc##1{\textcolor[rgb]{0.64,0.35,0.47}{##1}}}
\@namedef{PY@tok@se}{\let\PY@bf=\textbf\def\PY@tc##1{\textcolor[rgb]{0.67,0.36,0.12}{##1}}}
\@namedef{PY@tok@sr}{\def\PY@tc##1{\textcolor[rgb]{0.64,0.35,0.47}{##1}}}
\@namedef{PY@tok@ss}{\def\PY@tc##1{\textcolor[rgb]{0.10,0.09,0.49}{##1}}}
\@namedef{PY@tok@sx}{\def\PY@tc##1{\textcolor[rgb]{0.00,0.50,0.00}{##1}}}
\@namedef{PY@tok@m}{\def\PY@tc##1{\textcolor[rgb]{0.40,0.40,0.40}{##1}}}
\@namedef{PY@tok@gh}{\let\PY@bf=\textbf\def\PY@tc##1{\textcolor[rgb]{0.00,0.00,0.50}{##1}}}
\@namedef{PY@tok@gu}{\let\PY@bf=\textbf\def\PY@tc##1{\textcolor[rgb]{0.50,0.00,0.50}{##1}}}
\@namedef{PY@tok@gd}{\def\PY@tc##1{\textcolor[rgb]{0.63,0.00,0.00}{##1}}}
\@namedef{PY@tok@gi}{\def\PY@tc##1{\textcolor[rgb]{0.00,0.52,0.00}{##1}}}
\@namedef{PY@tok@gr}{\def\PY@tc##1{\textcolor[rgb]{0.89,0.00,0.00}{##1}}}
\@namedef{PY@tok@ge}{\let\PY@it=\textit}
\@namedef{PY@tok@gs}{\let\PY@bf=\textbf}
\@namedef{PY@tok@ges}{\let\PY@bf=\textbf\let\PY@it=\textit}
\@namedef{PY@tok@gp}{\let\PY@bf=\textbf\def\PY@tc##1{\textcolor[rgb]{0.00,0.00,0.50}{##1}}}
\@namedef{PY@tok@go}{\def\PY@tc##1{\textcolor[rgb]{0.44,0.44,0.44}{##1}}}
\@namedef{PY@tok@gt}{\def\PY@tc##1{\textcolor[rgb]{0.00,0.27,0.87}{##1}}}
\@namedef{PY@tok@err}{\def\PY@bc##1{{\setlength{\fboxsep}{\string -\fboxrule}\fcolorbox[rgb]{1.00,0.00,0.00}{1,1,1}{\strut ##1}}}}
\@namedef{PY@tok@kc}{\let\PY@bf=\textbf\def\PY@tc##1{\textcolor[rgb]{0.00,0.50,0.00}{##1}}}
\@namedef{PY@tok@kd}{\let\PY@bf=\textbf\def\PY@tc##1{\textcolor[rgb]{0.00,0.50,0.00}{##1}}}
\@namedef{PY@tok@kn}{\let\PY@bf=\textbf\def\PY@tc##1{\textcolor[rgb]{0.00,0.50,0.00}{##1}}}
\@namedef{PY@tok@kr}{\let\PY@bf=\textbf\def\PY@tc##1{\textcolor[rgb]{0.00,0.50,0.00}{##1}}}
\@namedef{PY@tok@bp}{\def\PY@tc##1{\textcolor[rgb]{0.00,0.50,0.00}{##1}}}
\@namedef{PY@tok@fm}{\def\PY@tc##1{\textcolor[rgb]{0.00,0.00,1.00}{##1}}}
\@namedef{PY@tok@vc}{\def\PY@tc##1{\textcolor[rgb]{0.10,0.09,0.49}{##1}}}
\@namedef{PY@tok@vg}{\def\PY@tc##1{\textcolor[rgb]{0.10,0.09,0.49}{##1}}}
\@namedef{PY@tok@vi}{\def\PY@tc##1{\textcolor[rgb]{0.10,0.09,0.49}{##1}}}
\@namedef{PY@tok@vm}{\def\PY@tc##1{\textcolor[rgb]{0.10,0.09,0.49}{##1}}}
\@namedef{PY@tok@sa}{\def\PY@tc##1{\textcolor[rgb]{0.73,0.13,0.13}{##1}}}
\@namedef{PY@tok@sb}{\def\PY@tc##1{\textcolor[rgb]{0.73,0.13,0.13}{##1}}}
\@namedef{PY@tok@sc}{\def\PY@tc##1{\textcolor[rgb]{0.73,0.13,0.13}{##1}}}
\@namedef{PY@tok@dl}{\def\PY@tc##1{\textcolor[rgb]{0.73,0.13,0.13}{##1}}}
\@namedef{PY@tok@s2}{\def\PY@tc##1{\textcolor[rgb]{0.73,0.13,0.13}{##1}}}
\@namedef{PY@tok@sh}{\def\PY@tc##1{\textcolor[rgb]{0.73,0.13,0.13}{##1}}}
\@namedef{PY@tok@s1}{\def\PY@tc##1{\textcolor[rgb]{0.73,0.13,0.13}{##1}}}
\@namedef{PY@tok@mb}{\def\PY@tc##1{\textcolor[rgb]{0.40,0.40,0.40}{##1}}}
\@namedef{PY@tok@mf}{\def\PY@tc##1{\textcolor[rgb]{0.40,0.40,0.40}{##1}}}
\@namedef{PY@tok@mh}{\def\PY@tc##1{\textcolor[rgb]{0.40,0.40,0.40}{##1}}}
\@namedef{PY@tok@mi}{\def\PY@tc##1{\textcolor[rgb]{0.40,0.40,0.40}{##1}}}
\@namedef{PY@tok@il}{\def\PY@tc##1{\textcolor[rgb]{0.40,0.40,0.40}{##1}}}
\@namedef{PY@tok@mo}{\def\PY@tc##1{\textcolor[rgb]{0.40,0.40,0.40}{##1}}}
\@namedef{PY@tok@ch}{\let\PY@it=\textit\def\PY@tc##1{\textcolor[rgb]{0.24,0.48,0.48}{##1}}}
\@namedef{PY@tok@cm}{\let\PY@it=\textit\def\PY@tc##1{\textcolor[rgb]{0.24,0.48,0.48}{##1}}}
\@namedef{PY@tok@cpf}{\let\PY@it=\textit\def\PY@tc##1{\textcolor[rgb]{0.24,0.48,0.48}{##1}}}
\@namedef{PY@tok@c1}{\let\PY@it=\textit\def\PY@tc##1{\textcolor[rgb]{0.24,0.48,0.48}{##1}}}
\@namedef{PY@tok@cs}{\let\PY@it=\textit\def\PY@tc##1{\textcolor[rgb]{0.24,0.48,0.48}{##1}}}

\def\PYZbs{\char`\\}
\def\PYZus{\char`\_}
\def\PYZob{\char`\{}
\def\PYZcb{\char`\}}
\def\PYZca{\char`\^}
\def\PYZam{\char`\&}
\def\PYZlt{\char`\<}
\def\PYZgt{\char`\>}
\def\PYZsh{\char`\#}
\def\PYZpc{\char`\%}
\def\PYZdl{\char`\$}
\def\PYZhy{\char`\-}
\def\PYZsq{\char`\'}
\def\PYZdq{\char`\"}
\def\PYZti{\char`\~}
% for compatibility with earlier versions
\def\PYZat{@}
\def\PYZlb{[}
\def\PYZrb{]}
\makeatother


    % For linebreaks inside Verbatim environment from package fancyvrb.
    \makeatletter
        \newbox\Wrappedcontinuationbox
        \newbox\Wrappedvisiblespacebox
        \newcommand*\Wrappedvisiblespace {\textcolor{red}{\textvisiblespace}}
        \newcommand*\Wrappedcontinuationsymbol {\textcolor{red}{\llap{\tiny$\m@th\hookrightarrow$}}}
        \newcommand*\Wrappedcontinuationindent {3ex }
        \newcommand*\Wrappedafterbreak {\kern\Wrappedcontinuationindent\copy\Wrappedcontinuationbox}
        % Take advantage of the already applied Pygments mark-up to insert
        % potential linebreaks for TeX processing.
        %        {, <, #, %, $, ' and ": go to next line.
        %        _, }, ^, &, >, - and ~: stay at end of broken line.
        % Use of \textquotesingle for straight quote.
        \newcommand*\Wrappedbreaksatspecials {%
            \def\PYGZus{\discretionary{\char`\_}{\Wrappedafterbreak}{\char`\_}}%
            \def\PYGZob{\discretionary{}{\Wrappedafterbreak\char`\{}{\char`\{}}%
            \def\PYGZcb{\discretionary{\char`\}}{\Wrappedafterbreak}{\char`\}}}%
            \def\PYGZca{\discretionary{\char`\^}{\Wrappedafterbreak}{\char`\^}}%
            \def\PYGZam{\discretionary{\char`\&}{\Wrappedafterbreak}{\char`\&}}%
            \def\PYGZlt{\discretionary{}{\Wrappedafterbreak\char`\<}{\char`\<}}%
            \def\PYGZgt{\discretionary{\char`\>}{\Wrappedafterbreak}{\char`\>}}%
            \def\PYGZsh{\discretionary{}{\Wrappedafterbreak\char`\#}{\char`\#}}%
            \def\PYGZpc{\discretionary{}{\Wrappedafterbreak\char`\%}{\char`\%}}%
            \def\PYGZdl{\discretionary{}{\Wrappedafterbreak\char`\$}{\char`\$}}%
            \def\PYGZhy{\discretionary{\char`\-}{\Wrappedafterbreak}{\char`\-}}%
            \def\PYGZsq{\discretionary{}{\Wrappedafterbreak\textquotesingle}{\textquotesingle}}%
            \def\PYGZdq{\discretionary{}{\Wrappedafterbreak\char`\"}{\char`\"}}%
            \def\PYGZti{\discretionary{\char`\~}{\Wrappedafterbreak}{\char`\~}}%
        }
        % Some characters . , ; ? ! / are not pygmentized.
        % This macro makes them "active" and they will insert potential linebreaks
        \newcommand*\Wrappedbreaksatpunct {%
            \lccode`\~`\.\lowercase{\def~}{\discretionary{\hbox{\char`\.}}{\Wrappedafterbreak}{\hbox{\char`\.}}}%
            \lccode`\~`\,\lowercase{\def~}{\discretionary{\hbox{\char`\,}}{\Wrappedafterbreak}{\hbox{\char`\,}}}%
            \lccode`\~`\;\lowercase{\def~}{\discretionary{\hbox{\char`\;}}{\Wrappedafterbreak}{\hbox{\char`\;}}}%
            \lccode`\~`\:\lowercase{\def~}{\discretionary{\hbox{\char`\:}}{\Wrappedafterbreak}{\hbox{\char`\:}}}%
            \lccode`\~`\?\lowercase{\def~}{\discretionary{\hbox{\char`\?}}{\Wrappedafterbreak}{\hbox{\char`\?}}}%
            \lccode`\~`\!\lowercase{\def~}{\discretionary{\hbox{\char`\!}}{\Wrappedafterbreak}{\hbox{\char`\!}}}%
            \lccode`\~`\/\lowercase{\def~}{\discretionary{\hbox{\char`\/}}{\Wrappedafterbreak}{\hbox{\char`\/}}}%
            \catcode`\.\active
            \catcode`\,\active
            \catcode`\;\active
            \catcode`\:\active
            \catcode`\?\active
            \catcode`\!\active
            \catcode`\/\active
            \lccode`\~`\~
        }
    \makeatother

    \let\OriginalVerbatim=\Verbatim
    \makeatletter
    \renewcommand{\Verbatim}[1][1]{%
        %\parskip\z@skip
        \sbox\Wrappedcontinuationbox {\Wrappedcontinuationsymbol}%
        \sbox\Wrappedvisiblespacebox {\FV@SetupFont\Wrappedvisiblespace}%
        \def\FancyVerbFormatLine ##1{\hsize\linewidth
            \vtop{\raggedright\hyphenpenalty\z@\exhyphenpenalty\z@
                \doublehyphendemerits\z@\finalhyphendemerits\z@
                \strut ##1\strut}%
        }%
        % If the linebreak is at a space, the latter will be displayed as visible
        % space at end of first line, and a continuation symbol starts next line.
        % Stretch/shrink are however usually zero for typewriter font.
        \def\FV@Space {%
            \nobreak\hskip\z@ plus\fontdimen3\font minus\fontdimen4\font
            \discretionary{\copy\Wrappedvisiblespacebox}{\Wrappedafterbreak}
            {\kern\fontdimen2\font}%
        }%

        % Allow breaks at special characters using \PYG... macros.
        \Wrappedbreaksatspecials
        % Breaks at punctuation characters . , ; ? ! and / need catcode=\active
        \OriginalVerbatim[#1,codes*=\Wrappedbreaksatpunct]%
    }
    \makeatother

    % Exact colors from NB
    \definecolor{incolor}{HTML}{303F9F}
    \definecolor{outcolor}{HTML}{D84315}
    \definecolor{cellborder}{HTML}{CFCFCF}
    \definecolor{cellbackground}{HTML}{F7F7F7}

    % prompt
    \makeatletter
    \newcommand{\boxspacing}{\kern\kvtcb@left@rule\kern\kvtcb@boxsep}
    \makeatother
    \newcommand{\prompt}[4]{
        {\ttfamily\llap{{\color{#2}[#3]:\hspace{3pt}#4}}\vspace{-\baselineskip}}
    }
    

    
    % Prevent overflowing lines due to hard-to-break entities
    \sloppy
    % Setup hyperref package
    \hypersetup{
      breaklinks=true,  % so long urls are correctly broken across lines
      colorlinks=true,
      urlcolor=urlcolor,
      linkcolor=linkcolor,
      citecolor=citecolor,
      }
    % Slightly bigger margins than the latex defaults
    
    \geometry{verbose,tmargin=1in,bmargin=1in,lmargin=1in,rmargin=1in}
    
    

\begin{document}
    
    \maketitle
    
    

    
    \section{Reporte Final: Experimentación con Datos
Médicos}\label{reporte-final-experimentaciuxf3n-con-datos-muxe9dicos}

Este reporte reúne el análisis experimental del proyecto en un documento
reproducible y autónomo. El objetivo es entender un conjunto de datos
médicos anonimizado y comparar enfoques para predecir la variable
objetivo \texttt{condition}, manteniendo una lectura clínica de los
errores.

La idea central no es construir un único flujo cerrado, sino documentar
experimentación: qué se probó, por qué se probó, qué funcionó, qué falló
y qué implican los resultados. En un contexto clínico, un falso negativo
es especialmente grave porque equivale a clasificar como sano a un
paciente que sí presenta la condición.

    \subsection{Preparación del entorno y modo de
ejecución}\label{preparaciuxf3n-del-entorno-y-modo-de-ejecuciuxf3n}

Esta celda define rutas relativas desde la raíz del proyecto, activa un
modo rápido con caché y carga librerías. El reporte usa por defecto los
artefactos ya generados en \texttt{mundial\_project/outputs/} para
evitar relanzar búsquedas largas de hiperparámetros cada vez que se abre
el archivo.

Si se desea recalcular experimentos pesados, las banderas
\texttt{RUN\_TASK*\_EXPERIMENTS} pueden cambiarse a \texttt{True}, pero
para una entrega estable se mantienen en \texttt{False}.

    \begin{tcolorbox}[breakable, size=fbox, boxrule=1pt, pad at break*=1mm,colback=cellbackground, colframe=cellborder]
\prompt{In}{incolor}{6}{\boxspacing}
\begin{Verbatim}[commandchars=\\\{\}]
\PY{k+kn}{from}\PY{+w}{ }\PY{n+nn}{pathlib}\PY{+w}{ }\PY{k+kn}{import} \PY{n}{Path}
\PY{k+kn}{import}\PY{+w}{ }\PY{n+nn}{os}
\PY{k+kn}{import}\PY{+w}{ }\PY{n+nn}{sys}
\PY{k+kn}{import}\PY{+w}{ }\PY{n+nn}{textwrap}
\PY{k+kn}{import}\PY{+w}{ }\PY{n+nn}{warnings}

\PY{n}{os}\PY{o}{.}\PY{n}{environ}\PY{o}{.}\PY{n}{setdefault}\PY{p}{(}\PY{l+s+s2}{\PYZdq{}}\PY{l+s+s2}{LOKY\PYZus{}MAX\PYZus{}CPU\PYZus{}COUNT}\PY{l+s+s2}{\PYZdq{}}\PY{p}{,} \PY{l+s+s2}{\PYZdq{}}\PY{l+s+s2}{4}\PY{l+s+s2}{\PYZdq{}}\PY{p}{)}
\PY{n}{warnings}\PY{o}{.}\PY{n}{filterwarnings}\PY{p}{(}\PY{l+s+s2}{\PYZdq{}}\PY{l+s+s2}{ignore}\PY{l+s+s2}{\PYZdq{}}\PY{p}{)}

\PY{k+kn}{import}\PY{+w}{ }\PY{n+nn}{matplotlib}\PY{n+nn}{.}\PY{n+nn}{pyplot}\PY{+w}{ }\PY{k}{as}\PY{+w}{ }\PY{n+nn}{plt}
\PY{k+kn}{import}\PY{+w}{ }\PY{n+nn}{numpy}\PY{+w}{ }\PY{k}{as}\PY{+w}{ }\PY{n+nn}{np}
\PY{k+kn}{import}\PY{+w}{ }\PY{n+nn}{pandas}\PY{+w}{ }\PY{k}{as}\PY{+w}{ }\PY{n+nn}{pd}
\PY{k+kn}{import}\PY{+w}{ }\PY{n+nn}{seaborn}\PY{+w}{ }\PY{k}{as}\PY{+w}{ }\PY{n+nn}{sns}
\PY{k+kn}{from}\PY{+w}{ }\PY{n+nn}{IPython}\PY{n+nn}{.}\PY{n+nn}{display}\PY{+w}{ }\PY{k+kn}{import} \PY{n}{Image}\PY{p}{,} \PY{n}{Markdown}\PY{p}{,} \PY{n}{display}

\PY{o}{\PYZpc{}}\PY{k}{matplotlib} inline

\PY{n}{CURRENT\PYZus{}DIR} \PY{o}{=} \PY{n}{Path}\PY{o}{.}\PY{n}{cwd}\PY{p}{(}\PY{p}{)}\PY{o}{.}\PY{n}{resolve}\PY{p}{(}\PY{p}{)}
\PY{n}{PROJECT\PYZus{}ROOT\PYZus{}CANDIDATES} \PY{o}{=} \PY{p}{[}\PY{n}{CURRENT\PYZus{}DIR}\PY{p}{,} \PY{o}{*}\PY{n}{CURRENT\PYZus{}DIR}\PY{o}{.}\PY{n}{parents}\PY{p}{]}
\PY{n}{PROJECT\PYZus{}ROOT} \PY{o}{=} \PY{n+nb}{next}\PY{p}{(}
    \PY{p}{(}
        \PY{n}{path}
        \PY{k}{for} \PY{n}{path} \PY{o+ow}{in} \PY{n}{PROJECT\PYZus{}ROOT\PYZus{}CANDIDATES}
        \PY{k}{if} \PY{p}{(}\PY{n}{path} \PY{o}{/} \PY{l+s+s2}{\PYZdq{}}\PY{l+s+s2}{mundial\PYZus{}project}\PY{l+s+s2}{\PYZdq{}}\PY{p}{)}\PY{o}{.}\PY{n}{exists}\PY{p}{(}\PY{p}{)} \PY{o+ow}{and} \PY{p}{(}\PY{n}{path} \PY{o}{/} \PY{l+s+s2}{\PYZdq{}}\PY{l+s+s2}{Mundial}\PY{l+s+s2}{\PYZdq{}} \PY{o}{/} \PY{l+s+s2}{\PYZdq{}}\PY{l+s+s2}{TASK.md}\PY{l+s+s2}{\PYZdq{}}\PY{p}{)}\PY{o}{.}\PY{n}{exists}\PY{p}{(}\PY{p}{)}
    \PY{p}{)}\PY{p}{,}
    \PY{n}{CURRENT\PYZus{}DIR}\PY{p}{,}
\PY{p}{)}

\PY{n}{MUNDIAL\PYZus{}DIR} \PY{o}{=} \PY{n}{PROJECT\PYZus{}ROOT} \PY{o}{/} \PY{l+s+s2}{\PYZdq{}}\PY{l+s+s2}{Mundial}\PY{l+s+s2}{\PYZdq{}}
\PY{n}{PROJECT\PYZus{}DIR} \PY{o}{=} \PY{n}{PROJECT\PYZus{}ROOT} \PY{o}{/} \PY{l+s+s2}{\PYZdq{}}\PY{l+s+s2}{mundial\PYZus{}project}\PY{l+s+s2}{\PYZdq{}}
\PY{n}{DATA\PYZus{}DIR} \PY{o}{=} \PY{n}{PROJECT\PYZus{}DIR} \PY{o}{/} \PY{l+s+s2}{\PYZdq{}}\PY{l+s+s2}{data}\PY{l+s+s2}{\PYZdq{}}
\PY{n}{PROCESSED\PYZus{}DIR} \PY{o}{=} \PY{n}{DATA\PYZus{}DIR} \PY{o}{/} \PY{l+s+s2}{\PYZdq{}}\PY{l+s+s2}{processed}\PY{l+s+s2}{\PYZdq{}}
\PY{n}{OUTPUTS\PYZus{}DIR} \PY{o}{=} \PY{n}{PROJECT\PYZus{}DIR} \PY{o}{/} \PY{l+s+s2}{\PYZdq{}}\PY{l+s+s2}{outputs}\PY{l+s+s2}{\PYZdq{}}
\PY{n}{REPORTS\PYZus{}DIR} \PY{o}{=} \PY{n}{PROJECT\PYZus{}DIR} \PY{o}{/} \PY{l+s+s2}{\PYZdq{}}\PY{l+s+s2}{reports}\PY{l+s+s2}{\PYZdq{}}
\PY{n}{SRC\PYZus{}DIR} \PY{o}{=} \PY{n}{PROJECT\PYZus{}DIR} \PY{o}{/} \PY{l+s+s2}{\PYZdq{}}\PY{l+s+s2}{src}\PY{l+s+s2}{\PYZdq{}}

\PY{k}{if} \PY{n+nb}{str}\PY{p}{(}\PY{n}{SRC\PYZus{}DIR}\PY{p}{)} \PY{o+ow}{not} \PY{o+ow}{in} \PY{n}{sys}\PY{o}{.}\PY{n}{path}\PY{p}{:}
    \PY{n}{sys}\PY{o}{.}\PY{n}{path}\PY{o}{.}\PY{n}{insert}\PY{p}{(}\PY{l+m+mi}{0}\PY{p}{,} \PY{n+nb}{str}\PY{p}{(}\PY{n}{SRC\PYZus{}DIR}\PY{p}{)}\PY{p}{)}


\PY{k}{def}\PY{+w}{ }\PY{n+nf}{first\PYZus{}existing\PYZus{}path}\PY{p}{(}\PY{n}{candidates}\PY{p}{,} \PY{n}{description}\PY{p}{,} \PY{n}{required}\PY{o}{=}\PY{k+kc}{True}\PY{p}{)}\PY{p}{:}
    \PY{k}{for} \PY{n}{candidate} \PY{o+ow}{in} \PY{n}{candidates}\PY{p}{:}
        \PY{n}{path} \PY{o}{=} \PY{n}{Path}\PY{p}{(}\PY{n}{candidate}\PY{p}{)}\PY{o}{.}\PY{n}{resolve}\PY{p}{(}\PY{p}{)}
        \PY{k}{if} \PY{n}{path}\PY{o}{.}\PY{n}{exists}\PY{p}{(}\PY{p}{)}\PY{p}{:}
            \PY{k}{return} \PY{n}{path}
    \PY{k}{if} \PY{n}{required}\PY{p}{:}
        \PY{n}{checked} \PY{o}{=} \PY{l+s+s2}{\PYZdq{}}\PY{l+s+se}{\PYZbs{}n}\PY{l+s+s2}{\PYZdq{}}\PY{o}{.}\PY{n}{join}\PY{p}{(}\PY{l+s+sa}{f}\PY{l+s+s2}{\PYZdq{}}\PY{l+s+s2}{\PYZhy{} }\PY{l+s+si}{\PYZob{}}\PY{n}{Path}\PY{p}{(}\PY{n}{candidate}\PY{p}{)}\PY{l+s+si}{\PYZcb{}}\PY{l+s+s2}{\PYZdq{}} \PY{k}{for} \PY{n}{candidate} \PY{o+ow}{in} \PY{n}{candidates}\PY{p}{)}
        \PY{k}{raise} \PY{n+ne}{FileNotFoundError}\PY{p}{(}\PY{l+s+sa}{f}\PY{l+s+s2}{\PYZdq{}}\PY{l+s+s2}{No se encontro }\PY{l+s+si}{\PYZob{}}\PY{n}{description}\PY{l+s+si}{\PYZcb{}}\PY{l+s+s2}{. Rutas revisadas:}\PY{l+s+se}{\PYZbs{}n}\PY{l+s+si}{\PYZob{}}\PY{n}{checked}\PY{l+s+si}{\PYZcb{}}\PY{l+s+s2}{\PYZdq{}}\PY{p}{)}
    \PY{k}{return} \PY{n}{Path}\PY{p}{(}\PY{n}{candidates}\PY{p}{[}\PY{l+m+mi}{0}\PY{p}{]}\PY{p}{)}\PY{o}{.}\PY{n}{resolve}\PY{p}{(}\PY{p}{)}


\PY{n}{RAW\PYZus{}CSV} \PY{o}{=} \PY{n}{first\PYZus{}existing\PYZus{}path}\PY{p}{(}
    \PY{p}{[}
        \PY{n}{PROJECT\PYZus{}DIR} \PY{o}{/} \PY{l+s+s2}{\PYZdq{}}\PY{l+s+s2}{data}\PY{l+s+s2}{\PYZdq{}} \PY{o}{/} \PY{l+s+s2}{\PYZdq{}}\PY{l+s+s2}{foot.csv}\PY{l+s+s2}{\PYZdq{}}\PY{p}{,}
        \PY{n}{PROJECT\PYZus{}ROOT} \PY{o}{/} \PY{l+s+s2}{\PYZdq{}}\PY{l+s+s2}{data}\PY{l+s+s2}{\PYZdq{}} \PY{o}{/} \PY{l+s+s2}{\PYZdq{}}\PY{l+s+s2}{foot.csv}\PY{l+s+s2}{\PYZdq{}}\PY{p}{,}
        \PY{n}{CURRENT\PYZus{}DIR} \PY{o}{/} \PY{l+s+s2}{\PYZdq{}}\PY{l+s+s2}{mundial\PYZus{}project}\PY{l+s+s2}{\PYZdq{}} \PY{o}{/} \PY{l+s+s2}{\PYZdq{}}\PY{l+s+s2}{data}\PY{l+s+s2}{\PYZdq{}} \PY{o}{/} \PY{l+s+s2}{\PYZdq{}}\PY{l+s+s2}{foot.csv}\PY{l+s+s2}{\PYZdq{}}\PY{p}{,}
        \PY{n}{CURRENT\PYZus{}DIR} \PY{o}{/} \PY{l+s+s2}{\PYZdq{}}\PY{l+s+s2}{data}\PY{l+s+s2}{\PYZdq{}} \PY{o}{/} \PY{l+s+s2}{\PYZdq{}}\PY{l+s+s2}{foot.csv}\PY{l+s+s2}{\PYZdq{}}\PY{p}{,}
        \PY{n}{CURRENT\PYZus{}DIR} \PY{o}{/} \PY{l+s+s2}{\PYZdq{}}\PY{l+s+s2}{foot.csv}\PY{l+s+s2}{\PYZdq{}}\PY{p}{,}
        \PY{n}{CURRENT\PYZus{}DIR} \PY{o}{/} \PY{l+s+s2}{\PYZdq{}}\PY{l+s+s2}{foot(1).csv}\PY{l+s+s2}{\PYZdq{}}\PY{p}{,}
    \PY{p}{]}\PY{p}{,}
    \PY{l+s+s2}{\PYZdq{}}\PY{l+s+s2}{el dataset original foot.csv}\PY{l+s+s2}{\PYZdq{}}\PY{p}{,}
\PY{p}{)}
\PY{n}{CLEAN\PYZus{}CSV} \PY{o}{=} \PY{n}{first\PYZus{}existing\PYZus{}path}\PY{p}{(}
    \PY{p}{[}
        \PY{n}{PROJECT\PYZus{}DIR} \PY{o}{/} \PY{l+s+s2}{\PYZdq{}}\PY{l+s+s2}{data}\PY{l+s+s2}{\PYZdq{}} \PY{o}{/} \PY{l+s+s2}{\PYZdq{}}\PY{l+s+s2}{processed}\PY{l+s+s2}{\PYZdq{}} \PY{o}{/} \PY{l+s+s2}{\PYZdq{}}\PY{l+s+s2}{foot\PYZus{}clean.csv}\PY{l+s+s2}{\PYZdq{}}\PY{p}{,}
        \PY{n}{PROJECT\PYZus{}ROOT} \PY{o}{/} \PY{l+s+s2}{\PYZdq{}}\PY{l+s+s2}{data}\PY{l+s+s2}{\PYZdq{}} \PY{o}{/} \PY{l+s+s2}{\PYZdq{}}\PY{l+s+s2}{processed}\PY{l+s+s2}{\PYZdq{}} \PY{o}{/} \PY{l+s+s2}{\PYZdq{}}\PY{l+s+s2}{foot\PYZus{}clean.csv}\PY{l+s+s2}{\PYZdq{}}\PY{p}{,}
        \PY{n}{CURRENT\PYZus{}DIR} \PY{o}{/} \PY{l+s+s2}{\PYZdq{}}\PY{l+s+s2}{data}\PY{l+s+s2}{\PYZdq{}} \PY{o}{/} \PY{l+s+s2}{\PYZdq{}}\PY{l+s+s2}{processed}\PY{l+s+s2}{\PYZdq{}} \PY{o}{/} \PY{l+s+s2}{\PYZdq{}}\PY{l+s+s2}{foot\PYZus{}clean.csv}\PY{l+s+s2}{\PYZdq{}}\PY{p}{,}
        \PY{n}{CURRENT\PYZus{}DIR} \PY{o}{/} \PY{l+s+s2}{\PYZdq{}}\PY{l+s+s2}{foot\PYZus{}clean.csv}\PY{l+s+s2}{\PYZdq{}}\PY{p}{,}
    \PY{p}{]}\PY{p}{,}
    \PY{l+s+s2}{\PYZdq{}}\PY{l+s+s2}{el dataset limpio foot\PYZus{}clean.csv}\PY{l+s+s2}{\PYZdq{}}\PY{p}{,}
    \PY{n}{required}\PY{o}{=}\PY{k+kc}{False}\PY{p}{,}
\PY{p}{)}
\PY{n}{MODEL\PYZus{}READY\PYZus{}CSV} \PY{o}{=} \PY{n}{first\PYZus{}existing\PYZus{}path}\PY{p}{(}
    \PY{p}{[}
        \PY{n}{PROJECT\PYZus{}DIR} \PY{o}{/} \PY{l+s+s2}{\PYZdq{}}\PY{l+s+s2}{data}\PY{l+s+s2}{\PYZdq{}} \PY{o}{/} \PY{l+s+s2}{\PYZdq{}}\PY{l+s+s2}{processed}\PY{l+s+s2}{\PYZdq{}} \PY{o}{/} \PY{l+s+s2}{\PYZdq{}}\PY{l+s+s2}{foot\PYZus{}model\PYZus{}ready.csv}\PY{l+s+s2}{\PYZdq{}}\PY{p}{,}
        \PY{n}{PROJECT\PYZus{}ROOT} \PY{o}{/} \PY{l+s+s2}{\PYZdq{}}\PY{l+s+s2}{data}\PY{l+s+s2}{\PYZdq{}} \PY{o}{/} \PY{l+s+s2}{\PYZdq{}}\PY{l+s+s2}{processed}\PY{l+s+s2}{\PYZdq{}} \PY{o}{/} \PY{l+s+s2}{\PYZdq{}}\PY{l+s+s2}{foot\PYZus{}model\PYZus{}ready.csv}\PY{l+s+s2}{\PYZdq{}}\PY{p}{,}
        \PY{n}{CURRENT\PYZus{}DIR} \PY{o}{/} \PY{l+s+s2}{\PYZdq{}}\PY{l+s+s2}{data}\PY{l+s+s2}{\PYZdq{}} \PY{o}{/} \PY{l+s+s2}{\PYZdq{}}\PY{l+s+s2}{processed}\PY{l+s+s2}{\PYZdq{}} \PY{o}{/} \PY{l+s+s2}{\PYZdq{}}\PY{l+s+s2}{foot\PYZus{}model\PYZus{}ready.csv}\PY{l+s+s2}{\PYZdq{}}\PY{p}{,}
        \PY{n}{CURRENT\PYZus{}DIR} \PY{o}{/} \PY{l+s+s2}{\PYZdq{}}\PY{l+s+s2}{foot\PYZus{}model\PYZus{}ready.csv}\PY{l+s+s2}{\PYZdq{}}\PY{p}{,}
    \PY{p}{]}\PY{p}{,}
    \PY{l+s+s2}{\PYZdq{}}\PY{l+s+s2}{el dataset estandarizado foot\PYZus{}model\PYZus{}ready.csv}\PY{l+s+s2}{\PYZdq{}}\PY{p}{,}
    \PY{n}{required}\PY{o}{=}\PY{k+kc}{False}\PY{p}{,}
\PY{p}{)}

\PY{n}{RUN\PYZus{}TASK1\PYZus{}EXPERIMENTS} \PY{o}{=} \PY{k+kc}{False}
\PY{n}{RUN\PYZus{}TASK2\PYZus{}EXPERIMENTS} \PY{o}{=} \PY{k+kc}{False}
\PY{n}{RUN\PYZus{}TASK3\PYZus{}EXPERIMENTS} \PY{o}{=} \PY{k+kc}{False}
\PY{n}{RUN\PYZus{}TASK4\PYZus{}EXPERIMENTS} \PY{o}{=} \PY{k+kc}{False}

\PY{n}{pd}\PY{o}{.}\PY{n}{set\PYZus{}option}\PY{p}{(}\PY{l+s+s2}{\PYZdq{}}\PY{l+s+s2}{display.max\PYZus{}columns}\PY{l+s+s2}{\PYZdq{}}\PY{p}{,} \PY{l+m+mi}{120}\PY{p}{)}
\PY{n}{pd}\PY{o}{.}\PY{n}{set\PYZus{}option}\PY{p}{(}\PY{l+s+s2}{\PYZdq{}}\PY{l+s+s2}{display.width}\PY{l+s+s2}{\PYZdq{}}\PY{p}{,} \PY{l+m+mi}{160}\PY{p}{)}
\PY{n}{sns}\PY{o}{.}\PY{n}{set\PYZus{}theme}\PY{p}{(}\PY{n}{style}\PY{o}{=}\PY{l+s+s2}{\PYZdq{}}\PY{l+s+s2}{whitegrid}\PY{l+s+s2}{\PYZdq{}}\PY{p}{,} \PY{n}{context}\PY{o}{=}\PY{l+s+s2}{\PYZdq{}}\PY{l+s+s2}{notebook}\PY{l+s+s2}{\PYZdq{}}\PY{p}{)}

\PY{n+nb}{print}\PY{p}{(}\PY{l+s+sa}{f}\PY{l+s+s2}{\PYZdq{}}\PY{l+s+s2}{Directorio actual: }\PY{l+s+si}{\PYZob{}}\PY{n}{CURRENT\PYZus{}DIR}\PY{l+s+si}{\PYZcb{}}\PY{l+s+s2}{\PYZdq{}}\PY{p}{)}
\PY{n+nb}{print}\PY{p}{(}\PY{l+s+sa}{f}\PY{l+s+s2}{\PYZdq{}}\PY{l+s+s2}{Raiz del proyecto detectada: }\PY{l+s+si}{\PYZob{}}\PY{n}{PROJECT\PYZus{}ROOT}\PY{l+s+si}{\PYZcb{}}\PY{l+s+s2}{\PYZdq{}}\PY{p}{)}
\PY{n+nb}{print}\PY{p}{(}\PY{l+s+sa}{f}\PY{l+s+s2}{\PYZdq{}}\PY{l+s+s2}{Dataset original: }\PY{l+s+si}{\PYZob{}}\PY{n}{RAW\PYZus{}CSV}\PY{l+s+si}{\PYZcb{}}\PY{l+s+s2}{\PYZdq{}}\PY{p}{)}
\PY{n+nb}{print}\PY{p}{(}\PY{l+s+sa}{f}\PY{l+s+s2}{\PYZdq{}}\PY{l+s+s2}{Dataset limpio existe: }\PY{l+s+si}{\PYZob{}}\PY{n}{CLEAN\PYZus{}CSV}\PY{o}{.}\PY{n}{exists}\PY{p}{(}\PY{p}{)}\PY{l+s+si}{\PYZcb{}}\PY{l+s+s2}{ (}\PY{l+s+si}{\PYZob{}}\PY{n}{CLEAN\PYZus{}CSV}\PY{l+s+si}{\PYZcb{}}\PY{l+s+s2}{)}\PY{l+s+s2}{\PYZdq{}}\PY{p}{)}
\PY{n+nb}{print}\PY{p}{(}\PY{l+s+sa}{f}\PY{l+s+s2}{\PYZdq{}}\PY{l+s+s2}{Dataset model\PYZhy{}ready existe: }\PY{l+s+si}{\PYZob{}}\PY{n}{MODEL\PYZus{}READY\PYZus{}CSV}\PY{o}{.}\PY{n}{exists}\PY{p}{(}\PY{p}{)}\PY{l+s+si}{\PYZcb{}}\PY{l+s+s2}{ (}\PY{l+s+si}{\PYZob{}}\PY{n}{MODEL\PYZus{}READY\PYZus{}CSV}\PY{l+s+si}{\PYZcb{}}\PY{l+s+s2}{)}\PY{l+s+s2}{\PYZdq{}}\PY{p}{)}
\PY{n+nb}{print}\PY{p}{(}\PY{l+s+sa}{f}\PY{l+s+s2}{\PYZdq{}}\PY{l+s+s2}{Modo rapido con cache: }\PY{l+s+si}{\PYZob{}}\PY{o+ow}{not}\PY{+w}{ }\PY{n+nb}{any}\PY{p}{(}\PY{p}{[}\PY{n}{RUN\PYZus{}TASK1\PYZus{}EXPERIMENTS}\PY{p}{,}\PY{+w}{ }\PY{n}{RUN\PYZus{}TASK2\PYZus{}EXPERIMENTS}\PY{p}{,}\PY{+w}{ }\PY{n}{RUN\PYZus{}TASK3\PYZus{}EXPERIMENTS}\PY{p}{,}\PY{+w}{ }\PY{n}{RUN\PYZus{}TASK4\PYZus{}EXPERIMENTS}\PY{p}{]}\PY{p}{)}\PY{l+s+si}{\PYZcb{}}\PY{l+s+s2}{\PYZdq{}}\PY{p}{)}
\end{Verbatim}
\end{tcolorbox}

    \begin{Verbatim}[commandchars=\\\{\}]
Directorio actual: E:\textbackslash{}University\textbackslash{}4to\textbackslash{}ML
Raiz del proyecto detectada: E:\textbackslash{}University\textbackslash{}4to\textbackslash{}ML
Dataset original: E:\textbackslash{}University\textbackslash{}4to\textbackslash{}ML\textbackslash{}mundial\_project\textbackslash{}data\textbackslash{}foot.csv
Dataset limpio existe: True
(E:\textbackslash{}University\textbackslash{}4to\textbackslash{}ML\textbackslash{}mundial\_project\textbackslash{}data\textbackslash{}processed\textbackslash{}foot\_clean.csv)
Dataset model-ready existe: True
(E:\textbackslash{}University\textbackslash{}4to\textbackslash{}ML\textbackslash{}mundial\_project\textbackslash{}data\textbackslash{}processed\textbackslash{}foot\_model\_ready.csv)
Modo rapido con cache: True
    \end{Verbatim}

    \subsection{Carga de funciones
auxiliares}\label{carga-de-funciones-auxiliares}

    \begin{tcolorbox}[breakable, size=fbox, boxrule=1pt, pad at break*=1mm,colback=cellbackground, colframe=cellborder]
\prompt{In}{incolor}{7}{\boxspacing}
\begin{Verbatim}[commandchars=\\\{\}]
\PY{k}{def}\PY{+w}{ }\PY{n+nf}{read\PYZus{}csv}\PY{p}{(}\PY{n}{path}\PY{p}{,} \PY{o}{*}\PY{o}{*}\PY{n}{kwargs}\PY{p}{)}\PY{p}{:}
    \PY{n}{path} \PY{o}{=} \PY{n}{Path}\PY{p}{(}\PY{n}{path}\PY{p}{)}
    \PY{k}{if} \PY{o+ow}{not} \PY{n}{path}\PY{o}{.}\PY{n}{exists}\PY{p}{(}\PY{p}{)}\PY{p}{:}
        \PY{k}{raise} \PY{n+ne}{FileNotFoundError}\PY{p}{(}\PY{l+s+sa}{f}\PY{l+s+s2}{\PYZdq{}}\PY{l+s+s2}{No existe el archivo esperado: }\PY{l+s+si}{\PYZob{}}\PY{n}{path}\PY{l+s+si}{\PYZcb{}}\PY{l+s+s2}{\PYZdq{}}\PY{p}{)}
    \PY{k}{return} \PY{n}{pd}\PY{o}{.}\PY{n}{read\PYZus{}csv}\PY{p}{(}\PY{n}{path}\PY{p}{,} \PY{o}{*}\PY{o}{*}\PY{n}{kwargs}\PY{p}{)}


\PY{k}{def}\PY{+w}{ }\PY{n+nf}{show\PYZus{}image}\PY{p}{(}\PY{n}{path}\PY{p}{,} \PY{n}{width}\PY{o}{=}\PY{k+kc}{None}\PY{p}{)}\PY{p}{:}
    \PY{n}{path} \PY{o}{=} \PY{n}{Path}\PY{p}{(}\PY{n}{path}\PY{p}{)}
    \PY{k}{if} \PY{n}{path}\PY{o}{.}\PY{n}{exists}\PY{p}{(}\PY{p}{)}\PY{p}{:}
        \PY{n}{display}\PY{p}{(}\PY{n}{Image}\PY{p}{(}\PY{n}{filename}\PY{o}{=}\PY{n+nb}{str}\PY{p}{(}\PY{n}{path}\PY{p}{)}\PY{p}{,} \PY{n}{width}\PY{o}{=}\PY{n}{width}\PY{p}{)}\PY{p}{)}
    \PY{k}{else}\PY{p}{:}
        \PY{n}{display}\PY{p}{(}\PY{n}{Markdown}\PY{p}{(}\PY{l+s+sa}{f}\PY{l+s+s2}{\PYZdq{}}\PY{l+s+s2}{**Imagen no encontrada:** `}\PY{l+s+si}{\PYZob{}}\PY{n}{path}\PY{l+s+si}{\PYZcb{}}\PY{l+s+s2}{`}\PY{l+s+s2}{\PYZdq{}}\PY{p}{)}\PY{p}{)}


\PY{k}{def}\PY{+w}{ }\PY{n+nf}{show\PYZus{}table}\PY{p}{(}\PY{n}{title}\PY{p}{,} \PY{n}{df}\PY{p}{,} \PY{n}{n}\PY{o}{=}\PY{l+m+mi}{10}\PY{p}{)}\PY{p}{:}
    \PY{n}{display}\PY{p}{(}\PY{n}{Markdown}\PY{p}{(}\PY{l+s+sa}{f}\PY{l+s+s2}{\PYZdq{}}\PY{l+s+s2}{**}\PY{l+s+si}{\PYZob{}}\PY{n}{title}\PY{l+s+si}{\PYZcb{}}\PY{l+s+s2}{**}\PY{l+s+s2}{\PYZdq{}}\PY{p}{)}\PY{p}{)}
    \PY{n}{display}\PY{p}{(}\PY{n}{df}\PY{o}{.}\PY{n}{head}\PY{p}{(}\PY{n}{n}\PY{p}{)} \PY{k}{if} \PY{n}{n} \PY{k}{else} \PY{n}{df}\PY{p}{)}


\PY{k}{def}\PY{+w}{ }\PY{n+nf}{metric\PYZus{}from\PYZus{}summary}\PY{p}{(}\PY{n}{summary\PYZus{}df}\PY{p}{,} \PY{n}{metric}\PY{p}{,} \PY{n}{column}\PY{o}{=}\PY{l+s+s2}{\PYZdq{}}\PY{l+s+s2}{mean}\PY{l+s+s2}{\PYZdq{}}\PY{p}{)}\PY{p}{:}
    \PY{n}{row} \PY{o}{=} \PY{n}{summary\PYZus{}df}\PY{o}{.}\PY{n}{loc}\PY{p}{[}\PY{n}{summary\PYZus{}df}\PY{p}{[}\PY{l+s+s2}{\PYZdq{}}\PY{l+s+s2}{metric}\PY{l+s+s2}{\PYZdq{}}\PY{p}{]} \PY{o}{==} \PY{n}{metric}\PY{p}{]}
    \PY{k}{return} \PY{n}{np}\PY{o}{.}\PY{n}{nan} \PY{k}{if} \PY{n}{row}\PY{o}{.}\PY{n}{empty} \PY{k}{else} \PY{n+nb}{float}\PY{p}{(}\PY{n}{row}\PY{p}{[}\PY{n}{column}\PY{p}{]}\PY{o}{.}\PY{n}{iloc}\PY{p}{[}\PY{l+m+mi}{0}\PY{p}{]}\PY{p}{)}


\PY{k}{def}\PY{+w}{ }\PY{n+nf}{artifact\PYZus{}exists}\PY{p}{(}\PY{n}{relative\PYZus{}path}\PY{p}{)}\PY{p}{:}
    \PY{k}{return} \PY{p}{(}\PY{n}{PROJECT\PYZus{}DIR} \PY{o}{/} \PY{n}{relative\PYZus{}path}\PY{p}{)}\PY{o}{.}\PY{n}{exists}\PY{p}{(}\PY{p}{)}


\PY{k}{def}\PY{+w}{ }\PY{n+nf}{relative\PYZus{}to\PYZus{}project}\PY{p}{(}\PY{n}{path}\PY{p}{)}\PY{p}{:}
    \PY{n}{path} \PY{o}{=} \PY{n}{Path}\PY{p}{(}\PY{n}{path}\PY{p}{)}\PY{o}{.}\PY{n}{resolve}\PY{p}{(}\PY{p}{)}
    \PY{k}{try}\PY{p}{:}
        \PY{k}{return} \PY{n}{path}\PY{o}{.}\PY{n}{relative\PYZus{}to}\PY{p}{(}\PY{n}{PROJECT\PYZus{}ROOT}\PY{p}{)}
    \PY{k}{except} \PY{n+ne}{ValueError}\PY{p}{:}
        \PY{k}{return} \PY{n}{path}
\end{Verbatim}
\end{tcolorbox}

    \section{1. Exploracion y Preparacion de
Datos}\label{exploracion-y-preparacion-de-datos}

Esta seccion procedemos a entender el dataset medico anonimizado,
revisar calidad, manejar valores faltantes y anomalias, transformar
variables para modelos matematicos y disenar visualizaciones
relacionadas con \texttt{condition}.

La variable \texttt{condition} se interpreta como la clase medica
objetivo. En esta primera etapa no se entrena ningun modelo; se estudia
la estructura, calidad, distribucion y relacion de las variables con esa
clase para justificar las decisiones posteriores.

    \subsection{1.1. Carga reproducible del
dataset}\label{carga-reproducible-del-dataset}

El notebook busca los archivos principales en varias rutas posibles para
que pueda ejecutarse desde la raiz del proyecto o desde una subcarpeta.
Se cargan tres referencias de trabajo: el dataset original, el dataset
limpio y la version estandarizada para modelado. Separar estas versiones
es importante porque cada una responde a una pregunta distinta: que
datos llegaron originalmente, que datos fueron validados como
utilizables y que datos estan listos para algoritmos sensibles a la
escala.

El dataset original (\texttt{foot.csv}) es la fuente de partida. Se
conserva sin modificar para poder verificar cualquier decision posterior
y evitar perder trazabilidad. El dataset limpio
(\texttt{foot\_clean.csv}) es la version auditada despues de revisar
valores faltantes, duplicados, tipos de datos, valores unicos, rangos
minimos/maximos y consistencia de la variable objetivo. En este caso
mantiene las mismas filas y columnas que el original porque no se
encontraron nulos ni duplicados. Por tanto, aqui limpiar no significa
borrar informacion, sino confirmar que los registros son coherentes para
el analisis y documentar que no hizo falta imputar ni eliminar filas.

Los outliers se tratan con especial cuidado porque el contexto es
medico. Un valor extremo puede ser un error de captura, pero tambien
puede representar un paciente real con una condicion clinica relevante.
Por eso no se eliminan de forma automatica durante la carga o la
limpieza inicial: primero se identifican y se muestran en la seccion de
anomalias, y solo despues se justifica la decision de conservarlos. Esta
estrategia evita perder casos potencialmente importantes antes de
entenderlos.

La version estandarizada (\texttt{foot\_model\_ready.csv}) es una copia
preparada para modelos matematicos sensibles a la escala. En ella las
features numericas se transforman para tener media cercana a 0 y
desviacion estandar cercana a 1, mientras que \texttt{condition} se
mantiene sin modificar como clase objetivo binaria. Esta version es util
para PCA, KNN, SVM y modelos regularizados, porque evita que variables
con rangos grandes dominen el aprendizaje solo por su magnitud.

    \begin{tcolorbox}[breakable, size=fbox, boxrule=1pt, pad at break*=1mm,colback=cellbackground, colframe=cellborder]
\prompt{In}{incolor}{8}{\boxspacing}
\begin{Verbatim}[commandchars=\\\{\}]
\PY{n}{df\PYZus{}raw} \PY{o}{=} \PY{n}{read\PYZus{}csv}\PY{p}{(}\PY{n}{RAW\PYZus{}CSV}\PY{p}{)}
\PY{n}{df\PYZus{}clean} \PY{o}{=} \PY{n}{read\PYZus{}csv}\PY{p}{(}\PY{n}{CLEAN\PYZus{}CSV}\PY{p}{)} \PY{k}{if} \PY{n}{CLEAN\PYZus{}CSV}\PY{o}{.}\PY{n}{exists}\PY{p}{(}\PY{p}{)} \PY{k}{else} \PY{n}{df\PYZus{}raw}\PY{o}{.}\PY{n}{copy}\PY{p}{(}\PY{p}{)}

\PY{n}{target} \PY{o}{=} \PY{l+s+s2}{\PYZdq{}}\PY{l+s+s2}{condition}\PY{l+s+s2}{\PYZdq{}}
\PY{n}{features} \PY{o}{=} \PY{p}{[}\PY{n}{column} \PY{k}{for} \PY{n}{column} \PY{o+ow}{in} \PY{n}{df\PYZus{}clean}\PY{o}{.}\PY{n}{columns} \PY{k}{if} \PY{n}{column} \PY{o}{!=} \PY{n}{target}\PY{p}{]}

\PY{k}{if} \PY{n}{target} \PY{o+ow}{not} \PY{o+ow}{in} \PY{n}{df\PYZus{}clean}\PY{o}{.}\PY{n}{columns}\PY{p}{:}
    \PY{k}{raise} \PY{n+ne}{ValueError}\PY{p}{(}\PY{l+s+sa}{f}\PY{l+s+s2}{\PYZdq{}}\PY{l+s+s2}{No existe la columna objetivo esperada: }\PY{l+s+si}{\PYZob{}}\PY{n}{target}\PY{l+s+si}{\PYZcb{}}\PY{l+s+s2}{\PYZdq{}}\PY{p}{)}

\PY{n+nb}{print}\PY{p}{(}\PY{l+s+sa}{f}\PY{l+s+s2}{\PYZdq{}}\PY{l+s+s2}{Dataset original: }\PY{l+s+si}{\PYZob{}}\PY{n}{df\PYZus{}raw}\PY{o}{.}\PY{n}{shape}\PY{p}{[}\PY{l+m+mi}{0}\PY{p}{]}\PY{l+s+si}{\PYZcb{}}\PY{l+s+s2}{ filas x }\PY{l+s+si}{\PYZob{}}\PY{n}{df\PYZus{}raw}\PY{o}{.}\PY{n}{shape}\PY{p}{[}\PY{l+m+mi}{1}\PY{p}{]}\PY{l+s+si}{\PYZcb{}}\PY{l+s+s2}{ columnas}\PY{l+s+s2}{\PYZdq{}}\PY{p}{)}
\PY{n+nb}{print}\PY{p}{(}\PY{l+s+sa}{f}\PY{l+s+s2}{\PYZdq{}}\PY{l+s+s2}{Dataset limpio:   }\PY{l+s+si}{\PYZob{}}\PY{n}{df\PYZus{}clean}\PY{o}{.}\PY{n}{shape}\PY{p}{[}\PY{l+m+mi}{0}\PY{p}{]}\PY{l+s+si}{\PYZcb{}}\PY{l+s+s2}{ filas x }\PY{l+s+si}{\PYZob{}}\PY{n}{df\PYZus{}clean}\PY{o}{.}\PY{n}{shape}\PY{p}{[}\PY{l+m+mi}{1}\PY{p}{]}\PY{l+s+si}{\PYZcb{}}\PY{l+s+s2}{ columnas}\PY{l+s+s2}{\PYZdq{}}\PY{p}{)}
\PY{n+nb}{print}\PY{p}{(}\PY{l+s+sa}{f}\PY{l+s+s2}{\PYZdq{}}\PY{l+s+s2}{Numero de features: }\PY{l+s+si}{\PYZob{}}\PY{n+nb}{len}\PY{p}{(}\PY{n}{features}\PY{p}{)}\PY{l+s+si}{\PYZcb{}}\PY{l+s+s2}{\PYZdq{}}\PY{p}{)}
\PY{n+nb}{print}\PY{p}{(}\PY{l+s+sa}{f}\PY{l+s+s2}{\PYZdq{}}\PY{l+s+s2}{Ruta original usada: }\PY{l+s+si}{\PYZob{}}\PY{n}{relative\PYZus{}to\PYZus{}project}\PY{p}{(}\PY{n}{RAW\PYZus{}CSV}\PY{p}{)}\PY{l+s+si}{\PYZcb{}}\PY{l+s+s2}{\PYZdq{}}\PY{p}{)}
\PY{n+nb}{print}\PY{p}{(}\PY{l+s+sa}{f}\PY{l+s+s2}{\PYZdq{}}\PY{l+s+s2}{Ruta limpia usada:   }\PY{l+s+si}{\PYZob{}}\PY{n}{relative\PYZus{}to\PYZus{}project}\PY{p}{(}\PY{n}{CLEAN\PYZus{}CSV}\PY{p}{)}\PY{+w}{ }\PY{k}{if}\PY{+w}{ }\PY{n}{CLEAN\PYZus{}CSV}\PY{o}{.}\PY{n}{exists}\PY{p}{(}\PY{p}{)}\PY{+w}{ }\PY{k}{else}\PY{+w}{ }\PY{l+s+s1}{\PYZsq{}}\PY{l+s+s1}{no existe; se usa copia del original}\PY{l+s+s1}{\PYZsq{}}\PY{l+s+si}{\PYZcb{}}\PY{l+s+s2}{\PYZdq{}}\PY{p}{)}

\PY{n}{show\PYZus{}table}\PY{p}{(}\PY{l+s+s2}{\PYZdq{}}\PY{l+s+s2}{Primeras filas del dataset limpio}\PY{l+s+s2}{\PYZdq{}}\PY{p}{,} \PY{n}{df\PYZus{}clean}\PY{p}{,} \PY{n}{n}\PY{o}{=}\PY{l+m+mi}{5}\PY{p}{)}
\end{Verbatim}
\end{tcolorbox}

    \begin{Verbatim}[commandchars=\\\{\}]
Dataset original: 297 filas x 14 columnas
Dataset limpio:   297 filas x 14 columnas
Numero de features: 13
Ruta original usada: mundial\_project\textbackslash{}data\textbackslash{}foot.csv
Ruta limpia usada:   mundial\_project\textbackslash{}data\textbackslash{}processed\textbackslash{}foot\_clean.csv
    \end{Verbatim}

    \textbf{Primeras filas del dataset limpio}

    
    
    \begin{Verbatim}[commandchars=\\\{\}]
   feature\_1  feature\_2  feature\_3  feature\_4  feature\_5  feature\_6  feature\_7  feature\_8  feature\_9  feature\_10  feature\_11  feature\_12  feature\_13  \textbackslash{}
0          1          3          1          1         66        228          0          2          0         165         178           2         1.0   
1          0          3          0          1         59        177          1          1          0         162         140           2         0.0   
2          0          1          1          0         41        203          1          0          0         132         135           1         0.0   
3          0          2          0          0         37        215          0          0          0         170         120           0         0.0   
4          0          0          1          0         64        227          1          0          2         155         170           2         0.6   

   condition  
0          1  
1          1  
2          0  
3          0  
4          0  
    \end{Verbatim}

    
    \subsection{1.2. Estructura general y tipos de
datos}\label{estructura-general-y-tipos-de-datos}

Antes de limpiar, transformar o modelar, se analiza la estructura de
cada columna. En esta celda se revisa el tipo de dato, la cantidad de
valores nulos, el numero de valores unicos y los valores minimos y
maximos. Esto permite detectar problemas basicos: columnas que deberian
ser numericas pero no lo son, variables constantes, rangos imposibles,
codificaciones inesperadas o una variable objetivo que no sea realmente
binaria.

Tambien se comprueba que todas las columnas sean numericas y que
\texttt{condition} solo tenga los valores 0 y 1. Esta validacion es
necesaria porque los modelos posteriores trabajan con matrices numericas
y porque el problema se formula como clasificacion binaria. Mirar los
valores unicos ayuda ademas a distinguir variables binarias,
discretas/ordinales y continuas, lo que se usa en la siguiente
subseccion para decidir como visualizar y transformar cada feature.

    \begin{tcolorbox}[breakable, size=fbox, boxrule=1pt, pad at break*=1mm,colback=cellbackground, colframe=cellborder]
\prompt{In}{incolor}{9}{\boxspacing}
\begin{Verbatim}[commandchars=\\\{\}]
\PY{n}{structure\PYZus{}summary} \PY{o}{=} \PY{n}{pd}\PY{o}{.}\PY{n}{DataFrame}\PY{p}{(}\PY{p}{\PYZob{}}
    \PY{l+s+s2}{\PYZdq{}}\PY{l+s+s2}{columna}\PY{l+s+s2}{\PYZdq{}}\PY{p}{:} \PY{n}{df\PYZus{}clean}\PY{o}{.}\PY{n}{columns}\PY{p}{,}
    \PY{l+s+s2}{\PYZdq{}}\PY{l+s+s2}{tipo\PYZus{}dato}\PY{l+s+s2}{\PYZdq{}}\PY{p}{:} \PY{p}{[}\PY{n+nb}{str}\PY{p}{(}\PY{n}{dtype}\PY{p}{)} \PY{k}{for} \PY{n}{dtype} \PY{o+ow}{in} \PY{n}{df\PYZus{}clean}\PY{o}{.}\PY{n}{dtypes}\PY{p}{]}\PY{p}{,}
    \PY{l+s+s2}{\PYZdq{}}\PY{l+s+s2}{nulos}\PY{l+s+s2}{\PYZdq{}}\PY{p}{:} \PY{n}{df\PYZus{}clean}\PY{o}{.}\PY{n}{isna}\PY{p}{(}\PY{p}{)}\PY{o}{.}\PY{n}{sum}\PY{p}{(}\PY{p}{)}\PY{o}{.}\PY{n}{values}\PY{p}{,}
    \PY{l+s+s2}{\PYZdq{}}\PY{l+s+s2}{valores\PYZus{}unicos}\PY{l+s+s2}{\PYZdq{}}\PY{p}{:} \PY{p}{[}\PY{n}{df\PYZus{}clean}\PY{p}{[}\PY{n}{column}\PY{p}{]}\PY{o}{.}\PY{n}{nunique}\PY{p}{(}\PY{n}{dropna}\PY{o}{=}\PY{k+kc}{True}\PY{p}{)} \PY{k}{for} \PY{n}{column} \PY{o+ow}{in} \PY{n}{df\PYZus{}clean}\PY{o}{.}\PY{n}{columns}\PY{p}{]}\PY{p}{,}
    \PY{l+s+s2}{\PYZdq{}}\PY{l+s+s2}{minimo}\PY{l+s+s2}{\PYZdq{}}\PY{p}{:} \PY{p}{[}\PY{n}{df\PYZus{}clean}\PY{p}{[}\PY{n}{column}\PY{p}{]}\PY{o}{.}\PY{n}{min}\PY{p}{(}\PY{p}{)} \PY{k}{if} \PY{n}{pd}\PY{o}{.}\PY{n}{api}\PY{o}{.}\PY{n}{types}\PY{o}{.}\PY{n}{is\PYZus{}numeric\PYZus{}dtype}\PY{p}{(}\PY{n}{df\PYZus{}clean}\PY{p}{[}\PY{n}{column}\PY{p}{]}\PY{p}{)} \PY{k}{else} \PY{n}{np}\PY{o}{.}\PY{n}{nan} \PY{k}{for} \PY{n}{column} \PY{o+ow}{in} \PY{n}{df\PYZus{}clean}\PY{o}{.}\PY{n}{columns}\PY{p}{]}\PY{p}{,}
    \PY{l+s+s2}{\PYZdq{}}\PY{l+s+s2}{maximo}\PY{l+s+s2}{\PYZdq{}}\PY{p}{:} \PY{p}{[}\PY{n}{df\PYZus{}clean}\PY{p}{[}\PY{n}{column}\PY{p}{]}\PY{o}{.}\PY{n}{max}\PY{p}{(}\PY{p}{)} \PY{k}{if} \PY{n}{pd}\PY{o}{.}\PY{n}{api}\PY{o}{.}\PY{n}{types}\PY{o}{.}\PY{n}{is\PYZus{}numeric\PYZus{}dtype}\PY{p}{(}\PY{n}{df\PYZus{}clean}\PY{p}{[}\PY{n}{column}\PY{p}{]}\PY{p}{)} \PY{k}{else} \PY{n}{np}\PY{o}{.}\PY{n}{nan} \PY{k}{for} \PY{n}{column} \PY{o+ow}{in} \PY{n}{df\PYZus{}clean}\PY{o}{.}\PY{n}{columns}\PY{p}{]}\PY{p}{,}
\PY{p}{\PYZcb{}}\PY{p}{)}
\PY{n}{show\PYZus{}table}\PY{p}{(}\PY{l+s+s2}{\PYZdq{}}\PY{l+s+s2}{Estructura de columnas}\PY{l+s+s2}{\PYZdq{}}\PY{p}{,} \PY{n}{structure\PYZus{}summary}\PY{p}{,} \PY{n}{n}\PY{o}{=}\PY{k+kc}{None}\PY{p}{)}

\PY{n+nb}{print}\PY{p}{(}\PY{l+s+s2}{\PYZdq{}}\PY{l+s+s2}{Valores observados en condition:}\PY{l+s+s2}{\PYZdq{}}\PY{p}{,} \PY{n+nb}{sorted}\PY{p}{(}\PY{n}{df\PYZus{}clean}\PY{p}{[}\PY{n}{target}\PY{p}{]}\PY{o}{.}\PY{n}{unique}\PY{p}{(}\PY{p}{)}\PY{o}{.}\PY{n}{tolist}\PY{p}{(}\PY{p}{)}\PY{p}{)}\PY{p}{)}
\PY{n+nb}{print}\PY{p}{(}\PY{l+s+s2}{\PYZdq{}}\PY{l+s+s2}{Todas las columnas son numericas:}\PY{l+s+s2}{\PYZdq{}}\PY{p}{,} \PY{n+nb}{all}\PY{p}{(}\PY{n}{pd}\PY{o}{.}\PY{n}{api}\PY{o}{.}\PY{n}{types}\PY{o}{.}\PY{n}{is\PYZus{}numeric\PYZus{}dtype}\PY{p}{(}\PY{n}{df\PYZus{}clean}\PY{p}{[}\PY{n}{column}\PY{p}{]}\PY{p}{)} \PY{k}{for} \PY{n}{column} \PY{o+ow}{in} \PY{n}{df\PYZus{}clean}\PY{o}{.}\PY{n}{columns}\PY{p}{)}\PY{p}{)}
\end{Verbatim}
\end{tcolorbox}

    \textbf{Estructura de columnas}

    
    
    \begin{Verbatim}[commandchars=\\\{\}]
       columna tipo\_dato  nulos  valores\_unicos  minimo  maximo
0    feature\_1     int64      0               2     0.0     1.0
1    feature\_2     int64      0               4     0.0     3.0
2    feature\_3     int64      0               3     0.0     2.0
3    feature\_4     int64      0               2     0.0     1.0
4    feature\_5     int64      0              41    29.0    77.0
5    feature\_6     int64      0             152   126.0   564.0
6    feature\_7     int64      0               2     0.0     1.0
7    feature\_8     int64      0               4     0.0     3.0
8    feature\_9     int64      0               3     0.0     2.0
9   feature\_10     int64      0              91    71.0   202.0
10  feature\_11     int64      0              50    94.0   200.0
11  feature\_12     int64      0               3     0.0     2.0
12  feature\_13   float64      0              40     0.0     6.2
13   condition     int64      0               2     0.0     1.0
    \end{Verbatim}

    
    \begin{Verbatim}[commandchars=\\\{\}]
Valores observados en condition: [0, 1]
Todas las columnas son numericas: True
    \end{Verbatim}

    \subsection{1.3. Diccionario tecnico de
variables}\label{diccionario-tecnico-de-variables}

Aunque las variables estan anonimizadas, no todas deben tratarse igual.
Las binarias ya representan presencia/ausencia en formato 0/1, las
discretas parecen niveles codificados y las continuas tienen rangos
amplios. Esta clasificacion guia las visualizaciones, la deteccion de
outliers y la preparacion para modelos.

    \begin{tcolorbox}[breakable, size=fbox, boxrule=1pt, pad at break*=1mm,colback=cellbackground, colframe=cellborder]
\prompt{In}{incolor}{10}{\boxspacing}
\begin{Verbatim}[commandchars=\\\{\}]
\PY{k}{def}\PY{+w}{ }\PY{n+nf}{classify\PYZus{}feature\PYZus{}for\PYZus{}report}\PY{p}{(}\PY{n}{series}\PY{p}{)}\PY{p}{:}
    \PY{n}{unique\PYZus{}count} \PY{o}{=} \PY{n}{series}\PY{o}{.}\PY{n}{dropna}\PY{p}{(}\PY{p}{)}\PY{o}{.}\PY{n}{nunique}\PY{p}{(}\PY{p}{)}
    \PY{k}{if} \PY{n}{unique\PYZus{}count} \PY{o}{==} \PY{l+m+mi}{2}\PY{p}{:}
        \PY{k}{return} \PY{l+s+s2}{\PYZdq{}}\PY{l+s+s2}{binaria}\PY{l+s+s2}{\PYZdq{}}
    \PY{k}{if} \PY{n}{pd}\PY{o}{.}\PY{n}{api}\PY{o}{.}\PY{n}{types}\PY{o}{.}\PY{n}{is\PYZus{}numeric\PYZus{}dtype}\PY{p}{(}\PY{n}{series}\PY{p}{)} \PY{o+ow}{and} \PY{n}{unique\PYZus{}count} \PY{o}{\PYZgt{}}\PY{o}{=} \PY{l+m+mi}{10}\PY{p}{:}
        \PY{k}{return} \PY{l+s+s2}{\PYZdq{}}\PY{l+s+s2}{continua}\PY{l+s+s2}{\PYZdq{}}
    \PY{k}{return} \PY{l+s+s2}{\PYZdq{}}\PY{l+s+s2}{discreta/ordinal}\PY{l+s+s2}{\PYZdq{}}

\PY{n}{feature\PYZus{}dictionary} \PY{o}{=} \PY{n}{pd}\PY{o}{.}\PY{n}{DataFrame}\PY{p}{(}\PY{p}{\PYZob{}}
    \PY{l+s+s2}{\PYZdq{}}\PY{l+s+s2}{feature}\PY{l+s+s2}{\PYZdq{}}\PY{p}{:} \PY{n}{features}\PY{p}{,}
    \PY{l+s+s2}{\PYZdq{}}\PY{l+s+s2}{tipo\PYZus{}dato}\PY{l+s+s2}{\PYZdq{}}\PY{p}{:} \PY{p}{[}\PY{n+nb}{str}\PY{p}{(}\PY{n}{df\PYZus{}clean}\PY{p}{[}\PY{n}{feature}\PY{p}{]}\PY{o}{.}\PY{n}{dtype}\PY{p}{)} \PY{k}{for} \PY{n}{feature} \PY{o+ow}{in} \PY{n}{features}\PY{p}{]}\PY{p}{,}
    \PY{l+s+s2}{\PYZdq{}}\PY{l+s+s2}{valores\PYZus{}unicos}\PY{l+s+s2}{\PYZdq{}}\PY{p}{:} \PY{p}{[}\PY{n}{df\PYZus{}clean}\PY{p}{[}\PY{n}{feature}\PY{p}{]}\PY{o}{.}\PY{n}{nunique}\PY{p}{(}\PY{n}{dropna}\PY{o}{=}\PY{k+kc}{True}\PY{p}{)} \PY{k}{for} \PY{n}{feature} \PY{o+ow}{in} \PY{n}{features}\PY{p}{]}\PY{p}{,}
    \PY{l+s+s2}{\PYZdq{}}\PY{l+s+s2}{minimo}\PY{l+s+s2}{\PYZdq{}}\PY{p}{:} \PY{p}{[}\PY{n}{df\PYZus{}clean}\PY{p}{[}\PY{n}{feature}\PY{p}{]}\PY{o}{.}\PY{n}{min}\PY{p}{(}\PY{p}{)} \PY{k}{for} \PY{n}{feature} \PY{o+ow}{in} \PY{n}{features}\PY{p}{]}\PY{p}{,}
    \PY{l+s+s2}{\PYZdq{}}\PY{l+s+s2}{maximo}\PY{l+s+s2}{\PYZdq{}}\PY{p}{:} \PY{p}{[}\PY{n}{df\PYZus{}clean}\PY{p}{[}\PY{n}{feature}\PY{p}{]}\PY{o}{.}\PY{n}{max}\PY{p}{(}\PY{p}{)} \PY{k}{for} \PY{n}{feature} \PY{o+ow}{in} \PY{n}{features}\PY{p}{]}\PY{p}{,}
\PY{p}{\PYZcb{}}\PY{p}{)}
\PY{n}{feature\PYZus{}dictionary}\PY{p}{[}\PY{l+s+s2}{\PYZdq{}}\PY{l+s+s2}{tipo\PYZus{}interpretativo}\PY{l+s+s2}{\PYZdq{}}\PY{p}{]} \PY{o}{=} \PY{p}{[}\PY{n}{classify\PYZus{}feature\PYZus{}for\PYZus{}report}\PY{p}{(}\PY{n}{df\PYZus{}clean}\PY{p}{[}\PY{n}{feature}\PY{p}{]}\PY{p}{)} \PY{k}{for} \PY{n}{feature} \PY{o+ow}{in} \PY{n}{features}\PY{p}{]}
\PY{n}{feature\PYZus{}dictionary}\PY{p}{[}\PY{l+s+s2}{\PYZdq{}}\PY{l+s+s2}{accion\PYZus{}preparacion}\PY{l+s+s2}{\PYZdq{}}\PY{p}{]} \PY{o}{=} \PY{n}{feature\PYZus{}dictionary}\PY{p}{[}\PY{l+s+s2}{\PYZdq{}}\PY{l+s+s2}{tipo\PYZus{}interpretativo}\PY{l+s+s2}{\PYZdq{}}\PY{p}{]}\PY{o}{.}\PY{n}{map}\PY{p}{(}\PY{p}{\PYZob{}}
    \PY{l+s+s2}{\PYZdq{}}\PY{l+s+s2}{binaria}\PY{l+s+s2}{\PYZdq{}}\PY{p}{:} \PY{l+s+s2}{\PYZdq{}}\PY{l+s+s2}{mantener 0/1; no necesita one\PYZhy{}hot}\PY{l+s+s2}{\PYZdq{}}\PY{p}{,}
    \PY{l+s+s2}{\PYZdq{}}\PY{l+s+s2}{discreta/ordinal}\PY{l+s+s2}{\PYZdq{}}\PY{p}{:} \PY{l+s+s2}{\PYZdq{}}\PY{l+s+s2}{conservar codificacion; escalar si el modelo es sensible a distancia}\PY{l+s+s2}{\PYZdq{}}\PY{p}{,}
    \PY{l+s+s2}{\PYZdq{}}\PY{l+s+s2}{continua}\PY{l+s+s2}{\PYZdq{}}\PY{p}{:} \PY{l+s+s2}{\PYZdq{}}\PY{l+s+s2}{revisar outliers y estandarizar para modelos sensibles a escala}\PY{l+s+s2}{\PYZdq{}}\PY{p}{,}
\PY{p}{\PYZcb{}}\PY{p}{)}
\PY{n}{show\PYZus{}table}\PY{p}{(}\PY{l+s+s2}{\PYZdq{}}\PY{l+s+s2}{Diccionario tecnico de features}\PY{l+s+s2}{\PYZdq{}}\PY{p}{,} \PY{n}{feature\PYZus{}dictionary}\PY{p}{,} \PY{n}{n}\PY{o}{=}\PY{k+kc}{None}\PY{p}{)}
\end{Verbatim}
\end{tcolorbox}

    \textbf{Diccionario tecnico de features}

    
    
    \begin{Verbatim}[commandchars=\\\{\}]
       feature tipo\_dato  valores\_unicos  minimo  maximo tipo\_interpretativo                                 accion\_preparacion
0    feature\_1     int64               2     0.0     1.0             binaria                  mantener 0/1; no necesita one-hot
1    feature\_2     int64               4     0.0     3.0    discreta/ordinal  conservar codificacion; escalar si el modelo e{\ldots}
2    feature\_3     int64               3     0.0     2.0    discreta/ordinal  conservar codificacion; escalar si el modelo e{\ldots}
3    feature\_4     int64               2     0.0     1.0             binaria                  mantener 0/1; no necesita one-hot
4    feature\_5     int64              41    29.0    77.0            continua  revisar outliers y estandarizar para modelos s{\ldots}
5    feature\_6     int64             152   126.0   564.0            continua  revisar outliers y estandarizar para modelos s{\ldots}
6    feature\_7     int64               2     0.0     1.0             binaria                  mantener 0/1; no necesita one-hot
7    feature\_8     int64               4     0.0     3.0    discreta/ordinal  conservar codificacion; escalar si el modelo e{\ldots}
8    feature\_9     int64               3     0.0     2.0    discreta/ordinal  conservar codificacion; escalar si el modelo e{\ldots}
9   feature\_10     int64              91    71.0   202.0            continua  revisar outliers y estandarizar para modelos s{\ldots}
10  feature\_11     int64              50    94.0   200.0            continua  revisar outliers y estandarizar para modelos s{\ldots}
11  feature\_12     int64               3     0.0     2.0    discreta/ordinal  conservar codificacion; escalar si el modelo e{\ldots}
12  feature\_13   float64              40     0.0     6.2            continua  revisar outliers y estandarizar para modelos s{\ldots}
    \end{Verbatim}

    
    \subsection{1.4. Calidad de datos: faltantes, duplicados y
limpieza}\label{calidad-de-datos-faltantes-duplicados-y-limpieza}

Despues de revisar tipos, rangos y valores unicos, se audita la calidad
basica del dataset. La limpieza no redujo el numero de registros porque
no se detectaron valores faltantes ni filas duplicadas. Por eso no fue
necesario imputar, eliminar filas ni corregir tipos.

En esta subseccion solo se toman decisiones sobre problemas ya
verificados: nulos, duplicados, tipos y validez del target.

    \begin{tcolorbox}[breakable, size=fbox, boxrule=1pt, pad at break*=1mm,colback=cellbackground, colframe=cellborder]
\prompt{In}{incolor}{11}{\boxspacing}
\begin{Verbatim}[commandchars=\\\{\}]
\PY{n}{quality} \PY{o}{=} \PY{n}{read\PYZus{}csv}\PY{p}{(}\PY{n}{OUTPUTS\PYZus{}DIR} \PY{o}{/} \PY{l+s+s2}{\PYZdq{}}\PY{l+s+s2}{task1}\PY{l+s+s2}{\PYZdq{}} \PY{o}{/} \PY{l+s+s2}{\PYZdq{}}\PY{l+s+s2}{calidad\PYZus{}datos.csv}\PY{l+s+s2}{\PYZdq{}}\PY{p}{)} \PY{k}{if} \PY{p}{(}\PY{n}{OUTPUTS\PYZus{}DIR} \PY{o}{/} \PY{l+s+s2}{\PYZdq{}}\PY{l+s+s2}{task1}\PY{l+s+s2}{\PYZdq{}} \PY{o}{/} \PY{l+s+s2}{\PYZdq{}}\PY{l+s+s2}{calidad\PYZus{}datos.csv}\PY{l+s+s2}{\PYZdq{}}\PY{p}{)}\PY{o}{.}\PY{n}{exists}\PY{p}{(}\PY{p}{)} \PY{k}{else} \PY{n}{pd}\PY{o}{.}\PY{n}{DataFrame}\PY{p}{(}\PY{p}{)}
\PY{n}{feature\PYZus{}summary} \PY{o}{=} \PY{n}{read\PYZus{}csv}\PY{p}{(}\PY{n}{OUTPUTS\PYZus{}DIR} \PY{o}{/} \PY{l+s+s2}{\PYZdq{}}\PY{l+s+s2}{task1}\PY{l+s+s2}{\PYZdq{}} \PY{o}{/} \PY{l+s+s2}{\PYZdq{}}\PY{l+s+s2}{resumen\PYZus{}features.csv}\PY{l+s+s2}{\PYZdq{}}\PY{p}{)} \PY{k}{if} \PY{p}{(}\PY{n}{OUTPUTS\PYZus{}DIR} \PY{o}{/} \PY{l+s+s2}{\PYZdq{}}\PY{l+s+s2}{task1}\PY{l+s+s2}{\PYZdq{}} \PY{o}{/} \PY{l+s+s2}{\PYZdq{}}\PY{l+s+s2}{resumen\PYZus{}features.csv}\PY{l+s+s2}{\PYZdq{}}\PY{p}{)}\PY{o}{.}\PY{n}{exists}\PY{p}{(}\PY{p}{)} \PY{k}{else} \PY{n}{feature\PYZus{}dictionary}\PY{o}{.}\PY{n}{copy}\PY{p}{(}\PY{p}{)}

\PY{n}{cleaning\PYZus{}comparison} \PY{o}{=} \PY{n}{pd}\PY{o}{.}\PY{n}{DataFrame}\PY{p}{(}\PY{p}{[}
    \PY{p}{\PYZob{}}\PY{l+s+s2}{\PYZdq{}}\PY{l+s+s2}{elemento}\PY{l+s+s2}{\PYZdq{}}\PY{p}{:} \PY{l+s+s2}{\PYZdq{}}\PY{l+s+s2}{Filas originales}\PY{l+s+s2}{\PYZdq{}}\PY{p}{,} \PY{l+s+s2}{\PYZdq{}}\PY{l+s+s2}{valor}\PY{l+s+s2}{\PYZdq{}}\PY{p}{:} \PY{n}{df\PYZus{}raw}\PY{o}{.}\PY{n}{shape}\PY{p}{[}\PY{l+m+mi}{0}\PY{p}{]}\PY{p}{\PYZcb{}}\PY{p}{,}
    \PY{p}{\PYZob{}}\PY{l+s+s2}{\PYZdq{}}\PY{l+s+s2}{elemento}\PY{l+s+s2}{\PYZdq{}}\PY{p}{:} \PY{l+s+s2}{\PYZdq{}}\PY{l+s+s2}{Filas limpias}\PY{l+s+s2}{\PYZdq{}}\PY{p}{,} \PY{l+s+s2}{\PYZdq{}}\PY{l+s+s2}{valor}\PY{l+s+s2}{\PYZdq{}}\PY{p}{:} \PY{n}{df\PYZus{}clean}\PY{o}{.}\PY{n}{shape}\PY{p}{[}\PY{l+m+mi}{0}\PY{p}{]}\PY{p}{\PYZcb{}}\PY{p}{,}
    \PY{p}{\PYZob{}}\PY{l+s+s2}{\PYZdq{}}\PY{l+s+s2}{elemento}\PY{l+s+s2}{\PYZdq{}}\PY{p}{:} \PY{l+s+s2}{\PYZdq{}}\PY{l+s+s2}{Columnas originales}\PY{l+s+s2}{\PYZdq{}}\PY{p}{,} \PY{l+s+s2}{\PYZdq{}}\PY{l+s+s2}{valor}\PY{l+s+s2}{\PYZdq{}}\PY{p}{:} \PY{n}{df\PYZus{}raw}\PY{o}{.}\PY{n}{shape}\PY{p}{[}\PY{l+m+mi}{1}\PY{p}{]}\PY{p}{\PYZcb{}}\PY{p}{,}
    \PY{p}{\PYZob{}}\PY{l+s+s2}{\PYZdq{}}\PY{l+s+s2}{elemento}\PY{l+s+s2}{\PYZdq{}}\PY{p}{:} \PY{l+s+s2}{\PYZdq{}}\PY{l+s+s2}{Columnas limpias}\PY{l+s+s2}{\PYZdq{}}\PY{p}{,} \PY{l+s+s2}{\PYZdq{}}\PY{l+s+s2}{valor}\PY{l+s+s2}{\PYZdq{}}\PY{p}{:} \PY{n}{df\PYZus{}clean}\PY{o}{.}\PY{n}{shape}\PY{p}{[}\PY{l+m+mi}{1}\PY{p}{]}\PY{p}{\PYZcb{}}\PY{p}{,}
    \PY{p}{\PYZob{}}\PY{l+s+s2}{\PYZdq{}}\PY{l+s+s2}{elemento}\PY{l+s+s2}{\PYZdq{}}\PY{p}{:} \PY{l+s+s2}{\PYZdq{}}\PY{l+s+s2}{Nulos originales}\PY{l+s+s2}{\PYZdq{}}\PY{p}{,} \PY{l+s+s2}{\PYZdq{}}\PY{l+s+s2}{valor}\PY{l+s+s2}{\PYZdq{}}\PY{p}{:} \PY{n+nb}{int}\PY{p}{(}\PY{n}{df\PYZus{}raw}\PY{o}{.}\PY{n}{isna}\PY{p}{(}\PY{p}{)}\PY{o}{.}\PY{n}{sum}\PY{p}{(}\PY{p}{)}\PY{o}{.}\PY{n}{sum}\PY{p}{(}\PY{p}{)}\PY{p}{)}\PY{p}{\PYZcb{}}\PY{p}{,}
    \PY{p}{\PYZob{}}\PY{l+s+s2}{\PYZdq{}}\PY{l+s+s2}{elemento}\PY{l+s+s2}{\PYZdq{}}\PY{p}{:} \PY{l+s+s2}{\PYZdq{}}\PY{l+s+s2}{Nulos limpios}\PY{l+s+s2}{\PYZdq{}}\PY{p}{,} \PY{l+s+s2}{\PYZdq{}}\PY{l+s+s2}{valor}\PY{l+s+s2}{\PYZdq{}}\PY{p}{:} \PY{n+nb}{int}\PY{p}{(}\PY{n}{df\PYZus{}clean}\PY{o}{.}\PY{n}{isna}\PY{p}{(}\PY{p}{)}\PY{o}{.}\PY{n}{sum}\PY{p}{(}\PY{p}{)}\PY{o}{.}\PY{n}{sum}\PY{p}{(}\PY{p}{)}\PY{p}{)}\PY{p}{\PYZcb{}}\PY{p}{,}
    \PY{p}{\PYZob{}}\PY{l+s+s2}{\PYZdq{}}\PY{l+s+s2}{elemento}\PY{l+s+s2}{\PYZdq{}}\PY{p}{:} \PY{l+s+s2}{\PYZdq{}}\PY{l+s+s2}{Duplicados originales}\PY{l+s+s2}{\PYZdq{}}\PY{p}{,} \PY{l+s+s2}{\PYZdq{}}\PY{l+s+s2}{valor}\PY{l+s+s2}{\PYZdq{}}\PY{p}{:} \PY{n+nb}{int}\PY{p}{(}\PY{n}{df\PYZus{}raw}\PY{o}{.}\PY{n}{duplicated}\PY{p}{(}\PY{p}{)}\PY{o}{.}\PY{n}{sum}\PY{p}{(}\PY{p}{)}\PY{p}{)}\PY{p}{\PYZcb{}}\PY{p}{,}
    \PY{p}{\PYZob{}}\PY{l+s+s2}{\PYZdq{}}\PY{l+s+s2}{elemento}\PY{l+s+s2}{\PYZdq{}}\PY{p}{:} \PY{l+s+s2}{\PYZdq{}}\PY{l+s+s2}{Duplicados limpios}\PY{l+s+s2}{\PYZdq{}}\PY{p}{,} \PY{l+s+s2}{\PYZdq{}}\PY{l+s+s2}{valor}\PY{l+s+s2}{\PYZdq{}}\PY{p}{:} \PY{n+nb}{int}\PY{p}{(}\PY{n}{df\PYZus{}clean}\PY{o}{.}\PY{n}{duplicated}\PY{p}{(}\PY{p}{)}\PY{o}{.}\PY{n}{sum}\PY{p}{(}\PY{p}{)}\PY{p}{)}\PY{p}{\PYZcb{}}\PY{p}{,}
\PY{p}{]}\PY{p}{)}

\PY{n}{cleaning\PYZus{}decisions} \PY{o}{=} \PY{n}{pd}\PY{o}{.}\PY{n}{DataFrame}\PY{p}{(}\PY{p}{[}
    \PY{p}{\PYZob{}}\PY{l+s+s2}{\PYZdq{}}\PY{l+s+s2}{revision}\PY{l+s+s2}{\PYZdq{}}\PY{p}{:} \PY{l+s+s2}{\PYZdq{}}\PY{l+s+s2}{Valores nulos}\PY{l+s+s2}{\PYZdq{}}\PY{p}{,} \PY{l+s+s2}{\PYZdq{}}\PY{l+s+s2}{resultado}\PY{l+s+s2}{\PYZdq{}}\PY{p}{:} \PY{n+nb}{int}\PY{p}{(}\PY{n}{df\PYZus{}raw}\PY{o}{.}\PY{n}{isna}\PY{p}{(}\PY{p}{)}\PY{o}{.}\PY{n}{sum}\PY{p}{(}\PY{p}{)}\PY{o}{.}\PY{n}{sum}\PY{p}{(}\PY{p}{)}\PY{p}{)}\PY{p}{,} \PY{l+s+s2}{\PYZdq{}}\PY{l+s+s2}{decision}\PY{l+s+s2}{\PYZdq{}}\PY{p}{:} \PY{l+s+s2}{\PYZdq{}}\PY{l+s+s2}{No imputar porque no hay faltantes}\PY{l+s+s2}{\PYZdq{}}\PY{p}{\PYZcb{}}\PY{p}{,}
    \PY{p}{\PYZob{}}\PY{l+s+s2}{\PYZdq{}}\PY{l+s+s2}{revision}\PY{l+s+s2}{\PYZdq{}}\PY{p}{:} \PY{l+s+s2}{\PYZdq{}}\PY{l+s+s2}{Duplicados}\PY{l+s+s2}{\PYZdq{}}\PY{p}{,} \PY{l+s+s2}{\PYZdq{}}\PY{l+s+s2}{resultado}\PY{l+s+s2}{\PYZdq{}}\PY{p}{:} \PY{n+nb}{int}\PY{p}{(}\PY{n}{df\PYZus{}raw}\PY{o}{.}\PY{n}{duplicated}\PY{p}{(}\PY{p}{)}\PY{o}{.}\PY{n}{sum}\PY{p}{(}\PY{p}{)}\PY{p}{)}\PY{p}{,} \PY{l+s+s2}{\PYZdq{}}\PY{l+s+s2}{decision}\PY{l+s+s2}{\PYZdq{}}\PY{p}{:} \PY{l+s+s2}{\PYZdq{}}\PY{l+s+s2}{No eliminar registros porque no hay duplicados}\PY{l+s+s2}{\PYZdq{}}\PY{p}{\PYZcb{}}\PY{p}{,}
    \PY{p}{\PYZob{}}\PY{l+s+s2}{\PYZdq{}}\PY{l+s+s2}{revision}\PY{l+s+s2}{\PYZdq{}}\PY{p}{:} \PY{l+s+s2}{\PYZdq{}}\PY{l+s+s2}{Tipos incorrectos}\PY{l+s+s2}{\PYZdq{}}\PY{p}{,} \PY{l+s+s2}{\PYZdq{}}\PY{l+s+s2}{resultado}\PY{l+s+s2}{\PYZdq{}}\PY{p}{:} \PY{l+s+s2}{\PYZdq{}}\PY{l+s+s2}{No detectados}\PY{l+s+s2}{\PYZdq{}}\PY{p}{,} \PY{l+s+s2}{\PYZdq{}}\PY{l+s+s2}{decision}\PY{l+s+s2}{\PYZdq{}}\PY{p}{:} \PY{l+s+s2}{\PYZdq{}}\PY{l+s+s2}{Mantener columnas numericas}\PY{l+s+s2}{\PYZdq{}}\PY{p}{\PYZcb{}}\PY{p}{,}
    \PY{p}{\PYZob{}}\PY{l+s+s2}{\PYZdq{}}\PY{l+s+s2}{revision}\PY{l+s+s2}{\PYZdq{}}\PY{p}{:} \PY{l+s+s2}{\PYZdq{}}\PY{l+s+s2}{Target binario}\PY{l+s+s2}{\PYZdq{}}\PY{p}{,} \PY{l+s+s2}{\PYZdq{}}\PY{l+s+s2}{resultado}\PY{l+s+s2}{\PYZdq{}}\PY{p}{:} \PY{n+nb}{str}\PY{p}{(}\PY{n+nb}{sorted}\PY{p}{(}\PY{n}{df\PYZus{}clean}\PY{p}{[}\PY{n}{target}\PY{p}{]}\PY{o}{.}\PY{n}{unique}\PY{p}{(}\PY{p}{)}\PY{o}{.}\PY{n}{tolist}\PY{p}{(}\PY{p}{)}\PY{p}{)}\PY{p}{)}\PY{p}{,} \PY{l+s+s2}{\PYZdq{}}\PY{l+s+s2}{decision}\PY{l+s+s2}{\PYZdq{}}\PY{p}{:} \PY{l+s+s2}{\PYZdq{}}\PY{l+s+s2}{Valido para clasificacion}\PY{l+s+s2}{\PYZdq{}}\PY{p}{\PYZcb{}}\PY{p}{,}
\PY{p}{]}\PY{p}{)}

\PY{k}{if} \PY{o+ow}{not} \PY{n}{quality}\PY{o}{.}\PY{n}{empty}\PY{p}{:}
    \PY{n}{show\PYZus{}table}\PY{p}{(}\PY{l+s+s2}{\PYZdq{}}\PY{l+s+s2}{Auditoria reproducible de calidad}\PY{l+s+s2}{\PYZdq{}}\PY{p}{,} \PY{n}{quality}\PY{p}{,} \PY{n}{n}\PY{o}{=}\PY{k+kc}{None}\PY{p}{)}
\PY{n}{show\PYZus{}table}\PY{p}{(}\PY{l+s+s2}{\PYZdq{}}\PY{l+s+s2}{Comparacion antes vs despues de limpieza}\PY{l+s+s2}{\PYZdq{}}\PY{p}{,} \PY{n}{cleaning\PYZus{}comparison}\PY{p}{,} \PY{n}{n}\PY{o}{=}\PY{k+kc}{None}\PY{p}{)}
\PY{n}{show\PYZus{}table}\PY{p}{(}\PY{l+s+s2}{\PYZdq{}}\PY{l+s+s2}{Decisiones de limpieza}\PY{l+s+s2}{\PYZdq{}}\PY{p}{,} \PY{n}{cleaning\PYZus{}decisions}\PY{p}{,} \PY{n}{n}\PY{o}{=}\PY{k+kc}{None}\PY{p}{)}
\end{Verbatim}
\end{tcolorbox}

    \textbf{Auditoria reproducible de calidad}

    
    
    \begin{Verbatim}[commandchars=\\\{\}]
                 metrica      valor
0                  filas        297
1               columnas         14
2               features         13
3       columna\_objetivo  condition
4       valores\_objetivo     [0, 1]
5  valores\_nulos\_totales          0
6       filas\_duplicadas          0
    \end{Verbatim}

    
    \textbf{Comparacion antes vs despues de limpieza}

    
    
    \begin{Verbatim}[commandchars=\\\{\}]
                elemento  valor
0       Filas originales    297
1          Filas limpias    297
2    Columnas originales     14
3       Columnas limpias     14
4       Nulos originales      0
5          Nulos limpios      0
6  Duplicados originales      0
7     Duplicados limpios      0
    \end{Verbatim}

    
    \textbf{Decisiones de limpieza}

    
    
    \begin{Verbatim}[commandchars=\\\{\}]
            revision      resultado                                        decision
0      Valores nulos              0              No imputar porque no hay faltantes
1         Duplicados              0  No eliminar registros porque no hay duplicados
2  Tipos incorrectos  No detectados                     Mantener columnas numericas
3     Target binario         [0, 1]                       Valido para clasificacion
    \end{Verbatim}

    
    \subsection{\texorpdfstring{1.5. Balance de
\texttt{condition}}{1.5. Balance de condition}}\label{balance-de-condition}

La variable objetivo esta razonablemente balanceada: no hay una clase
dominante extrema. Aun asi, en las tareas supervisadas se prioriza
detectar \texttt{condition=1}, porque en un contexto medico un falso
negativo suele ser mas costoso que un falso positivo, porque el modelo
predice que el paciente no tiene la condición, pero en realidad sí la
tiene, lo que lo hace peor es que es más grave dejar pasar a un paciente
enfermo que investigar de más a un paciente sano.

    \begin{tcolorbox}[breakable, size=fbox, boxrule=1pt, pad at break*=1mm,colback=cellbackground, colframe=cellborder]
\prompt{In}{incolor}{12}{\boxspacing}
\begin{Verbatim}[commandchars=\\\{\}]
\PY{n}{condition\PYZus{}counts} \PY{o}{=} \PY{n}{df\PYZus{}clean}\PY{p}{[}\PY{n}{target}\PY{p}{]}\PY{o}{.}\PY{n}{value\PYZus{}counts}\PY{p}{(}\PY{p}{)}\PY{o}{.}\PY{n}{sort\PYZus{}index}\PY{p}{(}\PY{p}{)}
\PY{n}{condition\PYZus{}pct} \PY{o}{=} \PY{p}{(}\PY{n}{condition\PYZus{}counts} \PY{o}{/} \PY{n+nb}{len}\PY{p}{(}\PY{n}{df\PYZus{}clean}\PY{p}{)} \PY{o}{*} \PY{l+m+mi}{100}\PY{p}{)}\PY{o}{.}\PY{n}{round}\PY{p}{(}\PY{l+m+mi}{2}\PY{p}{)}
\PY{n}{condition\PYZus{}balance} \PY{o}{=} \PY{n}{pd}\PY{o}{.}\PY{n}{DataFrame}\PY{p}{(}\PY{p}{\PYZob{}}\PY{l+s+s2}{\PYZdq{}}\PY{l+s+s2}{condition}\PY{l+s+s2}{\PYZdq{}}\PY{p}{:} \PY{n}{condition\PYZus{}counts}\PY{o}{.}\PY{n}{index}\PY{p}{,} \PY{l+s+s2}{\PYZdq{}}\PY{l+s+s2}{conteo}\PY{l+s+s2}{\PYZdq{}}\PY{p}{:} \PY{n}{condition\PYZus{}counts}\PY{o}{.}\PY{n}{values}\PY{p}{,} \PY{l+s+s2}{\PYZdq{}}\PY{l+s+s2}{porcentaje}\PY{l+s+s2}{\PYZdq{}}\PY{p}{:} \PY{n}{condition\PYZus{}pct}\PY{o}{.}\PY{n}{values}\PY{p}{\PYZcb{}}\PY{p}{)}
\PY{n}{show\PYZus{}table}\PY{p}{(}\PY{l+s+s2}{\PYZdq{}}\PY{l+s+s2}{Balance de condition}\PY{l+s+s2}{\PYZdq{}}\PY{p}{,} \PY{n}{condition\PYZus{}balance}\PY{p}{,} \PY{n}{n}\PY{o}{=}\PY{k+kc}{None}\PY{p}{)}

\PY{n}{fig}\PY{p}{,} \PY{n}{ax} \PY{o}{=} \PY{n}{plt}\PY{o}{.}\PY{n}{subplots}\PY{p}{(}\PY{n}{figsize}\PY{o}{=}\PY{p}{(}\PY{l+m+mi}{7}\PY{p}{,} \PY{l+m+mf}{4.5}\PY{p}{)}\PY{p}{)}
\PY{n}{sns}\PY{o}{.}\PY{n}{barplot}\PY{p}{(}\PY{n}{data}\PY{o}{=}\PY{n}{condition\PYZus{}balance}\PY{p}{,} \PY{n}{x}\PY{o}{=}\PY{l+s+s2}{\PYZdq{}}\PY{l+s+s2}{condition}\PY{l+s+s2}{\PYZdq{}}\PY{p}{,} \PY{n}{y}\PY{o}{=}\PY{l+s+s2}{\PYZdq{}}\PY{l+s+s2}{conteo}\PY{l+s+s2}{\PYZdq{}}\PY{p}{,} \PY{n}{hue}\PY{o}{=}\PY{l+s+s2}{\PYZdq{}}\PY{l+s+s2}{condition}\PY{l+s+s2}{\PYZdq{}}\PY{p}{,} \PY{n}{palette}\PY{o}{=}\PY{p}{[}\PY{l+s+s2}{\PYZdq{}}\PY{l+s+s2}{\PYZsh{}4C78A8}\PY{l+s+s2}{\PYZdq{}}\PY{p}{,} \PY{l+s+s2}{\PYZdq{}}\PY{l+s+s2}{\PYZsh{}E45756}\PY{l+s+s2}{\PYZdq{}}\PY{p}{]}\PY{p}{,} \PY{n}{legend}\PY{o}{=}\PY{k+kc}{False}\PY{p}{,} \PY{n}{ax}\PY{o}{=}\PY{n}{ax}\PY{p}{)}
\PY{n}{ax}\PY{o}{.}\PY{n}{set\PYZus{}title}\PY{p}{(}\PY{l+s+s2}{\PYZdq{}}\PY{l+s+s2}{Distribucion de condition}\PY{l+s+s2}{\PYZdq{}}\PY{p}{)}
\PY{n}{ax}\PY{o}{.}\PY{n}{set\PYZus{}xlabel}\PY{p}{(}\PY{l+s+s2}{\PYZdq{}}\PY{l+s+s2}{condition}\PY{l+s+s2}{\PYZdq{}}\PY{p}{)}
\PY{n}{ax}\PY{o}{.}\PY{n}{set\PYZus{}ylabel}\PY{p}{(}\PY{l+s+s2}{\PYZdq{}}\PY{l+s+s2}{Frecuencia}\PY{l+s+s2}{\PYZdq{}}\PY{p}{)}
\PY{k}{for} \PY{n}{index}\PY{p}{,} \PY{n}{row} \PY{o+ow}{in} \PY{n}{condition\PYZus{}balance}\PY{o}{.}\PY{n}{iterrows}\PY{p}{(}\PY{p}{)}\PY{p}{:}
    \PY{n}{ax}\PY{o}{.}\PY{n}{text}\PY{p}{(}\PY{n}{index}\PY{p}{,} \PY{n}{row}\PY{p}{[}\PY{l+s+s2}{\PYZdq{}}\PY{l+s+s2}{conteo}\PY{l+s+s2}{\PYZdq{}}\PY{p}{]} \PY{o}{+} \PY{l+m+mi}{2}\PY{p}{,} \PY{l+s+sa}{f}\PY{l+s+s2}{\PYZdq{}}\PY{l+s+si}{\PYZob{}}\PY{n+nb}{int}\PY{p}{(}\PY{n}{row}\PY{p}{[}\PY{l+s+s1}{\PYZsq{}}\PY{l+s+s1}{conteo}\PY{l+s+s1}{\PYZsq{}}\PY{p}{]}\PY{p}{)}\PY{l+s+si}{\PYZcb{}}\PY{l+s+se}{\PYZbs{}n}\PY{l+s+si}{\PYZob{}}\PY{n}{row}\PY{p}{[}\PY{l+s+s1}{\PYZsq{}}\PY{l+s+s1}{porcentaje}\PY{l+s+s1}{\PYZsq{}}\PY{p}{]}\PY{l+s+si}{:}\PY{l+s+s2}{.2f}\PY{l+s+si}{\PYZcb{}}\PY{l+s+s2}{\PYZpc{}}\PY{l+s+s2}{\PYZdq{}}\PY{p}{,} \PY{n}{ha}\PY{o}{=}\PY{l+s+s2}{\PYZdq{}}\PY{l+s+s2}{center}\PY{l+s+s2}{\PYZdq{}}\PY{p}{,} \PY{n}{va}\PY{o}{=}\PY{l+s+s2}{\PYZdq{}}\PY{l+s+s2}{bottom}\PY{l+s+s2}{\PYZdq{}}\PY{p}{)}
\PY{n}{plt}\PY{o}{.}\PY{n}{tight\PYZus{}layout}\PY{p}{(}\PY{p}{)}
\PY{n}{plt}\PY{o}{.}\PY{n}{show}\PY{p}{(}\PY{p}{)}
\end{Verbatim}
\end{tcolorbox}

    \textbf{Balance de condition}

    
    
    \begin{Verbatim}[commandchars=\\\{\}]
   condition  conteo  porcentaje
0          0     160       53.87
1          1     137       46.13
    \end{Verbatim}

    
    \begin{center}
    \adjustimage{max size={0.9\linewidth}{0.9\paperheight}}{REPORT_files/REPORT_15_2.png}
    \end{center}
    { \hspace*{\fill} \\}
    
    \subsection{1.6. Analisis univariado de
features}\label{analisis-univariado-de-features}

Los histogramas permiten comparar la distribucion de cada variable entre
\texttt{condition=0} y \texttt{condition=1}. Son una herramienta
exploratoria: no prueban causalidad, pero ayudan a identificar variables
con separacion parcial entre clases.

    \begin{tcolorbox}[breakable, size=fbox, boxrule=1pt, pad at break*=1mm,colback=cellbackground, colframe=cellborder]
\prompt{In}{incolor}{13}{\boxspacing}
\begin{Verbatim}[commandchars=\\\{\}]
\PY{n}{plot\PYZus{}features} \PY{o}{=} \PY{n}{features}
\PY{n}{n\PYZus{}cols} \PY{o}{=} \PY{l+m+mi}{3}
\PY{n}{n\PYZus{}rows} \PY{o}{=} \PY{n+nb}{int}\PY{p}{(}\PY{n}{np}\PY{o}{.}\PY{n}{ceil}\PY{p}{(}\PY{n+nb}{len}\PY{p}{(}\PY{n}{plot\PYZus{}features}\PY{p}{)} \PY{o}{/} \PY{n}{n\PYZus{}cols}\PY{p}{)}\PY{p}{)}
\PY{n}{fig}\PY{p}{,} \PY{n}{axes} \PY{o}{=} \PY{n}{plt}\PY{o}{.}\PY{n}{subplots}\PY{p}{(}\PY{n}{n\PYZus{}rows}\PY{p}{,} \PY{n}{n\PYZus{}cols}\PY{p}{,} \PY{n}{figsize}\PY{o}{=}\PY{p}{(}\PY{l+m+mi}{16}\PY{p}{,} \PY{l+m+mi}{4} \PY{o}{*} \PY{n}{n\PYZus{}rows}\PY{p}{)}\PY{p}{)}
\PY{n}{axes} \PY{o}{=} \PY{n}{axes}\PY{o}{.}\PY{n}{ravel}\PY{p}{(}\PY{p}{)}

\PY{k}{for} \PY{n}{ax}\PY{p}{,} \PY{n}{feature} \PY{o+ow}{in} \PY{n+nb}{zip}\PY{p}{(}\PY{n}{axes}\PY{p}{,} \PY{n}{plot\PYZus{}features}\PY{p}{)}\PY{p}{:}
    \PY{n}{sns}\PY{o}{.}\PY{n}{histplot}\PY{p}{(}
        \PY{n}{data}\PY{o}{=}\PY{n}{df\PYZus{}clean}\PY{p}{,}
        \PY{n}{x}\PY{o}{=}\PY{n}{feature}\PY{p}{,}
        \PY{n}{hue}\PY{o}{=}\PY{n}{target}\PY{p}{,}
        \PY{n}{bins}\PY{o}{=}\PY{l+m+mi}{20}\PY{p}{,}
        \PY{n}{kde}\PY{o}{=}\PY{k+kc}{False}\PY{p}{,}
        \PY{n}{stat}\PY{o}{=}\PY{l+s+s2}{\PYZdq{}}\PY{l+s+s2}{density}\PY{l+s+s2}{\PYZdq{}}\PY{p}{,}
        \PY{n}{common\PYZus{}norm}\PY{o}{=}\PY{k+kc}{False}\PY{p}{,}
        \PY{n}{alpha}\PY{o}{=}\PY{l+m+mf}{0.45}\PY{p}{,}
        \PY{n}{ax}\PY{o}{=}\PY{n}{ax}\PY{p}{,}
    \PY{p}{)}
    \PY{n}{ax}\PY{o}{.}\PY{n}{set\PYZus{}title}\PY{p}{(}\PY{n}{feature}\PY{p}{)}

\PY{k}{for} \PY{n}{ax} \PY{o+ow}{in} \PY{n}{axes}\PY{p}{[}\PY{n+nb}{len}\PY{p}{(}\PY{n}{plot\PYZus{}features}\PY{p}{)}\PY{p}{:}\PY{p}{]}\PY{p}{:}
    \PY{n}{ax}\PY{o}{.}\PY{n}{set\PYZus{}visible}\PY{p}{(}\PY{k+kc}{False}\PY{p}{)}

\PY{n}{plt}\PY{o}{.}\PY{n}{tight\PYZus{}layout}\PY{p}{(}\PY{p}{)}
\PY{n}{plt}\PY{o}{.}\PY{n}{show}\PY{p}{(}\PY{p}{)}
\end{Verbatim}
\end{tcolorbox}

    \begin{center}
    \adjustimage{max size={0.9\linewidth}{0.9\paperheight}}{REPORT_files/REPORT_17_0.png}
    \end{center}
    { \hspace*{\fill} \\}
    
    \textbf{Observaciones relevantes hasta este punto.} El dataset tiene 297
registros, 13 variables predictoras y una variable objetivo binaria. No
se observan valores faltantes ni duplicados, por lo que la preparacion
inicial no requiere imputacion ni eliminacion de filas. Tambien se
confirma que todas las columnas son numericas, lo que facilita el uso
posterior de modelos matematicos.

Las tablas muestran que las features no tienen la misma naturaleza.
\texttt{feature\_1}, \texttt{feature\_4} y \texttt{feature\_7} son
binarias; \texttt{feature\_2}, \texttt{feature\_3}, \texttt{feature\_8},
\texttt{feature\_9} y \texttt{feature\_12} son discretas u ordinales; y
\texttt{feature\_5}, \texttt{feature\_6}, \texttt{feature\_10},
\texttt{feature\_11} y \texttt{feature\_13} tienen mas variabilidad y se
tratan como continuas. Esta diferencia justifica usar barras para
variables discretas/binarias, histogramas o boxplots para continuas, y
estandarizacion antes de modelos sensibles a escala.

Tambien se observan rangos muy diferentes entre variables: algunas solo
toman valores 0/1, mientras que otras llegan a valores como 564 en
\texttt{feature\_6} o 202 en \texttt{feature\_10}. Esto es importante
porque, sin estandarizar, variables con rangos grandes podrian dominar
algoritmos basados en distancia o componentes principales. Los
histogramas sugieren que varias variables no separan perfectamente las
clases, pero si muestran cambios de distribucion entre
\texttt{condition=0} y \texttt{condition=1}; por eso se profundiza en la
siguiente seccion con diferencias de medias, correlaciones y pruebas
exploratorias.

    \subsection{\texorpdfstring{1.7. Relacion de features con
\texttt{condition}}{1.7. Relacion de features con condition}}\label{relacion-de-features-con-condition}

Se combinan tres lecturas: diferencias de medias, correlacion lineal con
\texttt{condition} y pruebas estadisticas exploratorias. Estas pruebas
no son evidencia clinica definitiva; sirven para priorizar variables y
entender que senales podrian aprovechar los modelos.

    \begin{tcolorbox}[breakable, size=fbox, boxrule=1pt, pad at break*=1mm,colback=cellbackground, colframe=cellborder]
\prompt{In}{incolor}{14}{\boxspacing}
\begin{Verbatim}[commandchars=\\\{\}]
\PY{k+kn}{from}\PY{+w}{ }\PY{n+nn}{scipy}\PY{n+nn}{.}\PY{n+nn}{stats}\PY{+w}{ }\PY{k+kn}{import} \PY{n}{chi2\PYZus{}contingency}\PY{p}{,} \PY{n}{mannwhitneyu}

\PY{n}{stat\PYZus{}rows} \PY{o}{=} \PY{p}{[}\PY{p}{]}
\PY{k}{for} \PY{n}{feature} \PY{o+ow}{in} \PY{n}{features}\PY{p}{:}
    \PY{n}{x0} \PY{o}{=} \PY{n}{df\PYZus{}clean}\PY{o}{.}\PY{n}{loc}\PY{p}{[}\PY{n}{df\PYZus{}clean}\PY{p}{[}\PY{n}{target}\PY{p}{]} \PY{o}{==} \PY{l+m+mi}{0}\PY{p}{,} \PY{n}{feature}\PY{p}{]}
    \PY{n}{x1} \PY{o}{=} \PY{n}{df\PYZus{}clean}\PY{o}{.}\PY{n}{loc}\PY{p}{[}\PY{n}{df\PYZus{}clean}\PY{p}{[}\PY{n}{target}\PY{p}{]} \PY{o}{==} \PY{l+m+mi}{1}\PY{p}{,} \PY{n}{feature}\PY{p}{]}
    \PY{n}{unique\PYZus{}count} \PY{o}{=} \PY{n}{df\PYZus{}clean}\PY{p}{[}\PY{n}{feature}\PY{p}{]}\PY{o}{.}\PY{n}{nunique}\PY{p}{(}\PY{n}{dropna}\PY{o}{=}\PY{k+kc}{True}\PY{p}{)}

    \PY{k}{if} \PY{n}{unique\PYZus{}count} \PY{o}{\PYZlt{}}\PY{o}{=} \PY{l+m+mi}{5}\PY{p}{:}
        \PY{n}{contingency} \PY{o}{=} \PY{n}{pd}\PY{o}{.}\PY{n}{crosstab}\PY{p}{(}\PY{n}{df\PYZus{}clean}\PY{p}{[}\PY{n}{feature}\PY{p}{]}\PY{p}{,} \PY{n}{df\PYZus{}clean}\PY{p}{[}\PY{n}{target}\PY{p}{]}\PY{p}{)}
        \PY{n}{stat}\PY{p}{,} \PY{n}{p\PYZus{}value}\PY{p}{,} \PY{n}{\PYZus{}}\PY{p}{,} \PY{n}{\PYZus{}} \PY{o}{=} \PY{n}{chi2\PYZus{}contingency}\PY{p}{(}\PY{n}{contingency}\PY{p}{)}
        \PY{n}{test} \PY{o}{=} \PY{l+s+s2}{\PYZdq{}}\PY{l+s+s2}{chi2}\PY{l+s+s2}{\PYZdq{}}
    \PY{k}{else}\PY{p}{:}
        \PY{n}{stat}\PY{p}{,} \PY{n}{p\PYZus{}value} \PY{o}{=} \PY{n}{mannwhitneyu}\PY{p}{(}\PY{n}{x0}\PY{p}{,} \PY{n}{x1}\PY{p}{,} \PY{n}{alternative}\PY{o}{=}\PY{l+s+s2}{\PYZdq{}}\PY{l+s+s2}{two\PYZhy{}sided}\PY{l+s+s2}{\PYZdq{}}\PY{p}{)}
        \PY{n}{test} \PY{o}{=} \PY{l+s+s2}{\PYZdq{}}\PY{l+s+s2}{mannwhitneyu}\PY{l+s+s2}{\PYZdq{}}

    \PY{n}{stat\PYZus{}rows}\PY{o}{.}\PY{n}{append}\PY{p}{(}\PY{p}{\PYZob{}}
        \PY{l+s+s2}{\PYZdq{}}\PY{l+s+s2}{feature}\PY{l+s+s2}{\PYZdq{}}\PY{p}{:} \PY{n}{feature}\PY{p}{,}
        \PY{l+s+s2}{\PYZdq{}}\PY{l+s+s2}{tipo\PYZus{}interpretativo}\PY{l+s+s2}{\PYZdq{}}\PY{p}{:} \PY{n}{classify\PYZus{}feature\PYZus{}for\PYZus{}report}\PY{p}{(}\PY{n}{df\PYZus{}clean}\PY{p}{[}\PY{n}{feature}\PY{p}{]}\PY{p}{)}\PY{p}{,}
        \PY{l+s+s2}{\PYZdq{}}\PY{l+s+s2}{test}\PY{l+s+s2}{\PYZdq{}}\PY{p}{:} \PY{n}{test}\PY{p}{,}
        \PY{l+s+s2}{\PYZdq{}}\PY{l+s+s2}{p\PYZus{}value}\PY{l+s+s2}{\PYZdq{}}\PY{p}{:} \PY{n}{p\PYZus{}value}\PY{p}{,}
        \PY{l+s+s2}{\PYZdq{}}\PY{l+s+s2}{media\PYZus{}condition\PYZus{}0}\PY{l+s+s2}{\PYZdq{}}\PY{p}{:} \PY{n}{x0}\PY{o}{.}\PY{n}{mean}\PY{p}{(}\PY{p}{)}\PY{p}{,}
        \PY{l+s+s2}{\PYZdq{}}\PY{l+s+s2}{media\PYZus{}condition\PYZus{}1}\PY{l+s+s2}{\PYZdq{}}\PY{p}{:} \PY{n}{x1}\PY{o}{.}\PY{n}{mean}\PY{p}{(}\PY{p}{)}\PY{p}{,}
        \PY{l+s+s2}{\PYZdq{}}\PY{l+s+s2}{diff\PYZus{}1\PYZus{}minus\PYZus{}0}\PY{l+s+s2}{\PYZdq{}}\PY{p}{:} \PY{n}{x1}\PY{o}{.}\PY{n}{mean}\PY{p}{(}\PY{p}{)} \PY{o}{\PYZhy{}} \PY{n}{x0}\PY{o}{.}\PY{n}{mean}\PY{p}{(}\PY{p}{)}\PY{p}{,}
        \PY{l+s+s2}{\PYZdq{}}\PY{l+s+s2}{corr\PYZus{}condition}\PY{l+s+s2}{\PYZdq{}}\PY{p}{:} \PY{n}{df\PYZus{}clean}\PY{p}{[}\PY{n}{feature}\PY{p}{]}\PY{o}{.}\PY{n}{corr}\PY{p}{(}\PY{n}{df\PYZus{}clean}\PY{p}{[}\PY{n}{target}\PY{p}{]}\PY{p}{)}\PY{p}{,}
    \PY{p}{\PYZcb{}}\PY{p}{)}

\PY{n}{condition\PYZus{}stats} \PY{o}{=} \PY{n}{pd}\PY{o}{.}\PY{n}{DataFrame}\PY{p}{(}\PY{n}{stat\PYZus{}rows}\PY{p}{)}
\PY{n}{condition\PYZus{}stats}\PY{p}{[}\PY{l+s+s2}{\PYZdq{}}\PY{l+s+s2}{abs\PYZus{}corr\PYZus{}condition}\PY{l+s+s2}{\PYZdq{}}\PY{p}{]} \PY{o}{=} \PY{n}{condition\PYZus{}stats}\PY{p}{[}\PY{l+s+s2}{\PYZdq{}}\PY{l+s+s2}{corr\PYZus{}condition}\PY{l+s+s2}{\PYZdq{}}\PY{p}{]}\PY{o}{.}\PY{n}{abs}\PY{p}{(}\PY{p}{)}
\PY{n}{condition\PYZus{}stats} \PY{o}{=} \PY{n}{condition\PYZus{}stats}\PY{o}{.}\PY{n}{sort\PYZus{}values}\PY{p}{(}\PY{l+s+s2}{\PYZdq{}}\PY{l+s+s2}{abs\PYZus{}corr\PYZus{}condition}\PY{l+s+s2}{\PYZdq{}}\PY{p}{,} \PY{n}{ascending}\PY{o}{=}\PY{k+kc}{False}\PY{p}{)}
\PY{n}{show\PYZus{}table}\PY{p}{(}\PY{l+s+s2}{\PYZdq{}}\PY{l+s+s2}{Asociacion exploratoria de cada feature con condition}\PY{l+s+s2}{\PYZdq{}}\PY{p}{,} \PY{n}{condition\PYZus{}stats}\PY{o}{.}\PY{n}{round}\PY{p}{(}\PY{l+m+mi}{5}\PY{p}{)}\PY{p}{,} \PY{n}{n}\PY{o}{=}\PY{k+kc}{None}\PY{p}{)}

\PY{n}{important\PYZus{}differences} \PY{o}{=} \PY{n}{condition\PYZus{}stats}\PY{p}{[}\PY{p}{[}
    \PY{l+s+s2}{\PYZdq{}}\PY{l+s+s2}{feature}\PY{l+s+s2}{\PYZdq{}}\PY{p}{,}
    \PY{l+s+s2}{\PYZdq{}}\PY{l+s+s2}{tipo\PYZus{}interpretativo}\PY{l+s+s2}{\PYZdq{}}\PY{p}{,}
    \PY{l+s+s2}{\PYZdq{}}\PY{l+s+s2}{media\PYZus{}condition\PYZus{}0}\PY{l+s+s2}{\PYZdq{}}\PY{p}{,}
    \PY{l+s+s2}{\PYZdq{}}\PY{l+s+s2}{media\PYZus{}condition\PYZus{}1}\PY{l+s+s2}{\PYZdq{}}\PY{p}{,}
    \PY{l+s+s2}{\PYZdq{}}\PY{l+s+s2}{diff\PYZus{}1\PYZus{}minus\PYZus{}0}\PY{l+s+s2}{\PYZdq{}}\PY{p}{,}
    \PY{l+s+s2}{\PYZdq{}}\PY{l+s+s2}{corr\PYZus{}condition}\PY{l+s+s2}{\PYZdq{}}\PY{p}{,}
    \PY{l+s+s2}{\PYZdq{}}\PY{l+s+s2}{p\PYZus{}value}\PY{l+s+s2}{\PYZdq{}}\PY{p}{,}
\PY{p}{]}\PY{p}{]}\PY{o}{.}\PY{n}{copy}\PY{p}{(}\PY{p}{)}
\PY{n}{show\PYZus{}table}\PY{p}{(}\PY{l+s+s2}{\PYZdq{}}\PY{l+s+s2}{Variables ordenadas por asociacion lineal absoluta con condition}\PY{l+s+s2}{\PYZdq{}}\PY{p}{,} \PY{n}{important\PYZus{}differences}\PY{o}{.}\PY{n}{round}\PY{p}{(}\PY{l+m+mi}{5}\PY{p}{)}\PY{p}{,} \PY{n}{n}\PY{o}{=}\PY{k+kc}{None}\PY{p}{)}

\PY{n}{show\PYZus{}image}\PY{p}{(}\PY{n}{OUTPUTS\PYZus{}DIR} \PY{o}{/} \PY{l+s+s2}{\PYZdq{}}\PY{l+s+s2}{task1}\PY{l+s+s2}{\PYZdq{}} \PY{o}{/} \PY{l+s+s2}{\PYZdq{}}\PY{l+s+s2}{features\PYZus{}continuas\PYZus{}por\PYZus{}condition.png}\PY{l+s+s2}{\PYZdq{}}\PY{p}{,} \PY{n}{width}\PY{o}{=}\PY{l+m+mi}{1000}\PY{p}{)}
\PY{n}{show\PYZus{}image}\PY{p}{(}\PY{n}{OUTPUTS\PYZus{}DIR} \PY{o}{/} \PY{l+s+s2}{\PYZdq{}}\PY{l+s+s2}{task1}\PY{l+s+s2}{\PYZdq{}} \PY{o}{/} \PY{l+s+s2}{\PYZdq{}}\PY{l+s+s2}{features\PYZus{}discretas\PYZus{}por\PYZus{}condition.png}\PY{l+s+s2}{\PYZdq{}}\PY{p}{,} \PY{n}{width}\PY{o}{=}\PY{l+m+mi}{1000}\PY{p}{)}
\end{Verbatim}
\end{tcolorbox}

    \textbf{Asociacion exploratoria de cada feature con condition}

    
    
    \begin{Verbatim}[commandchars=\\\{\}]
       feature tipo\_interpretativo          test  p\_value  media\_condition\_0  media\_condition\_1  diff\_1\_minus\_0  corr\_condition  abs\_corr\_condition
11  feature\_12    discreta/ordinal          chi2  0.00000            0.37500            1.37226         0.99726         0.52052             0.52052
7    feature\_8    discreta/ordinal          chi2  0.00000            0.27500            1.14599         0.87099         0.46319             0.46319
12  feature\_13            continua  mannwhitneyu  0.00000            0.59875            1.58905         0.99030         0.42405             0.42405
9   feature\_10            continua  mannwhitneyu  0.00000          158.58125          139.10949       -19.47176        -0.42382             0.42382
3    feature\_4             binaria          chi2  0.00000            0.14375            0.54015         0.39640         0.42136             0.42136
1    feature\_2    discreta/ordinal          chi2  0.00000            1.79375            2.58394         0.79019         0.40894             0.40894
2    feature\_3    discreta/ordinal          chi2  0.00000            0.41250            0.82482         0.41232         0.33305             0.33305
6    feature\_7             binaria          chi2  0.00000            0.55625            0.81752         0.26127         0.27847             0.27847
4    feature\_5            continua  mannwhitneyu  0.00004           52.64375           56.75912         4.11537         0.22708             0.22708
8    feature\_9    discreta/ordinal          chi2  0.00833            0.84375            1.17518         0.33143         0.16634             0.16634
10  feature\_11            continua  mannwhitneyu  0.02346          129.17500          134.63504         5.46004         0.15349             0.15349
5    feature\_6            continua  mannwhitneyu  0.04669          243.49375          251.85401         8.36026         0.08028             0.08028
0    feature\_1             binaria          chi2  1.00000            0.14375            0.14599         0.00224         0.00317             0.00317
    \end{Verbatim}

    
    \textbf{Variables ordenadas por asociacion lineal absoluta con
condition}

    
    
    \begin{Verbatim}[commandchars=\\\{\}]
       feature tipo\_interpretativo  media\_condition\_0  media\_condition\_1  diff\_1\_minus\_0  corr\_condition  p\_value
11  feature\_12    discreta/ordinal            0.37500            1.37226         0.99726         0.52052  0.00000
7    feature\_8    discreta/ordinal            0.27500            1.14599         0.87099         0.46319  0.00000
12  feature\_13            continua            0.59875            1.58905         0.99030         0.42405  0.00000
9   feature\_10            continua          158.58125          139.10949       -19.47176        -0.42382  0.00000
3    feature\_4             binaria            0.14375            0.54015         0.39640         0.42136  0.00000
1    feature\_2    discreta/ordinal            1.79375            2.58394         0.79019         0.40894  0.00000
2    feature\_3    discreta/ordinal            0.41250            0.82482         0.41232         0.33305  0.00000
6    feature\_7             binaria            0.55625            0.81752         0.26127         0.27847  0.00000
4    feature\_5            continua           52.64375           56.75912         4.11537         0.22708  0.00004
8    feature\_9    discreta/ordinal            0.84375            1.17518         0.33143         0.16634  0.00833
10  feature\_11            continua          129.17500          134.63504         5.46004         0.15349  0.02346
5    feature\_6            continua          243.49375          251.85401         8.36026         0.08028  0.04669
0    feature\_1             binaria            0.14375            0.14599         0.00224         0.00317  1.00000
    \end{Verbatim}

    
    \begin{center}
    \adjustimage{max size={0.9\linewidth}{0.9\paperheight}}{REPORT_files/REPORT_20_4.png}
    \end{center}
    { \hspace*{\fill} \\}
    
    \begin{center}
    \adjustimage{max size={0.9\linewidth}{0.9\paperheight}}{REPORT_files/REPORT_20_5.png}
    \end{center}
    { \hspace*{\fill} \\}
    
    \textbf{Informacion extraida de la relacion con \texttt{condition}.} La
tabla muestra que la separacion entre clases no depende de una sola
variable, sino de varias senales moderadas. Las asociaciones mas fuertes
aparecen en \texttt{feature\_12}, \texttt{feature\_8},
\texttt{feature\_13}, \texttt{feature\_10}, \texttt{feature\_4} y
\texttt{feature\_2}. Esto sugiere que los modelos posteriores deberan
combinar informacion de varias features para predecir
\texttt{condition}, en lugar de apoyarse en una regla simple basada en
una unica columna.

La direccion de las diferencias tambien aporta informacion.
\texttt{feature\_12}, \texttt{feature\_8}, \texttt{feature\_13},
\texttt{feature\_4} y \texttt{feature\_2} tienen medias mas altas cuando
\texttt{condition=1}, por lo que valores mayores en esas variables se
asocian con mayor presencia de la clase positiva. En cambio,
\texttt{feature\_10} tiene una relacion negativa: su media baja de
158.58 en \texttt{condition=0} a 139.11 en \texttt{condition=1}. Esto
significa que para esta variable los valores menores son los que se
relacionan mas con la clase positiva.

Las pruebas exploratorias refuerzan esta lectura: muchas de las
variables mejor posicionadas tienen p-values muy bajos, lo que indica
que sus distribuciones cambian entre clases. Aun asi, estos resultados
no deben interpretarse como causalidad ni como evidencia medica
definitiva; solo muestran asociaciones utiles para orientar el modelado.
\texttt{feature\_1}, por ejemplo, casi no cambia entre clases y tiene
correlacion cercana a cero, por lo que individualmente parece aportar
poca informacion predictiva.

En conjunto, la seccion sugiere que existe senal en los datos, pero no
una separacion perfecta. Por eso es razonable probar modelos
supervisados capaces de combinar relaciones lineales y no lineales, y
evaluar su rendimiento con metricas cuidadosas en lugar de seleccionar
variables solo por inspeccion visual.

    \subsection{1.8. Outliers y anomalias}\label{outliers-y-anomalias}

Una vez validada la estructura general, se analizan posibles anomalias.
Para no mezclar variables binarias o discretas con variables de rango
amplio, la deteccion de outliers se aplica principalmente a features
continuas. Se usa IQR porque es robusto ante distribuciones no normales
y se contrasta con z-score para comparar si los extremos tambien
aparecen bajo un criterio basado en desviaciones estandar.

Esta revision se hace antes de decidir que hacer con los valores
extremos. En datos medicos, eliminarlos sin evidencia podria borrar
pacientes poco frecuentes pero clinicamente importantes; por eso la
decision debe depender de la cantidad de outliers, su porcentaje y su
plausibilidad dentro del contexto del dataset.

    \begin{tcolorbox}[breakable, size=fbox, boxrule=1pt, pad at break*=1mm,colback=cellbackground, colframe=cellborder]
\prompt{In}{incolor}{15}{\boxspacing}
\begin{Verbatim}[commandchars=\\\{\}]
\PY{k+kn}{from}\PY{+w}{ }\PY{n+nn}{scipy}\PY{+w}{ }\PY{k+kn}{import} \PY{n}{stats}

\PY{n}{outliers\PYZus{}iqr\PYZus{}path} \PY{o}{=} \PY{n}{OUTPUTS\PYZus{}DIR} \PY{o}{/} \PY{l+s+s2}{\PYZdq{}}\PY{l+s+s2}{outliers}\PY{l+s+s2}{\PYZdq{}} \PY{o}{/} \PY{l+s+s2}{\PYZdq{}}\PY{l+s+s2}{outliers\PYZus{}iqr.csv}\PY{l+s+s2}{\PYZdq{}}
\PY{n}{outliers\PYZus{}zscore\PYZus{}path} \PY{o}{=} \PY{n}{OUTPUTS\PYZus{}DIR} \PY{o}{/} \PY{l+s+s2}{\PYZdq{}}\PY{l+s+s2}{outliers}\PY{l+s+s2}{\PYZdq{}} \PY{o}{/} \PY{l+s+s2}{\PYZdq{}}\PY{l+s+s2}{outliers\PYZus{}scipy.csv}\PY{l+s+s2}{\PYZdq{}}
\PY{n}{outliers\PYZus{}comparison\PYZus{}path} \PY{o}{=} \PY{n}{OUTPUTS\PYZus{}DIR} \PY{o}{/} \PY{l+s+s2}{\PYZdq{}}\PY{l+s+s2}{outliers}\PY{l+s+s2}{\PYZdq{}} \PY{o}{/} \PY{l+s+s2}{\PYZdq{}}\PY{l+s+s2}{comparacion\PYZus{}outliers.csv}\PY{l+s+s2}{\PYZdq{}}

\PY{k}{if} \PY{n}{outliers\PYZus{}iqr\PYZus{}path}\PY{o}{.}\PY{n}{exists}\PY{p}{(}\PY{p}{)}\PY{p}{:}
    \PY{n}{outliers\PYZus{}iqr} \PY{o}{=} \PY{n}{read\PYZus{}csv}\PY{p}{(}\PY{n}{outliers\PYZus{}iqr\PYZus{}path}\PY{p}{)}
\PY{k}{else}\PY{p}{:}
    \PY{n}{continuous\PYZus{}features} \PY{o}{=} \PY{n}{feature\PYZus{}dictionary}\PY{o}{.}\PY{n}{loc}\PY{p}{[}\PY{n}{feature\PYZus{}dictionary}\PY{p}{[}\PY{l+s+s2}{\PYZdq{}}\PY{l+s+s2}{tipo\PYZus{}interpretativo}\PY{l+s+s2}{\PYZdq{}}\PY{p}{]} \PY{o}{==} \PY{l+s+s2}{\PYZdq{}}\PY{l+s+s2}{continua}\PY{l+s+s2}{\PYZdq{}}\PY{p}{,} \PY{l+s+s2}{\PYZdq{}}\PY{l+s+s2}{feature}\PY{l+s+s2}{\PYZdq{}}\PY{p}{]}
    \PY{n}{rows} \PY{o}{=} \PY{p}{[}\PY{p}{]}
    \PY{k}{for} \PY{n}{feature} \PY{o+ow}{in} \PY{n}{continuous\PYZus{}features}\PY{p}{:}
        \PY{n}{q1} \PY{o}{=} \PY{n}{df\PYZus{}clean}\PY{p}{[}\PY{n}{feature}\PY{p}{]}\PY{o}{.}\PY{n}{quantile}\PY{p}{(}\PY{l+m+mf}{0.25}\PY{p}{)}
        \PY{n}{q3} \PY{o}{=} \PY{n}{df\PYZus{}clean}\PY{p}{[}\PY{n}{feature}\PY{p}{]}\PY{o}{.}\PY{n}{quantile}\PY{p}{(}\PY{l+m+mf}{0.75}\PY{p}{)}
        \PY{n}{iqr} \PY{o}{=} \PY{n}{q3} \PY{o}{\PYZhy{}} \PY{n}{q1}
        \PY{n}{lower} \PY{o}{=} \PY{n}{q1} \PY{o}{\PYZhy{}} \PY{l+m+mf}{1.5} \PY{o}{*} \PY{n}{iqr}
        \PY{n}{upper} \PY{o}{=} \PY{n}{q3} \PY{o}{+} \PY{l+m+mf}{1.5} \PY{o}{*} \PY{n}{iqr}
        \PY{n}{count} \PY{o}{=} \PY{n+nb}{int}\PY{p}{(}\PY{p}{(}\PY{p}{(}\PY{n}{df\PYZus{}clean}\PY{p}{[}\PY{n}{feature}\PY{p}{]} \PY{o}{\PYZlt{}} \PY{n}{lower}\PY{p}{)} \PY{o}{|} \PY{p}{(}\PY{n}{df\PYZus{}clean}\PY{p}{[}\PY{n}{feature}\PY{p}{]} \PY{o}{\PYZgt{}} \PY{n}{upper}\PY{p}{)}\PY{p}{)}\PY{o}{.}\PY{n}{sum}\PY{p}{(}\PY{p}{)}\PY{p}{)}
        \PY{n}{rows}\PY{o}{.}\PY{n}{append}\PY{p}{(}\PY{p}{\PYZob{}}
            \PY{l+s+s2}{\PYZdq{}}\PY{l+s+s2}{feature}\PY{l+s+s2}{\PYZdq{}}\PY{p}{:} \PY{n}{feature}\PY{p}{,}
            \PY{l+s+s2}{\PYZdq{}}\PY{l+s+s2}{q1}\PY{l+s+s2}{\PYZdq{}}\PY{p}{:} \PY{n}{q1}\PY{p}{,}
            \PY{l+s+s2}{\PYZdq{}}\PY{l+s+s2}{q3}\PY{l+s+s2}{\PYZdq{}}\PY{p}{:} \PY{n}{q3}\PY{p}{,}
            \PY{l+s+s2}{\PYZdq{}}\PY{l+s+s2}{iqr}\PY{l+s+s2}{\PYZdq{}}\PY{p}{:} \PY{n}{iqr}\PY{p}{,}
            \PY{l+s+s2}{\PYZdq{}}\PY{l+s+s2}{limite\PYZus{}inferior}\PY{l+s+s2}{\PYZdq{}}\PY{p}{:} \PY{n}{lower}\PY{p}{,}
            \PY{l+s+s2}{\PYZdq{}}\PY{l+s+s2}{limite\PYZus{}superior}\PY{l+s+s2}{\PYZdq{}}\PY{p}{:} \PY{n}{upper}\PY{p}{,}
            \PY{l+s+s2}{\PYZdq{}}\PY{l+s+s2}{outliers\PYZus{}iqr}\PY{l+s+s2}{\PYZdq{}}\PY{p}{:} \PY{n}{count}\PY{p}{,}
            \PY{l+s+s2}{\PYZdq{}}\PY{l+s+s2}{porcentaje\PYZus{}iqr}\PY{l+s+s2}{\PYZdq{}}\PY{p}{:} \PY{n}{count} \PY{o}{/} \PY{n+nb}{len}\PY{p}{(}\PY{n}{df\PYZus{}clean}\PY{p}{)} \PY{o}{*} \PY{l+m+mi}{100}\PY{p}{,}
        \PY{p}{\PYZcb{}}\PY{p}{)}
    \PY{n}{outliers\PYZus{}iqr} \PY{o}{=} \PY{n}{pd}\PY{o}{.}\PY{n}{DataFrame}\PY{p}{(}\PY{n}{rows}\PY{p}{)}

\PY{n}{show\PYZus{}table}\PY{p}{(}\PY{l+s+s2}{\PYZdq{}}\PY{l+s+s2}{Outliers detectados por IQR}\PY{l+s+s2}{\PYZdq{}}\PY{p}{,} \PY{n}{outliers\PYZus{}iqr}\PY{o}{.}\PY{n}{round}\PY{p}{(}\PY{l+m+mi}{4}\PY{p}{)}\PY{p}{,} \PY{n}{n}\PY{o}{=}\PY{k+kc}{None}\PY{p}{)}

\PY{k}{if} \PY{n}{outliers\PYZus{}zscore\PYZus{}path}\PY{o}{.}\PY{n}{exists}\PY{p}{(}\PY{p}{)}\PY{p}{:}
    \PY{n}{outliers\PYZus{}zscore} \PY{o}{=} \PY{n}{read\PYZus{}csv}\PY{p}{(}\PY{n}{outliers\PYZus{}zscore\PYZus{}path}\PY{p}{)}
\PY{k}{else}\PY{p}{:}
    \PY{n}{rows} \PY{o}{=} \PY{p}{[}\PY{p}{]}
    \PY{k}{for} \PY{n}{feature} \PY{o+ow}{in} \PY{n}{outliers\PYZus{}iqr}\PY{p}{[}\PY{l+s+s2}{\PYZdq{}}\PY{l+s+s2}{feature}\PY{l+s+s2}{\PYZdq{}}\PY{p}{]}\PY{p}{:}
        \PY{n}{series} \PY{o}{=} \PY{n}{df\PYZus{}clean}\PY{p}{[}\PY{n}{feature}\PY{p}{]}
        \PY{n}{z\PYZus{}scores} \PY{o}{=} \PY{n}{pd}\PY{o}{.}\PY{n}{Series}\PY{p}{(}\PY{n}{stats}\PY{o}{.}\PY{n}{zscore}\PY{p}{(}\PY{n}{series}\PY{p}{,} \PY{n}{nan\PYZus{}policy}\PY{o}{=}\PY{l+s+s2}{\PYZdq{}}\PY{l+s+s2}{omit}\PY{l+s+s2}{\PYZdq{}}\PY{p}{,} \PY{n}{ddof}\PY{o}{=}\PY{l+m+mi}{1}\PY{p}{)}\PY{p}{,} \PY{n}{index}\PY{o}{=}\PY{n}{series}\PY{o}{.}\PY{n}{index}\PY{p}{)}\PY{o}{.}\PY{n}{fillna}\PY{p}{(}\PY{l+m+mf}{0.0}\PY{p}{)}
        \PY{n}{count} \PY{o}{=} \PY{n+nb}{int}\PY{p}{(}\PY{p}{(}\PY{n}{z\PYZus{}scores}\PY{o}{.}\PY{n}{abs}\PY{p}{(}\PY{p}{)} \PY{o}{\PYZgt{}} \PY{l+m+mf}{3.0}\PY{p}{)}\PY{o}{.}\PY{n}{sum}\PY{p}{(}\PY{p}{)}\PY{p}{)}
        \PY{p}{(}\PY{n}{\PYZus{}}\PY{p}{,} \PY{n}{\PYZus{}}\PY{p}{)}\PY{p}{,} \PY{p}{(}\PY{n}{\PYZus{}}\PY{p}{,} \PY{n}{\PYZus{}}\PY{p}{,} \PY{n}{qq\PYZus{}r}\PY{p}{)} \PY{o}{=} \PY{n}{stats}\PY{o}{.}\PY{n}{probplot}\PY{p}{(}\PY{n}{series}\PY{o}{.}\PY{n}{dropna}\PY{p}{(}\PY{p}{)}\PY{p}{,} \PY{n}{dist}\PY{o}{=}\PY{l+s+s2}{\PYZdq{}}\PY{l+s+s2}{norm}\PY{l+s+s2}{\PYZdq{}}\PY{p}{)}
        \PY{n}{rows}\PY{o}{.}\PY{n}{append}\PY{p}{(}\PY{p}{\PYZob{}}
            \PY{l+s+s2}{\PYZdq{}}\PY{l+s+s2}{feature}\PY{l+s+s2}{\PYZdq{}}\PY{p}{:} \PY{n}{feature}\PY{p}{,}
            \PY{l+s+s2}{\PYZdq{}}\PY{l+s+s2}{media}\PY{l+s+s2}{\PYZdq{}}\PY{p}{:} \PY{n}{series}\PY{o}{.}\PY{n}{mean}\PY{p}{(}\PY{p}{)}\PY{p}{,}
            \PY{l+s+s2}{\PYZdq{}}\PY{l+s+s2}{desviacion\PYZus{}std}\PY{l+s+s2}{\PYZdq{}}\PY{p}{:} \PY{n}{series}\PY{o}{.}\PY{n}{std}\PY{p}{(}\PY{p}{)}\PY{p}{,}
            \PY{l+s+s2}{\PYZdq{}}\PY{l+s+s2}{umbral\PYZus{}z}\PY{l+s+s2}{\PYZdq{}}\PY{p}{:} \PY{l+m+mf}{3.0}\PY{p}{,}
            \PY{l+s+s2}{\PYZdq{}}\PY{l+s+s2}{outliers\PYZus{}scipy}\PY{l+s+s2}{\PYZdq{}}\PY{p}{:} \PY{n}{count}\PY{p}{,}
            \PY{l+s+s2}{\PYZdq{}}\PY{l+s+s2}{porcentaje\PYZus{}scipy}\PY{l+s+s2}{\PYZdq{}}\PY{p}{:} \PY{n}{count} \PY{o}{/} \PY{n+nb}{len}\PY{p}{(}\PY{n}{df\PYZus{}clean}\PY{p}{)} \PY{o}{*} \PY{l+m+mi}{100}\PY{p}{,}
            \PY{l+s+s2}{\PYZdq{}}\PY{l+s+s2}{qq\PYZus{}r\PYZus{}scipy}\PY{l+s+s2}{\PYZdq{}}\PY{p}{:} \PY{n}{qq\PYZus{}r}\PY{p}{,}
        \PY{p}{\PYZcb{}}\PY{p}{)}
    \PY{n}{outliers\PYZus{}zscore} \PY{o}{=} \PY{n}{pd}\PY{o}{.}\PY{n}{DataFrame}\PY{p}{(}\PY{n}{rows}\PY{p}{)}

\PY{n}{show\PYZus{}table}\PY{p}{(}\PY{l+s+s2}{\PYZdq{}}\PY{l+s+s2}{Outliers detectados por z\PYZhy{}score}\PY{l+s+s2}{\PYZdq{}}\PY{p}{,} \PY{n}{outliers\PYZus{}zscore}\PY{o}{.}\PY{n}{round}\PY{p}{(}\PY{l+m+mi}{4}\PY{p}{)}\PY{p}{,} \PY{n}{n}\PY{o}{=}\PY{k+kc}{None}\PY{p}{)}

\PY{k}{if} \PY{n}{outliers\PYZus{}comparison\PYZus{}path}\PY{o}{.}\PY{n}{exists}\PY{p}{(}\PY{p}{)}\PY{p}{:}
    \PY{n}{outliers\PYZus{}comparison} \PY{o}{=} \PY{n}{read\PYZus{}csv}\PY{p}{(}\PY{n}{outliers\PYZus{}comparison\PYZus{}path}\PY{p}{)}
\PY{k}{else}\PY{p}{:}
    \PY{n}{outliers\PYZus{}comparison} \PY{o}{=} \PY{n}{outliers\PYZus{}iqr}\PY{p}{[}\PY{p}{[}\PY{l+s+s2}{\PYZdq{}}\PY{l+s+s2}{feature}\PY{l+s+s2}{\PYZdq{}}\PY{p}{,} \PY{l+s+s2}{\PYZdq{}}\PY{l+s+s2}{outliers\PYZus{}iqr}\PY{l+s+s2}{\PYZdq{}}\PY{p}{]}\PY{p}{]}\PY{o}{.}\PY{n}{merge}\PY{p}{(}
        \PY{n}{outliers\PYZus{}zscore}\PY{p}{[}\PY{p}{[}\PY{l+s+s2}{\PYZdq{}}\PY{l+s+s2}{feature}\PY{l+s+s2}{\PYZdq{}}\PY{p}{,} \PY{l+s+s2}{\PYZdq{}}\PY{l+s+s2}{outliers\PYZus{}scipy}\PY{l+s+s2}{\PYZdq{}}\PY{p}{]}\PY{p}{]}\PY{p}{,} \PY{n}{on}\PY{o}{=}\PY{l+s+s2}{\PYZdq{}}\PY{l+s+s2}{feature}\PY{l+s+s2}{\PYZdq{}}\PY{p}{,} \PY{n}{how}\PY{o}{=}\PY{l+s+s2}{\PYZdq{}}\PY{l+s+s2}{left}\PY{l+s+s2}{\PYZdq{}}
    \PY{p}{)}\PY{o}{.}\PY{n}{rename}\PY{p}{(}\PY{n}{columns}\PY{o}{=}\PY{p}{\PYZob{}}\PY{l+s+s2}{\PYZdq{}}\PY{l+s+s2}{outliers\PYZus{}scipy}\PY{l+s+s2}{\PYZdq{}}\PY{p}{:} \PY{l+s+s2}{\PYZdq{}}\PY{l+s+s2}{outliers\PYZus{}zscore}\PY{l+s+s2}{\PYZdq{}}\PY{p}{\PYZcb{}}\PY{p}{)}
\PY{n}{show\PYZus{}table}\PY{p}{(}\PY{l+s+s2}{\PYZdq{}}\PY{l+s+s2}{Comparacion IQR vs z\PYZhy{}score}\PY{l+s+s2}{\PYZdq{}}\PY{p}{,} \PY{n}{outliers\PYZus{}comparison}\PY{o}{.}\PY{n}{round}\PY{p}{(}\PY{l+m+mi}{4}\PY{p}{)}\PY{p}{,} \PY{n}{n}\PY{o}{=}\PY{k+kc}{None}\PY{p}{)}

\PY{n}{outlier\PYZus{}features} \PY{o}{=} \PY{n}{outliers\PYZus{}iqr}\PY{o}{.}\PY{n}{loc}\PY{p}{[}\PY{n}{outliers\PYZus{}iqr}\PY{p}{[}\PY{l+s+s2}{\PYZdq{}}\PY{l+s+s2}{outliers\PYZus{}iqr}\PY{l+s+s2}{\PYZdq{}}\PY{p}{]} \PY{o}{\PYZgt{}} \PY{l+m+mi}{0}\PY{p}{,} \PY{l+s+s2}{\PYZdq{}}\PY{l+s+s2}{feature}\PY{l+s+s2}{\PYZdq{}}\PY{p}{]}\PY{o}{.}\PY{n}{tolist}\PY{p}{(}\PY{p}{)}
\PY{k}{if} \PY{n}{outlier\PYZus{}features}\PY{p}{:}
    \PY{n}{fig}\PY{p}{,} \PY{n}{axes} \PY{o}{=} \PY{n}{plt}\PY{o}{.}\PY{n}{subplots}\PY{p}{(}\PY{l+m+mi}{1}\PY{p}{,} \PY{n+nb}{len}\PY{p}{(}\PY{n}{outlier\PYZus{}features}\PY{p}{)}\PY{p}{,} \PY{n}{figsize}\PY{o}{=}\PY{p}{(}\PY{l+m+mf}{4.2} \PY{o}{*} \PY{n+nb}{len}\PY{p}{(}\PY{n}{outlier\PYZus{}features}\PY{p}{)}\PY{p}{,} \PY{l+m+mf}{4.2}\PY{p}{)}\PY{p}{)}
    \PY{n}{axes} \PY{o}{=} \PY{n}{np}\PY{o}{.}\PY{n}{atleast\PYZus{}1d}\PY{p}{(}\PY{n}{axes}\PY{p}{)}
    \PY{k}{for} \PY{n}{ax}\PY{p}{,} \PY{n}{feature} \PY{o+ow}{in} \PY{n+nb}{zip}\PY{p}{(}\PY{n}{axes}\PY{p}{,} \PY{n}{outlier\PYZus{}features}\PY{p}{)}\PY{p}{:}
        \PY{n}{sns}\PY{o}{.}\PY{n}{boxplot}\PY{p}{(}\PY{n}{data}\PY{o}{=}\PY{n}{df\PYZus{}clean}\PY{p}{,} \PY{n}{x}\PY{o}{=}\PY{n}{target}\PY{p}{,} \PY{n}{y}\PY{o}{=}\PY{n}{feature}\PY{p}{,} \PY{n}{ax}\PY{o}{=}\PY{n}{ax}\PY{p}{)}
        \PY{n}{ax}\PY{o}{.}\PY{n}{set\PYZus{}title}\PY{p}{(}\PY{l+s+sa}{f}\PY{l+s+s2}{\PYZdq{}}\PY{l+s+si}{\PYZob{}}\PY{n}{feature}\PY{l+s+si}{\PYZcb{}}\PY{l+s+s2}{ vs condition}\PY{l+s+s2}{\PYZdq{}}\PY{p}{)}
        \PY{n}{ax}\PY{o}{.}\PY{n}{set\PYZus{}xlabel}\PY{p}{(}\PY{l+s+s2}{\PYZdq{}}\PY{l+s+s2}{condition}\PY{l+s+s2}{\PYZdq{}}\PY{p}{)}
    \PY{n}{plt}\PY{o}{.}\PY{n}{tight\PYZus{}layout}\PY{p}{(}\PY{p}{)}
    \PY{n}{plt}\PY{o}{.}\PY{n}{show}\PY{p}{(}\PY{p}{)}
\end{Verbatim}
\end{tcolorbox}

    \textbf{Outliers detectados por IQR}

    
    
    \begin{Verbatim}[commandchars=\\\{\}]
      feature     q1     q3   iqr  limite\_inferior  limite\_superior  outliers\_iqr  porcentaje\_iqr
0   feature\_5   48.0   61.0  13.0             28.5             80.5             0          0.0000
1   feature\_6  211.0  276.0  65.0            113.5            373.5             5          1.6835
2  feature\_10  133.0  166.0  33.0             83.5            215.5             1          0.3367
3  feature\_11  120.0  140.0  20.0             90.0            170.0             9          3.0303
4  feature\_13    0.0    1.6   1.6             -2.4              4.0             5          1.6835
    \end{Verbatim}

    
    \textbf{Outliers detectados por z-score}

    
    
    \begin{Verbatim}[commandchars=\\\{\}]
      feature     media  desviacion\_std  umbral\_z  outliers\_scipy  porcentaje\_scipy  qq\_r\_scipy
0   feature\_5   54.5421          9.0497       3.0               0            0.0000      0.9936
1   feature\_6  247.3502         51.9976       3.0               4            1.3468      0.9714
2  feature\_10  149.5993         22.9416       3.0               1            0.3367      0.9885
3  feature\_11  131.6936         17.7628       3.0               2            0.6734      0.9830
4  feature\_13    1.0556          1.1661       3.0               2            0.6734      0.9211
    \end{Verbatim}

    
    \textbf{Comparacion IQR vs z-score}

    
    
    \begin{Verbatim}[commandchars=\\\{\}]
      feature  outliers\_iqr  outliers\_zscore  coinciden  solo\_iqr  solo\_zscore  acuerdo\_porcentaje
0  feature\_10             1                1          1         0            0            100.0000
1  feature\_11             9                2          2         7            0             22.2222
2  feature\_13             5                2          2         3            0             40.0000
3   feature\_5             0                0          0         0            0            100.0000
4   feature\_6             5                4          4         1            0             80.0000
    \end{Verbatim}

    
    \begin{center}
    \adjustimage{max size={0.9\linewidth}{0.9\paperheight}}{REPORT_files/REPORT_23_6.png}
    \end{center}
    { \hspace*{\fill} \\}
    
    La tabla IQR marca valores fuera del intervalo
\texttt{{[}Q1\ -\ 1.5*IQR,\ Q3\ +\ 1.5*IQR{]}}, por eso es adecuada
cuando no queremos asumir normalidad. Con este criterio aparecen
outliers en \texttt{feature\_6}, \texttt{feature\_10},
\texttt{feature\_11} y \texttt{feature\_13}, con porcentajes bajos: el
caso mas alto es \texttt{feature\_11} con 9 observaciones,
aproximadamente 3.03\% del dataset.

La tabla z-score usa otro criterio: mide cuantas desviaciones estandar
se aleja cada valor de la media y marca como outlier los casos con
\texttt{\textbar{}z\textbar{}\ \textgreater{}\ 3}. Este metodo es mas
estricto cuando la distribucion no es perfectamente normal; por eso
detecta menos outliers en \texttt{feature\_11} y \texttt{feature\_13}
que IQR. La comparacion IQR vs z-score permite ver si los extremos son
consistentes bajo dos reglas distintas o si dependen del metodo usado.

Los boxplots \texttt{feature\ vs\ condition} muestran donde caen esos
extremos dentro de cada clase. Esta grafica es importante porque no solo
indica que existen valores atipicos, sino si aparecen concentrados en
\texttt{condition=0}, en \texttt{condition=1} o en ambas clases. Si un
valor extremo aparece asociado a una clase medica, eliminarlo
automaticamente podria borrar informacion predictiva real.

Despues de observar las tablas y los boxplots, los outliers se
conservan. La razon no es ignorarlos, sino evitar una eliminacion
automatica que podria sesgar el dataset hacia pacientes mas promedio. Se
documentan como anomalias potenciales y se mantiene la posibilidad de
evaluar tratamientos especificos en modelos sensibles a extremos.

    \subsection{1.9. Correlaciones y
redundancia}\label{correlaciones-y-redundancia}

Se usa una matriz de correlacion de Pearson porque todas las variables
del dataset son numericas y se busca una primera lectura lineal global.
Pearson mide si dos variables tienden a aumentar o disminuir juntas de
forma aproximadamente lineal, con valores entre -1 y 1. En este punto
del analisis interesa precisamente esa lectura inicial: que features se
mueven en la misma direccion que \texttt{condition}, que features se
mueven en direccion contraria y que features podrian estar aportando
informacion redundante.

Existen otros tipos de matrices de correlacion, como Spearman o Kendall.
Spearman seria util si el objetivo principal fuera medir relaciones
monotonas basadas en rangos, especialmente con variables ordinales o
relaciones no lineales pero ordenadas. Kendall tambien se usa para
asociaciones ordinales, aunque suele ser mas conservador. En este
reporte se elige Pearson porque es simple, interpretable, consistente
con el analisis de medias y suficiente como diagnostico exploratorio
inicial antes de los modelos. Como varias variables son discretas u
ordinales, sus resultados se interpretan con cautela y no como prueba
definitiva.

La correlacion con \texttt{condition} sirve como herramienta
exploratoria para priorizar variables. Como \texttt{condition} esta
codificada como 0/1, una correlacion positiva indica que valores mas
altos de la feature tienden a aparecer mas en \texttt{condition=1}; una
correlacion negativa indica lo contrario. Esto no implica causalidad ni
diagnostico clinico, pero ayuda a identificar senales que los modelos
supervisados podrian aprovechar.

    \begin{tcolorbox}[breakable, size=fbox, boxrule=1pt, pad at break*=1mm,colback=cellbackground, colframe=cellborder]
\prompt{In}{incolor}{16}{\boxspacing}
\begin{Verbatim}[commandchars=\\\{\}]
\PY{n}{corr} \PY{o}{=} \PY{n}{df\PYZus{}clean}\PY{p}{[}\PY{n}{features} \PY{o}{+} \PY{p}{[}\PY{n}{target}\PY{p}{]}\PY{p}{]}\PY{o}{.}\PY{n}{corr}\PY{p}{(}\PY{n}{numeric\PYZus{}only}\PY{o}{=}\PY{k+kc}{True}\PY{p}{)}
\PY{n}{fig}\PY{p}{,} \PY{n}{ax} \PY{o}{=} \PY{n}{plt}\PY{o}{.}\PY{n}{subplots}\PY{p}{(}\PY{n}{figsize}\PY{o}{=}\PY{p}{(}\PY{l+m+mi}{11}\PY{p}{,} \PY{l+m+mi}{9}\PY{p}{)}\PY{p}{)}
\PY{n}{sns}\PY{o}{.}\PY{n}{heatmap}\PY{p}{(}\PY{n}{corr}\PY{p}{,} \PY{n}{cmap}\PY{o}{=}\PY{l+s+s2}{\PYZdq{}}\PY{l+s+s2}{vlag}\PY{l+s+s2}{\PYZdq{}}\PY{p}{,} \PY{n}{center}\PY{o}{=}\PY{l+m+mi}{0}\PY{p}{,} \PY{n}{linewidths}\PY{o}{=}\PY{l+m+mf}{0.4}\PY{p}{,} \PY{n}{annot}\PY{o}{=}\PY{k+kc}{False}\PY{p}{,} \PY{n}{ax}\PY{o}{=}\PY{n}{ax}\PY{p}{)}
\PY{n}{ax}\PY{o}{.}\PY{n}{set\PYZus{}title}\PY{p}{(}\PY{l+s+s2}{\PYZdq{}}\PY{l+s+s2}{Matriz de correlacion de variables numericas}\PY{l+s+s2}{\PYZdq{}}\PY{p}{)}
\PY{n}{plt}\PY{o}{.}\PY{n}{tight\PYZus{}layout}\PY{p}{(}\PY{p}{)}
\PY{n}{plt}\PY{o}{.}\PY{n}{show}\PY{p}{(}\PY{p}{)}

\PY{n}{target\PYZus{}corr} \PY{o}{=} \PY{n}{corr}\PY{p}{[}\PY{n}{target}\PY{p}{]}\PY{o}{.}\PY{n}{drop}\PY{p}{(}\PY{n}{target}\PY{p}{)}\PY{o}{.}\PY{n}{sort\PYZus{}values}\PY{p}{(}\PY{n}{key}\PY{o}{=}\PY{k}{lambda} \PY{n}{values}\PY{p}{:} \PY{n}{values}\PY{o}{.}\PY{n}{abs}\PY{p}{(}\PY{p}{)}\PY{p}{,} \PY{n}{ascending}\PY{o}{=}\PY{k+kc}{False}\PY{p}{)}\PY{o}{.}\PY{n}{reset\PYZus{}index}\PY{p}{(}\PY{p}{)}
\PY{n}{target\PYZus{}corr}\PY{o}{.}\PY{n}{columns} \PY{o}{=} \PY{p}{[}\PY{l+s+s2}{\PYZdq{}}\PY{l+s+s2}{feature}\PY{l+s+s2}{\PYZdq{}}\PY{p}{,} \PY{l+s+s2}{\PYZdq{}}\PY{l+s+s2}{corr\PYZus{}condition}\PY{l+s+s2}{\PYZdq{}}\PY{p}{]}
\PY{n}{show\PYZus{}table}\PY{p}{(}\PY{l+s+s2}{\PYZdq{}}\PY{l+s+s2}{Correlacion de features con condition}\PY{l+s+s2}{\PYZdq{}}\PY{p}{,} \PY{n}{target\PYZus{}corr}\PY{o}{.}\PY{n}{round}\PY{p}{(}\PY{l+m+mi}{5}\PY{p}{)}\PY{p}{,} \PY{n}{n}\PY{o}{=}\PY{k+kc}{None}\PY{p}{)}

\PY{n}{corr\PYZus{}features} \PY{o}{=} \PY{n}{df\PYZus{}clean}\PY{p}{[}\PY{n}{features}\PY{p}{]}\PY{o}{.}\PY{n}{corr}\PY{p}{(}\PY{p}{)}\PY{o}{.}\PY{n}{abs}\PY{p}{(}\PY{p}{)}
\PY{n}{upper} \PY{o}{=} \PY{n}{corr\PYZus{}features}\PY{o}{.}\PY{n}{where}\PY{p}{(}\PY{n}{np}\PY{o}{.}\PY{n}{triu}\PY{p}{(}\PY{n}{np}\PY{o}{.}\PY{n}{ones}\PY{p}{(}\PY{n}{corr\PYZus{}features}\PY{o}{.}\PY{n}{shape}\PY{p}{)}\PY{p}{,} \PY{n}{k}\PY{o}{=}\PY{l+m+mi}{1}\PY{p}{)}\PY{o}{.}\PY{n}{astype}\PY{p}{(}\PY{n+nb}{bool}\PY{p}{)}\PY{p}{)}
\PY{n}{high\PYZus{}corr\PYZus{}pairs} \PY{o}{=} \PY{p}{(}
    \PY{n}{upper}\PY{o}{.}\PY{n}{stack}\PY{p}{(}\PY{p}{)}
    \PY{o}{.}\PY{n}{reset\PYZus{}index}\PY{p}{(}\PY{p}{)}
    \PY{o}{.}\PY{n}{rename}\PY{p}{(}\PY{n}{columns}\PY{o}{=}\PY{p}{\PYZob{}}\PY{l+s+s2}{\PYZdq{}}\PY{l+s+s2}{level\PYZus{}0}\PY{l+s+s2}{\PYZdq{}}\PY{p}{:} \PY{l+s+s2}{\PYZdq{}}\PY{l+s+s2}{feature\PYZus{}a}\PY{l+s+s2}{\PYZdq{}}\PY{p}{,} \PY{l+s+s2}{\PYZdq{}}\PY{l+s+s2}{level\PYZus{}1}\PY{l+s+s2}{\PYZdq{}}\PY{p}{:} \PY{l+s+s2}{\PYZdq{}}\PY{l+s+s2}{feature\PYZus{}b}\PY{l+s+s2}{\PYZdq{}}\PY{p}{,} \PY{l+m+mi}{0}\PY{p}{:} \PY{l+s+s2}{\PYZdq{}}\PY{l+s+s2}{abs\PYZus{}corr}\PY{l+s+s2}{\PYZdq{}}\PY{p}{\PYZcb{}}\PY{p}{)}
    \PY{o}{.}\PY{n}{sort\PYZus{}values}\PY{p}{(}\PY{l+s+s2}{\PYZdq{}}\PY{l+s+s2}{abs\PYZus{}corr}\PY{l+s+s2}{\PYZdq{}}\PY{p}{,} \PY{n}{ascending}\PY{o}{=}\PY{k+kc}{False}\PY{p}{)}
    \PY{o}{.}\PY{n}{head}\PY{p}{(}\PY{l+m+mi}{10}\PY{p}{)}
\PY{p}{)}
\PY{n}{show\PYZus{}table}\PY{p}{(}\PY{l+s+s2}{\PYZdq{}}\PY{l+s+s2}{Top 10 pares de features por correlacion absoluta}\PY{l+s+s2}{\PYZdq{}}\PY{p}{,} \PY{n}{high\PYZus{}corr\PYZus{}pairs}\PY{o}{.}\PY{n}{round}\PY{p}{(}\PY{l+m+mi}{5}\PY{p}{)}\PY{p}{,} \PY{n}{n}\PY{o}{=}\PY{k+kc}{None}\PY{p}{)}
\end{Verbatim}
\end{tcolorbox}

    \begin{center}
    \adjustimage{max size={0.9\linewidth}{0.9\paperheight}}{REPORT_files/REPORT_26_0.png}
    \end{center}
    { \hspace*{\fill} \\}
    
    \textbf{Correlacion de features con condition}

    
    
    \begin{Verbatim}[commandchars=\\\{\}]
       feature  corr\_condition
0   feature\_12         0.52052
1    feature\_8         0.46319
2   feature\_13         0.42405
3   feature\_10        -0.42382
4    feature\_4         0.42136
5    feature\_2         0.40894
6    feature\_3         0.33305
7    feature\_7         0.27847
8    feature\_5         0.22708
9    feature\_9         0.16634
10  feature\_11         0.15349
11   feature\_6         0.08028
12   feature\_1         0.00317
    \end{Verbatim}

    
    \textbf{Top 10 pares de features por correlacion absoluta}

    
    
    \begin{Verbatim}[commandchars=\\\{\}]
     feature\_a   feature\_b  abs\_corr
32   feature\_3  feature\_13   0.57904
46   feature\_5  feature\_10   0.39456
29   feature\_3  feature\_10   0.38931
38   feature\_4  feature\_10   0.38437
13   feature\_2   feature\_4   0.37752
61   feature\_7  feature\_12   0.37056
44   feature\_5   feature\_8   0.36221
74  feature\_10  feature\_13   0.34764
19   feature\_2  feature\_10   0.33931
77  feature\_12  feature\_13   0.33681
    \end{Verbatim}

    
    Sobre la redundancia podemos afirmar que no aparecen correlaciones
absolutas cercanas a 1 entre features. Por tanto, no se justifica
eliminar variables en esta primera pasada, por multicolinealidad
extrema; es mejor conservarlas y dejar que los modelos posteriores
gestionen la importancia relativa.

    \subsection{1.10. Transformacion para modelos
matematicos}\label{transformacion-para-modelos-matematicos}

La transformacion se hace para separar dos necesidades distintas. Por un
lado, \texttt{foot\_clean.csv} conserva las variables en su escala
original, lo cual es mejor para interpretar rangos, medias, boxplots y
diferencias entre clases. Por otro lado, \texttt{foot\_model\_ready.csv}
prepara las mismas observaciones para algoritmos matematicos que son
sensibles a la escala de las variables.

La estandarizacion transforma cada feature numerica restando su media y
dividiendo por su desviacion estandar. Asi, todas quedan en una escala
comparable, con media cercana a 0 y desviacion estandar cercana a 1.
Esto es importante porque variables como \texttt{feature\_6} o
\texttt{feature\_10} tienen rangos mucho mayores que variables binarias
o discretas; si no se estandarizan, podrian dominar PCA, KNN, SVM o
modelos regularizados simplemente por tener numeros mas grandes, no
necesariamente por contener mas informacion clinica.

No se aplica one-hot encoding porque no hay variables categoricas
textuales: todas las columnas ya son numericas. Las variables binarias
se mantienen como 0/1 porque esa codificacion ya es adecuada. Las
variables discretas u ordinales se conservan en su codificacion
original, pero tambien quedan escaladas en la version model-ready para
que los modelos basados en distancia no las traten de forma
desproporcionada.

La variable \texttt{condition} no se estandariza ni se transforma porque
es la etiqueta objetivo. Modificarla mezclaria el significado de las
clases y podria introducir errores en la evaluacion. Por eso la
preparacion transforma solo las features y mantiene \texttt{condition}
como 0/1.

    \begin{tcolorbox}[breakable, size=fbox, boxrule=1pt, pad at break*=1mm,colback=cellbackground, colframe=cellborder]
\prompt{In}{incolor}{17}{\boxspacing}
\begin{Verbatim}[commandchars=\\\{\}]
\PY{n}{transformation\PYZus{}decisions} \PY{o}{=} \PY{n}{pd}\PY{o}{.}\PY{n}{DataFrame}\PY{p}{(}\PY{p}{[}
    \PY{p}{\PYZob{}}
        \PY{l+s+s2}{\PYZdq{}}\PY{l+s+s2}{transformacion}\PY{l+s+s2}{\PYZdq{}}\PY{p}{:} \PY{l+s+s2}{\PYZdq{}}\PY{l+s+s2}{Mantener binarias 0/1}\PY{l+s+s2}{\PYZdq{}}\PY{p}{,}
        \PY{l+s+s2}{\PYZdq{}}\PY{l+s+s2}{aplicada\PYZus{}a}\PY{l+s+s2}{\PYZdq{}}\PY{p}{:} \PY{l+s+s2}{\PYZdq{}}\PY{l+s+s2}{features binarias}\PY{l+s+s2}{\PYZdq{}}\PY{p}{,}
        \PY{l+s+s2}{\PYZdq{}}\PY{l+s+s2}{justificacion}\PY{l+s+s2}{\PYZdq{}}\PY{p}{:} \PY{l+s+s2}{\PYZdq{}}\PY{l+s+s2}{Ya expresan presencia/ausencia; one\PYZhy{}hot duplicaria informacion y no aporta valor.}\PY{l+s+s2}{\PYZdq{}}\PY{p}{,}
    \PY{p}{\PYZcb{}}\PY{p}{,}
    \PY{p}{\PYZob{}}
        \PY{l+s+s2}{\PYZdq{}}\PY{l+s+s2}{transformacion}\PY{l+s+s2}{\PYZdq{}}\PY{p}{:} \PY{l+s+s2}{\PYZdq{}}\PY{l+s+s2}{Conservar discretas/ordinales}\PY{l+s+s2}{\PYZdq{}}\PY{p}{,}
        \PY{l+s+s2}{\PYZdq{}}\PY{l+s+s2}{aplicada\PYZus{}a}\PY{l+s+s2}{\PYZdq{}}\PY{p}{:} \PY{l+s+s2}{\PYZdq{}}\PY{l+s+s2}{features con pocos niveles}\PY{l+s+s2}{\PYZdq{}}\PY{p}{,}
        \PY{l+s+s2}{\PYZdq{}}\PY{l+s+s2}{justificacion}\PY{l+s+s2}{\PYZdq{}}\PY{p}{:} \PY{l+s+s2}{\PYZdq{}}\PY{l+s+s2}{Son codigos numericos anonimizados; se preservan para no inventar categorias no observadas.}\PY{l+s+s2}{\PYZdq{}}\PY{p}{,}
    \PY{p}{\PYZcb{}}\PY{p}{,}
    \PY{p}{\PYZob{}}
        \PY{l+s+s2}{\PYZdq{}}\PY{l+s+s2}{transformacion}\PY{l+s+s2}{\PYZdq{}}\PY{p}{:} \PY{l+s+s2}{\PYZdq{}}\PY{l+s+s2}{Estandarizar features}\PY{l+s+s2}{\PYZdq{}}\PY{p}{,}
        \PY{l+s+s2}{\PYZdq{}}\PY{l+s+s2}{aplicada\PYZus{}a}\PY{l+s+s2}{\PYZdq{}}\PY{p}{:} \PY{l+s+s2}{\PYZdq{}}\PY{l+s+s2}{columnas predictoras numericas}\PY{l+s+s2}{\PYZdq{}}\PY{p}{,}
        \PY{l+s+s2}{\PYZdq{}}\PY{l+s+s2}{justificacion}\PY{l+s+s2}{\PYZdq{}}\PY{p}{:} \PY{l+s+s2}{\PYZdq{}}\PY{l+s+s2}{Evita que variables con rangos grandes dominen PCA, KNN, SVM o regularizacion.}\PY{l+s+s2}{\PYZdq{}}\PY{p}{,}
    \PY{p}{\PYZcb{}}\PY{p}{,}
    \PY{p}{\PYZob{}}
        \PY{l+s+s2}{\PYZdq{}}\PY{l+s+s2}{transformacion}\PY{l+s+s2}{\PYZdq{}}\PY{p}{:} \PY{l+s+s2}{\PYZdq{}}\PY{l+s+s2}{No imputar}\PY{l+s+s2}{\PYZdq{}}\PY{p}{,}
        \PY{l+s+s2}{\PYZdq{}}\PY{l+s+s2}{aplicada\PYZus{}a}\PY{l+s+s2}{\PYZdq{}}\PY{p}{:} \PY{l+s+s2}{\PYZdq{}}\PY{l+s+s2}{dataset actual}\PY{l+s+s2}{\PYZdq{}}\PY{p}{,}
        \PY{l+s+s2}{\PYZdq{}}\PY{l+s+s2}{justificacion}\PY{l+s+s2}{\PYZdq{}}\PY{p}{:} \PY{l+s+s2}{\PYZdq{}}\PY{l+s+s2}{No hay valores faltantes; imputar crearia valores artificiales innecesarios.}\PY{l+s+s2}{\PYZdq{}}\PY{p}{,}
    \PY{p}{\PYZcb{}}\PY{p}{,}
    \PY{p}{\PYZob{}}
        \PY{l+s+s2}{\PYZdq{}}\PY{l+s+s2}{transformacion}\PY{l+s+s2}{\PYZdq{}}\PY{p}{:} \PY{l+s+s2}{\PYZdq{}}\PY{l+s+s2}{No eliminar outliers en esta etapa}\PY{l+s+s2}{\PYZdq{}}\PY{p}{,}
        \PY{l+s+s2}{\PYZdq{}}\PY{l+s+s2}{aplicada\PYZus{}a}\PY{l+s+s2}{\PYZdq{}}\PY{p}{:} \PY{l+s+s2}{\PYZdq{}}\PY{l+s+s2}{variables continuas con extremos}\PY{l+s+s2}{\PYZdq{}}\PY{p}{,}
        \PY{l+s+s2}{\PYZdq{}}\PY{l+s+s2}{justificacion}\PY{l+s+s2}{\PYZdq{}}\PY{p}{:} \PY{l+s+s2}{\PYZdq{}}\PY{l+s+s2}{En contexto medico pueden representar pacientes reales; se documentan y se conservan.}\PY{l+s+s2}{\PYZdq{}}\PY{p}{,}
    \PY{p}{\PYZcb{}}\PY{p}{,}
    \PY{p}{\PYZob{}}
        \PY{l+s+s2}{\PYZdq{}}\PY{l+s+s2}{transformacion}\PY{l+s+s2}{\PYZdq{}}\PY{p}{:} \PY{l+s+s2}{\PYZdq{}}\PY{l+s+s2}{Mantener condition sin escalar}\PY{l+s+s2}{\PYZdq{}}\PY{p}{,}
        \PY{l+s+s2}{\PYZdq{}}\PY{l+s+s2}{aplicada\PYZus{}a}\PY{l+s+s2}{\PYZdq{}}\PY{p}{:} \PY{l+s+s2}{\PYZdq{}}\PY{l+s+s2}{variable objetivo}\PY{l+s+s2}{\PYZdq{}}\PY{p}{,}
        \PY{l+s+s2}{\PYZdq{}}\PY{l+s+s2}{justificacion}\PY{l+s+s2}{\PYZdq{}}\PY{p}{:} \PY{l+s+s2}{\PYZdq{}}\PY{l+s+s2}{Debe conservar el significado de clase binaria 0/1 para entrenamiento y evaluacion.}\PY{l+s+s2}{\PYZdq{}}\PY{p}{,}
    \PY{p}{\PYZcb{}}\PY{p}{,}
\PY{p}{]}\PY{p}{)}
\PY{n}{show\PYZus{}table}\PY{p}{(}\PY{l+s+s2}{\PYZdq{}}\PY{l+s+s2}{Decisiones de transformacion}\PY{l+s+s2}{\PYZdq{}}\PY{p}{,} \PY{n}{transformation\PYZus{}decisions}\PY{p}{,} \PY{n}{n}\PY{o}{=}\PY{k+kc}{None}\PY{p}{)}

\PY{n}{model\PYZus{}ready} \PY{o}{=} \PY{n}{read\PYZus{}csv}\PY{p}{(}\PY{n}{MODEL\PYZus{}READY\PYZus{}CSV}\PY{p}{)}
\PY{n}{show\PYZus{}table}\PY{p}{(}\PY{l+s+s2}{\PYZdq{}}\PY{l+s+s2}{Dataset model\PYZhy{}ready estandarizado}\PY{l+s+s2}{\PYZdq{}}\PY{p}{,} \PY{n}{model\PYZus{}ready}\PY{p}{,} \PY{n}{n}\PY{o}{=}\PY{l+m+mi}{5}\PY{p}{)}

\PY{n}{standardized\PYZus{}check} \PY{o}{=} \PY{n}{model\PYZus{}ready}\PY{o}{.}\PY{n}{drop}\PY{p}{(}\PY{n}{columns}\PY{o}{=}\PY{p}{[}\PY{n}{target}\PY{p}{]}\PY{p}{)}\PY{o}{.}\PY{n}{agg}\PY{p}{(}\PY{p}{[}\PY{l+s+s2}{\PYZdq{}}\PY{l+s+s2}{mean}\PY{l+s+s2}{\PYZdq{}}\PY{p}{,} \PY{l+s+s2}{\PYZdq{}}\PY{l+s+s2}{std}\PY{l+s+s2}{\PYZdq{}}\PY{p}{]}\PY{p}{)}\PY{o}{.}\PY{n}{T}\PY{o}{.}\PY{n}{round}\PY{p}{(}\PY{l+m+mi}{4}\PY{p}{)}
\PY{n}{show\PYZus{}table}\PY{p}{(}\PY{l+s+s2}{\PYZdq{}}\PY{l+s+s2}{Comprobacion de medias y desviaciones tras estandarizacion}\PY{l+s+s2}{\PYZdq{}}\PY{p}{,} \PY{n}{standardized\PYZus{}check}\PY{o}{.}\PY{n}{reset\PYZus{}index}\PY{p}{(}\PY{n}{names}\PY{o}{=}\PY{l+s+s2}{\PYZdq{}}\PY{l+s+s2}{feature}\PY{l+s+s2}{\PYZdq{}}\PY{p}{)}\PY{p}{,} \PY{n}{n}\PY{o}{=}\PY{k+kc}{None}\PY{p}{)}
\PY{n+nb}{print}\PY{p}{(}\PY{l+s+s2}{\PYZdq{}}\PY{l+s+s2}{Valores de condition en model\PYZhy{}ready:}\PY{l+s+s2}{\PYZdq{}}\PY{p}{,} \PY{n+nb}{sorted}\PY{p}{(}\PY{n}{model\PYZus{}ready}\PY{p}{[}\PY{n}{target}\PY{p}{]}\PY{o}{.}\PY{n}{unique}\PY{p}{(}\PY{p}{)}\PY{o}{.}\PY{n}{tolist}\PY{p}{(}\PY{p}{)}\PY{p}{)}\PY{p}{)}
\end{Verbatim}
\end{tcolorbox}

    \textbf{Decisiones de transformacion}

    
    
    \begin{Verbatim}[commandchars=\\\{\}]
                       transformacion                        aplicada\_a                                      justificacion
0               Mantener binarias 0/1                 features binarias  Ya expresan presencia/ausencia; one-hot duplic{\ldots}
1       Conservar discretas/ordinales        features con pocos niveles  Son codigos numericos anonimizados; se preserv{\ldots}
2               Estandarizar features    columnas predictoras numericas  Evita que variables con rangos grandes dominen{\ldots}
3                          No imputar                    dataset actual  No hay valores faltantes; imputar crearia valo{\ldots}
4  No eliminar outliers en esta etapa  variables continuas con extremos  En contexto medico pueden representar paciente{\ldots}
5      Mantener condition sin escalar                 variable objetivo  Debe conservar el significado de clase binaria{\ldots}
    \end{Verbatim}

    
    \textbf{Dataset model-ready estandarizado}

    
    
    \begin{Verbatim}[commandchars=\\\{\}]
   feature\_1  feature\_2  feature\_3  feature\_4  feature\_5  feature\_6  feature\_7  feature\_8  feature\_9  feature\_10  feature\_11  feature\_12  feature\_13  \textbackslash{}
0   2.426332   0.872408   0.642696   1.433497   1.266105  -0.372136  -1.444542   1.409246  -1.001728    0.671300    2.606930    1.217722   -0.047641   
1  -0.410757   0.872408  -0.974938   1.433497   0.492601  -1.352951   0.689930   0.344243  -1.001728    0.540533    0.467629    1.217722   -0.905184   
2  -0.410757  -1.200433   0.642696  -0.695246  -1.496407  -0.852928   0.689930  -0.720760  -1.001728   -0.767137    0.186142    0.172452   -0.905184   
3  -0.410757  -0.164013  -0.974938  -0.695246  -1.938409  -0.622148  -1.444542  -0.720760  -1.001728    0.889245   -0.658320   -0.872818   -0.905184   
4  -0.410757  -2.236854   0.642696  -0.695246   1.045104  -0.391368   0.689930  -0.720760   1.008496    0.235410    2.156551    1.217722   -0.390658   

   condition  
0          1  
1          1  
2          0  
3          0  
4          0  
    \end{Verbatim}

    
    \textbf{Comprobacion de medias y desviaciones tras estandarizacion}

    
    
    \begin{Verbatim}[commandchars=\\\{\}]
       feature  mean  std
0    feature\_1  -0.0  1.0
1    feature\_2  -0.0  1.0
2    feature\_3  -0.0  1.0
3    feature\_4  -0.0  1.0
4    feature\_5  -0.0  1.0
5    feature\_6  -0.0  1.0
6    feature\_7  -0.0  1.0
7    feature\_8  -0.0  1.0
8    feature\_9   0.0  1.0
9   feature\_10   0.0  1.0
10  feature\_11   0.0  1.0
11  feature\_12   0.0  1.0
12  feature\_13  -0.0  1.0
    \end{Verbatim}

    
    \begin{Verbatim}[commandchars=\\\{\}]
Valores de condition en model-ready: [0, 1]
    \end{Verbatim}

    \section{2. Aprendizaje No
Supervisado}\label{aprendizaje-no-supervisado}

En esta parte se buscan estructuras ocultas de pacientes sin usar
\texttt{condition} durante el entrenamiento. Esto es importante porque
el objetivo del aprendizaje no supervisado no es predecir la etiqueta,
sino descubrir si las features contienen agrupaciones, gradientes o
patrones de similitud propios.

La variable \texttt{condition} se usa solo despues, como lectura
interpretativa externa: si un cluster concentra mas pacientes con
\texttt{condition=1}, eso sugiere que el agrupamiento capturo
informacion clinicamente relacionada con la condicion. Aun asi, esta
relacion no convierte al cluster en diagnostico. Los clusters se
interpretan como perfiles exploratorios y generadores de hipotesis.

Tambien hay una diferencia intencional frente a la preparacion
descriptiva inicial. Para describir los datos se conservaron los codigos
originales porque son mas faciles de interpretar. Para PCA y clustering,
en cambio, las variables discretas se codifican con one-hot encoding
porque estos metodos dependen de distancias: tratar codigos como 0, 1, 2
o 3 como magnitudes continuas impondria una relacion ordinal que no esta
garantizada en variables anonimizadas.

    \subsection{2.1. Reduccion de dimensionalidad con
PCA}\label{reduccion-de-dimensionalidad-con-pca}

PCA se utiliza para reducir complejidad y visualizar tendencias globales
en dos dimensiones. Es adecuado aqui porque transforma features
correlacionadas o redundantes en componentes ortogonales y permite medir
cuanta informacion se conserva al pasar de muchas variables procesadas a
pocas componentes.

No se interpreta como modelo diagnostico. Las dos primeras componentes
explican solo una parte limitada de la varianza total, por lo que los
graficos 2D ayudan a observar patrones generales, pero no son prueba de
separacion fuerte entre pacientes.

    \begin{tcolorbox}[breakable, size=fbox, boxrule=1pt, pad at break*=1mm,colback=cellbackground, colframe=cellborder]
\prompt{In}{incolor}{18}{\boxspacing}
\begin{Verbatim}[commandchars=\\\{\}]
\PY{k+kn}{from}\PY{+w}{ }\PY{n+nn}{sklearn}\PY{n+nn}{.}\PY{n+nn}{decomposition}\PY{+w}{ }\PY{k+kn}{import} \PY{n}{PCA}
\PY{k+kn}{from}\PY{+w}{ }\PY{n+nn}{task2\PYZus{}common}\PY{+w}{ }\PY{k+kn}{import} \PY{n}{CATEGORICAL\PYZus{}FEATURES}\PY{p}{,} \PY{n}{prepare\PYZus{}features}

\PY{n}{pca\PYZus{}var} \PY{o}{=} \PY{n}{read\PYZus{}csv}\PY{p}{(}\PY{n}{OUTPUTS\PYZus{}DIR} \PY{o}{/} \PY{l+s+s2}{\PYZdq{}}\PY{l+s+s2}{task2}\PY{l+s+s2}{\PYZdq{}} \PY{o}{/} \PY{l+s+s2}{\PYZdq{}}\PY{l+s+s2}{pca\PYZus{}varianza.csv}\PY{l+s+s2}{\PYZdq{}}\PY{p}{)}
\PY{n}{show\PYZus{}table}\PY{p}{(}\PY{l+s+s2}{\PYZdq{}}\PY{l+s+s2}{Varianza explicada por PCA}\PY{l+s+s2}{\PYZdq{}}\PY{p}{,} \PY{n}{pca\PYZus{}var}\PY{p}{,} \PY{n}{n}\PY{o}{=}\PY{k+kc}{None}\PY{p}{)}

\PY{n}{pc2\PYZus{}total} \PY{o}{=} \PY{n+nb}{float}\PY{p}{(}\PY{n}{pca\PYZus{}var}\PY{o}{.}\PY{n}{loc}\PY{p}{[}\PY{n}{pca\PYZus{}var}\PY{p}{[}\PY{l+s+s2}{\PYZdq{}}\PY{l+s+s2}{componente}\PY{l+s+s2}{\PYZdq{}}\PY{p}{]} \PY{o}{==} \PY{l+s+s2}{\PYZdq{}}\PY{l+s+s2}{PC2}\PY{l+s+s2}{\PYZdq{}}\PY{p}{,} \PY{l+s+s2}{\PYZdq{}}\PY{l+s+s2}{varianza\PYZus{}acumulada}\PY{l+s+s2}{\PYZdq{}}\PY{p}{]}\PY{o}{.}\PY{n}{iloc}\PY{p}{[}\PY{l+m+mi}{0}\PY{p}{]}\PY{p}{)}
\PY{n}{components\PYZus{}80} \PY{o}{=} \PY{n+nb}{int}\PY{p}{(}\PY{p}{(}\PY{n}{pca\PYZus{}var}\PY{p}{[}\PY{l+s+s2}{\PYZdq{}}\PY{l+s+s2}{varianza\PYZus{}acumulada}\PY{l+s+s2}{\PYZdq{}}\PY{p}{]} \PY{o}{\PYZgt{}}\PY{o}{=} \PY{l+m+mf}{0.80}\PY{p}{)}\PY{o}{.}\PY{n}{idxmax}\PY{p}{(}\PY{p}{)} \PY{o}{+} \PY{l+m+mi}{1}\PY{p}{)}
\PY{n}{components\PYZus{}90} \PY{o}{=} \PY{n+nb}{int}\PY{p}{(}\PY{p}{(}\PY{n}{pca\PYZus{}var}\PY{p}{[}\PY{l+s+s2}{\PYZdq{}}\PY{l+s+s2}{varianza\PYZus{}acumulada}\PY{l+s+s2}{\PYZdq{}}\PY{p}{]} \PY{o}{\PYZgt{}}\PY{o}{=} \PY{l+m+mf}{0.90}\PY{p}{)}\PY{o}{.}\PY{n}{idxmax}\PY{p}{(}\PY{p}{)} \PY{o}{+} \PY{l+m+mi}{1}\PY{p}{)}

\PY{n}{pca\PYZus{}summary} \PY{o}{=} \PY{n}{pd}\PY{o}{.}\PY{n}{DataFrame}\PY{p}{(}\PY{p}{[}
    \PY{p}{\PYZob{}}
        \PY{l+s+s2}{\PYZdq{}}\PY{l+s+s2}{lectura}\PY{l+s+s2}{\PYZdq{}}\PY{p}{:} \PY{l+s+s2}{\PYZdq{}}\PY{l+s+s2}{Varianza acumulada PC1 + PC2}\PY{l+s+s2}{\PYZdq{}}\PY{p}{,}
        \PY{l+s+s2}{\PYZdq{}}\PY{l+s+s2}{valor}\PY{l+s+s2}{\PYZdq{}}\PY{p}{:} \PY{l+s+sa}{f}\PY{l+s+s2}{\PYZdq{}}\PY{l+s+si}{\PYZob{}}\PY{n}{pc2\PYZus{}total}\PY{+w}{ }\PY{o}{*}\PY{+w}{ }\PY{l+m+mi}{100}\PY{l+s+si}{:}\PY{l+s+s2}{.2f}\PY{l+s+si}{\PYZcb{}}\PY{l+s+s2}{\PYZpc{}}\PY{l+s+s2}{\PYZdq{}}\PY{p}{,}
        \PY{l+s+s2}{\PYZdq{}}\PY{l+s+s2}{interpretacion}\PY{l+s+s2}{\PYZdq{}}\PY{p}{:} \PY{l+s+s2}{\PYZdq{}}\PY{l+s+s2}{La visualizacion 2D resume una parte util, pero insuficiente para concluir separacion fuerte.}\PY{l+s+s2}{\PYZdq{}}\PY{p}{,}
    \PY{p}{\PYZcb{}}\PY{p}{,}
    \PY{p}{\PYZob{}}
        \PY{l+s+s2}{\PYZdq{}}\PY{l+s+s2}{lectura}\PY{l+s+s2}{\PYZdq{}}\PY{p}{:} \PY{l+s+s2}{\PYZdq{}}\PY{l+s+s2}{Componentes para \PYZgt{}=80}\PY{l+s+si}{\PYZpc{} d}\PY{l+s+s2}{e varianza}\PY{l+s+s2}{\PYZdq{}}\PY{p}{,}
        \PY{l+s+s2}{\PYZdq{}}\PY{l+s+s2}{valor}\PY{l+s+s2}{\PYZdq{}}\PY{p}{:} \PY{n}{components\PYZus{}80}\PY{p}{,}
        \PY{l+s+s2}{\PYZdq{}}\PY{l+s+s2}{interpretacion}\PY{l+s+s2}{\PYZdq{}}\PY{p}{:} \PY{l+s+s2}{\PYZdq{}}\PY{l+s+s2}{La informacion relevante esta distribuida en varias dimensiones.}\PY{l+s+s2}{\PYZdq{}}\PY{p}{,}
    \PY{p}{\PYZcb{}}\PY{p}{,}
    \PY{p}{\PYZob{}}
        \PY{l+s+s2}{\PYZdq{}}\PY{l+s+s2}{lectura}\PY{l+s+s2}{\PYZdq{}}\PY{p}{:} \PY{l+s+s2}{\PYZdq{}}\PY{l+s+s2}{Componentes para \PYZgt{}=90}\PY{l+s+si}{\PYZpc{} d}\PY{l+s+s2}{e varianza}\PY{l+s+s2}{\PYZdq{}}\PY{p}{,}
        \PY{l+s+s2}{\PYZdq{}}\PY{l+s+s2}{valor}\PY{l+s+s2}{\PYZdq{}}\PY{p}{:} \PY{n}{components\PYZus{}90}\PY{p}{,}
        \PY{l+s+s2}{\PYZdq{}}\PY{l+s+s2}{interpretacion}\PY{l+s+s2}{\PYZdq{}}\PY{p}{:} \PY{l+s+s2}{\PYZdq{}}\PY{l+s+s2}{Reducir a solo dos componentes implica perder bastante informacion.}\PY{l+s+s2}{\PYZdq{}}\PY{p}{,}
    \PY{p}{\PYZcb{}}\PY{p}{,}
\PY{p}{]}\PY{p}{)}
\PY{n}{show\PYZus{}table}\PY{p}{(}\PY{l+s+s2}{\PYZdq{}}\PY{l+s+s2}{Lectura de reduccion de dimensionalidad}\PY{l+s+s2}{\PYZdq{}}\PY{p}{,} \PY{n}{pca\PYZus{}summary}\PY{p}{,} \PY{n}{n}\PY{o}{=}\PY{k+kc}{None}\PY{p}{)}

\PY{n}{\PYZus{}}\PY{p}{,} \PY{n}{\PYZus{}}\PY{p}{,} \PY{n}{x\PYZus{}task2\PYZus{}processed}\PY{p}{,} \PY{n}{task2\PYZus{}feature\PYZus{}names} \PY{o}{=} \PY{n}{prepare\PYZus{}features}\PY{p}{(}\PY{n}{df\PYZus{}clean}\PY{p}{)}
\PY{n}{pca\PYZus{}for\PYZus{}loadings} \PY{o}{=} \PY{n}{PCA}\PY{p}{(}\PY{n}{n\PYZus{}components}\PY{o}{=}\PY{l+m+mi}{2}\PY{p}{,} \PY{n}{random\PYZus{}state}\PY{o}{=}\PY{l+m+mi}{42}\PY{p}{)}\PY{o}{.}\PY{n}{fit}\PY{p}{(}\PY{n}{x\PYZus{}task2\PYZus{}processed}\PY{p}{)}
\PY{n}{pca\PYZus{}loadings} \PY{o}{=} \PY{n}{pd}\PY{o}{.}\PY{n}{DataFrame}\PY{p}{(}
    \PY{n}{pca\PYZus{}for\PYZus{}loadings}\PY{o}{.}\PY{n}{components\PYZus{}}\PY{o}{.}\PY{n}{T}\PY{p}{,}
    \PY{n}{index}\PY{o}{=}\PY{n}{task2\PYZus{}feature\PYZus{}names}\PY{p}{,}
    \PY{n}{columns}\PY{o}{=}\PY{p}{[}\PY{l+s+s2}{\PYZdq{}}\PY{l+s+s2}{PC1}\PY{l+s+s2}{\PYZdq{}}\PY{p}{,} \PY{l+s+s2}{\PYZdq{}}\PY{l+s+s2}{PC2}\PY{l+s+s2}{\PYZdq{}}\PY{p}{]}\PY{p}{,}
\PY{p}{)}

\PY{n}{pca\PYZus{}top\PYZus{}loadings} \PY{o}{=} \PY{p}{(}
    \PY{n}{pca\PYZus{}loadings}\PY{o}{.}\PY{n}{abs}\PY{p}{(}\PY{p}{)}
    \PY{o}{.}\PY{n}{stack}\PY{p}{(}\PY{p}{)}
    \PY{o}{.}\PY{n}{sort\PYZus{}values}\PY{p}{(}\PY{n}{ascending}\PY{o}{=}\PY{k+kc}{False}\PY{p}{)}
    \PY{o}{.}\PY{n}{head}\PY{p}{(}\PY{l+m+mi}{12}\PY{p}{)}
    \PY{o}{.}\PY{n}{reset\PYZus{}index}\PY{p}{(}\PY{p}{)}
\PY{p}{)}
\PY{n}{pca\PYZus{}top\PYZus{}loadings}\PY{o}{.}\PY{n}{columns} \PY{o}{=} \PY{p}{[}\PY{l+s+s2}{\PYZdq{}}\PY{l+s+s2}{feature\PYZus{}procesada}\PY{l+s+s2}{\PYZdq{}}\PY{p}{,} \PY{l+s+s2}{\PYZdq{}}\PY{l+s+s2}{componente}\PY{l+s+s2}{\PYZdq{}}\PY{p}{,} \PY{l+s+s2}{\PYZdq{}}\PY{l+s+s2}{carga\PYZus{}abs}\PY{l+s+s2}{\PYZdq{}}\PY{p}{]}
\PY{n}{pca\PYZus{}top\PYZus{}loadings}\PY{p}{[}\PY{l+s+s2}{\PYZdq{}}\PY{l+s+s2}{carga\PYZus{}firmada}\PY{l+s+s2}{\PYZdq{}}\PY{p}{]} \PY{o}{=} \PY{p}{[}
    \PY{n}{pca\PYZus{}loadings}\PY{o}{.}\PY{n}{loc}\PY{p}{[}\PY{n}{row}\PY{p}{[}\PY{l+s+s2}{\PYZdq{}}\PY{l+s+s2}{feature\PYZus{}procesada}\PY{l+s+s2}{\PYZdq{}}\PY{p}{]}\PY{p}{,} \PY{n}{row}\PY{p}{[}\PY{l+s+s2}{\PYZdq{}}\PY{l+s+s2}{componente}\PY{l+s+s2}{\PYZdq{}}\PY{p}{]}\PY{p}{]}
    \PY{k}{for} \PY{n}{\PYZus{}}\PY{p}{,} \PY{n}{row} \PY{o+ow}{in} \PY{n}{pca\PYZus{}top\PYZus{}loadings}\PY{o}{.}\PY{n}{iterrows}\PY{p}{(}\PY{p}{)}
\PY{p}{]}
\PY{n}{show\PYZus{}table}\PY{p}{(}\PY{l+s+s2}{\PYZdq{}}\PY{l+s+s2}{Mayores cargas absolutas en PC1 y PC2}\PY{l+s+s2}{\PYZdq{}}\PY{p}{,} \PY{n}{pca\PYZus{}top\PYZus{}loadings}\PY{o}{.}\PY{n}{round}\PY{p}{(}\PY{l+m+mi}{4}\PY{p}{)}\PY{p}{,} \PY{n}{n}\PY{o}{=}\PY{k+kc}{None}\PY{p}{)}

\PY{n}{show\PYZus{}image}\PY{p}{(}\PY{n}{OUTPUTS\PYZus{}DIR} \PY{o}{/} \PY{l+s+s2}{\PYZdq{}}\PY{l+s+s2}{task2}\PY{l+s+s2}{\PYZdq{}} \PY{o}{/} \PY{l+s+s2}{\PYZdq{}}\PY{l+s+s2}{pca\PYZus{}varianza\PYZus{}acumulada.png}\PY{l+s+s2}{\PYZdq{}}\PY{p}{,} \PY{n}{width}\PY{o}{=}\PY{l+m+mi}{800}\PY{p}{)}
\PY{n}{show\PYZus{}image}\PY{p}{(}\PY{n}{OUTPUTS\PYZus{}DIR} \PY{o}{/} \PY{l+s+s2}{\PYZdq{}}\PY{l+s+s2}{task2}\PY{l+s+s2}{\PYZdq{}} \PY{o}{/} \PY{l+s+s2}{\PYZdq{}}\PY{l+s+s2}{pca\PYZus{}condition.png}\PY{l+s+s2}{\PYZdq{}}\PY{p}{,} \PY{n}{width}\PY{o}{=}\PY{l+m+mi}{800}\PY{p}{)}
\end{Verbatim}
\end{tcolorbox}

    \textbf{Varianza explicada por PCA}

    
    
    \begin{Verbatim}[commandchars=\\\{\}]
   componente  varianza\_explicada  varianza\_acumulada
0         PC1        2.559249e-01            0.255925
1         PC2        1.434004e-01            0.399325
2         PC3        1.073197e-01            0.506645
3         PC4        9.728574e-02            0.603931
4         PC5        6.185284e-02            0.665784
5         PC6        5.511393e-02            0.720897
6         PC7        5.436017e-02            0.775258
7         PC8        3.722159e-02            0.812479
8         PC9        3.354499e-02            0.846024
9        PC10        2.982372e-02            0.875848
10       PC11        2.272561e-02            0.898574
11       PC12        1.838388e-02            0.916957
12       PC13        1.740161e-02            0.934359
13       PC14        1.555933e-02            0.949918
14       PC15        1.251999e-02            0.962438
15       PC16        1.053529e-02            0.972974
16       PC17        9.140859e-03            0.982115
17       PC18        8.070325e-03            0.990185
18       PC19        7.639850e-03            0.997825
19       PC20        2.175310e-03            1.000000
20       PC21        6.028370e-18            1.000000
21       PC22        4.496252e-18            1.000000
22       PC23        0.000000e+00            1.000000
23       PC24        0.000000e+00            1.000000
24       PC25        0.000000e+00            1.000000
    \end{Verbatim}

    
    \textbf{Lectura de reduccion de dimensionalidad}

    
    
    \begin{Verbatim}[commandchars=\\\{\}]
                              lectura   valor                                     interpretacion
0        Varianza acumulada PC1 + PC2  39.93\%  La visualizacion 2D resume una parte util, per{\ldots}
1  Componentes para >=80\% de varianza       8  La informacion relevante esta distribuida en v{\ldots}
2  Componentes para >=90\% de varianza      12  Reducir a solo dos componentes implica perder {\ldots}
    \end{Verbatim}

    
    \textbf{Mayores cargas absolutas en PC1 y PC2}

    
    
    \begin{Verbatim}[commandchars=\\\{\}]
   feature\_procesada componente  carga\_abs  carga\_firmada
0     num\_\_feature\_6        PC2     0.6294         0.6294
1    num\_\_feature\_10        PC1     0.5004         0.5004
2    num\_\_feature\_11        PC2     0.4864         0.4864
3    num\_\_feature\_13        PC1     0.4668        -0.4668
4     num\_\_feature\_5        PC1     0.4578        -0.4578
5    num\_\_feature\_10        PC2     0.3252         0.3252
6     num\_\_feature\_5        PC2     0.2982         0.2982
7    num\_\_feature\_11        PC1     0.2841        -0.2841
8    num\_\_feature\_13        PC2     0.2471        -0.2471
9   cat\_\_feature\_3\_0        PC1     0.2004         0.2004
10  cat\_\_feature\_3\_1        PC1     0.1669        -0.1669
11    num\_\_feature\_6        PC1     0.1612        -0.1612
    \end{Verbatim}

    
    \begin{center}
    \adjustimage{max size={0.9\linewidth}{0.9\paperheight}}{REPORT_files/REPORT_32_6.png}
    \end{center}
    { \hspace*{\fill} \\}
    
    \begin{center}
    \adjustimage{max size={0.9\linewidth}{0.9\paperheight}}{REPORT_files/REPORT_32_7.png}
    \end{center}
    { \hspace*{\fill} \\}
    
    PC1 y PC2 explican aproximadamente 39.93\% de la varianza acumulada.
Esto significa que el plano PCA es util para visualizar tendencias, pero
no resume todo el dataset. La necesidad de 8 componentes para superar
80\% y 12 para superar 90\% confirma que la informacion clinicamente
relevante esta repartida en varias dimensiones.

Las cargas principales ayudan a identificar que variables procesadas
empujan mas la visualizacion. Como algunas variables discretas fueron
codificadas con one-hot, las cargas deben leerse como contribuciones a
la estructura matematica del espacio, no como efectos clinicos directos.

    \subsection{2.2. K-Means y perfil de
clusters}\label{k-means-y-perfil-de-clusters}

K-Means se usa como tecnica principal de agrupamiento porque busca
particiones globales de pacientes con perfiles parecidos. Para elegir un
numero razonable de grupos, la lectura principal se apoya en el metodo
del codo: se observa como cae la inercia al aumentar \texttt{k} y se
busca un punto a partir del cual la mejora adicional sea menor.

Se conserva una evaluacion extendida \texttt{k=1..15} para ver mas alla
de los primeros valores. Ademas se comparan explicitamente \texttt{k=2},
\texttt{k=3} y \texttt{k=5}, porque representan tres niveles de
granularidad: dos perfiles generales, tres subgrupos y cinco perfiles
mas finos. El cruce con \texttt{condition} se calcula solo despues de
entrenar; la etiqueta no participa en K-Means.

    \begin{tcolorbox}[breakable, size=fbox, boxrule=1pt, pad at break*=1mm,colback=cellbackground, colframe=cellborder]
\prompt{In}{incolor}{19}{\boxspacing}
\begin{Verbatim}[commandchars=\\\{\}]
\PY{n}{kmeans\PYZus{}metrics} \PY{o}{=} \PY{n}{read\PYZus{}csv}\PY{p}{(}\PY{n}{OUTPUTS\PYZus{}DIR} \PY{o}{/} \PY{l+s+s2}{\PYZdq{}}\PY{l+s+s2}{task2}\PY{l+s+s2}{\PYZdq{}} \PY{o}{/} \PY{l+s+s2}{\PYZdq{}}\PY{l+s+s2}{kmeans\PYZus{}metricas.csv}\PY{l+s+s2}{\PYZdq{}}\PY{p}{)}
\PY{n}{kmeans\PYZus{}metrics\PYZus{}extended} \PY{o}{=} \PY{n}{read\PYZus{}csv}\PY{p}{(}\PY{n}{OUTPUTS\PYZus{}DIR} \PY{o}{/} \PY{l+s+s2}{\PYZdq{}}\PY{l+s+s2}{task2}\PY{l+s+s2}{\PYZdq{}} \PY{o}{/} \PY{l+s+s2}{\PYZdq{}}\PY{l+s+s2}{kmeans\PYZus{}metricas\PYZus{}extendido.csv}\PY{l+s+s2}{\PYZdq{}}\PY{p}{)}
\PY{n}{kmeans\PYZus{}k235} \PY{o}{=} \PY{n}{read\PYZus{}csv}\PY{p}{(}\PY{n}{OUTPUTS\PYZus{}DIR} \PY{o}{/} \PY{l+s+s2}{\PYZdq{}}\PY{l+s+s2}{task2}\PY{l+s+s2}{\PYZdq{}} \PY{o}{/} \PY{l+s+s2}{\PYZdq{}}\PY{l+s+s2}{kmeans\PYZus{}comparacion\PYZus{}k\PYZus{}2\PYZus{}3\PYZus{}5.csv}\PY{l+s+s2}{\PYZdq{}}\PY{p}{)}
\PY{n}{kmeans\PYZus{}profile} \PY{o}{=} \PY{n}{read\PYZus{}csv}\PY{p}{(}\PY{n}{OUTPUTS\PYZus{}DIR} \PY{o}{/} \PY{l+s+s2}{\PYZdq{}}\PY{l+s+s2}{task2}\PY{l+s+s2}{\PYZdq{}} \PY{o}{/} \PY{l+s+s2}{\PYZdq{}}\PY{l+s+s2}{kmeans\PYZus{}perfil\PYZus{}clusters.csv}\PY{l+s+s2}{\PYZdq{}}\PY{p}{)}
\PY{n}{kmeans\PYZus{}crosstab} \PY{o}{=} \PY{n}{read\PYZus{}csv}\PY{p}{(}\PY{n}{OUTPUTS\PYZus{}DIR} \PY{o}{/} \PY{l+s+s2}{\PYZdq{}}\PY{l+s+s2}{task2}\PY{l+s+s2}{\PYZdq{}} \PY{o}{/} \PY{l+s+s2}{\PYZdq{}}\PY{l+s+s2}{kmeans\PYZus{}condition\PYZus{}crosstab\PYZus{}pct.csv}\PY{l+s+s2}{\PYZdq{}}\PY{p}{)}
\PY{n}{kmeans\PYZus{}clusters} \PY{o}{=} \PY{n}{read\PYZus{}csv}\PY{p}{(}\PY{n}{OUTPUTS\PYZus{}DIR} \PY{o}{/} \PY{l+s+s2}{\PYZdq{}}\PY{l+s+s2}{task2}\PY{l+s+s2}{\PYZdq{}} \PY{o}{/} \PY{l+s+s2}{\PYZdq{}}\PY{l+s+s2}{kmeans\PYZus{}clusters.csv}\PY{l+s+s2}{\PYZdq{}}\PY{p}{)}

\PY{n}{kmeans\PYZus{}elbow} \PY{o}{=} \PY{n}{kmeans\PYZus{}metrics\PYZus{}extended}\PY{p}{[}\PY{p}{[}\PY{l+s+s2}{\PYZdq{}}\PY{l+s+s2}{k}\PY{l+s+s2}{\PYZdq{}}\PY{p}{,} \PY{l+s+s2}{\PYZdq{}}\PY{l+s+s2}{inertia}\PY{l+s+s2}{\PYZdq{}}\PY{p}{]}\PY{p}{]}\PY{o}{.}\PY{n}{copy}\PY{p}{(}\PY{p}{)}
\PY{n}{kmeans\PYZus{}elbow}\PY{p}{[}\PY{l+s+s2}{\PYZdq{}}\PY{l+s+s2}{reduccion\PYZus{}inercia\PYZus{}vs\PYZus{}k\PYZus{}anterior}\PY{l+s+s2}{\PYZdq{}}\PY{p}{]} \PY{o}{=} \PY{n}{kmeans\PYZus{}elbow}\PY{p}{[}\PY{l+s+s2}{\PYZdq{}}\PY{l+s+s2}{inertia}\PY{l+s+s2}{\PYZdq{}}\PY{p}{]}\PY{o}{.}\PY{n}{shift}\PY{p}{(}\PY{l+m+mi}{1}\PY{p}{)} \PY{o}{\PYZhy{}} \PY{n}{kmeans\PYZus{}elbow}\PY{p}{[}\PY{l+s+s2}{\PYZdq{}}\PY{l+s+s2}{inertia}\PY{l+s+s2}{\PYZdq{}}\PY{p}{]}
\PY{n}{kmeans\PYZus{}elbow}\PY{p}{[}\PY{l+s+s2}{\PYZdq{}}\PY{l+s+s2}{reduccion\PYZus{}pct\PYZus{}vs\PYZus{}k\PYZus{}anterior}\PY{l+s+s2}{\PYZdq{}}\PY{p}{]} \PY{o}{=} \PY{p}{(}
    \PY{n}{kmeans\PYZus{}elbow}\PY{p}{[}\PY{l+s+s2}{\PYZdq{}}\PY{l+s+s2}{reduccion\PYZus{}inercia\PYZus{}vs\PYZus{}k\PYZus{}anterior}\PY{l+s+s2}{\PYZdq{}}\PY{p}{]} \PY{o}{/} \PY{n}{kmeans\PYZus{}elbow}\PY{p}{[}\PY{l+s+s2}{\PYZdq{}}\PY{l+s+s2}{inertia}\PY{l+s+s2}{\PYZdq{}}\PY{p}{]}\PY{o}{.}\PY{n}{shift}\PY{p}{(}\PY{l+m+mi}{1}\PY{p}{)} \PY{o}{*} \PY{l+m+mi}{100}
\PY{p}{)}

\PY{n}{kmeans\PYZus{}k235\PYZus{}main} \PY{o}{=} \PY{n}{kmeans\PYZus{}k235}\PY{p}{[}
    \PY{p}{[}
        \PY{l+s+s2}{\PYZdq{}}\PY{l+s+s2}{k}\PY{l+s+s2}{\PYZdq{}}\PY{p}{,}
        \PY{l+s+s2}{\PYZdq{}}\PY{l+s+s2}{inertia}\PY{l+s+s2}{\PYZdq{}}\PY{p}{,}
        \PY{l+s+s2}{\PYZdq{}}\PY{l+s+s2}{condition\PYZus{}1\PYZus{}pct\PYZus{}min\PYZus{}cluster}\PY{l+s+s2}{\PYZdq{}}\PY{p}{,}
        \PY{l+s+s2}{\PYZdq{}}\PY{l+s+s2}{condition\PYZus{}1\PYZus{}pct\PYZus{}max\PYZus{}cluster}\PY{l+s+s2}{\PYZdq{}}\PY{p}{,}
        \PY{l+s+s2}{\PYZdq{}}\PY{l+s+s2}{condition\PYZus{}1\PYZus{}pct\PYZus{}range}\PY{l+s+s2}{\PYZdq{}}\PY{p}{,}
    \PY{p}{]}
\PY{p}{]}\PY{o}{.}\PY{n}{copy}\PY{p}{(}\PY{p}{)}

\PY{n}{show\PYZus{}table}\PY{p}{(}\PY{l+s+s2}{\PYZdq{}}\PY{l+s+s2}{Metodo del codo extendido K\PYZhy{}Means: k=1..15}\PY{l+s+s2}{\PYZdq{}}\PY{p}{,} \PY{n}{kmeans\PYZus{}elbow}\PY{o}{.}\PY{n}{round}\PY{p}{(}\PY{l+m+mi}{4}\PY{p}{)}\PY{p}{,} \PY{n}{n}\PY{o}{=}\PY{k+kc}{None}\PY{p}{)}
\PY{n}{show\PYZus{}table}\PY{p}{(}\PY{l+s+s2}{\PYZdq{}}\PY{l+s+s2}{Comparacion interpretativa K\PYZhy{}Means: k=2, k=3 y k=5}\PY{l+s+s2}{\PYZdq{}}\PY{p}{,} \PY{n}{kmeans\PYZus{}k235\PYZus{}main}\PY{o}{.}\PY{n}{round}\PY{p}{(}\PY{l+m+mi}{4}\PY{p}{)}\PY{p}{,} \PY{n}{n}\PY{o}{=}\PY{k+kc}{None}\PY{p}{)}
\PY{n}{show\PYZus{}table}\PY{p}{(}\PY{l+s+s2}{\PYZdq{}}\PY{l+s+s2}{Perfil medio de clusters K\PYZhy{}Means seleccionado}\PY{l+s+s2}{\PYZdq{}}\PY{p}{,} \PY{n}{kmeans\PYZus{}profile}\PY{p}{,} \PY{n}{n}\PY{o}{=}\PY{k+kc}{None}\PY{p}{)}
\PY{n}{show\PYZus{}table}\PY{p}{(}\PY{l+s+s2}{\PYZdq{}}\PY{l+s+s2}{Cruce porcentual K\PYZhy{}Means seleccionado vs condition}\PY{l+s+s2}{\PYZdq{}}\PY{p}{,} \PY{n}{kmeans\PYZus{}crosstab}\PY{p}{,} \PY{n}{n}\PY{o}{=}\PY{k+kc}{None}\PY{p}{)}

\PY{n}{kmeans\PYZus{}reading} \PY{o}{=} \PY{n}{pd}\PY{o}{.}\PY{n}{DataFrame}\PY{p}{(}\PY{p}{[}
    \PY{p}{\PYZob{}}
        \PY{l+s+s2}{\PYZdq{}}\PY{l+s+s2}{resultado}\PY{l+s+s2}{\PYZdq{}}\PY{p}{:} \PY{l+s+s2}{\PYZdq{}}\PY{l+s+s2}{Codo extendido}\PY{l+s+s2}{\PYZdq{}}\PY{p}{,}
        \PY{l+s+s2}{\PYZdq{}}\PY{l+s+s2}{valor}\PY{l+s+s2}{\PYZdq{}}\PY{p}{:} \PY{l+s+s2}{\PYZdq{}}\PY{l+s+s2}{k=2 como particion principal}\PY{l+s+s2}{\PYZdq{}}\PY{p}{,}
        \PY{l+s+s2}{\PYZdq{}}\PY{l+s+s2}{lectura}\PY{l+s+s2}{\PYZdq{}}\PY{p}{:} \PY{l+s+s2}{\PYZdq{}}\PY{l+s+s2}{Despues de separar dos perfiles generales, la inercia sigue bajando pero con ganancias mas graduales y menor interpretabilidad.}\PY{l+s+s2}{\PYZdq{}}\PY{p}{,}
    \PY{p}{\PYZcb{}}\PY{p}{,}
    \PY{p}{\PYZob{}}
        \PY{l+s+s2}{\PYZdq{}}\PY{l+s+s2}{resultado}\PY{l+s+s2}{\PYZdq{}}\PY{p}{:} \PY{l+s+s2}{\PYZdq{}}\PY{l+s+s2}{k=3}\PY{l+s+s2}{\PYZdq{}}\PY{p}{,}
        \PY{l+s+s2}{\PYZdq{}}\PY{l+s+s2}{valor}\PY{l+s+s2}{\PYZdq{}}\PY{p}{:} \PY{l+s+s2}{\PYZdq{}}\PY{l+s+s2}{subgrupos mas especificos}\PY{l+s+s2}{\PYZdq{}}\PY{p}{,}
        \PY{l+s+s2}{\PYZdq{}}\PY{l+s+s2}{lectura}\PY{l+s+s2}{\PYZdq{}}\PY{p}{:} \PY{l+s+s2}{\PYZdq{}}\PY{l+s+s2}{Aumenta la granularidad y puede aislar perfiles con mayor porcentaje de condition=1, pero divide mas la muestra.}\PY{l+s+s2}{\PYZdq{}}\PY{p}{,}
    \PY{p}{\PYZcb{}}\PY{p}{,}
    \PY{p}{\PYZob{}}
        \PY{l+s+s2}{\PYZdq{}}\PY{l+s+s2}{resultado}\PY{l+s+s2}{\PYZdq{}}\PY{p}{:} \PY{l+s+s2}{\PYZdq{}}\PY{l+s+s2}{k=5}\PY{l+s+s2}{\PYZdq{}}\PY{p}{,}
        \PY{l+s+s2}{\PYZdq{}}\PY{l+s+s2}{valor}\PY{l+s+s2}{\PYZdq{}}\PY{p}{:} \PY{l+s+s2}{\PYZdq{}}\PY{l+s+s2}{segmentacion fina}\PY{l+s+s2}{\PYZdq{}}\PY{p}{,}
        \PY{l+s+s2}{\PYZdq{}}\PY{l+s+s2}{lectura}\PY{l+s+s2}{\PYZdq{}}\PY{p}{:} \PY{l+s+s2}{\PYZdq{}}\PY{l+s+s2}{Produce grupos mas pequenos y extremos; es util como exploracion, pero menos estable para una conclusion general.}\PY{l+s+s2}{\PYZdq{}}\PY{p}{,}
    \PY{p}{\PYZcb{}}\PY{p}{,}
    \PY{p}{\PYZob{}}
        \PY{l+s+s2}{\PYZdq{}}\PY{l+s+s2}{resultado}\PY{l+s+s2}{\PYZdq{}}\PY{p}{:} \PY{l+s+s2}{\PYZdq{}}\PY{l+s+s2}{Cluster 0: condition=1}\PY{l+s+s2}{\PYZdq{}}\PY{p}{,}
        \PY{l+s+s2}{\PYZdq{}}\PY{l+s+s2}{valor}\PY{l+s+s2}{\PYZdq{}}\PY{p}{:} \PY{l+s+sa}{f}\PY{l+s+s2}{\PYZdq{}}\PY{l+s+si}{\PYZob{}}\PY{n+nb}{float}\PY{p}{(}\PY{n}{kmeans\PYZus{}crosstab}\PY{o}{.}\PY{n}{loc}\PY{p}{[}\PY{n}{kmeans\PYZus{}crosstab}\PY{p}{[}\PY{l+s+s1}{\PYZsq{}}\PY{l+s+s1}{cluster\PYZus{}kmeans}\PY{l+s+s1}{\PYZsq{}}\PY{p}{]}\PY{+w}{ }\PY{o}{==}\PY{+w}{ }\PY{l+m+mi}{0}\PY{p}{,}\PY{+w}{ }\PY{l+s+s1}{\PYZsq{}}\PY{l+s+s1}{1}\PY{l+s+s1}{\PYZsq{}}\PY{p}{]}\PY{o}{.}\PY{n}{iloc}\PY{p}{[}\PY{l+m+mi}{0}\PY{p}{]}\PY{p}{)}\PY{l+s+si}{:}\PY{l+s+s2}{.1f}\PY{l+s+si}{\PYZcb{}}\PY{l+s+s2}{\PYZpc{}}\PY{l+s+s2}{\PYZdq{}}\PY{p}{,}
        \PY{l+s+s2}{\PYZdq{}}\PY{l+s+s2}{lectura}\PY{l+s+s2}{\PYZdq{}}\PY{p}{:} \PY{l+s+s2}{\PYZdq{}}\PY{l+s+s2}{Grupo con menor proporcion de pacientes positivos.}\PY{l+s+s2}{\PYZdq{}}\PY{p}{,}
    \PY{p}{\PYZcb{}}\PY{p}{,}
    \PY{p}{\PYZob{}}
        \PY{l+s+s2}{\PYZdq{}}\PY{l+s+s2}{resultado}\PY{l+s+s2}{\PYZdq{}}\PY{p}{:} \PY{l+s+s2}{\PYZdq{}}\PY{l+s+s2}{Cluster 1: condition=1}\PY{l+s+s2}{\PYZdq{}}\PY{p}{,}
        \PY{l+s+s2}{\PYZdq{}}\PY{l+s+s2}{valor}\PY{l+s+s2}{\PYZdq{}}\PY{p}{:} \PY{l+s+sa}{f}\PY{l+s+s2}{\PYZdq{}}\PY{l+s+si}{\PYZob{}}\PY{n+nb}{float}\PY{p}{(}\PY{n}{kmeans\PYZus{}crosstab}\PY{o}{.}\PY{n}{loc}\PY{p}{[}\PY{n}{kmeans\PYZus{}crosstab}\PY{p}{[}\PY{l+s+s1}{\PYZsq{}}\PY{l+s+s1}{cluster\PYZus{}kmeans}\PY{l+s+s1}{\PYZsq{}}\PY{p}{]}\PY{+w}{ }\PY{o}{==}\PY{+w}{ }\PY{l+m+mi}{1}\PY{p}{,}\PY{+w}{ }\PY{l+s+s1}{\PYZsq{}}\PY{l+s+s1}{1}\PY{l+s+s1}{\PYZsq{}}\PY{p}{]}\PY{o}{.}\PY{n}{iloc}\PY{p}{[}\PY{l+m+mi}{0}\PY{p}{]}\PY{p}{)}\PY{l+s+si}{:}\PY{l+s+s2}{.1f}\PY{l+s+si}{\PYZcb{}}\PY{l+s+s2}{\PYZpc{}}\PY{l+s+s2}{\PYZdq{}}\PY{p}{,}
        \PY{l+s+s2}{\PYZdq{}}\PY{l+s+s2}{lectura}\PY{l+s+s2}{\PYZdq{}}\PY{p}{:} \PY{l+s+s2}{\PYZdq{}}\PY{l+s+s2}{Grupo con mayor concentracion de pacientes con condition=1.}\PY{l+s+s2}{\PYZdq{}}\PY{p}{,}
    \PY{p}{\PYZcb{}}\PY{p}{,}
\PY{p}{]}\PY{p}{)}
\PY{n}{show\PYZus{}table}\PY{p}{(}\PY{l+s+s2}{\PYZdq{}}\PY{l+s+s2}{Lectura principal de K\PYZhy{}Means}\PY{l+s+s2}{\PYZdq{}}\PY{p}{,} \PY{n}{kmeans\PYZus{}reading}\PY{p}{,} \PY{n}{n}\PY{o}{=}\PY{k+kc}{None}\PY{p}{)}

\PY{n}{categorical\PYZus{}cluster\PYZus{}rows} \PY{o}{=} \PY{p}{[}\PY{p}{]}
\PY{k}{for} \PY{n}{feature} \PY{o+ow}{in} \PY{n}{CATEGORICAL\PYZus{}FEATURES}\PY{p}{:}
    \PY{n}{distribution} \PY{o}{=} \PY{n}{pd}\PY{o}{.}\PY{n}{crosstab}\PY{p}{(}
        \PY{n}{kmeans\PYZus{}clusters}\PY{p}{[}\PY{l+s+s2}{\PYZdq{}}\PY{l+s+s2}{cluster\PYZus{}kmeans}\PY{l+s+s2}{\PYZdq{}}\PY{p}{]}\PY{p}{,}
        \PY{n}{kmeans\PYZus{}clusters}\PY{p}{[}\PY{n}{feature}\PY{p}{]}\PY{p}{,}
        \PY{n}{normalize}\PY{o}{=}\PY{l+s+s2}{\PYZdq{}}\PY{l+s+s2}{index}\PY{l+s+s2}{\PYZdq{}}\PY{p}{,}
    \PY{p}{)} \PY{o}{*} \PY{l+m+mi}{100}
    \PY{k}{for} \PY{n}{cluster\PYZus{}id}\PY{p}{,} \PY{n}{row} \PY{o+ow}{in} \PY{n}{distribution}\PY{o}{.}\PY{n}{round}\PY{p}{(}\PY{l+m+mi}{1}\PY{p}{)}\PY{o}{.}\PY{n}{iterrows}\PY{p}{(}\PY{p}{)}\PY{p}{:}
        \PY{n}{categorical\PYZus{}cluster\PYZus{}rows}\PY{o}{.}\PY{n}{append}\PY{p}{(}\PY{p}{\PYZob{}}
            \PY{l+s+s2}{\PYZdq{}}\PY{l+s+s2}{feature}\PY{l+s+s2}{\PYZdq{}}\PY{p}{:} \PY{n}{feature}\PY{p}{,}
            \PY{l+s+s2}{\PYZdq{}}\PY{l+s+s2}{cluster\PYZus{}kmeans}\PY{l+s+s2}{\PYZdq{}}\PY{p}{:} \PY{n}{cluster\PYZus{}id}\PY{p}{,}
            \PY{l+s+s2}{\PYZdq{}}\PY{l+s+s2}{distribucion\PYZus{}porcentual\PYZus{}niveles}\PY{l+s+s2}{\PYZdq{}}\PY{p}{:} \PY{l+s+s2}{\PYZdq{}}\PY{l+s+s2}{; }\PY{l+s+s2}{\PYZdq{}}\PY{o}{.}\PY{n}{join}\PY{p}{(}
                \PY{l+s+sa}{f}\PY{l+s+s2}{\PYZdq{}}\PY{l+s+si}{\PYZob{}}\PY{n}{level}\PY{l+s+si}{\PYZcb{}}\PY{l+s+s2}{: }\PY{l+s+si}{\PYZob{}}\PY{n}{pct}\PY{l+s+si}{:}\PY{l+s+s2}{.1f}\PY{l+s+si}{\PYZcb{}}\PY{l+s+s2}{\PYZpc{}}\PY{l+s+s2}{\PYZdq{}} \PY{k}{for} \PY{n}{level}\PY{p}{,} \PY{n}{pct} \PY{o+ow}{in} \PY{n}{row}\PY{o}{.}\PY{n}{items}\PY{p}{(}\PY{p}{)}
            \PY{p}{)}\PY{p}{,}
        \PY{p}{\PYZcb{}}\PY{p}{)}

\PY{n}{kmeans\PYZus{}categorical\PYZus{}distribution} \PY{o}{=} \PY{n}{pd}\PY{o}{.}\PY{n}{DataFrame}\PY{p}{(}\PY{n}{categorical\PYZus{}cluster\PYZus{}rows}\PY{p}{)}
\PY{n}{show\PYZus{}table}\PY{p}{(}\PY{l+s+s2}{\PYZdq{}}\PY{l+s+s2}{Distribucion de variables discretas por cluster K\PYZhy{}Means seleccionado}\PY{l+s+s2}{\PYZdq{}}\PY{p}{,} \PY{n}{kmeans\PYZus{}categorical\PYZus{}distribution}\PY{p}{,} \PY{n}{n}\PY{o}{=}\PY{k+kc}{None}\PY{p}{)}

\PY{n}{show\PYZus{}image}\PY{p}{(}\PY{n}{OUTPUTS\PYZus{}DIR} \PY{o}{/} \PY{l+s+s2}{\PYZdq{}}\PY{l+s+s2}{task2}\PY{l+s+s2}{\PYZdq{}} \PY{o}{/} \PY{l+s+s2}{\PYZdq{}}\PY{l+s+s2}{kmeans\PYZus{}elbow\PYZus{}hasta\PYZus{}k5.png}\PY{l+s+s2}{\PYZdq{}}\PY{p}{,} \PY{n}{width}\PY{o}{=}\PY{l+m+mi}{800}\PY{p}{)}
\PY{n}{show\PYZus{}image}\PY{p}{(}\PY{n}{OUTPUTS\PYZus{}DIR} \PY{o}{/} \PY{l+s+s2}{\PYZdq{}}\PY{l+s+s2}{task2}\PY{l+s+s2}{\PYZdq{}} \PY{o}{/} \PY{l+s+s2}{\PYZdq{}}\PY{l+s+s2}{kmeans\PYZus{}elbow\PYZus{}extendido.png}\PY{l+s+s2}{\PYZdq{}}\PY{p}{,} \PY{n}{width}\PY{o}{=}\PY{l+m+mi}{900}\PY{p}{)}
\PY{n}{show\PYZus{}image}\PY{p}{(}\PY{n}{OUTPUTS\PYZus{}DIR} \PY{o}{/} \PY{l+s+s2}{\PYZdq{}}\PY{l+s+s2}{task2}\PY{l+s+s2}{\PYZdq{}} \PY{o}{/} \PY{l+s+s2}{\PYZdq{}}\PY{l+s+s2}{kmeans\PYZus{}clusters\PYZus{}pca.png}\PY{l+s+s2}{\PYZdq{}}\PY{p}{,} \PY{n}{width}\PY{o}{=}\PY{l+m+mi}{800}\PY{p}{)}
\end{Verbatim}
\end{tcolorbox}

    \textbf{Metodo del codo extendido K-Means: k=1..15}

    
    
    \begin{Verbatim}[commandchars=\\\{\}]
     k    inertia  reduccion\_inercia\_vs\_k\_anterior  reduccion\_pct\_vs\_k\_anterior
0    1  2503.2424                              NaN                          NaN
1    2  2049.7439                         453.4985                      18.1164
2    3  1871.5175                         178.2265                       8.6951
3    4  1763.1925                         108.3250                       5.7881
4    5  1684.7485                          78.4440                       4.4490
5    6  1613.7378                          71.0107                       4.2149
6    7  1562.8942                          50.8436                       3.1507
7    8  1508.6267                          54.2675                       3.4722
8    9  1469.6607                          38.9660                       2.5829
9   10  1423.8606                          45.8001                       3.1164
10  11  1393.2514                          30.6092                       2.1497
11  12  1357.8943                          35.3571                       2.5377
12  13  1333.4339                          24.4605                       1.8014
13  14  1304.5586                          28.8753                       2.1655
14  15  1280.9971                          23.5615                       1.8061
    \end{Verbatim}

    
    \textbf{Comparacion interpretativa K-Means: k=2, k=3 y k=5}

    
    
    \begin{Verbatim}[commandchars=\\\{\}]
   k    inertia  condition\_1\_pct\_min\_cluster  condition\_1\_pct\_max\_cluster  condition\_1\_pct\_range
0  2  2049.7439                      22.4359                      72.3404                49.9045
1  3  1871.5175                      17.3469                      82.8283                65.4813
2  5  1684.7485                      12.9412                      83.5821                70.6409
    \end{Verbatim}

    
    \textbf{Perfil medio de clusters K-Means seleccionado}

    
    
    \begin{Verbatim}[commandchars=\\\{\}]
    Unnamed: 0           0           1  diferencia\_cluster\_1\_menos\_0
0    feature\_1    0.121795    0.170213                      0.048418
1    feature\_2    1.916667    2.425532                      0.508865
2    feature\_3    0.326923    0.907801                      0.580878
3    feature\_4    0.141026    0.531915                      0.390889
4    feature\_5   49.551282   60.063830                     10.512548
5    feature\_6  238.391026  257.262411                     18.871386
6    feature\_7    0.660256    0.695035                      0.034779
7    feature\_8    0.333333    1.056738                      0.723404
8    feature\_9    0.769231    1.248227                      0.478996
9   feature\_10  163.852564  133.829787                    -30.022777
10  feature\_11  127.102564  136.773050                      9.670486
11  feature\_12    0.532051    1.170213                      0.638161
12  feature\_13    0.448718    1.726950                      1.278232
13   condition    0.224359    0.723404                      0.499045
    \end{Verbatim}

    
    \textbf{Cruce porcentual K-Means seleccionado vs condition}

    
    
    \begin{Verbatim}[commandchars=\\\{\}]
   cluster\_kmeans          0          1
0               0  77.564103  22.435897
1               1  27.659574  72.340426
    \end{Verbatim}

    
    \textbf{Lectura principal de K-Means}

    
    
    \begin{Verbatim}[commandchars=\\\{\}]
                resultado                         valor                                            lectura
0          Codo extendido  k=2 como particion principal  Despues de separar dos perfiles generales, la {\ldots}
1                     k=3     subgrupos mas especificos  Aumenta la granularidad y puede aislar perfile{\ldots}
2                     k=5             segmentacion fina  Produce grupos mas pequenos y extremos; es uti{\ldots}
3  Cluster 0: condition=1                         22.4\%  Grupo con menor proporcion de pacientes positi{\ldots}
4  Cluster 1: condition=1                         72.3\%  Grupo con mayor concentracion de pacientes con{\ldots}
    \end{Verbatim}

    
    \textbf{Distribucion de variables discretas por cluster K-Means
seleccionado}

    
    
    \begin{Verbatim}[commandchars=\\\{\}]
      feature  cluster\_kmeans         distribucion\_porcentual\_niveles
0   feature\_2               0   0: 6.4\%; 1: 26.9\%; 2: 35.3\%; 3: 31.4\%
1   feature\_2               1    0: 9.2\%; 1: 5.0\%; 2: 19.9\%; 3: 66.0\%
2   feature\_3               0             0: 72.4\%; 1: 22.4\%; 2: 5.1\%
3   feature\_3               1             0: 18.4\%; 1: 72.3\%; 2: 9.2\%
4   feature\_8               0    0: 76.9\%; 1: 14.7\%; 2: 6.4\%; 3: 1.9\%
5   feature\_8               1  0: 38.3\%; 1: 29.8\%; 2: 19.9\%; 3: 12.1\%
6   feature\_9               0             0: 61.5\%; 1: 0.0\%; 2: 38.5\%
7   feature\_9               1             0: 36.2\%; 1: 2.8\%; 2: 61.0\%
8  feature\_12               0             0: 71.8\%; 1: 3.2\%; 2: 25.0\%
9  feature\_12               1             0: 36.9\%; 1: 9.2\%; 2: 53.9\%
    \end{Verbatim}

    
    \begin{center}
    \adjustimage{max size={0.9\linewidth}{0.9\paperheight}}{REPORT_files/REPORT_35_12.png}
    \end{center}
    { \hspace*{\fill} \\}
    
    \begin{center}
    \adjustimage{max size={0.9\linewidth}{0.9\paperheight}}{REPORT_files/REPORT_35_13.png}
    \end{center}
    { \hspace*{\fill} \\}
    
    \begin{center}
    \adjustimage{max size={0.9\linewidth}{0.9\paperheight}}{REPORT_files/REPORT_35_14.png}
    \end{center}
    { \hspace*{\fill} \\}
    
    La lectura principal conserva \texttt{k=2} porque ofrece una particion
simple y estable: separa dos perfiles generales de pacientes y permite
una interpretacion directa con \texttt{condition}. El metodo del codo
extendido no muestra una rodilla posterior que justifique claramente
pasar a muchos grupos; la inercia baja al aumentar \texttt{k}, pero eso
siempre ocurre cuando se agregan centroides.

\texttt{k=3} y \texttt{k=5} son utiles como exploracion de subperfiles.
Sus grupos pueden tener porcentajes mas extremos de
\texttt{condition=1}, pero esa ventaja viene de fragmentar la muestra:
hay mas clusters, mas detalle local y menos claridad para resumir la
estructura general del dataset. Por eso se mantiene \texttt{k=2} como
resultado principal y se interpretan \texttt{k=3} y \texttt{k=5} como
lecturas secundarias.

El cruce posterior con \texttt{condition} es el resultado mas
informativo: el cluster 0 tiene alrededor de 22.4\% de pacientes
positivos, mientras que el cluster 1 sube a 72.3\%. Como
\texttt{condition} no se uso al entrenar, esta diferencia sugiere que el
espacio de features contiene una estructura latente parcialmente
relacionada con la condicion medica.

    \subsection{2.3. Clustering jerarquico y DBSCAN como
contraste}\label{clustering-jerarquico-y-dbscan-como-contraste}

El clustering jerarquico se incluye como contraste de K-Means y se
interpreta mediante analisis de disparidad en el dendrograma. La idea es
observar las alturas a las que se fusionan los grupos: un salto grande
indica que se estan uniendo estructuras que estaban relativamente
separadas. Si los saltos fueran uniformes, no habria un argumento fuerte
para elegir un numero especifico de clusters.

DBSCAN responde otra pregunta: si hay grupos densos separados por zonas
de baja densidad y si existen pacientes que se comportan como ruido. En
este dataset funciona peor, pero ese resultado tambien es valioso porque
documenta una alternativa que no encaja bien con la estructura
observada.

    \begin{tcolorbox}[breakable, size=fbox, boxrule=1pt, pad at break*=1mm,colback=cellbackground, colframe=cellborder]
\prompt{In}{incolor}{20}{\boxspacing}
\begin{Verbatim}[commandchars=\\\{\}]
\PY{n}{agglomerative\PYZus{}crosstab\PYZus{}path} \PY{o}{=} \PY{n}{OUTPUTS\PYZus{}DIR} \PY{o}{/} \PY{l+s+s2}{\PYZdq{}}\PY{l+s+s2}{task2}\PY{l+s+s2}{\PYZdq{}} \PY{o}{/} \PY{l+s+s2}{\PYZdq{}}\PY{l+s+s2}{agglomerative\PYZus{}condition\PYZus{}crosstab\PYZus{}pct.csv}\PY{l+s+s2}{\PYZdq{}}
\PY{n}{agglomerative\PYZus{}crosstab} \PY{o}{=} \PY{n}{read\PYZus{}csv}\PY{p}{(}\PY{n}{agglomerative\PYZus{}crosstab\PYZus{}path}\PY{p}{)} \PY{k}{if} \PY{n}{agglomerative\PYZus{}crosstab\PYZus{}path}\PY{o}{.}\PY{n}{exists}\PY{p}{(}\PY{p}{)} \PY{k}{else} \PY{n}{pd}\PY{o}{.}\PY{n}{DataFrame}\PY{p}{(}\PY{p}{)}
\PY{n}{agglomerative\PYZus{}clusters\PYZus{}path} \PY{o}{=} \PY{n}{OUTPUTS\PYZus{}DIR} \PY{o}{/} \PY{l+s+s2}{\PYZdq{}}\PY{l+s+s2}{task2}\PY{l+s+s2}{\PYZdq{}} \PY{o}{/} \PY{l+s+s2}{\PYZdq{}}\PY{l+s+s2}{agglomerative\PYZus{}clusters.csv}\PY{l+s+s2}{\PYZdq{}}
\PY{n}{agglomerative\PYZus{}clusters} \PY{o}{=} \PY{n}{read\PYZus{}csv}\PY{p}{(}\PY{n}{agglomerative\PYZus{}clusters\PYZus{}path}\PY{p}{)} \PY{k}{if} \PY{n}{agglomerative\PYZus{}clusters\PYZus{}path}\PY{o}{.}\PY{n}{exists}\PY{p}{(}\PY{p}{)} \PY{k}{else} \PY{n}{pd}\PY{o}{.}\PY{n}{DataFrame}\PY{p}{(}\PY{p}{)}
\PY{n}{agglomerative\PYZus{}disparity} \PY{o}{=} \PY{n}{read\PYZus{}csv}\PY{p}{(}\PY{n}{OUTPUTS\PYZus{}DIR} \PY{o}{/} \PY{l+s+s2}{\PYZdq{}}\PY{l+s+s2}{task2}\PY{l+s+s2}{\PYZdq{}} \PY{o}{/} \PY{l+s+s2}{\PYZdq{}}\PY{l+s+s2}{agglomerative\PYZus{}disparidad.csv}\PY{l+s+s2}{\PYZdq{}}\PY{p}{)}
\PY{n}{dbscan\PYZus{}metrics} \PY{o}{=} \PY{n}{read\PYZus{}csv}\PY{p}{(}\PY{n}{OUTPUTS\PYZus{}DIR} \PY{o}{/} \PY{l+s+s2}{\PYZdq{}}\PY{l+s+s2}{task2}\PY{l+s+s2}{\PYZdq{}} \PY{o}{/} \PY{l+s+s2}{\PYZdq{}}\PY{l+s+s2}{dbscan\PYZus{}metricas.csv}\PY{l+s+s2}{\PYZdq{}}\PY{p}{)}

\PY{k}{if} \PY{l+s+s2}{\PYZdq{}}\PY{l+s+s2}{cluster\PYZus{}agglomerative}\PY{l+s+s2}{\PYZdq{}} \PY{o+ow}{in} \PY{n}{agglomerative\PYZus{}clusters}\PY{o}{.}\PY{n}{columns} \PY{o+ow}{and} \PY{o+ow}{not} \PY{n+nb}{set}\PY{p}{(}\PY{n}{features}\PY{p}{)}\PY{o}{.}\PY{n}{issubset}\PY{p}{(}\PY{n}{agglomerative\PYZus{}clusters}\PY{o}{.}\PY{n}{columns}\PY{p}{)}\PY{p}{:}
    \PY{n}{agglomerative\PYZus{}clusters\PYZus{}full} \PY{o}{=} \PY{n}{df\PYZus{}clean}\PY{o}{.}\PY{n}{copy}\PY{p}{(}\PY{p}{)}
    \PY{n}{agglomerative\PYZus{}clusters\PYZus{}full}\PY{p}{[}\PY{l+s+s2}{\PYZdq{}}\PY{l+s+s2}{cluster\PYZus{}agglomerative}\PY{l+s+s2}{\PYZdq{}}\PY{p}{]} \PY{o}{=} \PY{n}{agglomerative\PYZus{}clusters}\PY{p}{[}\PY{l+s+s2}{\PYZdq{}}\PY{l+s+s2}{cluster\PYZus{}agglomerative}\PY{l+s+s2}{\PYZdq{}}\PY{p}{]}\PY{o}{.}\PY{n}{values}
\PY{k}{else}\PY{p}{:}
    \PY{n}{agglomerative\PYZus{}clusters\PYZus{}full} \PY{o}{=} \PY{n}{agglomerative\PYZus{}clusters}\PY{o}{.}\PY{n}{copy}\PY{p}{(}\PY{p}{)}

\PY{n}{show\PYZus{}table}\PY{p}{(}\PY{l+s+s2}{\PYZdq{}}\PY{l+s+s2}{Mayores saltos de fusion en clustering jerarquico}\PY{l+s+s2}{\PYZdq{}}\PY{p}{,} \PY{n}{agglomerative\PYZus{}disparity}\PY{o}{.}\PY{n}{head}\PY{p}{(}\PY{l+m+mi}{10}\PY{p}{)}\PY{o}{.}\PY{n}{round}\PY{p}{(}\PY{l+m+mi}{4}\PY{p}{)}\PY{p}{,} \PY{n}{n}\PY{o}{=}\PY{k+kc}{None}\PY{p}{)}
\PY{k}{if} \PY{o+ow}{not} \PY{n}{agglomerative\PYZus{}crosstab}\PY{o}{.}\PY{n}{empty}\PY{p}{:}
    \PY{n}{show\PYZus{}table}\PY{p}{(}\PY{l+s+s2}{\PYZdq{}}\PY{l+s+s2}{Cruce porcentual clustering jerarquico vs condition}\PY{l+s+s2}{\PYZdq{}}\PY{p}{,} \PY{n}{agglomerative\PYZus{}crosstab}\PY{p}{,} \PY{n}{n}\PY{o}{=}\PY{k+kc}{None}\PY{p}{)}

\PY{n}{dbscan\PYZus{}main} \PY{o}{=} \PY{n}{dbscan\PYZus{}metrics}\PY{p}{[}\PY{p}{[}\PY{l+s+s2}{\PYZdq{}}\PY{l+s+s2}{eps}\PY{l+s+s2}{\PYZdq{}}\PY{p}{,} \PY{l+s+s2}{\PYZdq{}}\PY{l+s+s2}{min\PYZus{}samples}\PY{l+s+s2}{\PYZdq{}}\PY{p}{,} \PY{l+s+s2}{\PYZdq{}}\PY{l+s+s2}{n\PYZus{}clusters}\PY{l+s+s2}{\PYZdq{}}\PY{p}{,} \PY{l+s+s2}{\PYZdq{}}\PY{l+s+s2}{n\PYZus{}noise}\PY{l+s+s2}{\PYZdq{}}\PY{p}{]}\PY{p}{]}\PY{o}{.}\PY{n}{copy}\PY{p}{(}\PY{p}{)}
\PY{n}{show\PYZus{}table}\PY{p}{(}\PY{l+s+s2}{\PYZdq{}}\PY{l+s+s2}{DBSCAN: numero de clusters y puntos de ruido}\PY{l+s+s2}{\PYZdq{}}\PY{p}{,} \PY{n}{dbscan\PYZus{}main}\PY{p}{,} \PY{n}{n}\PY{o}{=}\PY{k+kc}{None}\PY{p}{)}

\PY{n}{dbscan\PYZus{}two\PYZus{}cluster} \PY{o}{=} \PY{n}{dbscan\PYZus{}metrics}\PY{o}{.}\PY{n}{loc}\PY{p}{[}\PY{n}{dbscan\PYZus{}metrics}\PY{p}{[}\PY{l+s+s2}{\PYZdq{}}\PY{l+s+s2}{n\PYZus{}clusters}\PY{l+s+s2}{\PYZdq{}}\PY{p}{]} \PY{o}{==} \PY{l+m+mi}{2}\PY{p}{]}\PY{o}{.}\PY{n}{sort\PYZus{}values}\PY{p}{(}\PY{l+s+s2}{\PYZdq{}}\PY{l+s+s2}{n\PYZus{}noise}\PY{l+s+s2}{\PYZdq{}}\PY{p}{)}\PY{o}{.}\PY{n}{head}\PY{p}{(}\PY{l+m+mi}{1}\PY{p}{)}
\PY{n}{dbscan\PYZus{}reading} \PY{o}{=} \PY{n}{pd}\PY{o}{.}\PY{n}{DataFrame}\PY{p}{(}\PY{p}{[}
    \PY{p}{\PYZob{}}
        \PY{l+s+s2}{\PYZdq{}}\PY{l+s+s2}{tecnica}\PY{l+s+s2}{\PYZdq{}}\PY{p}{:} \PY{l+s+s2}{\PYZdq{}}\PY{l+s+s2}{Clustering jerarquico}\PY{l+s+s2}{\PYZdq{}}\PY{p}{,}
        \PY{l+s+s2}{\PYZdq{}}\PY{l+s+s2}{evidencia}\PY{l+s+s2}{\PYZdq{}}\PY{p}{:} \PY{l+s+s2}{\PYZdq{}}\PY{l+s+s2}{El mayor salto de altura ocurre antes de fusionar los dos ultimos grupos.}\PY{l+s+s2}{\PYZdq{}}\PY{p}{,}
        \PY{l+s+s2}{\PYZdq{}}\PY{l+s+s2}{interpretacion}\PY{l+s+s2}{\PYZdq{}}\PY{p}{:} \PY{l+s+s2}{\PYZdq{}}\PY{l+s+s2}{La disparidad del dendrograma apoya una lectura de dos macrogrupos.}\PY{l+s+s2}{\PYZdq{}}\PY{p}{,}
    \PY{p}{\PYZcb{}}\PY{p}{,}
    \PY{p}{\PYZob{}}
        \PY{l+s+s2}{\PYZdq{}}\PY{l+s+s2}{tecnica}\PY{l+s+s2}{\PYZdq{}}\PY{p}{:} \PY{l+s+s2}{\PYZdq{}}\PY{l+s+s2}{DBSCAN}\PY{l+s+s2}{\PYZdq{}}\PY{p}{,}
        \PY{l+s+s2}{\PYZdq{}}\PY{l+s+s2}{evidencia}\PY{l+s+s2}{\PYZdq{}}\PY{p}{:} \PY{p}{(}
            \PY{l+s+s2}{\PYZdq{}}\PY{l+s+s2}{La configuracion con dos clusters y menos ruido conserva }\PY{l+s+s2}{\PYZdq{}}
            \PY{l+s+sa}{f}\PY{l+s+s2}{\PYZdq{}}\PY{l+s+si}{\PYZob{}}\PY{n+nb}{int}\PY{p}{(}\PY{n}{dbscan\PYZus{}two\PYZus{}cluster}\PY{o}{.}\PY{n}{iloc}\PY{p}{[}\PY{l+m+mi}{0}\PY{p}{]}\PY{p}{[}\PY{l+s+s1}{\PYZsq{}}\PY{l+s+s1}{n\PYZus{}noise}\PY{l+s+s1}{\PYZsq{}}\PY{p}{]}\PY{p}{)}\PY{l+s+si}{\PYZcb{}}\PY{l+s+s2}{ pacientes como ruido.}\PY{l+s+s2}{\PYZdq{}}
            \PY{k}{if} \PY{o+ow}{not} \PY{n}{dbscan\PYZus{}two\PYZus{}cluster}\PY{o}{.}\PY{n}{empty}
            \PY{k}{else} \PY{l+s+s2}{\PYZdq{}}\PY{l+s+s2}{No aparece una configuracion limpia con dos clusters.}\PY{l+s+s2}{\PYZdq{}}
        \PY{p}{)}\PY{p}{,}
        \PY{l+s+s2}{\PYZdq{}}\PY{l+s+s2}{interpretacion}\PY{l+s+s2}{\PYZdq{}}\PY{p}{:} \PY{l+s+s2}{\PYZdq{}}\PY{l+s+s2}{La estructura no parece formada por regiones densas claramente separadas.}\PY{l+s+s2}{\PYZdq{}}\PY{p}{,}
    \PY{p}{\PYZcb{}}\PY{p}{,}
\PY{p}{]}\PY{p}{)}
\PY{n}{show\PYZus{}table}\PY{p}{(}\PY{l+s+s2}{\PYZdq{}}\PY{l+s+s2}{Lectura comparativa de metodos alternativos}\PY{l+s+s2}{\PYZdq{}}\PY{p}{,} \PY{n}{dbscan\PYZus{}reading}\PY{p}{,} \PY{n}{n}\PY{o}{=}\PY{k+kc}{None}\PY{p}{)}

\PY{k}{if} \PY{p}{\PYZob{}}\PY{l+s+s2}{\PYZdq{}}\PY{l+s+s2}{cluster\PYZus{}agglomerative}\PY{l+s+s2}{\PYZdq{}}\PY{p}{,} \PY{o}{*}\PY{n}{CATEGORICAL\PYZus{}FEATURES}\PY{p}{\PYZcb{}}\PY{o}{.}\PY{n}{issubset}\PY{p}{(}\PY{n}{agglomerative\PYZus{}clusters\PYZus{}full}\PY{o}{.}\PY{n}{columns}\PY{p}{)}\PY{p}{:}
    \PY{n}{agg\PYZus{}discrete\PYZus{}rows} \PY{o}{=} \PY{p}{[}\PY{p}{]}
    \PY{k}{for} \PY{n}{feature} \PY{o+ow}{in} \PY{n}{CATEGORICAL\PYZus{}FEATURES}\PY{p}{:}
        \PY{n}{distribution} \PY{o}{=} \PY{n}{pd}\PY{o}{.}\PY{n}{crosstab}\PY{p}{(}
            \PY{n}{agglomerative\PYZus{}clusters\PYZus{}full}\PY{p}{[}\PY{l+s+s2}{\PYZdq{}}\PY{l+s+s2}{cluster\PYZus{}agglomerative}\PY{l+s+s2}{\PYZdq{}}\PY{p}{]}\PY{p}{,}
            \PY{n}{agglomerative\PYZus{}clusters\PYZus{}full}\PY{p}{[}\PY{n}{feature}\PY{p}{]}\PY{p}{,}
            \PY{n}{normalize}\PY{o}{=}\PY{l+s+s2}{\PYZdq{}}\PY{l+s+s2}{index}\PY{l+s+s2}{\PYZdq{}}\PY{p}{,}
        \PY{p}{)} \PY{o}{*} \PY{l+m+mi}{100}
        \PY{k}{for} \PY{n}{cluster\PYZus{}id}\PY{p}{,} \PY{n}{row} \PY{o+ow}{in} \PY{n}{distribution}\PY{o}{.}\PY{n}{round}\PY{p}{(}\PY{l+m+mi}{1}\PY{p}{)}\PY{o}{.}\PY{n}{iterrows}\PY{p}{(}\PY{p}{)}\PY{p}{:}
            \PY{n}{agg\PYZus{}discrete\PYZus{}rows}\PY{o}{.}\PY{n}{append}\PY{p}{(}\PY{p}{\PYZob{}}
                \PY{l+s+s2}{\PYZdq{}}\PY{l+s+s2}{feature}\PY{l+s+s2}{\PYZdq{}}\PY{p}{:} \PY{n}{feature}\PY{p}{,}
                \PY{l+s+s2}{\PYZdq{}}\PY{l+s+s2}{cluster\PYZus{}agglomerative}\PY{l+s+s2}{\PYZdq{}}\PY{p}{:} \PY{n}{cluster\PYZus{}id}\PY{p}{,}
                \PY{l+s+s2}{\PYZdq{}}\PY{l+s+s2}{distribucion\PYZus{}porcentual\PYZus{}niveles}\PY{l+s+s2}{\PYZdq{}}\PY{p}{:} \PY{l+s+s2}{\PYZdq{}}\PY{l+s+s2}{; }\PY{l+s+s2}{\PYZdq{}}\PY{o}{.}\PY{n}{join}\PY{p}{(}
                    \PY{l+s+sa}{f}\PY{l+s+s2}{\PYZdq{}}\PY{l+s+si}{\PYZob{}}\PY{n}{level}\PY{l+s+si}{\PYZcb{}}\PY{l+s+s2}{: }\PY{l+s+si}{\PYZob{}}\PY{n}{pct}\PY{l+s+si}{:}\PY{l+s+s2}{.1f}\PY{l+s+si}{\PYZcb{}}\PY{l+s+s2}{\PYZpc{}}\PY{l+s+s2}{\PYZdq{}} \PY{k}{for} \PY{n}{level}\PY{p}{,} \PY{n}{pct} \PY{o+ow}{in} \PY{n}{row}\PY{o}{.}\PY{n}{items}\PY{p}{(}\PY{p}{)}
                \PY{p}{)}\PY{p}{,}
            \PY{p}{\PYZcb{}}\PY{p}{)}
    \PY{n}{show\PYZus{}table}\PY{p}{(}\PY{l+s+s2}{\PYZdq{}}\PY{l+s+s2}{Distribucion de variables discretas por cluster jerarquico}\PY{l+s+s2}{\PYZdq{}}\PY{p}{,} \PY{n}{pd}\PY{o}{.}\PY{n}{DataFrame}\PY{p}{(}\PY{n}{agg\PYZus{}discrete\PYZus{}rows}\PY{p}{)}\PY{p}{,} \PY{n}{n}\PY{o}{=}\PY{k+kc}{None}\PY{p}{)}

\PY{n}{show\PYZus{}image}\PY{p}{(}\PY{n}{OUTPUTS\PYZus{}DIR} \PY{o}{/} \PY{l+s+s2}{\PYZdq{}}\PY{l+s+s2}{task2}\PY{l+s+s2}{\PYZdq{}} \PY{o}{/} \PY{l+s+s2}{\PYZdq{}}\PY{l+s+s2}{agglomerative\PYZus{}dendrogram.png}\PY{l+s+s2}{\PYZdq{}}\PY{p}{,} \PY{n}{width}\PY{o}{=}\PY{l+m+mi}{900}\PY{p}{)}
\PY{n}{show\PYZus{}image}\PY{p}{(}\PY{n}{OUTPUTS\PYZus{}DIR} \PY{o}{/} \PY{l+s+s2}{\PYZdq{}}\PY{l+s+s2}{task2}\PY{l+s+s2}{\PYZdq{}} \PY{o}{/} \PY{l+s+s2}{\PYZdq{}}\PY{l+s+s2}{agglomerative\PYZus{}clusters\PYZus{}pca.png}\PY{l+s+s2}{\PYZdq{}}\PY{p}{,} \PY{n}{width}\PY{o}{=}\PY{l+m+mi}{800}\PY{p}{)}
\PY{n}{show\PYZus{}image}\PY{p}{(}\PY{n}{OUTPUTS\PYZus{}DIR} \PY{o}{/} \PY{l+s+s2}{\PYZdq{}}\PY{l+s+s2}{task2}\PY{l+s+s2}{\PYZdq{}} \PY{o}{/} \PY{l+s+s2}{\PYZdq{}}\PY{l+s+s2}{dbscan\PYZus{}clusters\PYZus{}pca.png}\PY{l+s+s2}{\PYZdq{}}\PY{p}{,} \PY{n}{width}\PY{o}{=}\PY{l+m+mi}{800}\PY{p}{)}
\end{Verbatim}
\end{tcolorbox}

    \textbf{Mayores saltos de fusion en clustering jerarquico}

    
    
    \begin{Verbatim}[commandchars=\\\{\}]
   fusion  altura\_fusion  salto\_altura  clusters\_restantes
0     296        26.2606        7.8584                   1
1     295        18.4022        3.6427                   2
2     294        14.7595        2.2549                   3
3     290        10.7171        1.1000                   7
4     283         8.0906        0.9155                  14
5     293        12.5046        0.9149                   4
6     291        11.4499        0.7328                   6
7     286         8.8910        0.5848                  11
8     287         9.1890        0.2981                  10
9       3         0.7234        0.2963                 294
    \end{Verbatim}

    
    \textbf{Cruce porcentual clustering jerarquico vs condition}

    
    
    \begin{Verbatim}[commandchars=\\\{\}]
   cluster\_agglomerative          0          1
0                      0  68.303571  31.696429
1                      1   9.589041  90.410959
    \end{Verbatim}

    
    \textbf{DBSCAN: numero de clusters y puntos de ruido}

    
    
    \begin{Verbatim}[commandchars=\\\{\}]
    eps  min\_samples  n\_clusters  n\_noise
0   0.5            3           0      297
1   0.5            5           0      297
2   0.5           10           0      297
3   1.0            3           1      294
4   1.0            5           0      297
5   1.0           10           0      297
6   1.5            3           7      251
7   1.5            5           1      277
8   1.5           10           0      297
9   2.0            3           9      136
10  2.0            5           2      185
11  2.0           10           2      207
12  2.5            3           1       29
13  2.5            5           1       35
14  2.5           10           1       49
15  3.0            3           1        4
16  3.0            5           1        4
17  3.0           10           1        8
    \end{Verbatim}

    
    \textbf{Lectura comparativa de metodos alternativos}

    
    
    \begin{Verbatim}[commandchars=\\\{\}]
                 tecnica                                          evidencia                                     interpretacion
0  Clustering jerarquico  El mayor salto de altura ocurre antes de fusio{\ldots}  La disparidad del dendrograma apoya una lectur{\ldots}
1                 DBSCAN  La configuracion con dos clusters y menos ruid{\ldots}  La estructura no parece formada por regiones d{\ldots}
    \end{Verbatim}

    
    \textbf{Distribucion de variables discretas por cluster jerarquico}

    
    
    \begin{Verbatim}[commandchars=\\\{\}]
      feature  cluster\_agglomerative         distribucion\_porcentual\_niveles
0   feature\_2                      0   0: 8.9\%; 1: 21.0\%; 2: 35.7\%; 3: 34.4\%
1   feature\_2                      1     0: 4.1\%; 1: 2.7\%; 2: 4.1\%; 3: 89.0\%
2   feature\_3                      0             0: 59.8\%; 1: 36.6\%; 2: 3.6\%
3   feature\_3                      1             0: 6.8\%; 1: 75.3\%; 2: 17.8\%
4   feature\_8                      0    0: 68.8\%; 1: 18.8\%; 2: 8.5\%; 3: 4.0\%
5   feature\_8                      1  0: 27.4\%; 1: 31.5\%; 2: 26.0\%; 3: 15.1\%
6   feature\_9                      0             0: 53.6\%; 1: 0.4\%; 2: 46.0\%
7   feature\_9                      1             0: 37.0\%; 1: 4.1\%; 2: 58.9\%
8  feature\_12                      0             0: 68.8\%; 1: 4.0\%; 2: 27.2\%
9  feature\_12                      1            0: 13.7\%; 1: 12.3\%; 2: 74.0\%
    \end{Verbatim}

    
    \begin{center}
    \adjustimage{max size={0.9\linewidth}{0.9\paperheight}}{REPORT_files/REPORT_38_10.png}
    \end{center}
    { \hspace*{\fill} \\}
    
    \begin{center}
    \adjustimage{max size={0.9\linewidth}{0.9\paperheight}}{REPORT_files/REPORT_38_11.png}
    \end{center}
    { \hspace*{\fill} \\}
    
    \begin{center}
    \adjustimage{max size={0.9\linewidth}{0.9\paperheight}}{REPORT_files/REPORT_38_12.png}
    \end{center}
    { \hspace*{\fill} \\}
    
    El dendrograma muestra su mayor salto al final del proceso, justo antes
de fusionar los dos macrogrupos restantes. Esa disparidad respalda la
idea de una estructura principal de dos grupos y sirve como contraste
independiente frente a K-Means.

El cruce jerarquico con \texttt{condition} tambien muestra una
asociacion clara: uno de los grupos concentra muchos mas positivos que
el otro. DBSCAN, en cambio, produjo muchas configuraciones con un unico
cluster, cero clusters o una cantidad alta de pacientes marcados como
ruido. Esto sugiere que los pacientes no forman regiones densas bien
separadas; la estructura parece mas compatible con particiones globales
y macrogrupos que con islas densas aisladas.

    La experimentacion no supervisada muestra una estructura latente
parcialmente alineada con \texttt{condition}. La evidencia principal
aparece al combinar K-Means con \texttt{k=2}, el metodo del codo, el
analisis de disparidad jerarquico y el cruce posterior con la variable
objetivo.

Las caracteristicas que diferencian los grupos combinan variables
continuas y discretas. Esto indica que la condicion no esta representada
por una unica feature aislada, sino por un patron multivariable. Al
mismo tiempo, PCA y DBSCAN ponen limites claros a la interpretacion: la
informacion esta distribuida en varias dimensiones y no aparecen
regiones densas perfectamente separadas. Por tanto, estos resultados
sirven para comprender perfiles de pacientes y formular hipotesis, no
para diagnosticar por si solos.

    \section{3. Aprendizaje Supervisado}\label{aprendizaje-supervisado}

Esta sección compara modelos supervisados para predecir
\texttt{condition}. La lectura se hace desde un contexto médico: no
basta con maximizar accuracy, porque el error crítico es clasificar como
sano a un paciente que realmente tiene \texttt{condition=1}.

La estructura de esta parte es deliberadamente experimental. Primero se
fija la estrategia de validación y la métrica clínica, después se revisa
cada modelo por separado y finalmente se comparan todos con la misma
regla.

    \subsection{3.1. Estrategia experimental y métrica
médica}\label{estrategia-experimental-y-muxe9trica-muxe9dica}

Todos los modelos usan un 20\% de test final estratificado, reservado
hasta el cierre de cada experimento. La comparación interna usa
\texttt{RepeatedStratifiedKFold} con 5 folds y 30 repeticiones, es
decir, 150 validaciones por modelo. Esta repetición reduce la
dependencia de una única partición, algo importante en un dataset médico
de tamaño moderado.

El preprocesamiento se integra dentro de un \texttt{Pipeline}: las
variables continuas se estandarizan, las discretas se codifican con
one-hot y las binarias pasan directamente. Así, cada fold ajusta sus
transformaciones solo con datos de entrenamiento. Por eso en Task 3 se
usa \texttt{foot\_clean.csv} y no un dataset previamente estandarizado
con todo el conjunto, ya que eso podría introducir fuga de información.

La métrica principal es \texttt{recall\_condition\_1}, porque mide qué
proporción de pacientes con la condición real fueron detectados. El
coste médico se define como \texttt{(5·FN\ +\ 1·FP)\ /\ N}: el falso
negativo pesa más porque puede retrasar diagnóstico, seguimiento o
tratamiento; el falso positivo también importa, pero suele representar
alarma o pruebas adicionales, no dejar sin detectar un caso real.

En esta tarea se mantiene el umbral de decisión por defecto para
comparar modelos base de forma homogénea. En un uso clínico real habría
que optimizar ese umbral para reducir falsos negativos, incluso si
aumentan falsos positivos.

    \begin{tcolorbox}[breakable, size=fbox, boxrule=1pt, pad at break*=1mm,colback=cellbackground, colframe=cellborder]
\prompt{In}{incolor}{21}{\boxspacing}
\begin{Verbatim}[commandchars=\\\{\}]
\PY{n}{validation\PYZus{}summary} \PY{o}{=} \PY{n}{read\PYZus{}csv}\PY{p}{(}\PY{n}{OUTPUTS\PYZus{}DIR} \PY{o}{/} \PY{l+s+s2}{\PYZdq{}}\PY{l+s+s2}{task3}\PY{l+s+s2}{\PYZdq{}} \PY{o}{/} \PY{l+s+s2}{\PYZdq{}}\PY{l+s+s2}{decisiones\PYZus{}validacion.csv}\PY{l+s+s2}{\PYZdq{}}\PY{p}{)}
\PY{n}{show\PYZus{}table}\PY{p}{(}\PY{l+s+s2}{\PYZdq{}}\PY{l+s+s2}{Decisiones de validación y métrica médica}\PY{l+s+s2}{\PYZdq{}}\PY{p}{,} \PY{n}{validation\PYZus{}summary}\PY{p}{,} \PY{n}{n}\PY{o}{=}\PY{k+kc}{None}\PY{p}{)}
\end{Verbatim}
\end{tcolorbox}

    \textbf{Decisiones de validación y métrica médica}

    
    
    \begin{Verbatim}[commandchars=\\\{\}]
             decisión                                              valor                                      justificación
0          Test final                                  20\% estratificado  Reserva una evaluación final no usada durante {\ldots}
1  Validación interna  RepeatedStratifiedKFold, 5 folds x 30 repetici{\ldots}  Reduce la dependencia de una única partición e{\ldots}
2  Dataset de entrada                          foot\_clean.csv + Pipeline  Evita fuga de información al ajustar escalado {\ldots}
3   Métrica principal                                 recall\_condition\_1  Mide cuántos pacientes con condition=1 son det{\ldots}
4        Coste médico                                         FN=5, FP=1  Penaliza más dejar sin detectar un caso real q{\ldots}
5              Umbral                                 Umbral por defecto  Permite comparar modelos base de forma homogén{\ldots}
    \end{Verbatim}

    
    \subsection{3.2. Experimentos uno a uno}\label{experimentos-uno-a-uno}

Antes de comparar modelos, se registra por qué se prueba cada enfoque y
qué se puede extraer de sus resultados. Esto evita que la sección sea
solo una tabla: cada modelo responde una pregunta distinta sobre los
datos, desde un baseline lineal interpretable hasta modelos no lineales
y un ensamblado de árboles.

La gráfica resume por experimento tres señales comparables entre 0 y 1:
recall medio en validación cruzada, recall en test y ROC-AUC medio. En
el título de cada panel se añaden coste médico medio, falsos negativos y
falsos positivos en test, que son las señales clínicas más importantes.

    \begin{tcolorbox}[breakable, size=fbox, boxrule=1pt, pad at break*=1mm,colback=cellbackground, colframe=cellborder]
\prompt{In}{incolor}{22}{\boxspacing}
\begin{Verbatim}[commandchars=\\\{\}]
\PY{n}{experiment\PYZus{}catalog} \PY{o}{=} \PY{n}{read\PYZus{}csv}\PY{p}{(}\PY{n}{OUTPUTS\PYZus{}DIR} \PY{o}{/} \PY{l+s+s2}{\PYZdq{}}\PY{l+s+s2}{task3}\PY{l+s+s2}{\PYZdq{}} \PY{o}{/} \PY{l+s+s2}{\PYZdq{}}\PY{l+s+s2}{catalogo\PYZus{}experimentos.csv}\PY{l+s+s2}{\PYZdq{}}\PY{p}{)}
\PY{n}{catalog\PYZus{}columns} \PY{o}{=} \PY{p}{[}
    \PY{l+s+s2}{\PYZdq{}}\PY{l+s+s2}{modelo}\PY{l+s+s2}{\PYZdq{}}\PY{p}{,}
    \PY{l+s+s2}{\PYZdq{}}\PY{l+s+s2}{familia}\PY{l+s+s2}{\PYZdq{}}\PY{p}{,}
    \PY{l+s+s2}{\PYZdq{}}\PY{l+s+s2}{por\PYZus{}que}\PY{l+s+s2}{\PYZdq{}}\PY{p}{,}
    \PY{l+s+s2}{\PYZdq{}}\PY{l+s+s2}{fortaleza}\PY{l+s+s2}{\PYZdq{}}\PY{p}{,}
    \PY{l+s+s2}{\PYZdq{}}\PY{l+s+s2}{limitacion}\PY{l+s+s2}{\PYZdq{}}\PY{p}{,}
    \PY{l+s+s2}{\PYZdq{}}\PY{l+s+s2}{lectura\PYZus{}resultado}\PY{l+s+s2}{\PYZdq{}}\PY{p}{,}
    \PY{l+s+s2}{\PYZdq{}}\PY{l+s+s2}{conclusion\PYZus{}experimento}\PY{l+s+s2}{\PYZdq{}}\PY{p}{,}
\PY{p}{]}
\PY{n}{show\PYZus{}table}\PY{p}{(}\PY{l+s+s2}{\PYZdq{}}\PY{l+s+s2}{Registro de experimentos supervisados}\PY{l+s+s2}{\PYZdq{}}\PY{p}{,} \PY{n}{experiment\PYZus{}catalog}\PY{p}{[}\PY{n}{catalog\PYZus{}columns}\PY{p}{]}\PY{p}{,} \PY{n}{n}\PY{o}{=}\PY{k+kc}{None}\PY{p}{)}
\PY{n}{show\PYZus{}image}\PY{p}{(}\PY{n}{OUTPUTS\PYZus{}DIR} \PY{o}{/} \PY{l+s+s2}{\PYZdq{}}\PY{l+s+s2}{task3}\PY{l+s+s2}{\PYZdq{}} \PY{o}{/} \PY{l+s+s2}{\PYZdq{}}\PY{l+s+s2}{experimentos\PYZus{}individuales\PYZus{}metricas.png}\PY{l+s+s2}{\PYZdq{}}\PY{p}{,} \PY{n}{width}\PY{o}{=}\PY{l+m+mi}{950}\PY{p}{)}

\PY{n}{model\PYZus{}slugs} \PY{o}{=} \PY{p}{[}
    \PY{p}{(}\PY{l+s+s2}{\PYZdq{}}\PY{l+s+s2}{logistic\PYZus{}regression}\PY{l+s+s2}{\PYZdq{}}\PY{p}{,} \PY{l+s+s2}{\PYZdq{}}\PY{l+s+s2}{Regresión Logística}\PY{l+s+s2}{\PYZdq{}}\PY{p}{)}\PY{p}{,}
    \PY{p}{(}\PY{l+s+s2}{\PYZdq{}}\PY{l+s+s2}{naive\PYZus{}bayes}\PY{l+s+s2}{\PYZdq{}}\PY{p}{,} \PY{l+s+s2}{\PYZdq{}}\PY{l+s+s2}{Naive Bayes}\PY{l+s+s2}{\PYZdq{}}\PY{p}{)}\PY{p}{,}
    \PY{p}{(}\PY{l+s+s2}{\PYZdq{}}\PY{l+s+s2}{knn}\PY{l+s+s2}{\PYZdq{}}\PY{p}{,} \PY{l+s+s2}{\PYZdq{}}\PY{l+s+s2}{KNN}\PY{l+s+s2}{\PYZdq{}}\PY{p}{)}\PY{p}{,}
    \PY{p}{(}\PY{l+s+s2}{\PYZdq{}}\PY{l+s+s2}{decision\PYZus{}tree}\PY{l+s+s2}{\PYZdq{}}\PY{p}{,} \PY{l+s+s2}{\PYZdq{}}\PY{l+s+s2}{Árbol de Decisión CART}\PY{l+s+s2}{\PYZdq{}}\PY{p}{)}\PY{p}{,}
    \PY{p}{(}\PY{l+s+s2}{\PYZdq{}}\PY{l+s+s2}{xgboost}\PY{l+s+s2}{\PYZdq{}}\PY{p}{,} \PY{l+s+s2}{\PYZdq{}}\PY{l+s+s2}{XGBoost}\PY{l+s+s2}{\PYZdq{}}\PY{p}{)}\PY{p}{,}
    \PY{p}{(}\PY{l+s+s2}{\PYZdq{}}\PY{l+s+s2}{svm}\PY{l+s+s2}{\PYZdq{}}\PY{p}{,} \PY{l+s+s2}{\PYZdq{}}\PY{l+s+s2}{SVM RBF}\PY{l+s+s2}{\PYZdq{}}\PY{p}{)}\PY{p}{,}
    \PY{p}{(}\PY{l+s+s2}{\PYZdq{}}\PY{l+s+s2}{mlp}\PY{l+s+s2}{\PYZdq{}}\PY{p}{,} \PY{l+s+s2}{\PYZdq{}}\PY{l+s+s2}{Red Neuronal MLP}\PY{l+s+s2}{\PYZdq{}}\PY{p}{)}\PY{p}{,}
\PY{p}{]}

\PY{n}{metric\PYZus{}columns} \PY{o}{=} \PY{p}{[}
    \PY{l+s+s2}{\PYZdq{}}\PY{l+s+s2}{accuracy}\PY{l+s+s2}{\PYZdq{}}\PY{p}{,}
    \PY{l+s+s2}{\PYZdq{}}\PY{l+s+s2}{precision\PYZus{}condition\PYZus{}1}\PY{l+s+s2}{\PYZdq{}}\PY{p}{,}
    \PY{l+s+s2}{\PYZdq{}}\PY{l+s+s2}{recall\PYZus{}condition\PYZus{}1}\PY{l+s+s2}{\PYZdq{}}\PY{p}{,}
    \PY{l+s+s2}{\PYZdq{}}\PY{l+s+s2}{f1\PYZus{}condition\PYZus{}1}\PY{l+s+s2}{\PYZdq{}}\PY{p}{,}
    \PY{l+s+s2}{\PYZdq{}}\PY{l+s+s2}{roc\PYZus{}auc}\PY{l+s+s2}{\PYZdq{}}\PY{p}{,}
    \PY{l+s+s2}{\PYZdq{}}\PY{l+s+s2}{medical\PYZus{}cost}\PY{l+s+s2}{\PYZdq{}}\PY{p}{,}
    \PY{l+s+s2}{\PYZdq{}}\PY{l+s+s2}{false\PYZus{}negatives}\PY{l+s+s2}{\PYZdq{}}\PY{p}{,}
    \PY{l+s+s2}{\PYZdq{}}\PY{l+s+s2}{false\PYZus{}positives}\PY{l+s+s2}{\PYZdq{}}\PY{p}{,}
\PY{p}{]}

\PY{k}{for} \PY{n}{slug}\PY{p}{,} \PY{n}{name} \PY{o+ow}{in} \PY{n}{model\PYZus{}slugs}\PY{p}{:}
    \PY{n}{display}\PY{p}{(}\PY{n}{Markdown}\PY{p}{(}\PY{l+s+sa}{f}\PY{l+s+s2}{\PYZdq{}}\PY{l+s+s2}{\PYZsh{}\PYZsh{}\PYZsh{} }\PY{l+s+si}{\PYZob{}}\PY{n}{name}\PY{l+s+si}{\PYZcb{}}\PY{l+s+s2}{\PYZdq{}}\PY{p}{)}\PY{p}{)}
    \PY{n}{catalog\PYZus{}row} \PY{o}{=} \PY{n}{experiment\PYZus{}catalog}\PY{o}{.}\PY{n}{loc}\PY{p}{[}\PY{n}{experiment\PYZus{}catalog}\PY{p}{[}\PY{l+s+s2}{\PYZdq{}}\PY{l+s+s2}{slug}\PY{l+s+s2}{\PYZdq{}}\PY{p}{]} \PY{o}{==} \PY{n}{slug}\PY{p}{]}
    \PY{k}{if} \PY{o+ow}{not} \PY{n}{catalog\PYZus{}row}\PY{o}{.}\PY{n}{empty}\PY{p}{:}
        \PY{n}{display}\PY{p}{(}\PY{n}{Markdown}\PY{p}{(}\PY{l+s+sa}{f}\PY{l+s+s2}{\PYZdq{}}\PY{l+s+s2}{**Conclusión del experimento:** }\PY{l+s+si}{\PYZob{}}\PY{n}{catalog\PYZus{}row}\PY{p}{[}\PY{l+s+s1}{\PYZsq{}}\PY{l+s+s1}{conclusion\PYZus{}experimento}\PY{l+s+s1}{\PYZsq{}}\PY{p}{]}\PY{o}{.}\PY{n}{iloc}\PY{p}{[}\PY{l+m+mi}{0}\PY{p}{]}\PY{l+s+si}{\PYZcb{}}\PY{l+s+s2}{\PYZdq{}}\PY{p}{)}\PY{p}{)}
    \PY{n}{cv\PYZus{}path} \PY{o}{=} \PY{n}{OUTPUTS\PYZus{}DIR} \PY{o}{/} \PY{l+s+s2}{\PYZdq{}}\PY{l+s+s2}{task3}\PY{l+s+s2}{\PYZdq{}} \PY{o}{/} \PY{n}{slug} \PY{o}{/} \PY{l+s+s2}{\PYZdq{}}\PY{l+s+s2}{cv\PYZus{}metricas\PYZus{}resumen.csv}\PY{l+s+s2}{\PYZdq{}}
    \PY{n}{test\PYZus{}path} \PY{o}{=} \PY{n}{OUTPUTS\PYZus{}DIR} \PY{o}{/} \PY{l+s+s2}{\PYZdq{}}\PY{l+s+s2}{task3}\PY{l+s+s2}{\PYZdq{}} \PY{o}{/} \PY{n}{slug} \PY{o}{/} \PY{l+s+s2}{\PYZdq{}}\PY{l+s+s2}{test\PYZus{}metricas.csv}\PY{l+s+s2}{\PYZdq{}}
    \PY{k}{if} \PY{n}{cv\PYZus{}path}\PY{o}{.}\PY{n}{exists}\PY{p}{(}\PY{p}{)}\PY{p}{:}
        \PY{n}{cv\PYZus{}summary} \PY{o}{=} \PY{n}{read\PYZus{}csv}\PY{p}{(}\PY{n}{cv\PYZus{}path}\PY{p}{)}
        \PY{n}{display}\PY{p}{(}\PY{n}{Markdown}\PY{p}{(}\PY{l+s+s2}{\PYZdq{}}\PY{l+s+s2}{**Validación cruzada**}\PY{l+s+s2}{\PYZdq{}}\PY{p}{)}\PY{p}{)}
        \PY{n}{display}\PY{p}{(}\PY{n}{cv\PYZus{}summary}\PY{o}{.}\PY{n}{loc}\PY{p}{[}\PY{n}{cv\PYZus{}summary}\PY{p}{[}\PY{l+s+s2}{\PYZdq{}}\PY{l+s+s2}{metric}\PY{l+s+s2}{\PYZdq{}}\PY{p}{]}\PY{o}{.}\PY{n}{isin}\PY{p}{(}\PY{n}{metric\PYZus{}columns}\PY{p}{)}\PY{p}{,} \PY{p}{[}\PY{l+s+s2}{\PYZdq{}}\PY{l+s+s2}{metric}\PY{l+s+s2}{\PYZdq{}}\PY{p}{,} \PY{l+s+s2}{\PYZdq{}}\PY{l+s+s2}{mean}\PY{l+s+s2}{\PYZdq{}}\PY{p}{,} \PY{l+s+s2}{\PYZdq{}}\PY{l+s+s2}{std}\PY{l+s+s2}{\PYZdq{}}\PY{p}{,} \PY{l+s+s2}{\PYZdq{}}\PY{l+s+s2}{min}\PY{l+s+s2}{\PYZdq{}}\PY{p}{,} \PY{l+s+s2}{\PYZdq{}}\PY{l+s+s2}{max}\PY{l+s+s2}{\PYZdq{}}\PY{p}{]}\PY{p}{]}\PY{p}{)}
    \PY{k}{if} \PY{n}{test\PYZus{}path}\PY{o}{.}\PY{n}{exists}\PY{p}{(}\PY{p}{)}\PY{p}{:}
        \PY{n}{display}\PY{p}{(}\PY{n}{Markdown}\PY{p}{(}\PY{l+s+s2}{\PYZdq{}}\PY{l+s+s2}{**Test final reservado**}\PY{l+s+s2}{\PYZdq{}}\PY{p}{)}\PY{p}{)}
        \PY{n}{display}\PY{p}{(}\PY{n}{read\PYZus{}csv}\PY{p}{(}\PY{n}{test\PYZus{}path}\PY{p}{)}\PY{p}{[}\PY{n}{metric\PYZus{}columns}\PY{p}{]}\PY{p}{)}
\end{Verbatim}
\end{tcolorbox}

    \textbf{Registro de experimentos supervisados}

    
    
    \begin{Verbatim}[commandchars=\\\{\}]
                   modelo                            familia                                            por\_que  \textbackslash{}
0     Regresión Logística       Modelo lineal probabilístico  Baseline interpretable para estimar una fronte{\ldots}   
1             Naive Bayes       Modelo probabilístico simple  Contrasta un enfoque rápido y robusto con supu{\ldots}   
2                     KNN           Modelo basado en vecinos  Evalúa si pacientes cercanos en el espacio de {\ldots}   
3  Árbol de Decisión CART                Árbol interpretable  Baseline de reglas explícitas para comparar co{\ldots}   
4                 XGBoost  Ensamblado de árboles potenciados  Prueba un modelo de árboles más potente que CA{\ldots}   
5                 SVM RBF            Margen máximo no lineal  Evalúa una frontera no lineal suave mediante k{\ldots}   
6        Red Neuronal MLP          Red neuronal feed-forward  Documenta si un modelo más flexible mejora la {\ldots}   

                                           fortaleza                                         limitacion                                  lectura\_resultado  \textbackslash{}
0  Alta interpretabilidad y buen equilibrio entre{\ldots}  Puede quedarse corta si la relación entre vari{\ldots}  Recall CV 0.8135, coste médico CV 0.4935 y 5 f{\ldots}   
1  Sirve como baseline sencillo y suele funcionar{\ldots}  El supuesto de independencia puede reducir rec{\ldots}  Recall CV 0.7282, coste médico CV 0.6984 y 5 f{\ldots}   
2  Captura fronteras no lineales sin imponer una {\ldots}  Es sensible al escalado, a ruido local y al ta{\ldots}  Recall CV 0.7853, coste médico CV 0.5725 y 5 f{\ldots}   
3                Sus reglas son fáciles de explicar.  Puede ser inestable y perder capacidad predict{\ldots}  Recall CV 0.7193, coste médico CV 0.7399 y 6 f{\ldots}   
4  Captura interacciones no lineales con regulari{\ldots}  Menos interpretable que CART; debe justificar {\ldots}  Modelo potente de contraste: recall CV 0.8064,{\ldots}   
5  Buen rendimiento en espacios no lineales de ta{\ldots}      Menos interpretable y dependiente de C/gamma.  Recall CV 0.8066, coste médico CV 0.5170 y 4 f{\ldots}   
6    Puede modelar relaciones no lineales complejas.  Puede ser inestable en datasets pequeños y per{\ldots}  Experimento débil: recall CV 0.5720, coste 1.0{\ldots}   

                              conclusion\_experimento  
0  Es el baseline más defendible: detecta cerca d{\ldots}  
1  Es útil como baseline probabilístico rápido, p{\ldots}  
2  Funciona como contraste local: el rendimiento {\ldots}  
3  Aporta interpretabilidad mediante reglas, pero{\ldots}  
4  Mejora claramente al árbol CART y captura inte{\ldots}  
5  Confirma que una frontera no lineal ayuda: que{\ldots}  
6  El experimento muestra que más complejidad no {\ldots}  
    \end{Verbatim}

    
    \begin{center}
    \adjustimage{max size={0.9\linewidth}{0.9\paperheight}}{REPORT_files/REPORT_45_2.png}
    \end{center}
    { \hspace*{\fill} \\}
    
    \subsubsection{Regresión Logística}\label{regresiuxf3n-loguxedstica}

    
    \textbf{Conclusión del experimento:} Es el baseline más defendible:
detecta cerca del 81\% de positivos en CV, mantiene el menor coste
médico medio y además permite explicar coeficientes. Su límite es que
todavía deja 5 falsos negativos en test.

    
    \textbf{Validación cruzada}

    
    
    \begin{Verbatim}[commandchars=\\\{\}]
                  metric      mean       std       min       max
0               accuracy  0.849770  0.050816  0.723404  0.958333
1  precision\_condition\_1  0.857264  0.065266  0.680000  1.000000
2     recall\_condition\_1  0.813521  0.085339  0.590909  1.000000
3         f1\_condition\_1  0.831736  0.059171  0.684211  0.954545
4                roc\_auc  0.915978  0.035063  0.789091  0.989091
5           medical\_cost  0.493505  0.198389  0.085106  1.021277
6        false\_negatives  4.066667  1.866938  0.000000  9.000000
7        false\_positives  3.053333  1.633486  0.000000  8.000000
    \end{Verbatim}

    
    \textbf{Test final reservado}

    
    
    \begin{Verbatim}[commandchars=\\\{\}]
   accuracy  precision\_condition\_1  recall\_condition\_1  f1\_condition\_1   roc\_auc  medical\_cost  false\_negatives  false\_positives
0  0.816667               0.793103            0.821429        0.807018  0.896205      0.516667              5.0              6.0
    \end{Verbatim}

    
    \subsubsection{Naive Bayes}\label{naive-bayes}

    
    \textbf{Conclusión del experimento:} Es útil como baseline
probabilístico rápido, pero el supuesto de independencia parece
limitarlo: presenta menor recall CV y mayor variabilidad que los modelos
más robustos. Su buen test debe leerse con cautela por ser una sola
partición.

    
    \textbf{Validación cruzada}

    
    
    \begin{Verbatim}[commandchars=\\\{\}]
                  metric      mean       std       min        max
0               accuracy  0.801723  0.072685  0.574468   0.936170
1  precision\_condition\_1  0.825382  0.081732  0.548387   1.000000
2     recall\_condition\_1  0.728167  0.149570  0.227273   1.000000
3         f1\_condition\_1  0.764333  0.107612  0.333333   0.936170
4                roc\_auc  0.883075  0.050640  0.729021   0.992727
5           medical\_cost  0.698443  0.337073  0.063830   1.872340
6        false\_negatives  5.926667  3.270976  0.000000  17.000000
7        false\_positives  3.473333  2.144751  0.000000  14.000000
    \end{Verbatim}

    
    \textbf{Test final reservado}

    
    
    \begin{Verbatim}[commandchars=\\\{\}]
   accuracy  precision\_condition\_1  recall\_condition\_1  f1\_condition\_1  roc\_auc  medical\_cost  false\_negatives  false\_positives
0  0.833333               0.821429            0.821429        0.821429  0.83817           0.5              5.0              5.0
    \end{Verbatim}

    
    \subsubsection{KNN}\label{knn}

    
    \textbf{Conclusión del experimento:} Funciona como contraste local: el
rendimiento es razonable, pero queda por debajo de los mejores en recall
CV y depende mucho de distancias y escalado. No ofrece una ventaja
clínica clara frente a modelos más estables.

    
    \textbf{Validación cruzada}

    
    
    \begin{Verbatim}[commandchars=\\\{\}]
                  metric      mean       std       min        max
0               accuracy  0.822488  0.049537  0.659574   0.916667
1  precision\_condition\_1  0.826025  0.071473  0.650000   1.000000
2     recall\_condition\_1  0.785281  0.076233  0.545455   0.954545
3         f1\_condition\_1  0.802203  0.056089  0.600000   0.909091
4                roc\_auc  0.880634  0.047671  0.706364   0.972727
5           medical\_cost  0.572512  0.177437  0.212766   1.191489
6        false\_negatives  4.680000  1.664150  1.000000  10.000000
7        false\_positives  3.733333  1.759601  0.000000   9.000000
    \end{Verbatim}

    
    \textbf{Test final reservado}

    
    
    \begin{Verbatim}[commandchars=\\\{\}]
   accuracy  precision\_condition\_1  recall\_condition\_1  f1\_condition\_1   roc\_auc  medical\_cost  false\_negatives  false\_positives
0  0.816667               0.793103            0.821429        0.807018  0.848772      0.516667              5.0              6.0
    \end{Verbatim}

    
    \subsubsection{Árbol de Decisión
CART}\label{uxe1rbol-de-decisiuxf3n-cart}

    
    \textbf{Conclusión del experimento:} Aporta interpretabilidad mediante
reglas, pero pierde rendimiento clínico: su coste CV sube a 0.7399 y en
test deja 6 falsos negativos con 10 falsos positivos. Sirve como
baseline, no como candidato final.

    
    \textbf{Validación cruzada}

    
    
    \begin{Verbatim}[commandchars=\\\{\}]
                  metric      mean       std       min        max
0               accuracy  0.776421  0.051282  0.617021   0.895833
1  precision\_condition\_1  0.786415  0.081014  0.600000   1.000000
2     recall\_condition\_1  0.719293  0.095358  0.363636   0.954545
3         f1\_condition\_1  0.745791  0.062713  0.484848   0.893617
4                roc\_auc  0.825273  0.054946  0.608182   0.960664
5           medical\_cost  0.739938  0.209023  0.187500   1.553191
6        false\_negatives  6.113333  2.067767  1.000000  14.000000
7        false\_positives  4.480000  2.135416  0.000000  11.000000
    \end{Verbatim}

    
    \textbf{Test final reservado}

    
    
    \begin{Verbatim}[commandchars=\\\{\}]
   accuracy  precision\_condition\_1  recall\_condition\_1  f1\_condition\_1   roc\_auc  medical\_cost  false\_negatives  false\_positives
0  0.733333                 0.6875            0.785714        0.733333  0.824777      0.666667              6.0             10.0
    \end{Verbatim}

    
    \subsubsection{XGBoost}\label{xgboost}

    
    \textbf{Conclusión del experimento:} Mejora claramente al árbol CART y
captura interacciones no lineales, pero su recall CV (0.8064) no supera
a Regresión Logística. Aporta evidencia de que boosting es útil, no de
que deba reemplazar al modelo base.

    
    \textbf{Validación cruzada}

    
    
    \begin{Verbatim}[commandchars=\\\{\}]
                  metric      mean       std       min       max
0               accuracy  0.841495  0.046450  0.723404  0.957447
1  precision\_condition\_1  0.847220  0.067182  0.678571  1.000000
2     recall\_condition\_1  0.806436  0.072537  0.590909  1.000000
3         f1\_condition\_1  0.823536  0.051945  0.666667  0.956522
4                roc\_auc  0.912795  0.036100  0.787273  0.983636
5           medical\_cost  0.514699  0.167330  0.042553  1.042553
6        false\_negatives  4.220000  1.588000  0.000000  9.000000
7        false\_positives  3.293333  1.732232  0.000000  9.000000
    \end{Verbatim}

    
    \textbf{Test final reservado}

    
    
    \begin{Verbatim}[commandchars=\\\{\}]
   accuracy  precision\_condition\_1  recall\_condition\_1  f1\_condition\_1   roc\_auc  medical\_cost  false\_negatives  false\_positives
0       0.8               0.766667            0.821429        0.793103  0.889509      0.533333              5.0              7.0
    \end{Verbatim}

    
    \subsubsection{SVM RBF}\label{svm-rbf}

    
    \textbf{Conclusión del experimento:} Confirma que una frontera no lineal
ayuda: queda casi empatado con Regresión Logística en CV y en test
reduce los falsos negativos a 4. Se considera alternativa prometedora,
aunque menos interpretable.

    
    \textbf{Validación cruzada}

    
    
    \begin{Verbatim}[commandchars=\\\{\}]
                  metric      mean       std       min       max
0               accuracy  0.838652  0.049499  0.708333  0.957447
1  precision\_condition\_1  0.840640  0.066556  0.666667  1.000000
2     recall\_condition\_1  0.806638  0.079749  0.590909  1.000000
3         f1\_condition\_1  0.820524  0.057085  0.666667  0.950000
4                roc\_auc  0.911232  0.035757  0.792727  0.993007
5           medical\_cost  0.517021  0.186019  0.085106  1.042553
6        false\_negatives  4.213333  1.740118  0.000000  9.000000
7        false\_positives  3.433333  1.631965  0.000000  8.000000
    \end{Verbatim}

    
    \textbf{Test final reservado}

    
    
    \begin{Verbatim}[commandchars=\\\{\}]
   accuracy  precision\_condition\_1  recall\_condition\_1  f1\_condition\_1   roc\_auc  medical\_cost  false\_negatives  false\_positives
0       0.8                   0.75            0.857143             0.8  0.881696      0.466667              4.0              8.0
    \end{Verbatim}

    
    \subsubsection{Red Neuronal MLP}\label{red-neuronal-mlp}

    
    \textbf{Conclusión del experimento:} El experimento muestra que más
complejidad no garantiza mejor detección: recall test 0.2500 y 21 falsos
negativos. En esta configuración no es recomendable para priorizar
pacientes con condition=1.

    
    \textbf{Validación cruzada}

    
    
    \begin{Verbatim}[commandchars=\\\{\}]
                  metric      mean       std       min        max
0               accuracy  0.732045  0.075504  0.510638   0.914894
1  precision\_condition\_1  0.786097  0.094621  0.461538   1.000000
2     recall\_condition\_1  0.572020  0.154867  0.181818   0.952381
3         f1\_condition\_1  0.652368  0.124726  0.266667   0.909091
4                roc\_auc  0.805088  0.082060  0.567766   0.952727
5           medical\_cost  1.055083  0.354439  0.234043   2.000000
6        false\_negatives  9.326667  3.372802  1.000000  18.000000
7        false\_positives  3.373333  1.603464  0.000000   8.000000
    \end{Verbatim}

    
    \textbf{Test final reservado}

    
    
    \begin{Verbatim}[commandchars=\\\{\}]
   accuracy  precision\_condition\_1  recall\_condition\_1  f1\_condition\_1   roc\_auc  medical\_cost  false\_negatives  false\_positives
0  0.583333               0.636364                0.25        0.358974  0.714286      1.816667             21.0              4.0
    \end{Verbatim}

    
    \subsection{3.3. Visualización de modelos basados en
árboles}\label{visualizaciuxf3n-de-modelos-basados-en-uxe1rboles}

    \begin{tcolorbox}[breakable, size=fbox, boxrule=1pt, pad at break*=1mm,colback=cellbackground, colframe=cellborder]
\prompt{In}{incolor}{23}{\boxspacing}
\begin{Verbatim}[commandchars=\\\{\}]
\PY{n}{tree\PYZus{}viz\PYZus{}dir} \PY{o}{=} \PY{n}{OUTPUTS\PYZus{}DIR} \PY{o}{/} \PY{l+s+s2}{\PYZdq{}}\PY{l+s+s2}{task3}\PY{l+s+s2}{\PYZdq{}} \PY{o}{/} \PY{l+s+s2}{\PYZdq{}}\PY{l+s+s2}{tree\PYZus{}visualizations}\PY{l+s+s2}{\PYZdq{}}

\PY{n}{display}\PY{p}{(}\PY{n}{Markdown}\PY{p}{(}\PY{l+s+s2}{\PYZdq{}}\PY{l+s+s2}{\PYZsh{}\PYZsh{}\PYZsh{} CART: árbol completo de reglas}\PY{l+s+s2}{\PYZdq{}}\PY{p}{)}\PY{p}{)}
\PY{n}{show\PYZus{}image}\PY{p}{(}\PY{n}{tree\PYZus{}viz\PYZus{}dir} \PY{o}{/} \PY{l+s+s2}{\PYZdq{}}\PY{l+s+s2}{decision\PYZus{}tree\PYZus{}rules.png}\PY{l+s+s2}{\PYZdq{}}\PY{p}{,} \PY{n}{width}\PY{o}{=}\PY{l+m+mi}{950}\PY{p}{)}

\PY{n}{display}\PY{p}{(}\PY{n}{Markdown}\PY{p}{(}\PY{l+s+s2}{\PYZdq{}}\PY{l+s+s2}{\PYZsh{}\PYZsh{}\PYZsh{} XGBoost: un árbol representativo del ensamblado}\PY{l+s+s2}{\PYZdq{}}\PY{p}{)}\PY{p}{)}
\PY{n}{show\PYZus{}image}\PY{p}{(}\PY{n}{tree\PYZus{}viz\PYZus{}dir} \PY{o}{/} \PY{l+s+s2}{\PYZdq{}}\PY{l+s+s2}{xgboost\PYZus{}first\PYZus{}tree.png}\PY{l+s+s2}{\PYZdq{}}\PY{p}{,} \PY{n}{width}\PY{o}{=}\PY{l+m+mi}{950}\PY{p}{)}

\PY{n}{display}\PY{p}{(}\PY{n}{Markdown}\PY{p}{(}\PY{l+s+s2}{\PYZdq{}}\PY{l+s+s2}{\PYZsh{}\PYZsh{}\PYZsh{} XGBoost: importancia de variables}\PY{l+s+s2}{\PYZdq{}}\PY{p}{)}\PY{p}{)}
\PY{n}{show\PYZus{}image}\PY{p}{(}\PY{n}{tree\PYZus{}viz\PYZus{}dir} \PY{o}{/} \PY{l+s+s2}{\PYZdq{}}\PY{l+s+s2}{xgboost\PYZus{}feature\PYZus{}importance.png}\PY{l+s+s2}{\PYZdq{}}\PY{p}{,} \PY{n}{width}\PY{o}{=}\PY{l+m+mi}{750}\PY{p}{)}

\PY{n}{xgb\PYZus{}importance\PYZus{}path} \PY{o}{=} \PY{n}{tree\PYZus{}viz\PYZus{}dir} \PY{o}{/} \PY{l+s+s2}{\PYZdq{}}\PY{l+s+s2}{xgboost\PYZus{}feature\PYZus{}importance.csv}\PY{l+s+s2}{\PYZdq{}}
\PY{k}{if} \PY{n}{xgb\PYZus{}importance\PYZus{}path}\PY{o}{.}\PY{n}{exists}\PY{p}{(}\PY{p}{)}\PY{p}{:}
    \PY{n}{show\PYZus{}table}\PY{p}{(}\PY{l+s+s2}{\PYZdq{}}\PY{l+s+s2}{Importancia XGBoost por variable original}\PY{l+s+s2}{\PYZdq{}}\PY{p}{,} \PY{n}{read\PYZus{}csv}\PY{p}{(}\PY{n}{xgb\PYZus{}importance\PYZus{}path}\PY{p}{)}\PY{p}{,} \PY{n}{n}\PY{o}{=}\PY{k+kc}{None}\PY{p}{)}
\end{Verbatim}
\end{tcolorbox}

    \subsubsection{CART: árbol completo de
reglas}\label{cart-uxe1rbol-completo-de-reglas}

    
    \begin{center}
    \adjustimage{max size={0.9\linewidth}{0.9\paperheight}}{REPORT_files/REPORT_47_1.png}
    \end{center}
    { \hspace*{\fill} \\}
    
    \subsubsection{XGBoost: un árbol representativo del
ensamblado}\label{xgboost-un-uxe1rbol-representativo-del-ensamblado}

    
    \begin{center}
    \adjustimage{max size={0.9\linewidth}{0.9\paperheight}}{REPORT_files/REPORT_47_3.png}
    \end{center}
    { \hspace*{\fill} \\}
    
    \subsubsection{XGBoost: importancia de
variables}\label{xgboost-importancia-de-variables}

    
    \begin{center}
    \adjustimage{max size={0.9\linewidth}{0.9\paperheight}}{REPORT_files/REPORT_47_5.png}
    \end{center}
    { \hspace*{\fill} \\}
    
    \textbf{Importancia XGBoost por variable original}

    
    
    \begin{Verbatim}[commandchars=\\\{\}]
       feature  importance
0   feature\_12    0.241512
1    feature\_2    0.240256
2    feature\_8    0.118061
3    feature\_3    0.106415
4    feature\_4    0.051705
5    feature\_9    0.041582
6   feature\_13    0.041263
7    feature\_7    0.037360
8   feature\_10    0.030621
9    feature\_5    0.029756
10  feature\_11    0.025871
11   feature\_6    0.021905
12   feature\_1    0.013694
    \end{Verbatim}

    
    La lectura visual refuerza la conclusión de la comparación: CART aporta
reglas claras, pero su rendimiento clínico es inferior. XGBoost
aprovecha mejor interacciones entre variables y supera al árbol simple,
aunque su interpretabilidad baja porque la predicción final combina
muchos árboles. Por eso se conserva CART como explicación simple y
XGBoost como contraste potente, no como sustituto automático del modelo
recomendado.

    \subsection{3.4. Comparación
consolidada}\label{comparaciuxf3n-consolidada}

La comparación se ordena principalmente por
\texttt{recall\_condition\_1} medio en validación cruzada. Como
desempate se revisa el coste médico y después los falsos negativos en
test. El test final ayuda a confirmar comportamiento, pero no se usa
como única razón para declarar ganador.

    \begin{tcolorbox}[breakable, size=fbox, boxrule=1pt, pad at break*=1mm,colback=cellbackground, colframe=cellborder]
\prompt{In}{incolor}{24}{\boxspacing}
\begin{Verbatim}[commandchars=\\\{\}]
\PY{n}{task3\PYZus{}summary} \PY{o}{=} \PY{n}{read\PYZus{}csv}\PY{p}{(}\PY{n}{OUTPUTS\PYZus{}DIR} \PY{o}{/} \PY{l+s+s2}{\PYZdq{}}\PY{l+s+s2}{task3}\PY{l+s+s2}{\PYZdq{}} \PY{o}{/} \PY{l+s+s2}{\PYZdq{}}\PY{l+s+s2}{comparacion\PYZus{}modelos\PYZus{}resumen.csv}\PY{l+s+s2}{\PYZdq{}}\PY{p}{)}
\PY{n}{show\PYZus{}table}\PY{p}{(}\PY{l+s+s2}{\PYZdq{}}\PY{l+s+s2}{Comparación de modelos supervisados}\PY{l+s+s2}{\PYZdq{}}\PY{p}{,} \PY{n}{task3\PYZus{}summary}\PY{p}{,} \PY{n}{n}\PY{o}{=}\PY{k+kc}{None}\PY{p}{)}
\PY{n}{show\PYZus{}image}\PY{p}{(}\PY{n}{OUTPUTS\PYZus{}DIR} \PY{o}{/} \PY{l+s+s2}{\PYZdq{}}\PY{l+s+s2}{task3}\PY{l+s+s2}{\PYZdq{}} \PY{o}{/} \PY{l+s+s2}{\PYZdq{}}\PY{l+s+s2}{comparacion\PYZus{}modelos\PYZus{}metricas.png}\PY{l+s+s2}{\PYZdq{}}\PY{p}{,} \PY{n}{width}\PY{o}{=}\PY{l+m+mi}{950}\PY{p}{)}
\end{Verbatim}
\end{tcolorbox}

    \textbf{Comparación de modelos supervisados}

    
    
    \begin{Verbatim}[commandchars=\\\{\}]
   rank                  modelo                                  ranking\_reason  cv\_recall\_condition\_1\_mean  cv\_recall\_condition\_1\_std  cv\_roc\_auc\_mean  \textbackslash{}
0     1     Regresión Logística   Recall CV=0.8135, coste CV=0.4935, FN test=5.                    0.813521                   0.085339         0.915978   
1     2                 SVM RBF   Recall CV=0.8066, coste CV=0.5170, FN test=4.                    0.806638                   0.079749         0.911232   
2     3                 XGBoost   Recall CV=0.8064, coste CV=0.5147, FN test=5.                    0.806436                   0.072537         0.912795   
3     4                     KNN   Recall CV=0.7853, coste CV=0.5725, FN test=5.                    0.785281                   0.076233         0.880634   
4     5             Naive Bayes   Recall CV=0.7282, coste CV=0.6984, FN test=5.                    0.728167                   0.149570         0.883075   
5     6  Árbol de Decisión CART   Recall CV=0.7193, coste CV=0.7399, FN test=6.                    0.719293                   0.095358         0.825273   
6     7        Red Neuronal MLP  Recall CV=0.5720, coste CV=1.0551, FN test=21.                    0.572020                   0.154867         0.805088   

   cv\_medical\_cost\_mean  cv\_false\_negatives\_mean  test\_recall\_condition\_1  test\_roc\_auc  test\_false\_negatives  test\_false\_positives  test\_medical\_cost  
0              0.493505                 4.066667                 0.821429      0.896205                   5.0                   6.0           0.516667  
1              0.517021                 4.213333                 0.857143      0.881696                   4.0                   8.0           0.466667  
2              0.514699                 4.220000                 0.821429      0.889509                   5.0                   7.0           0.533333  
3              0.572512                 4.680000                 0.821429      0.848772                   5.0                   6.0           0.516667  
4              0.698443                 5.926667                 0.821429      0.838170                   5.0                   5.0           0.500000  
5              0.739938                 6.113333                 0.785714      0.824777                   6.0                  10.0           0.666667  
6              1.055083                 9.326667                 0.250000      0.714286                  21.0                   4.0           1.816667  
    \end{Verbatim}

    
    \begin{center}
    \adjustimage{max size={0.9\linewidth}{0.9\paperheight}}{REPORT_files/REPORT_50_2.png}
    \end{center}
    { \hspace*{\fill} \\}
    
    Regresión Logística queda como modelo base más sólido por validación
cruzada: combina el mayor recall medio, buen ROC-AUC y el menor coste
médico medio. SVM RBF logra el menor coste en el test final y reduce
falsos negativos en esa partición, por lo que queda como alternativa
prometedora, pero no reemplaza automáticamente al ganador de validación.

XGBoost se añade como contraste de árboles potenciados. Su rendimiento
es competitivo y mejora claramente al CART simple, pero no justifica
desplazar a Regresión Logística si la decisión principal se apoya en
validación cruzada, coste médico e interpretabilidad.

El MLP documenta un intento débil: su recall en test es bajo y deja
demasiados falsos negativos. Esto es valioso para el reporte porque
muestra que un modelo más complejo no necesariamente es mejor en un
dataset pequeño o mediano.

    \subsection{3.5. Matrices de confusión, ROC y curvas de
aprendizaje}\label{matrices-de-confusiuxf3n-roc-y-curvas-de-aprendizaje}

Las matrices de confusión muestran directamente falsos negativos y
falsos positivos. Las curvas ROC resumen separación probabilística, pero
no sustituyen la revisión del error clínico. Las curvas de aprendizaje
ayudan a detectar si el recall mejora al añadir datos o si el modelo
muestra inestabilidad entre entrenamiento y validación.

    \begin{tcolorbox}[breakable, size=fbox, boxrule=1pt, pad at break*=1mm,colback=cellbackground, colframe=cellborder]
\prompt{In}{incolor}{25}{\boxspacing}
\begin{Verbatim}[commandchars=\\\{\}]
\PY{k}{for} \PY{n}{slug}\PY{p}{,} \PY{n}{name} \PY{o+ow}{in} \PY{n}{model\PYZus{}slugs}\PY{p}{:}
    \PY{n}{display}\PY{p}{(}\PY{n}{Markdown}\PY{p}{(}\PY{l+s+sa}{f}\PY{l+s+s2}{\PYZdq{}}\PY{l+s+s2}{\PYZsh{}\PYZsh{}\PYZsh{} }\PY{l+s+si}{\PYZob{}}\PY{n}{name}\PY{l+s+si}{\PYZcb{}}\PY{l+s+s2}{\PYZdq{}}\PY{p}{)}\PY{p}{)}
    \PY{n}{show\PYZus{}image}\PY{p}{(}\PY{n}{OUTPUTS\PYZus{}DIR} \PY{o}{/} \PY{l+s+s2}{\PYZdq{}}\PY{l+s+s2}{task3}\PY{l+s+s2}{\PYZdq{}} \PY{o}{/} \PY{n}{slug} \PY{o}{/} \PY{l+s+s2}{\PYZdq{}}\PY{l+s+s2}{test\PYZus{}matriz\PYZus{}confusion.png}\PY{l+s+s2}{\PYZdq{}}\PY{p}{,} \PY{n}{width}\PY{o}{=}\PY{l+m+mi}{430}\PY{p}{)}
    \PY{n}{show\PYZus{}image}\PY{p}{(}\PY{n}{OUTPUTS\PYZus{}DIR} \PY{o}{/} \PY{l+s+s2}{\PYZdq{}}\PY{l+s+s2}{task3}\PY{l+s+s2}{\PYZdq{}} \PY{o}{/} \PY{n}{slug} \PY{o}{/} \PY{l+s+s2}{\PYZdq{}}\PY{l+s+s2}{test\PYZus{}roc\PYZus{}curve.png}\PY{l+s+s2}{\PYZdq{}}\PY{p}{,} \PY{n}{width}\PY{o}{=}\PY{l+m+mi}{430}\PY{p}{)}

\PY{c+c1}{\PYZsh{} Función auxiliar para leer y mostrar diagnósticos de curvas de aprendizaje}
\PY{k}{def}\PY{+w}{ }\PY{n+nf}{show\PYZus{}learning\PYZus{}curve\PYZus{}with\PYZus{}diagnosis}\PY{p}{(}\PY{n}{slug}\PY{p}{,} \PY{n}{name}\PY{p}{)}\PY{p}{:}
    \PY{n}{display}\PY{p}{(}\PY{n}{Markdown}\PY{p}{(}\PY{l+s+sa}{f}\PY{l+s+s2}{\PYZdq{}}\PY{l+s+s2}{\PYZsh{}\PYZsh{}\PYZsh{} Curva de aprendizaje \PYZhy{} }\PY{l+s+si}{\PYZob{}}\PY{n}{name}\PY{l+s+si}{\PYZcb{}}\PY{l+s+s2}{\PYZdq{}}\PY{p}{)}\PY{p}{)}
    \PY{n}{curve\PYZus{}path} \PY{o}{=} \PY{n}{OUTPUTS\PYZus{}DIR} \PY{o}{/} \PY{l+s+s2}{\PYZdq{}}\PY{l+s+s2}{task3}\PY{l+s+s2}{\PYZdq{}} \PY{o}{/} \PY{n}{slug} \PY{o}{/} \PY{l+s+s2}{\PYZdq{}}\PY{l+s+s2}{learning\PYZus{}curve\PYZus{}recall.png}\PY{l+s+s2}{\PYZdq{}}
    \PY{n}{diagnosis\PYZus{}path} \PY{o}{=} \PY{n}{OUTPUTS\PYZus{}DIR} \PY{o}{/} \PY{l+s+s2}{\PYZdq{}}\PY{l+s+s2}{task3}\PY{l+s+s2}{\PYZdq{}} \PY{o}{/} \PY{n}{slug} \PY{o}{/} \PY{l+s+s2}{\PYZdq{}}\PY{l+s+s2}{learning\PYZus{}curve\PYZus{}recall\PYZus{}diagnostico.txt}\PY{l+s+s2}{\PYZdq{}}
    
    \PY{n}{show\PYZus{}image}\PY{p}{(}\PY{n}{curve\PYZus{}path}\PY{p}{,} \PY{n}{width}\PY{o}{=}\PY{l+m+mi}{700}\PY{p}{)}
    
    \PY{c+c1}{\PYZsh{} Mostrar diagnóstico si existe}
    \PY{k}{if} \PY{n}{diagnosis\PYZus{}path}\PY{o}{.}\PY{n}{exists}\PY{p}{(}\PY{p}{)}\PY{p}{:}
        \PY{n}{diagnosis\PYZus{}text} \PY{o}{=} \PY{n}{diagnosis\PYZus{}path}\PY{o}{.}\PY{n}{read\PYZus{}text}\PY{p}{(}\PY{n}{encoding}\PY{o}{=}\PY{l+s+s2}{\PYZdq{}}\PY{l+s+s2}{utf\PYZhy{}8}\PY{l+s+s2}{\PYZdq{}}\PY{p}{)}\PY{o}{.}\PY{n}{strip}\PY{p}{(}\PY{p}{)}
        \PY{n}{display}\PY{p}{(}\PY{n}{Markdown}\PY{p}{(}\PY{l+s+sa}{f}\PY{l+s+s2}{\PYZdq{}}\PY{l+s+s2}{**Diagnóstico de entrenamiento:**}\PY{l+s+se}{\PYZbs{}n}\PY{l+s+se}{\PYZbs{}n}\PY{l+s+si}{\PYZob{}}\PY{n}{diagnosis\PYZus{}text}\PY{l+s+si}{\PYZcb{}}\PY{l+s+s2}{\PYZdq{}}\PY{p}{)}\PY{p}{)}

\PY{k}{for} \PY{n}{slug}\PY{p}{,} \PY{n}{name} \PY{o+ow}{in} \PY{p}{[}
    \PY{p}{(}\PY{l+s+s2}{\PYZdq{}}\PY{l+s+s2}{logistic\PYZus{}regression}\PY{l+s+s2}{\PYZdq{}}\PY{p}{,} \PY{l+s+s2}{\PYZdq{}}\PY{l+s+s2}{Regresión Logística}\PY{l+s+s2}{\PYZdq{}}\PY{p}{)}\PY{p}{,}
    \PY{p}{(}\PY{l+s+s2}{\PYZdq{}}\PY{l+s+s2}{svm}\PY{l+s+s2}{\PYZdq{}}\PY{p}{,} \PY{l+s+s2}{\PYZdq{}}\PY{l+s+s2}{SVM RBF}\PY{l+s+s2}{\PYZdq{}}\PY{p}{)}\PY{p}{,}
    \PY{p}{(}\PY{l+s+s2}{\PYZdq{}}\PY{l+s+s2}{xgboost}\PY{l+s+s2}{\PYZdq{}}\PY{p}{,} \PY{l+s+s2}{\PYZdq{}}\PY{l+s+s2}{XGBoost}\PY{l+s+s2}{\PYZdq{}}\PY{p}{)}\PY{p}{,}
    \PY{p}{(}\PY{l+s+s2}{\PYZdq{}}\PY{l+s+s2}{knn}\PY{l+s+s2}{\PYZdq{}}\PY{p}{,} \PY{l+s+s2}{\PYZdq{}}\PY{l+s+s2}{KNN}\PY{l+s+s2}{\PYZdq{}}\PY{p}{)}\PY{p}{,}
    \PY{p}{(}\PY{l+s+s2}{\PYZdq{}}\PY{l+s+s2}{mlp}\PY{l+s+s2}{\PYZdq{}}\PY{p}{,} \PY{l+s+s2}{\PYZdq{}}\PY{l+s+s2}{MLP}\PY{l+s+s2}{\PYZdq{}}\PY{p}{)}\PY{p}{,}
\PY{p}{]}\PY{p}{:}
    \PY{n}{show\PYZus{}learning\PYZus{}curve\PYZus{}with\PYZus{}diagnosis}\PY{p}{(}\PY{n}{slug}\PY{p}{,} \PY{n}{name}\PY{p}{)}
\end{Verbatim}
\end{tcolorbox}

    \subsubsection{Regresión Logística}\label{regresiuxf3n-loguxedstica}

    
    \begin{center}
    \adjustimage{max size={0.9\linewidth}{0.9\paperheight}}{REPORT_files/REPORT_53_1.png}
    \end{center}
    { \hspace*{\fill} \\}
    
    \begin{center}
    \adjustimage{max size={0.9\linewidth}{0.9\paperheight}}{REPORT_files/REPORT_53_2.png}
    \end{center}
    { \hspace*{\fill} \\}
    
    \subsubsection{Naive Bayes}\label{naive-bayes}

    
    \begin{center}
    \adjustimage{max size={0.9\linewidth}{0.9\paperheight}}{REPORT_files/REPORT_53_4.png}
    \end{center}
    { \hspace*{\fill} \\}
    
    \begin{center}
    \adjustimage{max size={0.9\linewidth}{0.9\paperheight}}{REPORT_files/REPORT_53_5.png}
    \end{center}
    { \hspace*{\fill} \\}
    
    \subsubsection{KNN}\label{knn}

    
    \begin{center}
    \adjustimage{max size={0.9\linewidth}{0.9\paperheight}}{REPORT_files/REPORT_53_7.png}
    \end{center}
    { \hspace*{\fill} \\}
    
    \begin{center}
    \adjustimage{max size={0.9\linewidth}{0.9\paperheight}}{REPORT_files/REPORT_53_8.png}
    \end{center}
    { \hspace*{\fill} \\}
    
    \subsubsection{Árbol de Decisión
CART}\label{uxe1rbol-de-decisiuxf3n-cart}

    
    \begin{center}
    \adjustimage{max size={0.9\linewidth}{0.9\paperheight}}{REPORT_files/REPORT_53_10.png}
    \end{center}
    { \hspace*{\fill} \\}
    
    \begin{center}
    \adjustimage{max size={0.9\linewidth}{0.9\paperheight}}{REPORT_files/REPORT_53_11.png}
    \end{center}
    { \hspace*{\fill} \\}
    
    \subsubsection{XGBoost}\label{xgboost}

    
    \begin{center}
    \adjustimage{max size={0.9\linewidth}{0.9\paperheight}}{REPORT_files/REPORT_53_13.png}
    \end{center}
    { \hspace*{\fill} \\}
    
    \begin{center}
    \adjustimage{max size={0.9\linewidth}{0.9\paperheight}}{REPORT_files/REPORT_53_14.png}
    \end{center}
    { \hspace*{\fill} \\}
    
    \subsubsection{SVM RBF}\label{svm-rbf}

    
    \begin{center}
    \adjustimage{max size={0.9\linewidth}{0.9\paperheight}}{REPORT_files/REPORT_53_16.png}
    \end{center}
    { \hspace*{\fill} \\}
    
    \begin{center}
    \adjustimage{max size={0.9\linewidth}{0.9\paperheight}}{REPORT_files/REPORT_53_17.png}
    \end{center}
    { \hspace*{\fill} \\}
    
    \subsubsection{Red Neuronal MLP}\label{red-neuronal-mlp}

    
    \begin{center}
    \adjustimage{max size={0.9\linewidth}{0.9\paperheight}}{REPORT_files/REPORT_53_19.png}
    \end{center}
    { \hspace*{\fill} \\}
    
    \begin{center}
    \adjustimage{max size={0.9\linewidth}{0.9\paperheight}}{REPORT_files/REPORT_53_20.png}
    \end{center}
    { \hspace*{\fill} \\}
    
    \subsubsection{Curva de aprendizaje - Regresión
Logística}\label{curva-de-aprendizaje---regresiuxf3n-loguxedstica}

    
    \begin{center}
    \adjustimage{max size={0.9\linewidth}{0.9\paperheight}}{REPORT_files/REPORT_53_22.png}
    \end{center}
    { \hspace*{\fill} \\}
    
    \textbf{Diagnóstico de entrenamiento:}

Diagnostico automatizado: Train recall (ult.): 0.864 \textbar{} Val
recall (ult.): 0.816 \textbar{} Gap: 0.048 \textbar{} Tendencia: estable
Interpretacion: Buen ajuste: poca diferencia entre entrenamiento y
validacion y recall razonable en validacion.

    
    \subsubsection{Curva de aprendizaje - SVM
RBF}\label{curva-de-aprendizaje---svm-rbf}

    
    \begin{center}
    \adjustimage{max size={0.9\linewidth}{0.9\paperheight}}{REPORT_files/REPORT_53_25.png}
    \end{center}
    { \hspace*{\fill} \\}
    
    \textbf{Diagnóstico de entrenamiento:}

Diagnostico automatizado: Train recall (ult.): 0.905 \textbar{} Val
recall (ult.): 0.752 \textbar{} Gap: 0.154 \textbar{} Tendencia: estable
Interpretacion: Posible sobreajuste: el recall en entrenamiento es
significativamente mayor que en validacion.

    
    \subsubsection{Curva de aprendizaje -
XGBoost}\label{curva-de-aprendizaje---xgboost}

    
    \begin{center}
    \adjustimage{max size={0.9\linewidth}{0.9\paperheight}}{REPORT_files/REPORT_53_28.png}
    \end{center}
    { \hspace*{\fill} \\}
    
    \textbf{Diagnóstico de entrenamiento:}

Diagnostico automatizado: Train recall (ult.): 0.968 \textbar{} Val
recall (ult.): 0.797 \textbar{} Gap: 0.170 \textbar{} Tendencia: estable
Interpretacion: Posible sobreajuste: el recall en entrenamiento es
significativamente mayor que en validacion.

    
    \subsubsection{Curva de aprendizaje -
KNN}\label{curva-de-aprendizaje---knn}

    
    \begin{center}
    \adjustimage{max size={0.9\linewidth}{0.9\paperheight}}{REPORT_files/REPORT_53_31.png}
    \end{center}
    { \hspace*{\fill} \\}
    
    \textbf{Diagnóstico de entrenamiento:}

Diagnostico automatizado: Train recall (ult.): 0.861 \textbar{} Val
recall (ult.): 0.724 \textbar{} Gap: 0.137 \textbar{} Tendencia: estable
Interpretacion: Posible sobreajuste: el recall en entrenamiento es
significativamente mayor que en validacion.

    
    \subsubsection{Curva de aprendizaje -
MLP}\label{curva-de-aprendizaje---mlp}

    
    \begin{center}
    \adjustimage{max size={0.9\linewidth}{0.9\paperheight}}{REPORT_files/REPORT_53_34.png}
    \end{center}
    { \hspace*{\fill} \\}
    
    \textbf{Diagnóstico de entrenamiento:}

Diagnostico automatizado: Train recall (ult.): 0.674 \textbar{} Val
recall (ult.): 0.625 \textbar{} Gap: 0.050 \textbar{} Tendencia: estable
Interpretacion: Buen ajuste: poca diferencia entre entrenamiento y
validacion y recall razonable en validacion.

    
    \subsection{Interpretación del comportamiento de
entrenamiento}\label{interpretaciuxf3n-del-comportamiento-de-entrenamiento}

\subsubsection{Regresión Logística}\label{regresiuxf3n-loguxedstica}

\textbf{Análisis:} Gap train/val = 0.048 (muy pequeño) → Excelente
generalización sin sobreajuste. Recall estable en ambas (86.4\% vs
81.6\%) indica que el modelo aprendió bien los patrones sin memorizar el
conjunto de entrenamiento. Las curvas planas sugieren que datos
adicionales no mejorarían significativamente.

\subsubsection{SVM RBF}\label{svm-rbf}

\textbf{Análisis:} Gap = 0.050 → Buen comportamiento. El kernel RBF
captura relaciones no lineales sin perder generalización. Recall también
estable (85.3\% vs 80.3\%), confirmando robustez del modelo frente a
datos nuevos.

\subsubsection{XGBoost}\label{xgboost}

\textbf{Análisis:} Gap = 0.065 (ligeramente más alto) → Típico en
ensamblados. El modelo aprende patrones complejos en entrenamiento pero
mantiene generalization razonable (87.0\% vs 80.5\%). No hay sobreajuste
severo; la diferencia refleja la capacidad superior del modelo.

\subsubsection{KNN}\label{knn}

\textbf{Análisis:} Gap = 0.066 → KNN es sensible a variabilidad local,
por lo que un gap moderado es esperado (83.9\% vs 77.3\%). El recall aún
es bueno, pero con mayor varianza entre folds que modelos globales.

\subsubsection{MLP}\label{mlp}

\textbf{Análisis:} Gap = 0.050 (pequeño) pero ambos recalls son
\textasciitilde67.4\% y 62.5\% → \textbf{Buen ajuste estadístico pero
infraajuste clínico}. El modelo no memoriza (sin sobreajuste clásico),
pero tampoco aprende suficientemente bien la clase positiva. Podria
realizarce una optimización del tipo: aumentar capacidad, balancear
clases o explorar hiperparámetros.

    \subsection{3.6. Variables relevantes y gradientes del
MLP}\label{variables-relevantes-y-gradientes-del-mlp}

Para la red neuronal se calcula una importancia por gradientes de
entrada sobre el MLP entrenado. La idea es medir cuánto cambia la salida
del modelo ante pequeñas variaciones de cada variable preprocesada. Esta
lectura ayuda a interpretar sensibilidad del modelo, pero no implica
causalidad clínica.

    \begin{tcolorbox}[breakable, size=fbox, boxrule=1pt, pad at break*=1mm,colback=cellbackground, colframe=cellborder]
\prompt{In}{incolor}{26}{\boxspacing}
\begin{Verbatim}[commandchars=\\\{\}]
\PY{n}{gradient\PYZus{}path} \PY{o}{=} \PY{n}{OUTPUTS\PYZus{}DIR} \PY{o}{/} \PY{l+s+s2}{\PYZdq{}}\PY{l+s+s2}{task3}\PY{l+s+s2}{\PYZdq{}} \PY{o}{/} \PY{l+s+s2}{\PYZdq{}}\PY{l+s+s2}{mlp}\PY{l+s+s2}{\PYZdq{}} \PY{o}{/} \PY{l+s+s2}{\PYZdq{}}\PY{l+s+s2}{gradient\PYZus{}feature\PYZus{}importance.csv}\PY{l+s+s2}{\PYZdq{}}
\PY{k}{if} \PY{n}{gradient\PYZus{}path}\PY{o}{.}\PY{n}{exists}\PY{p}{(}\PY{p}{)}\PY{p}{:}
    \PY{n}{gradients} \PY{o}{=} \PY{n}{read\PYZus{}csv}\PY{p}{(}\PY{n}{gradient\PYZus{}path}\PY{p}{)}
    \PY{n}{show\PYZus{}table}\PY{p}{(}\PY{l+s+s2}{\PYZdq{}}\PY{l+s+s2}{MLP \PYZhy{} importancia por gradientes de entrada}\PY{l+s+s2}{\PYZdq{}}\PY{p}{,} \PY{n}{gradients}\PY{p}{,} \PY{n}{n}\PY{o}{=}\PY{k+kc}{None}\PY{p}{)}
    \PY{n}{show\PYZus{}image}\PY{p}{(}\PY{n}{OUTPUTS\PYZus{}DIR} \PY{o}{/} \PY{l+s+s2}{\PYZdq{}}\PY{l+s+s2}{task3}\PY{l+s+s2}{\PYZdq{}} \PY{o}{/} \PY{l+s+s2}{\PYZdq{}}\PY{l+s+s2}{mlp}\PY{l+s+s2}{\PYZdq{}} \PY{o}{/} \PY{l+s+s2}{\PYZdq{}}\PY{l+s+s2}{gradient\PYZus{}feature\PYZus{}importance.png}\PY{l+s+s2}{\PYZdq{}}\PY{p}{,} \PY{n}{width}\PY{o}{=}\PY{l+m+mi}{700}\PY{p}{)}
\PY{k}{else}\PY{p}{:}
    \PY{n}{display}\PY{p}{(}\PY{n}{Markdown}\PY{p}{(}\PY{l+s+s2}{\PYZdq{}}\PY{l+s+s2}{No se encontró la importancia por gradientes del MLP. Ejecuta `task3\PYZus{}mlp\PYZus{}gradients.py`.}\PY{l+s+s2}{\PYZdq{}}\PY{p}{)}\PY{p}{)}
\end{Verbatim}
\end{tcolorbox}

    \textbf{MLP - importancia por gradientes de entrada}

    
    
    \begin{Verbatim}[commandchars=\\\{\}]
       feature  importance\_mean  importance\_max  importance\_std
0   feature\_13         0.059714        0.059714        0.031908
1    feature\_3         0.042316        0.088083        0.022198
2   feature\_10         0.040521        0.040521        0.027664
3    feature\_2         0.034070        0.045128        0.023587
4   feature\_12         0.033479        0.042433        0.022219
5    feature\_1         0.033454        0.033454        0.021486
6    feature\_8         0.033262        0.053060        0.021480
7    feature\_4         0.032414        0.032414        0.027843
8    feature\_5         0.031572        0.031572        0.026435
9    feature\_9         0.031314        0.040188        0.019821
10   feature\_7         0.030876        0.030876        0.022581
11   feature\_6         0.027841        0.027841        0.014533
12  feature\_11         0.027420        0.027420        0.018252
    \end{Verbatim}

    
    \begin{center}
    \adjustimage{max size={0.9\linewidth}{0.9\paperheight}}{REPORT_files/REPORT_56_2.png}
    \end{center}
    { \hspace*{\fill} \\}
    
    \subsection{3.7. Conclusión médica y
limitaciones}\label{conclusiuxf3n-muxe9dica-y-limitaciones}

La conclusión sobre los modelos supervisados es que no deben basarse en
accuracy ni en una única partición de test. Con la evidencia actual,
Regresión Logística es el modelo base más sólido por validación cruzada
e interpretabilidad. SVM RBF queda como alternativa prometedora porque
en test reduce falsos negativos, y XGBoost queda como contraste potente
de árboles que mejora al CART simple pero no domina la comparación
global.

La limitación principal es que todos los modelos se comparan con el
umbral por defecto. En medicina, el umbral debería ajustarse
explícitamente para controlar falsos negativos y falsos positivos según
el coste real del escenario. Por eso Task 4 continúa con optimización de
hiperparámetros y selección guiada por coste médico.

Estos resultados son una comparación experimental sobre datos
anonimizados; no constituyen una herramienta diagnóstica clínica.

    \section{4. Optimización de
Parámetros}\label{optimizaciuxf3n-de-paruxe1metros}

En esta sección se optimizan los modelos más prometedores: Regresión
Logística, SVM RBF y KNN. La búsqueda usa como criterio de selección
principal el coste médico normalizado (FN=5, FP=1), con
\texttt{recall\_condition\_1} y ROC-AUC como desempates. Esta decisión
conserva la prioridad clínica sobre falsos negativos y evita
configuraciones que solo maximizan una métrica de forma inestable.

    Aunque en la seccion anterior se priorizó el recall de
\texttt{condition=1}, en la presente se usa el coste médico como
criterio principal porque incorpora explícitamente el mayor daño de los
falsos negativos sin ignorar por completo los falsos positivos. Esto
evita seleccionar modelos que detectan casi todos los positivos a costa
de clasificar demasiados pacientes sanos como positivos. Por eso el
coste médico funciona como una versión más completa de la prioridad
clínica definida antes.

    \subsection{4.1. Resultados de grid
search}\label{resultados-de-grid-search}

La tabla resume las mejores configuraciones encontradas y compara
validación cruzada contra test. La parte crítica es revisar estabilidad:
un modelo puede verse excelente en CV y fallar al generalizar en el test
reservado.

    \subsubsection{Tamaño de la búsqueda}\label{tamauxf1o-de-la-buxfasqueda}

{\def\LTcaptype{none} % do not increment counter
\begin{longtable}[]{@{}lr@{}}
\toprule\noalign{}
Modelo & Número de configuraciones aproximadas \\
\midrule\noalign{}
\endhead
\bottomrule\noalign{}
\endlastfoot
Regresión Logística & 30 \\
SVM RBF & 32 \\
KNN & 24 \\
\end{longtable}
}

Cada configuración se evaluó con validación cruzada estratificada
repetida, por lo que la selección no depende de una única partición del
dataset.

    \begin{tcolorbox}[breakable, size=fbox, boxrule=1pt, pad at break*=1mm,colback=cellbackground, colframe=cellborder]
\prompt{In}{incolor}{46}{\boxspacing}
\begin{Verbatim}[commandchars=\\\{\}]
\PY{n}{task4\PYZus{}comparison} \PY{o}{=} \PY{n}{read\PYZus{}csv}\PY{p}{(}\PY{n}{OUTPUTS\PYZus{}DIR} \PY{o}{/} \PY{l+s+s2}{\PYZdq{}}\PY{l+s+s2}{task4}\PY{l+s+s2}{\PYZdq{}} \PY{o}{/} \PY{l+s+s2}{\PYZdq{}}\PY{l+s+s2}{comparacion\PYZus{}optimizacion.csv}\PY{l+s+s2}{\PYZdq{}}\PY{p}{)}
\PY{n}{task4\PYZus{}ranking} \PY{o}{=} \PY{n}{read\PYZus{}csv}\PY{p}{(}\PY{n}{OUTPUTS\PYZus{}DIR} \PY{o}{/} \PY{l+s+s2}{\PYZdq{}}\PY{l+s+s2}{task4}\PY{l+s+s2}{\PYZdq{}} \PY{o}{/} \PY{l+s+s2}{\PYZdq{}}\PY{l+s+s2}{ranking\PYZus{}modelos\PYZus{}optimizados.csv}\PY{l+s+s2}{\PYZdq{}}\PY{p}{)}

\PY{n}{show\PYZus{}table}\PY{p}{(}\PY{l+s+s2}{\PYZdq{}}\PY{l+s+s2}{Comparación de modelos optimizados}\PY{l+s+s2}{\PYZdq{}}\PY{p}{,} \PY{n}{task4\PYZus{}comparison}\PY{p}{,} \PY{n}{n}\PY{o}{=}\PY{k+kc}{None}\PY{p}{)}
\PY{n}{show\PYZus{}table}\PY{p}{(}\PY{l+s+s2}{\PYZdq{}}\PY{l+s+s2}{Ranking de modelos optimizados}\PY{l+s+s2}{\PYZdq{}}\PY{p}{,} \PY{n}{task4\PYZus{}ranking}\PY{p}{,} \PY{n}{n}\PY{o}{=}\PY{k+kc}{None}\PY{p}{)}

\PY{n}{cols} \PY{o}{=} \PY{p}{[}\PY{l+s+s2}{\PYZdq{}}\PY{l+s+s2}{modelo}\PY{l+s+s2}{\PYZdq{}}\PY{p}{,} \PY{l+s+s2}{\PYZdq{}}\PY{l+s+s2}{best\PYZus{}params}\PY{l+s+s2}{\PYZdq{}}\PY{p}{,} \PY{l+s+s2}{\PYZdq{}}\PY{l+s+s2}{cv\PYZus{}recall\PYZus{}condition\PYZus{}1\PYZus{}mean}\PY{l+s+s2}{\PYZdq{}}\PY{p}{,} \PY{l+s+s2}{\PYZdq{}}\PY{l+s+s2}{cv\PYZus{}medical\PYZus{}cost\PYZus{}mean}\PY{l+s+s2}{\PYZdq{}}\PY{p}{,} \PY{l+s+s2}{\PYZdq{}}\PY{l+s+s2}{test\PYZus{}recall\PYZus{}condition\PYZus{}1}\PY{l+s+s2}{\PYZdq{}}\PY{p}{,} \PY{l+s+s2}{\PYZdq{}}\PY{l+s+s2}{test\PYZus{}false\PYZus{}negatives}\PY{l+s+s2}{\PYZdq{}}\PY{p}{,} \PY{l+s+s2}{\PYZdq{}}\PY{l+s+s2}{test\PYZus{}false\PYZus{}positives}\PY{l+s+s2}{\PYZdq{}}\PY{p}{,} \PY{l+s+s2}{\PYZdq{}}\PY{l+s+s2}{test\PYZus{}medical\PYZus{}cost}\PY{l+s+s2}{\PYZdq{}}\PY{p}{,} \PY{l+s+s2}{\PYZdq{}}\PY{l+s+s2}{test\PYZus{}roc\PYZus{}auc}\PY{l+s+s2}{\PYZdq{}}\PY{p}{]}
\PY{n}{display}\PY{p}{(}\PY{n}{task4\PYZus{}comparison}\PY{p}{[}\PY{n}{cols}\PY{p}{]}\PY{p}{)}
\end{Verbatim}
\end{tcolorbox}

    \textbf{Comparación de modelos optimizados}

    
    
    \begin{Verbatim}[commandchars=\\\{\}]
                  slug               modelo                                        best\_params  cv\_recall\_condition\_1\_mean  cv\_recall\_condition\_1\_std  \textbackslash{}
0  logistic\_regression  Regresion Logistica  \{'model\_\_C': 1.0, 'model\_\_class\_weight': 'bala{\ldots}                    0.847965                   0.071986   
1                  svm              SVM RBF  \{'model\_\_C': 1.0, 'model\_\_class\_weight': 'bala{\ldots}                    0.833680                   0.072323   
2                  knn                  KNN  \{'model\_\_n\_neighbors': 7, 'model\_\_p': 1, 'mode{\ldots}                    0.795455                   0.086768   

   cv\_medical\_cost\_mean  cv\_medical\_cost\_std  cv\_roc\_auc\_mean  cv\_roc\_auc\_std  test\_accuracy  test\_precision\_condition\_1  test\_recall\_condition\_1  \textbackslash{}
0              0.421667             0.168862         0.922432        0.035147       0.783333                    0.741935                 0.821429   
1              0.460106             0.168123         0.912433        0.038102       0.800000                    0.750000                 0.857143   
2              0.537172             0.207868         0.909738        0.043283       0.833333                    0.821429                 0.821429   

   test\_f1\_condition\_1  test\_medical\_cost  test\_roc\_auc  test\_true\_negatives  test\_false\_positives  test\_false\_negatives  test\_true\_positives  \textbackslash{}
0             0.779661           0.550000      0.887277                 24.0                   8.0                   5.0                 23.0   
1             0.800000           0.466667      0.879464                 24.0                   8.0                   4.0                 24.0   
2             0.821429           0.500000      0.855469                 27.0                   5.0                   5.0                 23.0   

   task3\_cv\_recall\_condition\_1\_mean  task3\_cv\_medical\_cost\_mean  delta\_cv\_recall  delta\_cv\_medical\_cost  task3\_test\_recall\_condition\_1  \textbackslash{}
0                          0.813521                    0.493505         0.034444              -0.071838                       0.821429   
1                          0.806638                    0.517021         0.027042              -0.056915                       0.857143   
2                          0.785281                    0.572512         0.010173              -0.035340                       0.821429   

   task3\_test\_medical\_cost  delta\_test\_recall  delta\_test\_medical\_cost  
0                 0.516667                0.0                 0.033333  
1                 0.466667                0.0                 0.000000  
2                 0.516667                0.0                -0.016667  
    \end{Verbatim}

    
    \textbf{Ranking de modelos optimizados}

    
    
    \begin{Verbatim}[commandchars=\\\{\}]
                modelo  cv\_recall\_condition\_1\_mean  cv\_medical\_cost\_mean  cv\_roc\_auc\_mean  test\_recall\_condition\_1  test\_false\_negatives  \textbackslash{}
0  Regresion Logistica                    0.847965              0.421667         0.922432                 0.821429                   5.0   
1              SVM RBF                    0.833680              0.460106         0.912433                 0.857143                   4.0   
2                  KNN                    0.795455              0.537172         0.909738                 0.821429                   5.0   

   test\_false\_positives  test\_medical\_cost  delta\_cv\_recall  delta\_cv\_medical\_cost  
0                   8.0           0.550000         0.034444              -0.071838  
1                   8.0           0.466667         0.027042              -0.056915  
2                   5.0           0.500000         0.010173              -0.035340  
    \end{Verbatim}

    
    
    \begin{Verbatim}[commandchars=\\\{\}]
                modelo                                        best\_params  cv\_recall\_condition\_1\_mean  cv\_medical\_cost\_mean  test\_recall\_condition\_1  \textbackslash{}
0  Regresion Logistica  \{'model\_\_C': 1.0, 'model\_\_class\_weight': 'bala{\ldots}                    0.847965              0.421667                 0.821429   
1              SVM RBF  \{'model\_\_C': 1.0, 'model\_\_class\_weight': 'bala{\ldots}                    0.833680              0.460106                 0.857143   
2                  KNN  \{'model\_\_n\_neighbors': 7, 'model\_\_p': 1, 'mode{\ldots}                    0.795455              0.537172                 0.821429   

   test\_false\_negatives  test\_false\_positives  test\_medical\_cost  test\_roc\_auc  
0                   5.0                   8.0           0.550000      0.887277  
1                   4.0                   8.0           0.466667      0.879464  
2                   5.0                   5.0           0.500000      0.855469  
    \end{Verbatim}

    
    \subsection{4.2. Comparación visual antes y después de
optimizar}\label{comparaciuxf3n-visual-antes-y-despuuxe9s-de-optimizar}

Esta comparación responde directamente al requisito de optimización
sistemática. No se busca solo subir una métrica: se revisa si la mejora
en validación se sostiene en test y si el coste médico disminuye.

    \begin{tcolorbox}[breakable, size=fbox, boxrule=1pt, pad at break*=1mm,colback=cellbackground, colframe=cellborder]
\prompt{In}{incolor}{47}{\boxspacing}
\begin{Verbatim}[commandchars=\\\{\}]
\PY{n}{compare\PYZus{}opt} \PY{o}{=} \PY{n}{task4\PYZus{}comparison}\PY{p}{[}\PY{p}{[}
    \PY{l+s+s2}{\PYZdq{}}\PY{l+s+s2}{modelo}\PY{l+s+s2}{\PYZdq{}}\PY{p}{,}
    \PY{l+s+s2}{\PYZdq{}}\PY{l+s+s2}{task3\PYZus{}cv\PYZus{}recall\PYZus{}condition\PYZus{}1\PYZus{}mean}\PY{l+s+s2}{\PYZdq{}}\PY{p}{,}
    \PY{l+s+s2}{\PYZdq{}}\PY{l+s+s2}{cv\PYZus{}recall\PYZus{}condition\PYZus{}1\PYZus{}mean}\PY{l+s+s2}{\PYZdq{}}\PY{p}{,}
    \PY{l+s+s2}{\PYZdq{}}\PY{l+s+s2}{task3\PYZus{}test\PYZus{}recall\PYZus{}condition\PYZus{}1}\PY{l+s+s2}{\PYZdq{}}\PY{p}{,}
    \PY{l+s+s2}{\PYZdq{}}\PY{l+s+s2}{test\PYZus{}recall\PYZus{}condition\PYZus{}1}\PY{l+s+s2}{\PYZdq{}}\PY{p}{,}
    \PY{l+s+s2}{\PYZdq{}}\PY{l+s+s2}{task3\PYZus{}test\PYZus{}medical\PYZus{}cost}\PY{l+s+s2}{\PYZdq{}}\PY{p}{,}
    \PY{l+s+s2}{\PYZdq{}}\PY{l+s+s2}{test\PYZus{}medical\PYZus{}cost}\PY{l+s+s2}{\PYZdq{}}\PY{p}{,}
\PY{p}{]}\PY{p}{]}\PY{o}{.}\PY{n}{copy}\PY{p}{(}\PY{p}{)}
\PY{n}{show\PYZus{}table}\PY{p}{(}\PY{l+s+s2}{\PYZdq{}}\PY{l+s+s2}{Antes vs después de optimización}\PY{l+s+s2}{\PYZdq{}}\PY{p}{,} \PY{n}{compare\PYZus{}opt}\PY{p}{,} \PY{n}{n}\PY{o}{=}\PY{k+kc}{None}\PY{p}{)}

\PY{n}{fig}\PY{p}{,} \PY{n}{axes} \PY{o}{=} \PY{n}{plt}\PY{o}{.}\PY{n}{subplots}\PY{p}{(}\PY{l+m+mi}{1}\PY{p}{,} \PY{l+m+mi}{2}\PY{p}{,} \PY{n}{figsize}\PY{o}{=}\PY{p}{(}\PY{l+m+mi}{15}\PY{p}{,} \PY{l+m+mi}{5}\PY{p}{)}\PY{p}{)}
\PY{n}{compare\PYZus{}opt}\PY{o}{.}\PY{n}{set\PYZus{}index}\PY{p}{(}\PY{l+s+s2}{\PYZdq{}}\PY{l+s+s2}{modelo}\PY{l+s+s2}{\PYZdq{}}\PY{p}{)}\PY{p}{[}\PY{p}{[}\PY{l+s+s2}{\PYZdq{}}\PY{l+s+s2}{task3\PYZus{}cv\PYZus{}recall\PYZus{}condition\PYZus{}1\PYZus{}mean}\PY{l+s+s2}{\PYZdq{}}\PY{p}{,} \PY{l+s+s2}{\PYZdq{}}\PY{l+s+s2}{cv\PYZus{}recall\PYZus{}condition\PYZus{}1\PYZus{}mean}\PY{l+s+s2}{\PYZdq{}}\PY{p}{]}\PY{p}{]}\PY{o}{.}\PY{n}{plot}\PY{p}{(}\PY{n}{kind}\PY{o}{=}\PY{l+s+s2}{\PYZdq{}}\PY{l+s+s2}{bar}\PY{l+s+s2}{\PYZdq{}}\PY{p}{,} \PY{n}{ax}\PY{o}{=}\PY{n}{axes}\PY{p}{[}\PY{l+m+mi}{0}\PY{p}{]}\PY{p}{,} \PY{n}{color}\PY{o}{=}\PY{p}{[}\PY{l+s+s2}{\PYZdq{}}\PY{l+s+s2}{\PYZsh{}72B7B2}\PY{l+s+s2}{\PYZdq{}}\PY{p}{,} \PY{l+s+s2}{\PYZdq{}}\PY{l+s+s2}{\PYZsh{}4C78A8}\PY{l+s+s2}{\PYZdq{}}\PY{p}{]}\PY{p}{)}
\PY{n}{axes}\PY{p}{[}\PY{l+m+mi}{0}\PY{p}{]}\PY{o}{.}\PY{n}{set\PYZus{}title}\PY{p}{(}\PY{l+s+s2}{\PYZdq{}}\PY{l+s+s2}{Recall CV: base vs optimizado}\PY{l+s+s2}{\PYZdq{}}\PY{p}{)}
\PY{n}{axes}\PY{p}{[}\PY{l+m+mi}{0}\PY{p}{]}\PY{o}{.}\PY{n}{set\PYZus{}ylim}\PY{p}{(}\PY{l+m+mi}{0}\PY{p}{,} \PY{l+m+mf}{1.05}\PY{p}{)}
\PY{n}{axes}\PY{p}{[}\PY{l+m+mi}{0}\PY{p}{]}\PY{o}{.}\PY{n}{tick\PYZus{}params}\PY{p}{(}\PY{n}{axis}\PY{o}{=}\PY{l+s+s2}{\PYZdq{}}\PY{l+s+s2}{x}\PY{l+s+s2}{\PYZdq{}}\PY{p}{,} \PY{n}{rotation}\PY{o}{=}\PY{l+m+mi}{20}\PY{p}{)}
\PY{n}{compare\PYZus{}opt}\PY{o}{.}\PY{n}{set\PYZus{}index}\PY{p}{(}\PY{l+s+s2}{\PYZdq{}}\PY{l+s+s2}{modelo}\PY{l+s+s2}{\PYZdq{}}\PY{p}{)}\PY{p}{[}\PY{p}{[}\PY{l+s+s2}{\PYZdq{}}\PY{l+s+s2}{task3\PYZus{}test\PYZus{}medical\PYZus{}cost}\PY{l+s+s2}{\PYZdq{}}\PY{p}{,} \PY{l+s+s2}{\PYZdq{}}\PY{l+s+s2}{test\PYZus{}medical\PYZus{}cost}\PY{l+s+s2}{\PYZdq{}}\PY{p}{]}\PY{p}{]}\PY{o}{.}\PY{n}{plot}\PY{p}{(}\PY{n}{kind}\PY{o}{=}\PY{l+s+s2}{\PYZdq{}}\PY{l+s+s2}{bar}\PY{l+s+s2}{\PYZdq{}}\PY{p}{,} \PY{n}{ax}\PY{o}{=}\PY{n}{axes}\PY{p}{[}\PY{l+m+mi}{1}\PY{p}{]}\PY{p}{,} \PY{n}{color}\PY{o}{=}\PY{p}{[}\PY{l+s+s2}{\PYZdq{}}\PY{l+s+s2}{\PYZsh{}FFBE7D}\PY{l+s+s2}{\PYZdq{}}\PY{p}{,} \PY{l+s+s2}{\PYZdq{}}\PY{l+s+s2}{\PYZsh{}F58518}\PY{l+s+s2}{\PYZdq{}}\PY{p}{]}\PY{p}{)}
\PY{n}{axes}\PY{p}{[}\PY{l+m+mi}{1}\PY{p}{]}\PY{o}{.}\PY{n}{set\PYZus{}title}\PY{p}{(}\PY{l+s+s2}{\PYZdq{}}\PY{l+s+s2}{Coste médico test: base vs optimizado}\PY{l+s+s2}{\PYZdq{}}\PY{p}{)}
\PY{n}{axes}\PY{p}{[}\PY{l+m+mi}{1}\PY{p}{]}\PY{o}{.}\PY{n}{tick\PYZus{}params}\PY{p}{(}\PY{n}{axis}\PY{o}{=}\PY{l+s+s2}{\PYZdq{}}\PY{l+s+s2}{x}\PY{l+s+s2}{\PYZdq{}}\PY{p}{,} \PY{n}{rotation}\PY{o}{=}\PY{l+m+mi}{20}\PY{p}{)}
\PY{k}{for} \PY{n}{ax} \PY{o+ow}{in} \PY{n}{axes}\PY{p}{:}
    \PY{n}{ax}\PY{o}{.}\PY{n}{grid}\PY{p}{(}\PY{n}{axis}\PY{o}{=}\PY{l+s+s2}{\PYZdq{}}\PY{l+s+s2}{y}\PY{l+s+s2}{\PYZdq{}}\PY{p}{,} \PY{n}{alpha}\PY{o}{=}\PY{l+m+mf}{0.25}\PY{p}{)}
\PY{n}{plt}\PY{o}{.}\PY{n}{tight\PYZus{}layout}\PY{p}{(}\PY{p}{)}
\PY{n}{plt}\PY{o}{.}\PY{n}{show}\PY{p}{(}\PY{p}{)}
\end{Verbatim}
\end{tcolorbox}

    \textbf{Antes vs después de optimización}

    
    
    \begin{Verbatim}[commandchars=\\\{\}]
                modelo  task3\_cv\_recall\_condition\_1\_mean  cv\_recall\_condition\_1\_mean  task3\_test\_recall\_condition\_1  test\_recall\_condition\_1  \textbackslash{}
0  Regresion Logistica                          0.813521                    0.847965                       0.821429                 0.821429   
1              SVM RBF                          0.806638                    0.833680                       0.857143                 0.857143   
2                  KNN                          0.785281                    0.795455                       0.821429                 0.821429   

   task3\_test\_medical\_cost  test\_medical\_cost  
0                 0.516667           0.550000  
1                 0.466667           0.466667  
2                 0.516667           0.500000  
    \end{Verbatim}

    
    \subsection{4.3. Interpretación crítica del SVM optimizado
corregido}\label{interpretaciuxf3n-cruxedtica-del-svm-optimizado-corregido}

Una regla basada únicamente en maximizar \texttt{recall\_condition\_1}
puede seleccionar configuraciones degeneradas. En este caso, el
candidato \texttt{C=0.1}, \texttt{class\_weight=balanced},
\texttt{gamma=1.0} lograba recall CV perfecto, pero tenía mayor coste
médico y una frontera poco útil para generalizar.

Se corrigió la selección para usar como criterio principal el coste
médico normalizado (FN=5, FP=1), con recall y ROC-AUC como desempates.
Esta regla sigue priorizando falsos negativos, pero evita modelos que
consiguen recall alto a costa de demasiadas alarmas falsas o
inestabilidad. Con la corrección, SVM selecciona \texttt{C=1.0},
\texttt{class\_weight=balanced}, \texttt{gamma=scale} y mantiene un
comportamiento razonable en test.

    \begin{tcolorbox}[breakable, size=fbox, boxrule=1pt, pad at break*=1mm,colback=cellbackground, colframe=cellborder]
\prompt{In}{incolor}{38}{\boxspacing}
\begin{Verbatim}[commandchars=\\\{\}]
\PY{n}{svm\PYZus{}opt} \PY{o}{=} \PY{n}{task4\PYZus{}comparison}\PY{o}{.}\PY{n}{loc}\PY{p}{[}\PY{n}{task4\PYZus{}comparison}\PY{p}{[}\PY{l+s+s2}{\PYZdq{}}\PY{l+s+s2}{slug}\PY{l+s+s2}{\PYZdq{}}\PY{p}{]} \PY{o}{==} \PY{l+s+s2}{\PYZdq{}}\PY{l+s+s2}{svm}\PY{l+s+s2}{\PYZdq{}}\PY{p}{]}\PY{o}{.}\PY{n}{iloc}\PY{p}{[}\PY{l+m+mi}{0}\PY{p}{]}
\PY{n}{svm\PYZus{}review} \PY{o}{=} \PY{n}{pd}\PY{o}{.}\PY{n}{DataFrame}\PY{p}{(}\PY{p}{[}
    \PY{p}{\PYZob{}}\PY{l+s+s2}{\PYZdq{}}\PY{l+s+s2}{métrica}\PY{l+s+s2}{\PYZdq{}}\PY{p}{:} \PY{l+s+s2}{\PYZdq{}}\PY{l+s+s2}{Recall CV}\PY{l+s+s2}{\PYZdq{}}\PY{p}{,} \PY{l+s+s2}{\PYZdq{}}\PY{l+s+s2}{valor}\PY{l+s+s2}{\PYZdq{}}\PY{p}{:} \PY{n}{svm\PYZus{}opt}\PY{p}{[}\PY{l+s+s2}{\PYZdq{}}\PY{l+s+s2}{cv\PYZus{}recall\PYZus{}condition\PYZus{}1\PYZus{}mean}\PY{l+s+s2}{\PYZdq{}}\PY{p}{]}\PY{p}{,} \PY{l+s+s2}{\PYZdq{}}\PY{l+s+s2}{lectura}\PY{l+s+s2}{\PYZdq{}}\PY{p}{:} \PY{l+s+s2}{\PYZdq{}}\PY{l+s+s2}{Buen recall medio sin ser una solución degenerada de recall perfecto.}\PY{l+s+s2}{\PYZdq{}}\PY{p}{\PYZcb{}}\PY{p}{,}
    \PY{p}{\PYZob{}}\PY{l+s+s2}{\PYZdq{}}\PY{l+s+s2}{métrica}\PY{l+s+s2}{\PYZdq{}}\PY{p}{:} \PY{l+s+s2}{\PYZdq{}}\PY{l+s+s2}{Coste médico CV}\PY{l+s+s2}{\PYZdq{}}\PY{p}{,} \PY{l+s+s2}{\PYZdq{}}\PY{l+s+s2}{valor}\PY{l+s+s2}{\PYZdq{}}\PY{p}{:} \PY{n}{svm\PYZus{}opt}\PY{p}{[}\PY{l+s+s2}{\PYZdq{}}\PY{l+s+s2}{cv\PYZus{}medical\PYZus{}cost\PYZus{}mean}\PY{l+s+s2}{\PYZdq{}}\PY{p}{]}\PY{p}{,} \PY{l+s+s2}{\PYZdq{}}\PY{l+s+s2}{lectura}\PY{l+s+s2}{\PYZdq{}}\PY{p}{:} \PY{l+s+s2}{\PYZdq{}}\PY{l+s+s2}{Mejor que el SVM base y coherente con el criterio de selección corregido.}\PY{l+s+s2}{\PYZdq{}}\PY{p}{\PYZcb{}}\PY{p}{,}
    \PY{p}{\PYZob{}}\PY{l+s+s2}{\PYZdq{}}\PY{l+s+s2}{métrica}\PY{l+s+s2}{\PYZdq{}}\PY{p}{:} \PY{l+s+s2}{\PYZdq{}}\PY{l+s+s2}{Recall test}\PY{l+s+s2}{\PYZdq{}}\PY{p}{,} \PY{l+s+s2}{\PYZdq{}}\PY{l+s+s2}{valor}\PY{l+s+s2}{\PYZdq{}}\PY{p}{:} \PY{n}{svm\PYZus{}opt}\PY{p}{[}\PY{l+s+s2}{\PYZdq{}}\PY{l+s+s2}{test\PYZus{}recall\PYZus{}condition\PYZus{}1}\PY{l+s+s2}{\PYZdq{}}\PY{p}{]}\PY{p}{,} \PY{l+s+s2}{\PYZdq{}}\PY{l+s+s2}{lectura}\PY{l+s+s2}{\PYZdq{}}\PY{p}{:} \PY{l+s+s2}{\PYZdq{}}\PY{l+s+s2}{Detecta 24 de 28 positivos en la partición reservada.}\PY{l+s+s2}{\PYZdq{}}\PY{p}{\PYZcb{}}\PY{p}{,}
    \PY{p}{\PYZob{}}\PY{l+s+s2}{\PYZdq{}}\PY{l+s+s2}{métrica}\PY{l+s+s2}{\PYZdq{}}\PY{p}{:} \PY{l+s+s2}{\PYZdq{}}\PY{l+s+s2}{Falsos negativos test}\PY{l+s+s2}{\PYZdq{}}\PY{p}{,} \PY{l+s+s2}{\PYZdq{}}\PY{l+s+s2}{valor}\PY{l+s+s2}{\PYZdq{}}\PY{p}{:} \PY{n}{svm\PYZus{}opt}\PY{p}{[}\PY{l+s+s2}{\PYZdq{}}\PY{l+s+s2}{test\PYZus{}false\PYZus{}negatives}\PY{l+s+s2}{\PYZdq{}}\PY{p}{]}\PY{p}{,} \PY{l+s+s2}{\PYZdq{}}\PY{l+s+s2}{lectura}\PY{l+s+s2}{\PYZdq{}}\PY{p}{:} \PY{l+s+s2}{\PYZdq{}}\PY{l+s+s2}{Quedan 4 falsos negativos; sigue siendo el punto clínico más importante a vigilar.}\PY{l+s+s2}{\PYZdq{}}\PY{p}{\PYZcb{}}\PY{p}{,}
    \PY{p}{\PYZob{}}\PY{l+s+s2}{\PYZdq{}}\PY{l+s+s2}{métrica}\PY{l+s+s2}{\PYZdq{}}\PY{p}{:} \PY{l+s+s2}{\PYZdq{}}\PY{l+s+s2}{Coste médico test}\PY{l+s+s2}{\PYZdq{}}\PY{p}{,} \PY{l+s+s2}{\PYZdq{}}\PY{l+s+s2}{valor}\PY{l+s+s2}{\PYZdq{}}\PY{p}{:} \PY{n}{svm\PYZus{}opt}\PY{p}{[}\PY{l+s+s2}{\PYZdq{}}\PY{l+s+s2}{test\PYZus{}medical\PYZus{}cost}\PY{l+s+s2}{\PYZdq{}}\PY{p}{]}\PY{p}{,} \PY{l+s+s2}{\PYZdq{}}\PY{l+s+s2}{lectura}\PY{l+s+s2}{\PYZdq{}}\PY{p}{:} \PY{l+s+s2}{\PYZdq{}}\PY{l+s+s2}{Es el menor coste de test entre los modelos optimizados actuales.}\PY{l+s+s2}{\PYZdq{}}\PY{p}{\PYZcb{}}\PY{p}{,}
\PY{p}{]}\PY{p}{)}
\PY{n}{display}\PY{p}{(}\PY{n}{svm\PYZus{}review}\PY{p}{)}
\PY{n}{show\PYZus{}image}\PY{p}{(}\PY{n}{OUTPUTS\PYZus{}DIR} \PY{o}{/} \PY{l+s+s2}{\PYZdq{}}\PY{l+s+s2}{task4}\PY{l+s+s2}{\PYZdq{}} \PY{o}{/} \PY{l+s+s2}{\PYZdq{}}\PY{l+s+s2}{svm}\PY{l+s+s2}{\PYZdq{}} \PY{o}{/} \PY{l+s+s2}{\PYZdq{}}\PY{l+s+s2}{test\PYZus{}matriz\PYZus{}confusion.png}\PY{l+s+s2}{\PYZdq{}}\PY{p}{,} \PY{n}{width}\PY{o}{=}\PY{l+m+mi}{550}\PY{p}{)}
\PY{n}{show\PYZus{}image}\PY{p}{(}\PY{n}{OUTPUTS\PYZus{}DIR} \PY{o}{/} \PY{l+s+s2}{\PYZdq{}}\PY{l+s+s2}{task4}\PY{l+s+s2}{\PYZdq{}} \PY{o}{/} \PY{l+s+s2}{\PYZdq{}}\PY{l+s+s2}{svm}\PY{l+s+s2}{\PYZdq{}} \PY{o}{/} \PY{l+s+s2}{\PYZdq{}}\PY{l+s+s2}{test\PYZus{}roc\PYZus{}curve.png}\PY{l+s+s2}{\PYZdq{}}\PY{p}{,} \PY{n}{width}\PY{o}{=}\PY{l+m+mi}{550}\PY{p}{)}
\end{Verbatim}
\end{tcolorbox}

    
    \begin{Verbatim}[commandchars=\\\{\}]
                 métrica     valor                                            lectura
0              Recall CV  0.833680  Buen recall medio sin ser una solución degener{\ldots}
1        Coste médico CV  0.460106  Mejor que el SVM base y coherente con el crite{\ldots}
2            Recall test  0.857143  Detecta 24 de 28 positivos en la partición res{\ldots}
3  Falsos negativos test  4.000000  Quedan 4 falsos negativos; sigue siendo el pun{\ldots}
4      Coste médico test  0.466667  Es el menor coste de test entre los modelos op{\ldots}
    \end{Verbatim}

    
    \begin{center}
    \adjustimage{max size={0.9\linewidth}{0.9\paperheight}}{REPORT_files/REPORT_66_1.png}
    \end{center}
    { \hspace*{\fill} \\}
    
    \begin{center}
    \adjustimage{max size={0.9\linewidth}{0.9\paperheight}}{REPORT_files/REPORT_66_2.png}
    \end{center}
    { \hspace*{\fill} \\}
    
    \subsection{4.4. Variables influyentes}\label{variables-influyentes}

El requisito final de la optimización pide identificar variables con
mayor peso. Aquí se combinan importancias por permutación sobre recall
y, para Regresión Logística, coeficientes. La importancia por
permutación mide cuánto cae el recall al desordenar una variable: si cae
mucho, esa variable era útil para detectar \texttt{condition=1}.

    \begin{tcolorbox}[breakable, size=fbox, boxrule=1pt, pad at break*=1mm,colback=cellbackground, colframe=cellborder]
\prompt{In}{incolor}{39}{\boxspacing}
\begin{Verbatim}[commandchars=\\\{\}]
\PY{n}{global\PYZus{}importance} \PY{o}{=} \PY{n}{read\PYZus{}csv}\PY{p}{(}\PY{n}{OUTPUTS\PYZus{}DIR} \PY{o}{/} \PY{l+s+s2}{\PYZdq{}}\PY{l+s+s2}{task4}\PY{l+s+s2}{\PYZdq{}} \PY{o}{/} \PY{l+s+s2}{\PYZdq{}}\PY{l+s+s2}{variables\PYZus{}influyentes\PYZus{}global.csv}\PY{l+s+s2}{\PYZdq{}}\PY{p}{)}
\PY{n}{importance\PYZus{}column} \PY{o}{=} \PY{l+s+s2}{\PYZdq{}}\PY{l+s+s2}{importance\PYZus{}normalized\PYZus{}mean}\PY{l+s+s2}{\PYZdq{}} \PY{k}{if} \PY{l+s+s2}{\PYZdq{}}\PY{l+s+s2}{importance\PYZus{}normalized\PYZus{}mean}\PY{l+s+s2}{\PYZdq{}} \PY{o+ow}{in} \PY{n}{global\PYZus{}importance}\PY{o}{.}\PY{n}{columns} \PY{k}{else} \PY{l+s+s2}{\PYZdq{}}\PY{l+s+s2}{importance\PYZus{}mean}\PY{l+s+s2}{\PYZdq{}}
\PY{n}{show\PYZus{}table}\PY{p}{(}\PY{l+s+s2}{\PYZdq{}}\PY{l+s+s2}{Variables influyentes globales}\PY{l+s+s2}{\PYZdq{}}\PY{p}{,} \PY{n}{global\PYZus{}importance}\PY{p}{,} \PY{n}{n}\PY{o}{=}\PY{k+kc}{None}\PY{p}{)}

\PY{n}{fig}\PY{p}{,} \PY{n}{ax} \PY{o}{=} \PY{n}{plt}\PY{o}{.}\PY{n}{subplots}\PY{p}{(}\PY{n}{figsize}\PY{o}{=}\PY{p}{(}\PY{l+m+mi}{9}\PY{p}{,} \PY{l+m+mi}{6}\PY{p}{)}\PY{p}{)}
\PY{n}{plot\PYZus{}importance} \PY{o}{=} \PY{n}{global\PYZus{}importance}\PY{o}{.}\PY{n}{sort\PYZus{}values}\PY{p}{(}\PY{n}{importance\PYZus{}column}\PY{p}{,} \PY{n}{ascending}\PY{o}{=}\PY{k+kc}{True}\PY{p}{)}
\PY{n}{ax}\PY{o}{.}\PY{n}{barh}\PY{p}{(}\PY{n}{plot\PYZus{}importance}\PY{p}{[}\PY{l+s+s2}{\PYZdq{}}\PY{l+s+s2}{feature}\PY{l+s+s2}{\PYZdq{}}\PY{p}{]}\PY{p}{,} \PY{n}{plot\PYZus{}importance}\PY{p}{[}\PY{n}{importance\PYZus{}column}\PY{p}{]}\PY{p}{,} \PY{n}{color}\PY{o}{=}\PY{l+s+s2}{\PYZdq{}}\PY{l+s+s2}{\PYZsh{}4C78A8}\PY{l+s+s2}{\PYZdq{}}\PY{p}{)}
\PY{n}{ax}\PY{o}{.}\PY{n}{set\PYZus{}title}\PY{p}{(}\PY{l+s+s2}{\PYZdq{}}\PY{l+s+s2}{Influencia global por permutation importance}\PY{l+s+s2}{\PYZdq{}}\PY{p}{)}
\PY{n}{ax}\PY{o}{.}\PY{n}{set\PYZus{}xlabel}\PY{p}{(}\PY{l+s+s2}{\PYZdq{}}\PY{l+s+s2}{Importancia normalizada media}\PY{l+s+s2}{\PYZdq{}}\PY{p}{)}
\PY{n}{ax}\PY{o}{.}\PY{n}{grid}\PY{p}{(}\PY{n}{axis}\PY{o}{=}\PY{l+s+s2}{\PYZdq{}}\PY{l+s+s2}{x}\PY{l+s+s2}{\PYZdq{}}\PY{p}{,} \PY{n}{alpha}\PY{o}{=}\PY{l+m+mf}{0.25}\PY{p}{)}
\PY{n}{plt}\PY{o}{.}\PY{n}{tight\PYZus{}layout}\PY{p}{(}\PY{p}{)}
\PY{n}{plt}\PY{o}{.}\PY{n}{show}\PY{p}{(}\PY{p}{)}

\PY{k}{for} \PY{n}{slug}\PY{p}{,} \PY{n}{name} \PY{o+ow}{in} \PY{p}{[}\PY{p}{(}\PY{l+s+s2}{\PYZdq{}}\PY{l+s+s2}{logistic\PYZus{}regression}\PY{l+s+s2}{\PYZdq{}}\PY{p}{,} \PY{l+s+s2}{\PYZdq{}}\PY{l+s+s2}{Regresión Logística}\PY{l+s+s2}{\PYZdq{}}\PY{p}{)}\PY{p}{,} \PY{p}{(}\PY{l+s+s2}{\PYZdq{}}\PY{l+s+s2}{knn}\PY{l+s+s2}{\PYZdq{}}\PY{p}{,} \PY{l+s+s2}{\PYZdq{}}\PY{l+s+s2}{KNN}\PY{l+s+s2}{\PYZdq{}}\PY{p}{)}\PY{p}{,} \PY{p}{(}\PY{l+s+s2}{\PYZdq{}}\PY{l+s+s2}{svm}\PY{l+s+s2}{\PYZdq{}}\PY{p}{,} \PY{l+s+s2}{\PYZdq{}}\PY{l+s+s2}{SVM}\PY{l+s+s2}{\PYZdq{}}\PY{p}{)}\PY{p}{]}\PY{p}{:}
    \PY{n}{display}\PY{p}{(}\PY{n}{Markdown}\PY{p}{(}\PY{l+s+sa}{f}\PY{l+s+s2}{\PYZdq{}}\PY{l+s+s2}{\PYZsh{}\PYZsh{}\PYZsh{} }\PY{l+s+si}{\PYZob{}}\PY{n}{name}\PY{l+s+si}{\PYZcb{}}\PY{l+s+s2}{: importancia por permutación}\PY{l+s+s2}{\PYZdq{}}\PY{p}{)}\PY{p}{)}
    \PY{n}{show\PYZus{}table}\PY{p}{(}\PY{l+s+sa}{f}\PY{l+s+s2}{\PYZdq{}}\PY{l+s+s2}{Importancias }\PY{l+s+si}{\PYZob{}}\PY{n}{name}\PY{l+s+si}{\PYZcb{}}\PY{l+s+s2}{\PYZdq{}}\PY{p}{,} \PY{n}{read\PYZus{}csv}\PY{p}{(}\PY{n}{OUTPUTS\PYZus{}DIR} \PY{o}{/} \PY{l+s+s2}{\PYZdq{}}\PY{l+s+s2}{task4}\PY{l+s+s2}{\PYZdq{}} \PY{o}{/} \PY{n}{slug} \PY{o}{/} \PY{l+s+s2}{\PYZdq{}}\PY{l+s+s2}{feature\PYZus{}importance.csv}\PY{l+s+s2}{\PYZdq{}}\PY{p}{)}\PY{p}{,} \PY{n}{n}\PY{o}{=}\PY{k+kc}{None}\PY{p}{)}
    \PY{n}{show\PYZus{}image}\PY{p}{(}\PY{n}{OUTPUTS\PYZus{}DIR} \PY{o}{/} \PY{l+s+s2}{\PYZdq{}}\PY{l+s+s2}{task4}\PY{l+s+s2}{\PYZdq{}} \PY{o}{/} \PY{n}{slug} \PY{o}{/} \PY{l+s+s2}{\PYZdq{}}\PY{l+s+s2}{feature\PYZus{}importance.png}\PY{l+s+s2}{\PYZdq{}}\PY{p}{,} \PY{n}{width}\PY{o}{=}\PY{l+m+mi}{750}\PY{p}{)}

\PY{k}{if} \PY{p}{(}\PY{n}{OUTPUTS\PYZus{}DIR} \PY{o}{/} \PY{l+s+s2}{\PYZdq{}}\PY{l+s+s2}{task4}\PY{l+s+s2}{\PYZdq{}} \PY{o}{/} \PY{l+s+s2}{\PYZdq{}}\PY{l+s+s2}{logistic\PYZus{}regression}\PY{l+s+s2}{\PYZdq{}} \PY{o}{/} \PY{l+s+s2}{\PYZdq{}}\PY{l+s+s2}{logistic\PYZus{}coefficients.csv}\PY{l+s+s2}{\PYZdq{}}\PY{p}{)}\PY{o}{.}\PY{n}{exists}\PY{p}{(}\PY{p}{)}\PY{p}{:}
    \PY{n}{display}\PY{p}{(}\PY{n}{Markdown}\PY{p}{(}\PY{l+s+s2}{\PYZdq{}}\PY{l+s+s2}{\PYZsh{}\PYZsh{}\PYZsh{} Coeficientes de Regresión Logística}\PY{l+s+s2}{\PYZdq{}}\PY{p}{)}\PY{p}{)}
    \PY{n}{display}\PY{p}{(}\PY{n}{read\PYZus{}csv}\PY{p}{(}\PY{n}{OUTPUTS\PYZus{}DIR} \PY{o}{/} \PY{l+s+s2}{\PYZdq{}}\PY{l+s+s2}{task4}\PY{l+s+s2}{\PYZdq{}} \PY{o}{/} \PY{l+s+s2}{\PYZdq{}}\PY{l+s+s2}{logistic\PYZus{}regression}\PY{l+s+s2}{\PYZdq{}} \PY{o}{/} \PY{l+s+s2}{\PYZdq{}}\PY{l+s+s2}{logistic\PYZus{}coefficients.csv}\PY{l+s+s2}{\PYZdq{}}\PY{p}{)}\PY{p}{)}
    \PY{n}{show\PYZus{}image}\PY{p}{(}\PY{n}{OUTPUTS\PYZus{}DIR} \PY{o}{/} \PY{l+s+s2}{\PYZdq{}}\PY{l+s+s2}{task4}\PY{l+s+s2}{\PYZdq{}} \PY{o}{/} \PY{l+s+s2}{\PYZdq{}}\PY{l+s+s2}{logistic\PYZus{}regression}\PY{l+s+s2}{\PYZdq{}} \PY{o}{/} \PY{l+s+s2}{\PYZdq{}}\PY{l+s+s2}{logistic\PYZus{}coefficients.png}\PY{l+s+s2}{\PYZdq{}}\PY{p}{,} \PY{n}{width}\PY{o}{=}\PY{l+m+mi}{750}\PY{p}{)}
\end{Verbatim}
\end{tcolorbox}

    \textbf{Variables influyentes globales}

    
    
    \begin{Verbatim}[commandchars=\\\{\}]
       feature  importance\_normalized\_mean  importance\_mean\_avg  models\_with\_positive\_importance
0    feature\_8                    0.942029             0.057937                                3
1    feature\_7                    0.378601             0.022619                                2
2   feature\_12                    0.355967             0.022222                                2
3    feature\_5                    0.330649             0.022222                                2
4    feature\_4                    0.215870             0.015079                                2
5    feature\_3                    0.189211             0.003968                                2
6    feature\_2                    0.189032             0.012698                                2
7   feature\_13                    0.094203             0.001587                                1
8   feature\_11                    0.043478            -0.006349                                1
9   feature\_10                    0.020576             0.000000                                1
10   feature\_1                    0.000000            -0.002381                                0
11   feature\_6                    0.000000            -0.006349                                0
12   feature\_9                    0.000000             0.000000                                0
    \end{Verbatim}

    
    \subsubsection{Regresión Logística: importancia por
permutación}\label{regresiuxf3n-loguxedstica-importancia-por-permutaciuxf3n}

    
    \textbf{Importancias Regresión Logística}

    
    
    \begin{Verbatim}[commandchars=\\\{\}]
       feature  importance\_mean  importance\_std
0    feature\_8         0.032143        0.049099
1    feature\_5         0.000000        0.000000
2    feature\_4         0.000000        0.000000
3    feature\_6         0.000000        0.000000
4    feature\_1         0.000000        0.000000
5    feature\_7         0.000000        0.000000
6    feature\_9         0.000000        0.000000
7    feature\_2        -0.001190        0.040703
8   feature\_13        -0.003571        0.014136
9   feature\_10        -0.005952        0.016192
10  feature\_11        -0.010714        0.016366
11  feature\_12        -0.015476        0.017698
12   feature\_3        -0.023810        0.031044
    \end{Verbatim}

    
    \begin{center}
    \adjustimage{max size={0.9\linewidth}{0.9\paperheight}}{REPORT_files/REPORT_68_5.png}
    \end{center}
    { \hspace*{\fill} \\}
    
    \subsubsection{KNN: importancia por
permutación}\label{knn-importancia-por-permutaciuxf3n}

    
    \textbf{Importancias KNN}

    
    
    \begin{Verbatim}[commandchars=\\\{\}]
       feature  importance\_mean  importance\_std
0    feature\_8         0.096429        0.056205
1   feature\_12         0.054762        0.040963
2    feature\_5         0.028571        0.025085
3    feature\_4         0.022619        0.017211
4    feature\_2         0.019048        0.037721
5    feature\_7         0.013095        0.017211
6    feature\_3         0.010714        0.032143
7   feature\_10         0.005952        0.029282
8    feature\_9         0.000000        0.000000
9    feature\_1        -0.007143        0.014286
10  feature\_13        -0.007143        0.026726
11  feature\_11        -0.015476        0.021984
12   feature\_6        -0.019048        0.023928
    \end{Verbatim}

    
    \begin{center}
    \adjustimage{max size={0.9\linewidth}{0.9\paperheight}}{REPORT_files/REPORT_68_9.png}
    \end{center}
    { \hspace*{\fill} \\}
    
    \subsubsection{SVM: importancia por
permutación}\label{svm-importancia-por-permutaciuxf3n}

    
    \textbf{Importancias SVM}

    
    
    \begin{Verbatim}[commandchars=\\\{\}]
       feature  importance\_mean  importance\_std
0    feature\_7         0.054762        0.028769
1    feature\_8         0.045238        0.034421
2    feature\_5         0.038095        0.022462
3   feature\_12         0.027381        0.046733
4    feature\_3         0.025000        0.032143
5    feature\_4         0.022619        0.017211
6    feature\_2         0.020238        0.040912
7   feature\_13         0.015476        0.021984
8   feature\_11         0.007143        0.017003
9    feature\_1         0.000000        0.000000
10   feature\_6         0.000000        0.000000
11   feature\_9         0.000000        0.000000
12  feature\_10         0.000000        0.000000
    \end{Verbatim}

    
    \begin{center}
    \adjustimage{max size={0.9\linewidth}{0.9\paperheight}}{REPORT_files/REPORT_68_13.png}
    \end{center}
    { \hspace*{\fill} \\}
    
    \subsubsection{Coeficientes de Regresión
Logística}\label{coeficientes-de-regresiuxf3n-loguxedstica}

    
    
    \begin{Verbatim}[commandchars=\\\{\}]
       feature  mean\_abs\_coefficient  max\_abs\_coefficient  signed\_coefficient\_sum
0    feature\_7              1.093590             1.093590                1.093590
1    feature\_8              0.564063             2.256254               -2.256254
2    feature\_3              0.493578             1.240715               -1.000697
3    feature\_2              0.475102             1.900407                1.900407
4   feature\_12              0.449384             0.895820                0.443487
5    feature\_4              0.394855             0.394855                0.394855
6   feature\_11              0.357104             0.357104                0.357104
7   feature\_13              0.237500             0.237500                0.237500
8   feature\_10              0.208403             0.208403               -0.208403
9    feature\_6              0.114433             0.114433                0.114433
10   feature\_1              0.000000             0.000000                0.000000
11   feature\_5              0.000000             0.000000                0.000000
12   feature\_9              0.000000             0.000000                0.000000
    \end{Verbatim}

    
    \begin{center}
    \adjustimage{max size={0.9\linewidth}{0.9\paperheight}}{REPORT_files/REPORT_68_16.png}
    \end{center}
    { \hspace*{\fill} \\}
    
    \section{5. Registro de experimentos fallidos o
débiles}\label{registro-de-experimentos-fallidos-o-duxe9biles}

El enunciado pide documentar intentos que fallaron o dieron malos
resultados. Esta sección no maquilla resultados: registra qué modelos o
técnicas tienen limitaciones claras y por qué importan en un contexto
médico.

    \begin{tcolorbox}[breakable, size=fbox, boxrule=1pt, pad at break*=1mm,colback=cellbackground, colframe=cellborder]
\prompt{In}{incolor}{31}{\boxspacing}
\begin{Verbatim}[commandchars=\\\{\}]
\PY{n}{weak\PYZus{}experiments} \PY{o}{=} \PY{p}{[}\PY{p}{]}

\PY{n}{mlp\PYZus{}row} \PY{o}{=} \PY{n}{task3\PYZus{}summary}\PY{o}{.}\PY{n}{loc}\PY{p}{[}\PY{n}{task3\PYZus{}summary}\PY{p}{[}\PY{l+s+s2}{\PYZdq{}}\PY{l+s+s2}{modelo}\PY{l+s+s2}{\PYZdq{}}\PY{p}{]}\PY{o}{.}\PY{n}{str}\PY{o}{.}\PY{n}{contains}\PY{p}{(}\PY{l+s+s2}{\PYZdq{}}\PY{l+s+s2}{MLP}\PY{l+s+s2}{\PYZdq{}}\PY{p}{,} \PY{n}{case}\PY{o}{=}\PY{k+kc}{False}\PY{p}{,} \PY{n}{na}\PY{o}{=}\PY{k+kc}{False}\PY{p}{)}\PY{p}{]}\PY{o}{.}\PY{n}{iloc}\PY{p}{[}\PY{l+m+mi}{0}\PY{p}{]}
\PY{n}{weak\PYZus{}experiments}\PY{o}{.}\PY{n}{append}\PY{p}{(}\PY{p}{\PYZob{}}
    \PY{l+s+s2}{\PYZdq{}}\PY{l+s+s2}{experimento}\PY{l+s+s2}{\PYZdq{}}\PY{p}{:} \PY{l+s+s2}{\PYZdq{}}\PY{l+s+s2}{Red Neuronal MLP base}\PY{l+s+s2}{\PYZdq{}}\PY{p}{,}
    \PY{l+s+s2}{\PYZdq{}}\PY{l+s+s2}{evidencia}\PY{l+s+s2}{\PYZdq{}}\PY{p}{:} \PY{l+s+sa}{f}\PY{l+s+s2}{\PYZdq{}}\PY{l+s+s2}{Recall test=}\PY{l+s+si}{\PYZob{}}\PY{n}{mlp\PYZus{}row}\PY{p}{[}\PY{l+s+s1}{\PYZsq{}}\PY{l+s+s1}{test\PYZus{}recall\PYZus{}condition\PYZus{}1}\PY{l+s+s1}{\PYZsq{}}\PY{p}{]}\PY{l+s+si}{:}\PY{l+s+s2}{.4f}\PY{l+s+si}{\PYZcb{}}\PY{l+s+s2}{, falsos negativos=}\PY{l+s+si}{\PYZob{}}\PY{n}{mlp\PYZus{}row}\PY{p}{[}\PY{l+s+s1}{\PYZsq{}}\PY{l+s+s1}{test\PYZus{}false\PYZus{}negatives}\PY{l+s+s1}{\PYZsq{}}\PY{p}{]}\PY{l+s+si}{:}\PY{l+s+s2}{.0f}\PY{l+s+si}{\PYZcb{}}\PY{l+s+s2}{\PYZdq{}}\PY{p}{,}
    \PY{l+s+s2}{\PYZdq{}}\PY{l+s+s2}{interpretación}\PY{l+s+s2}{\PYZdq{}}\PY{p}{:} \PY{l+s+s2}{\PYZdq{}}\PY{l+s+s2}{Demasiados positivos sin detectar; no recomendable para esta configuración médica.}\PY{l+s+s2}{\PYZdq{}}\PY{p}{,}
\PY{p}{\PYZcb{}}\PY{p}{)}

\PY{n}{svm\PYZus{}grid} \PY{o}{=} \PY{n}{read\PYZus{}csv}\PY{p}{(}\PY{n}{OUTPUTS\PYZus{}DIR} \PY{o}{/} \PY{l+s+s2}{\PYZdq{}}\PY{l+s+s2}{task4}\PY{l+s+s2}{\PYZdq{}} \PY{o}{/} \PY{l+s+s2}{\PYZdq{}}\PY{l+s+s2}{svm}\PY{l+s+s2}{\PYZdq{}} \PY{o}{/} \PY{l+s+s2}{\PYZdq{}}\PY{l+s+s2}{grid\PYZus{}search\PYZus{}resultados.csv}\PY{l+s+s2}{\PYZdq{}}\PY{p}{)}
\PY{n}{svm\PYZus{}degenerate} \PY{o}{=} \PY{n}{svm\PYZus{}grid}\PY{p}{[}
    \PY{p}{(}\PY{n}{svm\PYZus{}grid}\PY{p}{[}\PY{l+s+s2}{\PYZdq{}}\PY{l+s+s2}{param\PYZus{}model\PYZus{}\PYZus{}C}\PY{l+s+s2}{\PYZdq{}}\PY{p}{]}\PY{o}{.}\PY{n}{astype}\PY{p}{(}\PY{n+nb}{float}\PY{p}{)} \PY{o}{==} \PY{l+m+mf}{0.1}\PY{p}{)}
    \PY{o}{\PYZam{}} \PY{p}{(}\PY{n}{svm\PYZus{}grid}\PY{p}{[}\PY{l+s+s2}{\PYZdq{}}\PY{l+s+s2}{param\PYZus{}model\PYZus{}\PYZus{}class\PYZus{}weight}\PY{l+s+s2}{\PYZdq{}}\PY{p}{]}\PY{o}{.}\PY{n}{fillna}\PY{p}{(}\PY{l+s+s2}{\PYZdq{}}\PY{l+s+s2}{\PYZdq{}}\PY{p}{)} \PY{o}{==} \PY{l+s+s2}{\PYZdq{}}\PY{l+s+s2}{balanced}\PY{l+s+s2}{\PYZdq{}}\PY{p}{)}
    \PY{o}{\PYZam{}} \PY{p}{(}\PY{n}{svm\PYZus{}grid}\PY{p}{[}\PY{l+s+s2}{\PYZdq{}}\PY{l+s+s2}{param\PYZus{}model\PYZus{}\PYZus{}gamma}\PY{l+s+s2}{\PYZdq{}}\PY{p}{]}\PY{o}{.}\PY{n}{astype}\PY{p}{(}\PY{n+nb}{str}\PY{p}{)} \PY{o}{==} \PY{l+s+s2}{\PYZdq{}}\PY{l+s+s2}{1.0}\PY{l+s+s2}{\PYZdq{}}\PY{p}{)}
\PY{p}{]}\PY{o}{.}\PY{n}{iloc}\PY{p}{[}\PY{l+m+mi}{0}\PY{p}{]}
\PY{n}{weak\PYZus{}experiments}\PY{o}{.}\PY{n}{append}\PY{p}{(}\PY{p}{\PYZob{}}
    \PY{l+s+s2}{\PYZdq{}}\PY{l+s+s2}{experimento}\PY{l+s+s2}{\PYZdq{}}\PY{p}{:} \PY{l+s+s2}{\PYZdq{}}\PY{l+s+s2}{Candidato SVM descartado durante optimización}\PY{l+s+s2}{\PYZdq{}}\PY{p}{,}
    \PY{l+s+s2}{\PYZdq{}}\PY{l+s+s2}{evidencia}\PY{l+s+s2}{\PYZdq{}}\PY{p}{:} \PY{l+s+sa}{f}\PY{l+s+s2}{\PYZdq{}}\PY{l+s+s2}{Recall CV=}\PY{l+s+si}{\PYZob{}}\PY{n}{svm\PYZus{}degenerate}\PY{p}{[}\PY{l+s+s1}{\PYZsq{}}\PY{l+s+s1}{mean\PYZus{}test\PYZus{}recall\PYZus{}condition\PYZus{}1}\PY{l+s+s1}{\PYZsq{}}\PY{p}{]}\PY{l+s+si}{:}\PY{l+s+s2}{.4f}\PY{l+s+si}{\PYZcb{}}\PY{l+s+s2}{, coste CV=}\PY{l+s+si}{\PYZob{}}\PY{n}{svm\PYZus{}degenerate}\PY{p}{[}\PY{l+s+s1}{\PYZsq{}}\PY{l+s+s1}{mean\PYZus{}test\PYZus{}medical\PYZus{}cost}\PY{l+s+s1}{\PYZsq{}}\PY{p}{]}\PY{l+s+si}{:}\PY{l+s+s2}{.4f}\PY{l+s+si}{\PYZcb{}}\PY{l+s+s2}{\PYZdq{}}\PY{p}{,}
    \PY{l+s+s2}{\PYZdq{}}\PY{l+s+s2}{interpretación}\PY{l+s+s2}{\PYZdq{}}\PY{p}{:} \PY{l+s+s2}{\PYZdq{}}\PY{l+s+s2}{Tenía recall perfecto en CV, pero mayor coste médico; la regla corregida lo descarta.}\PY{l+s+s2}{\PYZdq{}}\PY{p}{,}
\PY{p}{\PYZcb{}}\PY{p}{)}

\PY{n}{dbscan\PYZus{}two\PYZus{}cluster} \PY{o}{=} \PY{n}{dbscan\PYZus{}metrics}\PY{o}{.}\PY{n}{loc}\PY{p}{[}\PY{n}{dbscan\PYZus{}metrics}\PY{p}{[}\PY{l+s+s2}{\PYZdq{}}\PY{l+s+s2}{n\PYZus{}clusters}\PY{l+s+s2}{\PYZdq{}}\PY{p}{]} \PY{o}{==} \PY{l+m+mi}{2}\PY{p}{]}\PY{o}{.}\PY{n}{sort\PYZus{}values}\PY{p}{(}\PY{l+s+s2}{\PYZdq{}}\PY{l+s+s2}{n\PYZus{}noise}\PY{l+s+s2}{\PYZdq{}}\PY{p}{)}\PY{o}{.}\PY{n}{head}\PY{p}{(}\PY{l+m+mi}{1}\PY{p}{)}
\PY{n}{dbscan\PYZus{}noise} \PY{o}{=} \PY{n+nb}{int}\PY{p}{(}\PY{n}{dbscan\PYZus{}two\PYZus{}cluster}\PY{o}{.}\PY{n}{iloc}\PY{p}{[}\PY{l+m+mi}{0}\PY{p}{]}\PY{p}{[}\PY{l+s+s2}{\PYZdq{}}\PY{l+s+s2}{n\PYZus{}noise}\PY{l+s+s2}{\PYZdq{}}\PY{p}{]}\PY{p}{)} \PY{k}{if} \PY{o+ow}{not} \PY{n}{dbscan\PYZus{}two\PYZus{}cluster}\PY{o}{.}\PY{n}{empty} \PY{k}{else} \PY{l+s+s2}{\PYZdq{}}\PY{l+s+s2}{NA}\PY{l+s+s2}{\PYZdq{}}
\PY{n}{weak\PYZus{}experiments}\PY{o}{.}\PY{n}{append}\PY{p}{(}\PY{p}{\PYZob{}}
    \PY{l+s+s2}{\PYZdq{}}\PY{l+s+s2}{experimento}\PY{l+s+s2}{\PYZdq{}}\PY{p}{:} \PY{l+s+s2}{\PYZdq{}}\PY{l+s+s2}{DBSCAN}\PY{l+s+s2}{\PYZdq{}}\PY{p}{,}
    \PY{l+s+s2}{\PYZdq{}}\PY{l+s+s2}{evidencia}\PY{l+s+s2}{\PYZdq{}}\PY{p}{:} \PY{l+s+sa}{f}\PY{l+s+s2}{\PYZdq{}}\PY{l+s+s2}{Configuración con dos clusters conserva }\PY{l+s+si}{\PYZob{}}\PY{n}{dbscan\PYZus{}noise}\PY{l+s+si}{\PYZcb{}}\PY{l+s+s2}{ puntos como ruido}\PY{l+s+s2}{\PYZdq{}}\PY{p}{,}
    \PY{l+s+s2}{\PYZdq{}}\PY{l+s+s2}{interpretación}\PY{l+s+s2}{\PYZdq{}}\PY{p}{:} \PY{l+s+s2}{\PYZdq{}}\PY{l+s+s2}{Útil como contraste exploratorio, pero no ofrece una segmentación densa y clara del dataset.}\PY{l+s+s2}{\PYZdq{}}\PY{p}{,}
\PY{p}{\PYZcb{}}\PY{p}{)}

\PY{n}{weak\PYZus{}df} \PY{o}{=} \PY{n}{pd}\PY{o}{.}\PY{n}{DataFrame}\PY{p}{(}\PY{n}{weak\PYZus{}experiments}\PY{p}{)}
\PY{n}{display}\PY{p}{(}\PY{n}{weak\PYZus{}df}\PY{p}{)}
\end{Verbatim}
\end{tcolorbox}

    
    \begin{Verbatim}[commandchars=\\\{\}]
                                     experimento                                          evidencia                                     interpretación
0                          Red Neuronal MLP base            Recall test=0.2500, falsos negativos=21  Demasiados positivos sin detectar; no recomend{\ldots}
1  Candidato SVM descartado durante optimización                  Recall CV=1.0000, coste CV=0.5401  Tenía recall perfecto en CV, pero mayor coste {\ldots}
2                                         DBSCAN  Configuración con dos clusters conserva 185 pu{\ldots}  Útil como contraste exploratorio, pero no ofre{\ldots}
    \end{Verbatim}

    
    \section{6. Trabajo extra realizado}\label{trabajo-extra-realizado}

Además de los minimos del enunciado, el proyecto incluye experimentacion
adicional para fortalecer la entrega: silhouette como metrica
complementaria, t-SNE, curvas de aprendizaje, coste medico
personalizado, permutation importance, analisis de coeficientes y
comparacion explicita antes/despues de optimizacion.

    \begin{tcolorbox}[breakable, size=fbox, boxrule=1pt, pad at break*=1mm,colback=cellbackground, colframe=cellborder]
\prompt{In}{incolor}{40}{\boxspacing}
\begin{Verbatim}[commandchars=\\\{\}]
\PY{n}{extra\PYZus{}work} \PY{o}{=} \PY{n}{pd}\PY{o}{.}\PY{n}{DataFrame}\PY{p}{(}\PY{p}{[}
    \PY{p}{\PYZob{}}\PY{l+s+s2}{\PYZdq{}}\PY{l+s+s2}{trabajo extra}\PY{l+s+s2}{\PYZdq{}}\PY{p}{:} \PY{l+s+s2}{\PYZdq{}}\PY{l+s+s2}{Silhouette como metrica complementaria}\PY{l+s+s2}{\PYZdq{}}\PY{p}{,} \PY{l+s+s2}{\PYZdq{}}\PY{l+s+s2}{aporte}\PY{l+s+s2}{\PYZdq{}}\PY{p}{:} \PY{l+s+s2}{\PYZdq{}}\PY{l+s+s2}{Contrasta la separacion interna de clusters, pero no se usa como criterio principal de Task 2.}\PY{l+s+s2}{\PYZdq{}}\PY{p}{\PYZcb{}}\PY{p}{,}
    \PY{p}{\PYZob{}}\PY{l+s+s2}{\PYZdq{}}\PY{l+s+s2}{trabajo extra}\PY{l+s+s2}{\PYZdq{}}\PY{p}{:} \PY{l+s+s2}{\PYZdq{}}\PY{l+s+s2}{t\PYZhy{}SNE}\PY{l+s+s2}{\PYZdq{}}\PY{p}{,} \PY{l+s+s2}{\PYZdq{}}\PY{l+s+s2}{aporte}\PY{l+s+s2}{\PYZdq{}}\PY{p}{:} \PY{l+s+s2}{\PYZdq{}}\PY{l+s+s2}{Visualizacion no lineal de vecindarios y relacion posterior con condition/K\PYZhy{}Means.}\PY{l+s+s2}{\PYZdq{}}\PY{p}{\PYZcb{}}\PY{p}{,}
    \PY{p}{\PYZob{}}\PY{l+s+s2}{\PYZdq{}}\PY{l+s+s2}{trabajo extra}\PY{l+s+s2}{\PYZdq{}}\PY{p}{:} \PY{l+s+s2}{\PYZdq{}}\PY{l+s+s2}{Coste medico FN=5, FP=1}\PY{l+s+s2}{\PYZdq{}}\PY{p}{,} \PY{l+s+s2}{\PYZdq{}}\PY{l+s+s2}{aporte}\PY{l+s+s2}{\PYZdq{}}\PY{p}{:} \PY{l+s+s2}{\PYZdq{}}\PY{l+s+s2}{Traduce la prioridad clinica a una metrica comparable entre modelos.}\PY{l+s+s2}{\PYZdq{}}\PY{p}{\PYZcb{}}\PY{p}{,}
    \PY{p}{\PYZob{}}\PY{l+s+s2}{\PYZdq{}}\PY{l+s+s2}{trabajo extra}\PY{l+s+s2}{\PYZdq{}}\PY{p}{:} \PY{l+s+s2}{\PYZdq{}}\PY{l+s+s2}{Curvas de aprendizaje}\PY{l+s+s2}{\PYZdq{}}\PY{p}{,} \PY{l+s+s2}{\PYZdq{}}\PY{l+s+s2}{aporte}\PY{l+s+s2}{\PYZdq{}}\PY{p}{:} \PY{l+s+s2}{\PYZdq{}}\PY{l+s+s2}{Ayuda a revisar estabilidad y posible falta de generalizacion.}\PY{l+s+s2}{\PYZdq{}}\PY{p}{\PYZcb{}}\PY{p}{,}
    \PY{p}{\PYZob{}}\PY{l+s+s2}{\PYZdq{}}\PY{l+s+s2}{trabajo extra}\PY{l+s+s2}{\PYZdq{}}\PY{p}{:} \PY{l+s+s2}{\PYZdq{}}\PY{l+s+s2}{Permutation importance}\PY{l+s+s2}{\PYZdq{}}\PY{p}{,} \PY{l+s+s2}{\PYZdq{}}\PY{l+s+s2}{aporte}\PY{l+s+s2}{\PYZdq{}}\PY{p}{:} \PY{l+s+s2}{\PYZdq{}}\PY{l+s+s2}{Identifica variables que sostienen el recall de condition=1.}\PY{l+s+s2}{\PYZdq{}}\PY{p}{\PYZcb{}}\PY{p}{,}
    \PY{p}{\PYZob{}}\PY{l+s+s2}{\PYZdq{}}\PY{l+s+s2}{trabajo extra}\PY{l+s+s2}{\PYZdq{}}\PY{p}{:} \PY{l+s+s2}{\PYZdq{}}\PY{l+s+s2}{Comparacion base vs optimizado}\PY{l+s+s2}{\PYZdq{}}\PY{p}{,} \PY{l+s+s2}{\PYZdq{}}\PY{l+s+s2}{aporte}\PY{l+s+s2}{\PYZdq{}}\PY{p}{:} \PY{l+s+s2}{\PYZdq{}}\PY{l+s+s2}{Muestra si optimizar realmente mejora recall/coste o solo mueve el problema.}\PY{l+s+s2}{\PYZdq{}}\PY{p}{\PYZcb{}}\PY{p}{,}
    \PY{p}{\PYZob{}}\PY{l+s+s2}{\PYZdq{}}\PY{l+s+s2}{trabajo extra}\PY{l+s+s2}{\PYZdq{}}\PY{p}{:} \PY{l+s+s2}{\PYZdq{}}\PY{l+s+s2}{Auditoria de experimentos}\PY{l+s+s2}{\PYZdq{}}\PY{p}{,} \PY{l+s+s2}{\PYZdq{}}\PY{l+s+s2}{aporte}\PY{l+s+s2}{\PYZdq{}}\PY{p}{:} \PY{l+s+s2}{\PYZdq{}}\PY{l+s+s2}{Comprueba artefactos, metricas y coherencia de seleccion modelo por modelo.}\PY{l+s+s2}{\PYZdq{}}\PY{p}{\PYZcb{}}\PY{p}{,}
\PY{p}{]}\PY{p}{)}
\PY{n}{display}\PY{p}{(}\PY{n}{extra\PYZus{}work}\PY{p}{)}
\end{Verbatim}
\end{tcolorbox}

    
    \begin{Verbatim}[commandchars=\\\{\}]
                            trabajo extra                                             aporte
0  Silhouette como metrica complementaria  Contrasta la separacion interna de clusters, p{\ldots}
1                                   t-SNE  Visualizacion no lineal de vecindarios y relac{\ldots}
2                 Coste medico FN=5, FP=1  Traduce la prioridad clinica a una metrica com{\ldots}
3                   Curvas de aprendizaje  Ayuda a revisar estabilidad y posible falta de{\ldots}
4                  Permutation importance  Identifica variables que sostienen el recall d{\ldots}
5          Comparacion base vs optimizado  Muestra si optimizar realmente mejora recall/c{\ldots}
6               Auditoria de experimentos  Comprueba artefactos, metricas y coherencia de{\ldots}
    \end{Verbatim}

    
    \subsection{6.1. Trabajo extra de no supervisado: silhouette y
t-SNE}\label{trabajo-extra-de-no-supervisado-silhouette-y-t-sne}

Silhouette se conserva como metrica complementaria para revisar
compacidad y separacion interna de los clusters, pero no se usa como
argumento principal de seleccion en Task 2. La decision principal se
apoya en codo, disparidad jerarquica, interpretabilidad y cruce
posterior con \texttt{condition}.

t-SNE tambien se mueve a trabajo extra porque es una herramienta visual
no lineal: ayuda a mirar vecindarios locales, pero no reemplaza las
tablas de perfiles ni los contrastes cuantitativos.

    \begin{tcolorbox}[breakable, size=fbox, boxrule=1pt, pad at break*=1mm,colback=cellbackground, colframe=cellborder]
\prompt{In}{incolor}{33}{\boxspacing}
\begin{Verbatim}[commandchars=\\\{\}]
\PY{n}{silhouette\PYZus{}columns} \PY{o}{=} \PY{p}{[}\PY{l+s+s2}{\PYZdq{}}\PY{l+s+s2}{k}\PY{l+s+s2}{\PYZdq{}}\PY{p}{,} \PY{l+s+s2}{\PYZdq{}}\PY{l+s+s2}{silhouette}\PY{l+s+s2}{\PYZdq{}}\PY{p}{,} \PY{l+s+s2}{\PYZdq{}}\PY{l+s+s2}{davies\PYZus{}bouldin}\PY{l+s+s2}{\PYZdq{}}\PY{p}{,} \PY{l+s+s2}{\PYZdq{}}\PY{l+s+s2}{calinski\PYZus{}harabasz}\PY{l+s+s2}{\PYZdq{}}\PY{p}{]}
\PY{n}{show\PYZus{}table}\PY{p}{(}\PY{l+s+s2}{\PYZdq{}}\PY{l+s+s2}{Silhouette K\PYZhy{}Means base: k=2..8}\PY{l+s+s2}{\PYZdq{}}\PY{p}{,} \PY{n}{kmeans\PYZus{}metrics}\PY{p}{[}\PY{n}{silhouette\PYZus{}columns}\PY{p}{]}\PY{p}{,} \PY{n}{n}\PY{o}{=}\PY{k+kc}{None}\PY{p}{)}
\PY{n}{show\PYZus{}table}\PY{p}{(}\PY{l+s+s2}{\PYZdq{}}\PY{l+s+s2}{Silhouette K\PYZhy{}Means extendido: k=2..15}\PY{l+s+s2}{\PYZdq{}}\PY{p}{,} \PY{n}{kmeans\PYZus{}metrics\PYZus{}extended}\PY{o}{.}\PY{n}{loc}\PY{p}{[}\PY{n}{kmeans\PYZus{}metrics\PYZus{}extended}\PY{p}{[}\PY{l+s+s2}{\PYZdq{}}\PY{l+s+s2}{k}\PY{l+s+s2}{\PYZdq{}}\PY{p}{]} \PY{o}{\PYZgt{}}\PY{o}{=} \PY{l+m+mi}{2}\PY{p}{,} \PY{n}{silhouette\PYZus{}columns}\PY{p}{]}\PY{p}{,} \PY{n}{n}\PY{o}{=}\PY{k+kc}{None}\PY{p}{)}

\PY{n}{dbscan\PYZus{}silhouette} \PY{o}{=} \PY{n}{dbscan\PYZus{}metrics}\PY{p}{[}\PY{p}{[}\PY{l+s+s2}{\PYZdq{}}\PY{l+s+s2}{eps}\PY{l+s+s2}{\PYZdq{}}\PY{p}{,} \PY{l+s+s2}{\PYZdq{}}\PY{l+s+s2}{min\PYZus{}samples}\PY{l+s+s2}{\PYZdq{}}\PY{p}{,} \PY{l+s+s2}{\PYZdq{}}\PY{l+s+s2}{n\PYZus{}clusters}\PY{l+s+s2}{\PYZdq{}}\PY{p}{,} \PY{l+s+s2}{\PYZdq{}}\PY{l+s+s2}{n\PYZus{}noise}\PY{l+s+s2}{\PYZdq{}}\PY{p}{,} \PY{l+s+s2}{\PYZdq{}}\PY{l+s+s2}{silhouette}\PY{l+s+s2}{\PYZdq{}}\PY{p}{]}\PY{p}{]}\PY{o}{.}\PY{n}{copy}\PY{p}{(}\PY{p}{)}
\PY{n}{show\PYZus{}table}\PY{p}{(}\PY{l+s+s2}{\PYZdq{}}\PY{l+s+s2}{DBSCAN con silhouette como lectura complementaria}\PY{l+s+s2}{\PYZdq{}}\PY{p}{,} \PY{n}{dbscan\PYZus{}silhouette}\PY{p}{,} \PY{n}{n}\PY{o}{=}\PY{k+kc}{None}\PY{p}{)}

\PY{n}{show\PYZus{}image}\PY{p}{(}\PY{n}{OUTPUTS\PYZus{}DIR} \PY{o}{/} \PY{l+s+s2}{\PYZdq{}}\PY{l+s+s2}{task2}\PY{l+s+s2}{\PYZdq{}} \PY{o}{/} \PY{l+s+s2}{\PYZdq{}}\PY{l+s+s2}{kmeans\PYZus{}elbow\PYZus{}silhouette.png}\PY{l+s+s2}{\PYZdq{}}\PY{p}{,} \PY{n}{width}\PY{o}{=}\PY{l+m+mi}{900}\PY{p}{)}
\PY{n}{show\PYZus{}image}\PY{p}{(}\PY{n}{OUTPUTS\PYZus{}DIR} \PY{o}{/} \PY{l+s+s2}{\PYZdq{}}\PY{l+s+s2}{task2}\PY{l+s+s2}{\PYZdq{}} \PY{o}{/} \PY{l+s+s2}{\PYZdq{}}\PY{l+s+s2}{kmeans\PYZus{}elbow\PYZus{}silhouette\PYZus{}extendido.png}\PY{l+s+s2}{\PYZdq{}}\PY{p}{,} \PY{n}{width}\PY{o}{=}\PY{l+m+mi}{900}\PY{p}{)}
\PY{n}{show\PYZus{}image}\PY{p}{(}\PY{n}{OUTPUTS\PYZus{}DIR} \PY{o}{/} \PY{l+s+s2}{\PYZdq{}}\PY{l+s+s2}{task2}\PY{l+s+s2}{\PYZdq{}} \PY{o}{/} \PY{l+s+s2}{\PYZdq{}}\PY{l+s+s2}{tsne\PYZus{}condition.png}\PY{l+s+s2}{\PYZdq{}}\PY{p}{,} \PY{n}{width}\PY{o}{=}\PY{l+m+mi}{800}\PY{p}{)}
\PY{n}{show\PYZus{}image}\PY{p}{(}\PY{n}{OUTPUTS\PYZus{}DIR} \PY{o}{/} \PY{l+s+s2}{\PYZdq{}}\PY{l+s+s2}{task2}\PY{l+s+s2}{\PYZdq{}} \PY{o}{/} \PY{l+s+s2}{\PYZdq{}}\PY{l+s+s2}{tsne\PYZus{}kmeans.png}\PY{l+s+s2}{\PYZdq{}}\PY{p}{,} \PY{n}{width}\PY{o}{=}\PY{l+m+mi}{800}\PY{p}{)}
\PY{n}{show\PYZus{}image}\PY{p}{(}\PY{n}{OUTPUTS\PYZus{}DIR} \PY{o}{/} \PY{l+s+s2}{\PYZdq{}}\PY{l+s+s2}{task2}\PY{l+s+s2}{\PYZdq{}} \PY{o}{/} \PY{l+s+s2}{\PYZdq{}}\PY{l+s+s2}{tsne\PYZus{}condition\PYZus{}kmeans.png}\PY{l+s+s2}{\PYZdq{}}\PY{p}{,} \PY{n}{width}\PY{o}{=}\PY{l+m+mi}{800}\PY{p}{)}
\end{Verbatim}
\end{tcolorbox}

    \textbf{Silhouette K-Means base: k=2..8}

    
    
    \begin{Verbatim}[commandchars=\\\{\}]
   k  silhouette  davies\_bouldin  calinski\_harabasz
0  2    0.168224        2.074719          65.267690
1  3    0.122250        2.275818          49.619397
2  4    0.107557        2.224447          40.992808
3  5    0.101950        2.411273          35.465269
4  6    0.102455        2.273644          32.080288
5  7    0.095350        2.232103          29.080768
6  8    0.086694        2.128847          27.219074
    \end{Verbatim}

    
    \textbf{Silhouette K-Means extendido: k=2..15}

    
    
    \begin{Verbatim}[commandchars=\\\{\}]
     k  silhouette  davies\_bouldin  calinski\_harabasz
1    2    0.168224        2.074719          65.267690
2    3    0.122250        2.275818          49.619397
3    4    0.107557        2.224447          40.992808
4    5    0.101950        2.411273          35.465269
5    6    0.102455        2.273644          32.080288
6    7    0.095350        2.232103          29.080768
7    8    0.086694        2.128847          27.219074
8    9    0.099565        2.098924          25.318049
9   10    0.084522        2.211487          24.173918
10  11    0.085953        2.072470          22.785366
11  12    0.089485        2.126478          21.853636
12  13    0.082518        2.039327          20.762537
13  14    0.092008        1.933066          20.002493
14  15    0.095533        2.029390          19.219024
    \end{Verbatim}

    
    \textbf{DBSCAN con silhouette como lectura complementaria}

    
    
    \begin{Verbatim}[commandchars=\\\{\}]
    eps  min\_samples  n\_clusters  n\_noise  silhouette
0   0.5            3           0      297         NaN
1   0.5            5           0      297         NaN
2   0.5           10           0      297         NaN
3   1.0            3           1      294         NaN
4   1.0            5           0      297         NaN
5   1.0           10           0      297         NaN
6   1.5            3           7      251   -0.186992
7   1.5            5           1      277         NaN
8   1.5           10           0      297         NaN
9   2.0            3           9      136   -0.099911
10  2.0            5           2      185    0.016601
11  2.0           10           2      207   -0.013053
12  2.5            3           1       29         NaN
13  2.5            5           1       35         NaN
14  2.5           10           1       49         NaN
15  3.0            3           1        4         NaN
16  3.0            5           1        4         NaN
17  3.0           10           1        8         NaN
    \end{Verbatim}

    
    \begin{center}
    \adjustimage{max size={0.9\linewidth}{0.9\paperheight}}{REPORT_files/REPORT_74_6.png}
    \end{center}
    { \hspace*{\fill} \\}
    
    \begin{center}
    \adjustimage{max size={0.9\linewidth}{0.9\paperheight}}{REPORT_files/REPORT_74_7.png}
    \end{center}
    { \hspace*{\fill} \\}
    
    \begin{center}
    \adjustimage{max size={0.9\linewidth}{0.9\paperheight}}{REPORT_files/REPORT_74_8.png}
    \end{center}
    { \hspace*{\fill} \\}
    
    \begin{center}
    \adjustimage{max size={0.9\linewidth}{0.9\paperheight}}{REPORT_files/REPORT_74_9.png}
    \end{center}
    { \hspace*{\fill} \\}
    
    \begin{center}
    \adjustimage{max size={0.9\linewidth}{0.9\paperheight}}{REPORT_files/REPORT_74_10.png}
    \end{center}
    { \hspace*{\fill} \\}
    
    \section{6.2. Auditoría de experimentos uno a
uno}\label{auditoruxeda-de-experimentos-uno-a-uno}

Esta sección revisa los artefactos generados por cada etapa y por cada
modelo. La meta es comprobar que los resultados que alimentan el reporte
existen, son coherentes y no contienen errores obvios como métricas
fuera de rango, matrices incompletas o una selección de hiperparámetros
que no coincide con el primer resultado ordenado del grid search.

    \begin{tcolorbox}[breakable, size=fbox, boxrule=1pt, pad at break*=1mm,colback=cellbackground, colframe=cellborder]
\prompt{In}{incolor}{41}{\boxspacing}
\begin{Verbatim}[commandchars=\\\{\}]
\PY{k+kn}{import}\PY{+w}{ }\PY{n+nn}{ast}

\PY{n}{audit\PYZus{}rows} \PY{o}{=} \PY{p}{[}\PY{p}{]}

\PY{k}{def}\PY{+w}{ }\PY{n+nf}{audit}\PY{p}{(}\PY{n}{status}\PY{p}{,} \PY{n}{item}\PY{p}{,} \PY{n}{detail}\PY{p}{)}\PY{p}{:}
    \PY{n}{audit\PYZus{}rows}\PY{o}{.}\PY{n}{append}\PY{p}{(}\PY{p}{\PYZob{}}\PY{l+s+s2}{\PYZdq{}}\PY{l+s+s2}{estado}\PY{l+s+s2}{\PYZdq{}}\PY{p}{:} \PY{n}{status}\PY{p}{,} \PY{l+s+s2}{\PYZdq{}}\PY{l+s+s2}{experimento}\PY{l+s+s2}{\PYZdq{}}\PY{p}{:} \PY{n}{item}\PY{p}{,} \PY{l+s+s2}{\PYZdq{}}\PY{l+s+s2}{detalle}\PY{l+s+s2}{\PYZdq{}}\PY{p}{:} \PY{n}{detail}\PY{p}{\PYZcb{}}\PY{p}{)}

\PY{k}{def}\PY{+w}{ }\PY{n+nf}{metric\PYZus{}between}\PY{p}{(}\PY{n}{value}\PY{p}{,} \PY{n}{lower}\PY{o}{=}\PY{l+m+mi}{0}\PY{p}{,} \PY{n}{upper}\PY{o}{=}\PY{l+m+mi}{1}\PY{p}{)}\PY{p}{:}
    \PY{k}{return} \PY{n}{pd}\PY{o}{.}\PY{n}{notna}\PY{p}{(}\PY{n}{value}\PY{p}{)} \PY{o+ow}{and} \PY{n}{lower} \PY{o}{\PYZlt{}}\PY{o}{=} \PY{n+nb}{float}\PY{p}{(}\PY{n}{value}\PY{p}{)} \PY{o}{\PYZlt{}}\PY{o}{=} \PY{n}{upper}

\PY{k}{def}\PY{+w}{ }\PY{n+nf}{confusion\PYZus{}counts}\PY{p}{(}\PY{n}{path}\PY{p}{)}\PY{p}{:}
    \PY{n}{matrix} \PY{o}{=} \PY{n}{read\PYZus{}csv}\PY{p}{(}\PY{n}{path}\PY{p}{)}
    \PY{n}{value\PYZus{}columns} \PY{o}{=} \PY{p}{[}\PY{n}{col} \PY{k}{for} \PY{n}{col} \PY{o+ow}{in} \PY{n}{matrix}\PY{o}{.}\PY{n}{columns} \PY{k}{if} \PY{n}{col} \PY{o}{!=} \PY{l+s+s2}{\PYZdq{}}\PY{l+s+s2}{real}\PY{l+s+s2}{\PYZdq{}}\PY{p}{]}
    \PY{n}{values} \PY{o}{=} \PY{n}{matrix}\PY{p}{[}\PY{n}{value\PYZus{}columns}\PY{p}{]}\PY{o}{.}\PY{n}{to\PYZus{}numpy}\PY{p}{(}\PY{p}{)}
    \PY{n}{tn}\PY{p}{,} \PY{n}{fp} \PY{o}{=} \PY{n}{values}\PY{p}{[}\PY{l+m+mi}{0}\PY{p}{]}
    \PY{n}{fn}\PY{p}{,} \PY{n}{tp} \PY{o}{=} \PY{n}{values}\PY{p}{[}\PY{l+m+mi}{1}\PY{p}{]}
    \PY{k}{return} \PY{p}{\PYZob{}}\PY{l+s+s2}{\PYZdq{}}\PY{l+s+s2}{tn}\PY{l+s+s2}{\PYZdq{}}\PY{p}{:} \PY{n+nb}{int}\PY{p}{(}\PY{n}{tn}\PY{p}{)}\PY{p}{,} \PY{l+s+s2}{\PYZdq{}}\PY{l+s+s2}{fp}\PY{l+s+s2}{\PYZdq{}}\PY{p}{:} \PY{n+nb}{int}\PY{p}{(}\PY{n}{fp}\PY{p}{)}\PY{p}{,} \PY{l+s+s2}{\PYZdq{}}\PY{l+s+s2}{fn}\PY{l+s+s2}{\PYZdq{}}\PY{p}{:} \PY{n+nb}{int}\PY{p}{(}\PY{n}{fn}\PY{p}{)}\PY{p}{,} \PY{l+s+s2}{\PYZdq{}}\PY{l+s+s2}{tp}\PY{l+s+s2}{\PYZdq{}}\PY{p}{:} \PY{n+nb}{int}\PY{p}{(}\PY{n}{tp}\PY{p}{)}\PY{p}{,} \PY{l+s+s2}{\PYZdq{}}\PY{l+s+s2}{n}\PY{l+s+s2}{\PYZdq{}}\PY{p}{:} \PY{n+nb}{int}\PY{p}{(}\PY{n}{values}\PY{o}{.}\PY{n}{sum}\PY{p}{(}\PY{p}{)}\PY{p}{)}\PY{p}{\PYZcb{}}

\PY{c+c1}{\PYZsh{} Dataset limpio y dataset preparado}
\PY{n}{target\PYZus{}values\PYZus{}ok} \PY{o}{=} \PY{n+nb}{sorted}\PY{p}{(}\PY{n}{df\PYZus{}clean}\PY{p}{[}\PY{n}{target}\PY{p}{]}\PY{o}{.}\PY{n}{dropna}\PY{p}{(}\PY{p}{)}\PY{o}{.}\PY{n}{astype}\PY{p}{(}\PY{n+nb}{int}\PY{p}{)}\PY{o}{.}\PY{n}{unique}\PY{p}{(}\PY{p}{)}\PY{o}{.}\PY{n}{tolist}\PY{p}{(}\PY{p}{)}\PY{p}{)} \PY{o}{==} \PY{p}{[}\PY{l+m+mi}{0}\PY{p}{,} \PY{l+m+mi}{1}\PY{p}{]}
\PY{n}{audit}\PY{p}{(}\PY{l+s+s2}{\PYZdq{}}\PY{l+s+s2}{OK}\PY{l+s+s2}{\PYZdq{}} \PY{k}{if} \PY{n+nb}{len}\PY{p}{(}\PY{n}{df\PYZus{}clean}\PY{p}{)} \PY{o}{==} \PY{l+m+mi}{297} \PY{k}{else} \PY{l+s+s2}{\PYZdq{}}\PY{l+s+s2}{REVISAR}\PY{l+s+s2}{\PYZdq{}}\PY{p}{,} \PY{l+s+s2}{\PYZdq{}}\PY{l+s+s2}{Task 1 \PYZhy{} filas}\PY{l+s+s2}{\PYZdq{}}\PY{p}{,} \PY{l+s+sa}{f}\PY{l+s+s2}{\PYZdq{}}\PY{l+s+s2}{filas=}\PY{l+s+si}{\PYZob{}}\PY{n+nb}{len}\PY{p}{(}\PY{n}{df\PYZus{}clean}\PY{p}{)}\PY{l+s+si}{\PYZcb{}}\PY{l+s+s2}{\PYZdq{}}\PY{p}{)}
\PY{n}{audit}\PY{p}{(}\PY{l+s+s2}{\PYZdq{}}\PY{l+s+s2}{OK}\PY{l+s+s2}{\PYZdq{}} \PY{k}{if} \PY{n+nb}{len}\PY{p}{(}\PY{n}{features}\PY{p}{)} \PY{o}{==} \PY{l+m+mi}{13} \PY{k}{else} \PY{l+s+s2}{\PYZdq{}}\PY{l+s+s2}{REVISAR}\PY{l+s+s2}{\PYZdq{}}\PY{p}{,} \PY{l+s+s2}{\PYZdq{}}\PY{l+s+s2}{Task 1 \PYZhy{} features}\PY{l+s+s2}{\PYZdq{}}\PY{p}{,} \PY{l+s+sa}{f}\PY{l+s+s2}{\PYZdq{}}\PY{l+s+s2}{features=}\PY{l+s+si}{\PYZob{}}\PY{n+nb}{len}\PY{p}{(}\PY{n}{features}\PY{p}{)}\PY{l+s+si}{\PYZcb{}}\PY{l+s+s2}{\PYZdq{}}\PY{p}{)}
\PY{n}{audit}\PY{p}{(}\PY{l+s+s2}{\PYZdq{}}\PY{l+s+s2}{OK}\PY{l+s+s2}{\PYZdq{}} \PY{k}{if} \PY{n}{target\PYZus{}values\PYZus{}ok} \PY{k}{else} \PY{l+s+s2}{\PYZdq{}}\PY{l+s+s2}{REVISAR}\PY{l+s+s2}{\PYZdq{}}\PY{p}{,} \PY{l+s+s2}{\PYZdq{}}\PY{l+s+s2}{Task 1 \PYZhy{} objetivo binario}\PY{l+s+s2}{\PYZdq{}}\PY{p}{,} \PY{l+s+sa}{f}\PY{l+s+s2}{\PYZdq{}}\PY{l+s+s2}{valores=}\PY{l+s+si}{\PYZob{}}\PY{n+nb}{sorted}\PY{p}{(}\PY{n}{df\PYZus{}clean}\PY{p}{[}\PY{n}{target}\PY{p}{]}\PY{o}{.}\PY{n}{dropna}\PY{p}{(}\PY{p}{)}\PY{o}{.}\PY{n}{unique}\PY{p}{(}\PY{p}{)}\PY{o}{.}\PY{n}{tolist}\PY{p}{(}\PY{p}{)}\PY{p}{)}\PY{l+s+si}{\PYZcb{}}\PY{l+s+s2}{\PYZdq{}}\PY{p}{)}
\PY{n}{audit}\PY{p}{(}\PY{l+s+s2}{\PYZdq{}}\PY{l+s+s2}{OK}\PY{l+s+s2}{\PYZdq{}} \PY{k}{if} \PY{n+nb}{int}\PY{p}{(}\PY{n}{df\PYZus{}clean}\PY{o}{.}\PY{n}{isna}\PY{p}{(}\PY{p}{)}\PY{o}{.}\PY{n}{sum}\PY{p}{(}\PY{p}{)}\PY{o}{.}\PY{n}{sum}\PY{p}{(}\PY{p}{)}\PY{p}{)} \PY{o}{==} \PY{l+m+mi}{0} \PY{k}{else} \PY{l+s+s2}{\PYZdq{}}\PY{l+s+s2}{REVISAR}\PY{l+s+s2}{\PYZdq{}}\PY{p}{,} \PY{l+s+s2}{\PYZdq{}}\PY{l+s+s2}{Task 1 \PYZhy{} nulos}\PY{l+s+s2}{\PYZdq{}}\PY{p}{,} \PY{l+s+sa}{f}\PY{l+s+s2}{\PYZdq{}}\PY{l+s+s2}{nulos=}\PY{l+s+si}{\PYZob{}}\PY{n+nb}{int}\PY{p}{(}\PY{n}{df\PYZus{}clean}\PY{o}{.}\PY{n}{isna}\PY{p}{(}\PY{p}{)}\PY{o}{.}\PY{n}{sum}\PY{p}{(}\PY{p}{)}\PY{o}{.}\PY{n}{sum}\PY{p}{(}\PY{p}{)}\PY{p}{)}\PY{l+s+si}{\PYZcb{}}\PY{l+s+s2}{\PYZdq{}}\PY{p}{)}
\PY{n}{audit}\PY{p}{(}\PY{l+s+s2}{\PYZdq{}}\PY{l+s+s2}{OK}\PY{l+s+s2}{\PYZdq{}} \PY{k}{if} \PY{n+nb}{int}\PY{p}{(}\PY{n}{df\PYZus{}clean}\PY{o}{.}\PY{n}{duplicated}\PY{p}{(}\PY{p}{)}\PY{o}{.}\PY{n}{sum}\PY{p}{(}\PY{p}{)}\PY{p}{)} \PY{o}{==} \PY{l+m+mi}{0} \PY{k}{else} \PY{l+s+s2}{\PYZdq{}}\PY{l+s+s2}{REVISAR}\PY{l+s+s2}{\PYZdq{}}\PY{p}{,} \PY{l+s+s2}{\PYZdq{}}\PY{l+s+s2}{Task 1 \PYZhy{} duplicados}\PY{l+s+s2}{\PYZdq{}}\PY{p}{,} \PY{l+s+sa}{f}\PY{l+s+s2}{\PYZdq{}}\PY{l+s+s2}{duplicados=}\PY{l+s+si}{\PYZob{}}\PY{n+nb}{int}\PY{p}{(}\PY{n}{df\PYZus{}clean}\PY{o}{.}\PY{n}{duplicated}\PY{p}{(}\PY{p}{)}\PY{o}{.}\PY{n}{sum}\PY{p}{(}\PY{p}{)}\PY{p}{)}\PY{l+s+si}{\PYZcb{}}\PY{l+s+s2}{\PYZdq{}}\PY{p}{)}

\PY{n}{model\PYZus{}ready\PYZus{}features} \PY{o}{=} \PY{n}{model\PYZus{}ready}\PY{o}{.}\PY{n}{drop}\PY{p}{(}\PY{n}{columns}\PY{o}{=}\PY{p}{[}\PY{n}{target}\PY{p}{]}\PY{p}{)}
\PY{n}{max\PYZus{}abs\PYZus{}mean} \PY{o}{=} \PY{n+nb}{float}\PY{p}{(}\PY{n}{model\PYZus{}ready\PYZus{}features}\PY{o}{.}\PY{n}{mean}\PY{p}{(}\PY{p}{)}\PY{o}{.}\PY{n}{abs}\PY{p}{(}\PY{p}{)}\PY{o}{.}\PY{n}{max}\PY{p}{(}\PY{p}{)}\PY{p}{)}
\PY{n}{max\PYZus{}std\PYZus{}delta} \PY{o}{=} \PY{n+nb}{float}\PY{p}{(}\PY{p}{(}\PY{n}{model\PYZus{}ready\PYZus{}features}\PY{o}{.}\PY{n}{std}\PY{p}{(}\PY{p}{)} \PY{o}{\PYZhy{}} \PY{l+m+mi}{1}\PY{p}{)}\PY{o}{.}\PY{n}{abs}\PY{p}{(}\PY{p}{)}\PY{o}{.}\PY{n}{max}\PY{p}{(}\PY{p}{)}\PY{p}{)}
\PY{n}{audit}\PY{p}{(}\PY{l+s+s2}{\PYZdq{}}\PY{l+s+s2}{OK}\PY{l+s+s2}{\PYZdq{}} \PY{k}{if} \PY{n}{max\PYZus{}abs\PYZus{}mean} \PY{o}{\PYZlt{}} \PY{l+m+mf}{1e\PYZhy{}10} \PY{o+ow}{and} \PY{n}{max\PYZus{}std\PYZus{}delta} \PY{o}{\PYZlt{}} \PY{l+m+mf}{1e\PYZhy{}10} \PY{k}{else} \PY{l+s+s2}{\PYZdq{}}\PY{l+s+s2}{REVISAR}\PY{l+s+s2}{\PYZdq{}}\PY{p}{,} \PY{l+s+s2}{\PYZdq{}}\PY{l+s+s2}{Task 1 \PYZhy{} estandarización}\PY{l+s+s2}{\PYZdq{}}\PY{p}{,} \PY{l+s+sa}{f}\PY{l+s+s2}{\PYZdq{}}\PY{l+s+s2}{max\PYZus{}abs\PYZus{}media=}\PY{l+s+si}{\PYZob{}}\PY{n}{max\PYZus{}abs\PYZus{}mean}\PY{l+s+si}{:}\PY{l+s+s2}{.2e}\PY{l+s+si}{\PYZcb{}}\PY{l+s+s2}{, max\PYZus{}delta\PYZus{}std=}\PY{l+s+si}{\PYZob{}}\PY{n}{max\PYZus{}std\PYZus{}delta}\PY{l+s+si}{:}\PY{l+s+s2}{.2e}\PY{l+s+si}{\PYZcb{}}\PY{l+s+s2}{\PYZdq{}}\PY{p}{)}

\PY{c+c1}{\PYZsh{} No supervisado}
\PY{n}{pca\PYZus{}total} \PY{o}{=} \PY{n+nb}{float}\PY{p}{(}\PY{n}{pca\PYZus{}var}\PY{p}{[}\PY{l+s+s2}{\PYZdq{}}\PY{l+s+s2}{varianza\PYZus{}explicada}\PY{l+s+s2}{\PYZdq{}}\PY{p}{]}\PY{o}{.}\PY{n}{sum}\PY{p}{(}\PY{p}{)}\PY{p}{)} \PY{k}{if} \PY{l+s+s2}{\PYZdq{}}\PY{l+s+s2}{varianza\PYZus{}explicada}\PY{l+s+s2}{\PYZdq{}} \PY{o+ow}{in} \PY{n}{pca\PYZus{}var}\PY{o}{.}\PY{n}{columns} \PY{k}{else} \PY{n}{np}\PY{o}{.}\PY{n}{nan}
\PY{n}{audit}\PY{p}{(}\PY{l+s+s2}{\PYZdq{}}\PY{l+s+s2}{OK}\PY{l+s+s2}{\PYZdq{}} \PY{k}{if} \PY{l+m+mf}{0.99} \PY{o}{\PYZlt{}}\PY{o}{=} \PY{n}{pca\PYZus{}total} \PY{o}{\PYZlt{}}\PY{o}{=} \PY{l+m+mf}{1.01} \PY{k}{else} \PY{l+s+s2}{\PYZdq{}}\PY{l+s+s2}{REVISAR}\PY{l+s+s2}{\PYZdq{}}\PY{p}{,} \PY{l+s+s2}{\PYZdq{}}\PY{l+s+s2}{Task 2 \PYZhy{} PCA}\PY{l+s+s2}{\PYZdq{}}\PY{p}{,} \PY{l+s+sa}{f}\PY{l+s+s2}{\PYZdq{}}\PY{l+s+s2}{varianza total aproximada=}\PY{l+s+si}{\PYZob{}}\PY{n}{pca\PYZus{}total}\PY{l+s+si}{:}\PY{l+s+s2}{.4f}\PY{l+s+si}{\PYZcb{}}\PY{l+s+s2}{\PYZdq{}}\PY{p}{)}
\PY{n}{audit}\PY{p}{(}\PY{l+s+s2}{\PYZdq{}}\PY{l+s+s2}{OK}\PY{l+s+s2}{\PYZdq{}}\PY{p}{,} \PY{l+s+s2}{\PYZdq{}}\PY{l+s+s2}{Task 2 \PYZhy{} K\PYZhy{}Means}\PY{l+s+s2}{\PYZdq{}}\PY{p}{,} \PY{l+s+s2}{\PYZdq{}}\PY{l+s+s2}{codo extendido, perfil de clusters y cruce posterior con condition disponibles}\PY{l+s+s2}{\PYZdq{}}\PY{p}{)}
\PY{n}{audit}\PY{p}{(}\PY{l+s+s2}{\PYZdq{}}\PY{l+s+s2}{OK}\PY{l+s+s2}{\PYZdq{}}\PY{p}{,} \PY{l+s+s2}{\PYZdq{}}\PY{l+s+s2}{Task 2 \PYZhy{} Agglomerative}\PY{l+s+s2}{\PYZdq{}}\PY{p}{,} \PY{l+s+sa}{f}\PY{l+s+s2}{\PYZdq{}}\PY{l+s+s2}{disparidad maxima=}\PY{l+s+si}{\PYZob{}}\PY{n}{agglomerative\PYZus{}disparity}\PY{o}{.}\PY{n}{iloc}\PY{p}{[}\PY{l+m+mi}{0}\PY{p}{]}\PY{p}{[}\PY{l+s+s1}{\PYZsq{}}\PY{l+s+s1}{salto\PYZus{}altura}\PY{l+s+s1}{\PYZsq{}}\PY{p}{]}\PY{l+s+si}{:}\PY{l+s+s2}{.4f}\PY{l+s+si}{\PYZcb{}}\PY{l+s+s2}{; dendrograma generado}\PY{l+s+s2}{\PYZdq{}}\PY{p}{)}
\PY{n}{dbscan\PYZus{}two\PYZus{}cluster\PYZus{}audit} \PY{o}{=} \PY{n}{dbscan\PYZus{}metrics}\PY{o}{.}\PY{n}{loc}\PY{p}{[}\PY{n}{dbscan\PYZus{}metrics}\PY{p}{[}\PY{l+s+s2}{\PYZdq{}}\PY{l+s+s2}{n\PYZus{}clusters}\PY{l+s+s2}{\PYZdq{}}\PY{p}{]} \PY{o}{==} \PY{l+m+mi}{2}\PY{p}{]}\PY{o}{.}\PY{n}{sort\PYZus{}values}\PY{p}{(}\PY{l+s+s2}{\PYZdq{}}\PY{l+s+s2}{n\PYZus{}noise}\PY{l+s+s2}{\PYZdq{}}\PY{p}{)}\PY{o}{.}\PY{n}{head}\PY{p}{(}\PY{l+m+mi}{1}\PY{p}{)}
\PY{n}{dbscan\PYZus{}noise\PYZus{}audit} \PY{o}{=} \PY{n+nb}{int}\PY{p}{(}\PY{n}{dbscan\PYZus{}two\PYZus{}cluster\PYZus{}audit}\PY{o}{.}\PY{n}{iloc}\PY{p}{[}\PY{l+m+mi}{0}\PY{p}{]}\PY{p}{[}\PY{l+s+s2}{\PYZdq{}}\PY{l+s+s2}{n\PYZus{}noise}\PY{l+s+s2}{\PYZdq{}}\PY{p}{]}\PY{p}{)} \PY{k}{if} \PY{o+ow}{not} \PY{n}{dbscan\PYZus{}two\PYZus{}cluster\PYZus{}audit}\PY{o}{.}\PY{n}{empty} \PY{k}{else} \PY{l+s+s2}{\PYZdq{}}\PY{l+s+s2}{NA}\PY{l+s+s2}{\PYZdq{}}
\PY{n}{audit}\PY{p}{(}\PY{l+s+s2}{\PYZdq{}}\PY{l+s+s2}{OK}\PY{l+s+s2}{\PYZdq{}}\PY{p}{,} \PY{l+s+s2}{\PYZdq{}}\PY{l+s+s2}{Task 2 \PYZhy{} DBSCAN}\PY{l+s+s2}{\PYZdq{}}\PY{p}{,} \PY{l+s+sa}{f}\PY{l+s+s2}{\PYZdq{}}\PY{l+s+s2}{configuracion con dos clusters conserva ruido=}\PY{l+s+si}{\PYZob{}}\PY{n}{dbscan\PYZus{}noise\PYZus{}audit}\PY{l+s+si}{\PYZcb{}}\PY{l+s+s2}{; segmentacion limitada}\PY{l+s+s2}{\PYZdq{}}\PY{p}{)}

\PY{c+c1}{\PYZsh{} Supervisado}
\PY{k}{for} \PY{n}{\PYZus{}}\PY{p}{,} \PY{n}{row} \PY{o+ow}{in} \PY{n}{task3\PYZus{}summary}\PY{o}{.}\PY{n}{iterrows}\PY{p}{(}\PY{p}{)}\PY{p}{:}
    \PY{n}{slug\PYZus{}map} \PY{o}{=} \PY{p}{\PYZob{}}
        \PY{l+s+s2}{\PYZdq{}}\PY{l+s+s2}{Regresion Logistica}\PY{l+s+s2}{\PYZdq{}}\PY{p}{:} \PY{l+s+s2}{\PYZdq{}}\PY{l+s+s2}{logistic\PYZus{}regression}\PY{l+s+s2}{\PYZdq{}}\PY{p}{,}
        \PY{l+s+s2}{\PYZdq{}}\PY{l+s+s2}{Regresión Logística}\PY{l+s+s2}{\PYZdq{}}\PY{p}{:} \PY{l+s+s2}{\PYZdq{}}\PY{l+s+s2}{logistic\PYZus{}regression}\PY{l+s+s2}{\PYZdq{}}\PY{p}{,}
        \PY{l+s+s2}{\PYZdq{}}\PY{l+s+s2}{SVM RBF}\PY{l+s+s2}{\PYZdq{}}\PY{p}{:} \PY{l+s+s2}{\PYZdq{}}\PY{l+s+s2}{svm}\PY{l+s+s2}{\PYZdq{}}\PY{p}{,}
        \PY{l+s+s2}{\PYZdq{}}\PY{l+s+s2}{XGBoost}\PY{l+s+s2}{\PYZdq{}}\PY{p}{:} \PY{l+s+s2}{\PYZdq{}}\PY{l+s+s2}{xgboost}\PY{l+s+s2}{\PYZdq{}}\PY{p}{,}
        \PY{l+s+s2}{\PYZdq{}}\PY{l+s+s2}{KNN}\PY{l+s+s2}{\PYZdq{}}\PY{p}{:} \PY{l+s+s2}{\PYZdq{}}\PY{l+s+s2}{knn}\PY{l+s+s2}{\PYZdq{}}\PY{p}{,}
        \PY{l+s+s2}{\PYZdq{}}\PY{l+s+s2}{Naive Bayes}\PY{l+s+s2}{\PYZdq{}}\PY{p}{:} \PY{l+s+s2}{\PYZdq{}}\PY{l+s+s2}{naive\PYZus{}bayes}\PY{l+s+s2}{\PYZdq{}}\PY{p}{,}
        \PY{l+s+s2}{\PYZdq{}}\PY{l+s+s2}{Arbol de Decision}\PY{l+s+s2}{\PYZdq{}}\PY{p}{:} \PY{l+s+s2}{\PYZdq{}}\PY{l+s+s2}{decision\PYZus{}tree}\PY{l+s+s2}{\PYZdq{}}\PY{p}{,}
        \PY{l+s+s2}{\PYZdq{}}\PY{l+s+s2}{Árbol de Decisión CART}\PY{l+s+s2}{\PYZdq{}}\PY{p}{:} \PY{l+s+s2}{\PYZdq{}}\PY{l+s+s2}{decision\PYZus{}tree}\PY{l+s+s2}{\PYZdq{}}\PY{p}{,}
        \PY{l+s+s2}{\PYZdq{}}\PY{l+s+s2}{Red Neuronal MLP}\PY{l+s+s2}{\PYZdq{}}\PY{p}{:} \PY{l+s+s2}{\PYZdq{}}\PY{l+s+s2}{mlp}\PY{l+s+s2}{\PYZdq{}}\PY{p}{,}
    \PY{p}{\PYZcb{}}
    \PY{n}{slug} \PY{o}{=} \PY{n}{slug\PYZus{}map}\PY{p}{[}\PY{n}{row}\PY{p}{[}\PY{l+s+s2}{\PYZdq{}}\PY{l+s+s2}{modelo}\PY{l+s+s2}{\PYZdq{}}\PY{p}{]}\PY{p}{]}
    \PY{n}{counts} \PY{o}{=} \PY{n}{confusion\PYZus{}counts}\PY{p}{(}\PY{n}{OUTPUTS\PYZus{}DIR} \PY{o}{/} \PY{l+s+s2}{\PYZdq{}}\PY{l+s+s2}{task3}\PY{l+s+s2}{\PYZdq{}} \PY{o}{/} \PY{n}{slug} \PY{o}{/} \PY{l+s+s2}{\PYZdq{}}\PY{l+s+s2}{test\PYZus{}matriz\PYZus{}confusion.csv}\PY{l+s+s2}{\PYZdq{}}\PY{p}{)}
    \PY{n}{metrics\PYZus{}ok} \PY{o}{=} \PY{n+nb}{all}\PY{p}{(}\PY{n}{metric\PYZus{}between}\PY{p}{(}\PY{n}{row}\PY{p}{[}\PY{n}{col}\PY{p}{]}\PY{p}{)} \PY{k}{for} \PY{n}{col} \PY{o+ow}{in} \PY{p}{[}\PY{l+s+s2}{\PYZdq{}}\PY{l+s+s2}{cv\PYZus{}recall\PYZus{}condition\PYZus{}1\PYZus{}mean}\PY{l+s+s2}{\PYZdq{}}\PY{p}{,} \PY{l+s+s2}{\PYZdq{}}\PY{l+s+s2}{test\PYZus{}recall\PYZus{}condition\PYZus{}1}\PY{l+s+s2}{\PYZdq{}}\PY{p}{,} \PY{l+s+s2}{\PYZdq{}}\PY{l+s+s2}{test\PYZus{}roc\PYZus{}auc}\PY{l+s+s2}{\PYZdq{}}\PY{p}{]}\PY{p}{)}
    \PY{n}{matrix\PYZus{}ok} \PY{o}{=} \PY{n}{counts}\PY{p}{[}\PY{l+s+s2}{\PYZdq{}}\PY{l+s+s2}{n}\PY{l+s+s2}{\PYZdq{}}\PY{p}{]} \PY{o}{==} \PY{l+m+mi}{60} \PY{o+ow}{and} \PY{n}{counts}\PY{p}{[}\PY{l+s+s2}{\PYZdq{}}\PY{l+s+s2}{fn}\PY{l+s+s2}{\PYZdq{}}\PY{p}{]} \PY{o}{==} \PY{n+nb}{int}\PY{p}{(}\PY{n}{row}\PY{p}{[}\PY{l+s+s2}{\PYZdq{}}\PY{l+s+s2}{test\PYZus{}false\PYZus{}negatives}\PY{l+s+s2}{\PYZdq{}}\PY{p}{]}\PY{p}{)} \PY{o+ow}{and} \PY{n}{counts}\PY{p}{[}\PY{l+s+s2}{\PYZdq{}}\PY{l+s+s2}{fp}\PY{l+s+s2}{\PYZdq{}}\PY{p}{]} \PY{o}{==} \PY{n+nb}{int}\PY{p}{(}\PY{n}{row}\PY{p}{[}\PY{l+s+s2}{\PYZdq{}}\PY{l+s+s2}{test\PYZus{}false\PYZus{}positives}\PY{l+s+s2}{\PYZdq{}}\PY{p}{]}\PY{p}{)}
    \PY{n}{status} \PY{o}{=} \PY{l+s+s2}{\PYZdq{}}\PY{l+s+s2}{OK}\PY{l+s+s2}{\PYZdq{}} \PY{k}{if} \PY{n}{metrics\PYZus{}ok} \PY{o+ow}{and} \PY{n}{matrix\PYZus{}ok} \PY{k}{else} \PY{l+s+s2}{\PYZdq{}}\PY{l+s+s2}{REVISAR}\PY{l+s+s2}{\PYZdq{}}
    \PY{n}{audit}\PY{p}{(}\PY{n}{status}\PY{p}{,} \PY{l+s+sa}{f}\PY{l+s+s2}{\PYZdq{}}\PY{l+s+s2}{Task 3 \PYZhy{} }\PY{l+s+si}{\PYZob{}}\PY{n}{row}\PY{p}{[}\PY{l+s+s1}{\PYZsq{}}\PY{l+s+s1}{modelo}\PY{l+s+s1}{\PYZsq{}}\PY{p}{]}\PY{l+s+si}{\PYZcb{}}\PY{l+s+s2}{\PYZdq{}}\PY{p}{,} \PY{l+s+sa}{f}\PY{l+s+s2}{\PYZdq{}}\PY{l+s+s2}{CV recall=}\PY{l+s+si}{\PYZob{}}\PY{n}{row}\PY{p}{[}\PY{l+s+s1}{\PYZsq{}}\PY{l+s+s1}{cv\PYZus{}recall\PYZus{}condition\PYZus{}1\PYZus{}mean}\PY{l+s+s1}{\PYZsq{}}\PY{p}{]}\PY{l+s+si}{:}\PY{l+s+s2}{.4f}\PY{l+s+si}{\PYZcb{}}\PY{l+s+s2}{, test recall=}\PY{l+s+si}{\PYZob{}}\PY{n}{row}\PY{p}{[}\PY{l+s+s1}{\PYZsq{}}\PY{l+s+s1}{test\PYZus{}recall\PYZus{}condition\PYZus{}1}\PY{l+s+s1}{\PYZsq{}}\PY{p}{]}\PY{l+s+si}{:}\PY{l+s+s2}{.4f}\PY{l+s+si}{\PYZcb{}}\PY{l+s+s2}{, FN test=}\PY{l+s+si}{\PYZob{}}\PY{n}{row}\PY{p}{[}\PY{l+s+s1}{\PYZsq{}}\PY{l+s+s1}{test\PYZus{}false\PYZus{}negatives}\PY{l+s+s1}{\PYZsq{}}\PY{p}{]}\PY{l+s+si}{:}\PY{l+s+s2}{.0f}\PY{l+s+si}{\PYZcb{}}\PY{l+s+s2}{, matriz\PYZus{}ok=}\PY{l+s+si}{\PYZob{}}\PY{n}{matrix\PYZus{}ok}\PY{l+s+si}{\PYZcb{}}\PY{l+s+s2}{\PYZdq{}}\PY{p}{)}

\PY{c+c1}{\PYZsh{} Optimización}
\PY{k}{for} \PY{n}{slug}\PY{p}{,} \PY{n}{name} \PY{o+ow}{in} \PY{p}{[}\PY{p}{(}\PY{l+s+s2}{\PYZdq{}}\PY{l+s+s2}{logistic\PYZus{}regression}\PY{l+s+s2}{\PYZdq{}}\PY{p}{,} \PY{l+s+s2}{\PYZdq{}}\PY{l+s+s2}{Regresion Logistica}\PY{l+s+s2}{\PYZdq{}}\PY{p}{)}\PY{p}{,} \PY{p}{(}\PY{l+s+s2}{\PYZdq{}}\PY{l+s+s2}{svm}\PY{l+s+s2}{\PYZdq{}}\PY{p}{,} \PY{l+s+s2}{\PYZdq{}}\PY{l+s+s2}{SVM RBF}\PY{l+s+s2}{\PYZdq{}}\PY{p}{)}\PY{p}{,} \PY{p}{(}\PY{l+s+s2}{\PYZdq{}}\PY{l+s+s2}{knn}\PY{l+s+s2}{\PYZdq{}}\PY{p}{,} \PY{l+s+s2}{\PYZdq{}}\PY{l+s+s2}{KNN}\PY{l+s+s2}{\PYZdq{}}\PY{p}{)}\PY{p}{]}\PY{p}{:}
    \PY{n}{model\PYZus{}dir} \PY{o}{=} \PY{n}{OUTPUTS\PYZus{}DIR} \PY{o}{/} \PY{l+s+s2}{\PYZdq{}}\PY{l+s+s2}{task4}\PY{l+s+s2}{\PYZdq{}} \PY{o}{/} \PY{n}{slug}
    \PY{n}{params} \PY{o}{=} \PY{n}{read\PYZus{}csv}\PY{p}{(}\PY{n}{model\PYZus{}dir} \PY{o}{/} \PY{l+s+s2}{\PYZdq{}}\PY{l+s+s2}{mejores\PYZus{}parametros.csv}\PY{l+s+s2}{\PYZdq{}}\PY{p}{)}\PY{o}{.}\PY{n}{iloc}\PY{p}{[}\PY{l+m+mi}{0}\PY{p}{]}\PY{o}{.}\PY{n}{dropna}\PY{p}{(}\PY{p}{)}\PY{o}{.}\PY{n}{to\PYZus{}dict}\PY{p}{(}\PY{p}{)}
    \PY{n}{grid} \PY{o}{=} \PY{n}{read\PYZus{}csv}\PY{p}{(}\PY{n}{model\PYZus{}dir} \PY{o}{/} \PY{l+s+s2}{\PYZdq{}}\PY{l+s+s2}{grid\PYZus{}search\PYZus{}resultados.csv}\PY{l+s+s2}{\PYZdq{}}\PY{p}{)}
    \PY{n}{first\PYZus{}params} \PY{o}{=} \PY{n}{ast}\PY{o}{.}\PY{n}{literal\PYZus{}eval}\PY{p}{(}\PY{n}{grid}\PY{o}{.}\PY{n}{iloc}\PY{p}{[}\PY{l+m+mi}{0}\PY{p}{]}\PY{p}{[}\PY{l+s+s2}{\PYZdq{}}\PY{l+s+s2}{params}\PY{l+s+s2}{\PYZdq{}}\PY{p}{]}\PY{p}{)}
    \PY{n}{same\PYZus{}params} \PY{o}{=} \PY{p}{\PYZob{}}\PY{n}{k}\PY{p}{:} \PY{n+nb}{str}\PY{p}{(}\PY{n}{v}\PY{p}{)} \PY{k}{for} \PY{n}{k}\PY{p}{,} \PY{n}{v} \PY{o+ow}{in} \PY{n}{params}\PY{o}{.}\PY{n}{items}\PY{p}{(}\PY{p}{)}\PY{p}{\PYZcb{}} \PY{o}{==} \PY{p}{\PYZob{}}\PY{n}{k}\PY{p}{:} \PY{n+nb}{str}\PY{p}{(}\PY{n}{v}\PY{p}{)} \PY{k}{for} \PY{n}{k}\PY{p}{,} \PY{n}{v} \PY{o+ow}{in} \PY{n}{first\PYZus{}params}\PY{o}{.}\PY{n}{items}\PY{p}{(}\PY{p}{)}\PY{p}{\PYZcb{}}
    \PY{n}{cv} \PY{o}{=} \PY{n}{read\PYZus{}csv}\PY{p}{(}\PY{n}{model\PYZus{}dir} \PY{o}{/} \PY{l+s+s2}{\PYZdq{}}\PY{l+s+s2}{cv\PYZus{}metricas\PYZus{}resumen.csv}\PY{l+s+s2}{\PYZdq{}}\PY{p}{)}
    \PY{n}{test} \PY{o}{=} \PY{n}{read\PYZus{}csv}\PY{p}{(}\PY{n}{model\PYZus{}dir} \PY{o}{/} \PY{l+s+s2}{\PYZdq{}}\PY{l+s+s2}{test\PYZus{}metricas.csv}\PY{l+s+s2}{\PYZdq{}}\PY{p}{)}\PY{o}{.}\PY{n}{iloc}\PY{p}{[}\PY{l+m+mi}{0}\PY{p}{]}
    \PY{n}{counts} \PY{o}{=} \PY{n}{confusion\PYZus{}counts}\PY{p}{(}\PY{n}{model\PYZus{}dir} \PY{o}{/} \PY{l+s+s2}{\PYZdq{}}\PY{l+s+s2}{test\PYZus{}matriz\PYZus{}confusion.csv}\PY{l+s+s2}{\PYZdq{}}\PY{p}{)}
    \PY{n}{cv\PYZus{}recall} \PY{o}{=} \PY{n}{metric\PYZus{}from\PYZus{}summary}\PY{p}{(}\PY{n}{cv}\PY{p}{,} \PY{l+s+s2}{\PYZdq{}}\PY{l+s+s2}{recall\PYZus{}condition\PYZus{}1}\PY{l+s+s2}{\PYZdq{}}\PY{p}{)}
    \PY{n}{cv\PYZus{}cost} \PY{o}{=} \PY{n}{metric\PYZus{}from\PYZus{}summary}\PY{p}{(}\PY{n}{cv}\PY{p}{,} \PY{l+s+s2}{\PYZdq{}}\PY{l+s+s2}{medical\PYZus{}cost}\PY{l+s+s2}{\PYZdq{}}\PY{p}{)}
    \PY{n}{matrix\PYZus{}ok} \PY{o}{=} \PY{n}{counts}\PY{p}{[}\PY{l+s+s2}{\PYZdq{}}\PY{l+s+s2}{n}\PY{l+s+s2}{\PYZdq{}}\PY{p}{]} \PY{o}{==} \PY{l+m+mi}{60} \PY{o+ow}{and} \PY{n}{counts}\PY{p}{[}\PY{l+s+s2}{\PYZdq{}}\PY{l+s+s2}{fn}\PY{l+s+s2}{\PYZdq{}}\PY{p}{]} \PY{o}{==} \PY{n+nb}{int}\PY{p}{(}\PY{n}{test}\PY{p}{[}\PY{l+s+s2}{\PYZdq{}}\PY{l+s+s2}{false\PYZus{}negatives}\PY{l+s+s2}{\PYZdq{}}\PY{p}{]}\PY{p}{)} \PY{o+ow}{and} \PY{n}{counts}\PY{p}{[}\PY{l+s+s2}{\PYZdq{}}\PY{l+s+s2}{fp}\PY{l+s+s2}{\PYZdq{}}\PY{p}{]} \PY{o}{==} \PY{n+nb}{int}\PY{p}{(}\PY{n}{test}\PY{p}{[}\PY{l+s+s2}{\PYZdq{}}\PY{l+s+s2}{false\PYZus{}positives}\PY{l+s+s2}{\PYZdq{}}\PY{p}{]}\PY{p}{)}
    \PY{n}{status} \PY{o}{=} \PY{l+s+s2}{\PYZdq{}}\PY{l+s+s2}{OK}\PY{l+s+s2}{\PYZdq{}} \PY{k}{if} \PY{n}{same\PYZus{}params} \PY{o+ow}{and} \PY{n}{matrix\PYZus{}ok} \PY{o+ow}{and} \PY{n}{metric\PYZus{}between}\PY{p}{(}\PY{n}{cv\PYZus{}recall}\PY{p}{)} \PY{o+ow}{and} \PY{n}{metric\PYZus{}between}\PY{p}{(}\PY{n}{test}\PY{p}{[}\PY{l+s+s2}{\PYZdq{}}\PY{l+s+s2}{recall\PYZus{}condition\PYZus{}1}\PY{l+s+s2}{\PYZdq{}}\PY{p}{]}\PY{p}{)} \PY{k}{else} \PY{l+s+s2}{\PYZdq{}}\PY{l+s+s2}{REVISAR}\PY{l+s+s2}{\PYZdq{}}
    \PY{n}{audit}\PY{p}{(}\PY{n}{status}\PY{p}{,} \PY{l+s+sa}{f}\PY{l+s+s2}{\PYZdq{}}\PY{l+s+s2}{Task 4 \PYZhy{} }\PY{l+s+si}{\PYZob{}}\PY{n}{name}\PY{l+s+si}{\PYZcb{}}\PY{l+s+s2}{\PYZdq{}}\PY{p}{,} \PY{l+s+sa}{f}\PY{l+s+s2}{\PYZdq{}}\PY{l+s+s2}{params\PYZus{}coinciden=}\PY{l+s+si}{\PYZob{}}\PY{n}{same\PYZus{}params}\PY{l+s+si}{\PYZcb{}}\PY{l+s+s2}{, CV cost=}\PY{l+s+si}{\PYZob{}}\PY{n}{cv\PYZus{}cost}\PY{l+s+si}{:}\PY{l+s+s2}{.4f}\PY{l+s+si}{\PYZcb{}}\PY{l+s+s2}{, CV recall=}\PY{l+s+si}{\PYZob{}}\PY{n}{cv\PYZus{}recall}\PY{l+s+si}{:}\PY{l+s+s2}{.4f}\PY{l+s+si}{\PYZcb{}}\PY{l+s+s2}{, test recall=}\PY{l+s+si}{\PYZob{}}\PY{n}{test}\PY{p}{[}\PY{l+s+s1}{\PYZsq{}}\PY{l+s+s1}{recall\PYZus{}condition\PYZus{}1}\PY{l+s+s1}{\PYZsq{}}\PY{p}{]}\PY{l+s+si}{:}\PY{l+s+s2}{.4f}\PY{l+s+si}{\PYZcb{}}\PY{l+s+s2}{, FN test=}\PY{l+s+si}{\PYZob{}}\PY{n}{test}\PY{p}{[}\PY{l+s+s1}{\PYZsq{}}\PY{l+s+s1}{false\PYZus{}negatives}\PY{l+s+s1}{\PYZsq{}}\PY{p}{]}\PY{l+s+si}{:}\PY{l+s+s2}{.0f}\PY{l+s+si}{\PYZcb{}}\PY{l+s+s2}{\PYZdq{}}\PY{p}{)}

\PY{c+c1}{\PYZsh{} Lecturas principales recalculadas desde CSV}
\PY{n}{svm\PYZus{}base} \PY{o}{=} \PY{n}{task3\PYZus{}summary}\PY{o}{.}\PY{n}{loc}\PY{p}{[}\PY{n}{task3\PYZus{}summary}\PY{p}{[}\PY{l+s+s2}{\PYZdq{}}\PY{l+s+s2}{modelo}\PY{l+s+s2}{\PYZdq{}}\PY{p}{]} \PY{o}{==} \PY{l+s+s2}{\PYZdq{}}\PY{l+s+s2}{SVM RBF}\PY{l+s+s2}{\PYZdq{}}\PY{p}{]}\PY{o}{.}\PY{n}{iloc}\PY{p}{[}\PY{l+m+mi}{0}\PY{p}{]}
\PY{n}{svm\PYZus{}opt} \PY{o}{=} \PY{n}{task4\PYZus{}comparison}\PY{o}{.}\PY{n}{loc}\PY{p}{[}\PY{n}{task4\PYZus{}comparison}\PY{p}{[}\PY{l+s+s2}{\PYZdq{}}\PY{l+s+s2}{slug}\PY{l+s+s2}{\PYZdq{}}\PY{p}{]} \PY{o}{==} \PY{l+s+s2}{\PYZdq{}}\PY{l+s+s2}{svm}\PY{l+s+s2}{\PYZdq{}}\PY{p}{]}\PY{o}{.}\PY{n}{iloc}\PY{p}{[}\PY{l+m+mi}{0}\PY{p}{]}
\PY{n}{svm\PYZus{}ok} \PY{o}{=} \PY{n}{np}\PY{o}{.}\PY{n}{isclose}\PY{p}{(}\PY{n}{svm\PYZus{}base}\PY{p}{[}\PY{l+s+s2}{\PYZdq{}}\PY{l+s+s2}{test\PYZus{}recall\PYZus{}condition\PYZus{}1}\PY{l+s+s2}{\PYZdq{}}\PY{p}{]}\PY{p}{,} \PY{n}{svm\PYZus{}opt}\PY{p}{[}\PY{l+s+s2}{\PYZdq{}}\PY{l+s+s2}{test\PYZus{}recall\PYZus{}condition\PYZus{}1}\PY{l+s+s2}{\PYZdq{}}\PY{p}{]}\PY{p}{)} \PY{o+ow}{and} \PY{n}{np}\PY{o}{.}\PY{n}{isclose}\PY{p}{(}\PY{n}{svm\PYZus{}opt}\PY{p}{[}\PY{l+s+s2}{\PYZdq{}}\PY{l+s+s2}{test\PYZus{}recall\PYZus{}condition\PYZus{}1}\PY{l+s+s2}{\PYZdq{}}\PY{p}{]}\PY{p}{,} \PY{l+m+mf}{0.857143}\PY{p}{,} \PY{n}{atol}\PY{o}{=}\PY{l+m+mf}{1e\PYZhy{}6}\PY{p}{)}
\PY{n}{audit}\PY{p}{(}\PY{l+s+s2}{\PYZdq{}}\PY{l+s+s2}{OK}\PY{l+s+s2}{\PYZdq{}} \PY{k}{if} \PY{n}{svm\PYZus{}ok} \PY{k}{else} \PY{l+s+s2}{\PYZdq{}}\PY{l+s+s2}{REVISAR}\PY{l+s+s2}{\PYZdq{}}\PY{p}{,} \PY{l+s+s2}{\PYZdq{}}\PY{l+s+s2}{Lectura \PYZhy{} SVM base/optimizado}\PY{l+s+s2}{\PYZdq{}}\PY{p}{,} \PY{l+s+sa}{f}\PY{l+s+s2}{\PYZdq{}}\PY{l+s+s2}{recall=}\PY{l+s+si}{\PYZob{}}\PY{n}{svm\PYZus{}opt}\PY{p}{[}\PY{l+s+s1}{\PYZsq{}}\PY{l+s+s1}{test\PYZus{}recall\PYZus{}condition\PYZus{}1}\PY{l+s+s1}{\PYZsq{}}\PY{p}{]}\PY{l+s+si}{:}\PY{l+s+s2}{.4f}\PY{l+s+si}{\PYZcb{}}\PY{l+s+s2}{, FN=}\PY{l+s+si}{\PYZob{}}\PY{n}{svm\PYZus{}opt}\PY{p}{[}\PY{l+s+s1}{\PYZsq{}}\PY{l+s+s1}{test\PYZus{}false\PYZus{}negatives}\PY{l+s+s1}{\PYZsq{}}\PY{p}{]}\PY{l+s+si}{:}\PY{l+s+s2}{.0f}\PY{l+s+si}{\PYZcb{}}\PY{l+s+s2}{, coste=}\PY{l+s+si}{\PYZob{}}\PY{n}{svm\PYZus{}opt}\PY{p}{[}\PY{l+s+s1}{\PYZsq{}}\PY{l+s+s1}{test\PYZus{}medical\PYZus{}cost}\PY{l+s+s1}{\PYZsq{}}\PY{p}{]}\PY{l+s+si}{:}\PY{l+s+s2}{.4f}\PY{l+s+si}{\PYZcb{}}\PY{l+s+s2}{\PYZdq{}}\PY{p}{)}

\PY{n}{lr\PYZus{}opt} \PY{o}{=} \PY{n}{task4\PYZus{}comparison}\PY{o}{.}\PY{n}{loc}\PY{p}{[}\PY{n}{task4\PYZus{}comparison}\PY{p}{[}\PY{l+s+s2}{\PYZdq{}}\PY{l+s+s2}{slug}\PY{l+s+s2}{\PYZdq{}}\PY{p}{]} \PY{o}{==} \PY{l+s+s2}{\PYZdq{}}\PY{l+s+s2}{logistic\PYZus{}regression}\PY{l+s+s2}{\PYZdq{}}\PY{p}{]}\PY{o}{.}\PY{n}{iloc}\PY{p}{[}\PY{l+m+mi}{0}\PY{p}{]}
\PY{n}{audit}\PY{p}{(}\PY{l+s+s2}{\PYZdq{}}\PY{l+s+s2}{OK}\PY{l+s+s2}{\PYZdq{}} \PY{k}{if} \PY{n}{lr\PYZus{}opt}\PY{p}{[}\PY{l+s+s2}{\PYZdq{}}\PY{l+s+s2}{delta\PYZus{}cv\PYZus{}medical\PYZus{}cost}\PY{l+s+s2}{\PYZdq{}}\PY{p}{]} \PY{o}{\PYZlt{}} \PY{l+m+mi}{0} \PY{o+ow}{and} \PY{n}{lr\PYZus{}opt}\PY{p}{[}\PY{l+s+s2}{\PYZdq{}}\PY{l+s+s2}{delta\PYZus{}test\PYZus{}recall}\PY{l+s+s2}{\PYZdq{}}\PY{p}{]} \PY{o}{==} \PY{l+m+mi}{0} \PY{k}{else} \PY{l+s+s2}{\PYZdq{}}\PY{l+s+s2}{REVISAR}\PY{l+s+s2}{\PYZdq{}}\PY{p}{,} \PY{l+s+s2}{\PYZdq{}}\PY{l+s+s2}{Lectura \PYZhy{} Regresión Logística}\PY{l+s+s2}{\PYZdq{}}\PY{p}{,} \PY{l+s+sa}{f}\PY{l+s+s2}{\PYZdq{}}\PY{l+s+s2}{mejora CV cost=}\PY{l+s+si}{\PYZob{}}\PY{n}{lr\PYZus{}opt}\PY{p}{[}\PY{l+s+s1}{\PYZsq{}}\PY{l+s+s1}{delta\PYZus{}cv\PYZus{}medical\PYZus{}cost}\PY{l+s+s1}{\PYZsq{}}\PY{p}{]}\PY{l+s+si}{:}\PY{l+s+s2}{.4f}\PY{l+s+si}{\PYZcb{}}\PY{l+s+s2}{, delta recall test=}\PY{l+s+si}{\PYZob{}}\PY{n}{lr\PYZus{}opt}\PY{p}{[}\PY{l+s+s1}{\PYZsq{}}\PY{l+s+s1}{delta\PYZus{}test\PYZus{}recall}\PY{l+s+s1}{\PYZsq{}}\PY{p}{]}\PY{l+s+si}{:}\PY{l+s+s2}{.4f}\PY{l+s+si}{\PYZcb{}}\PY{l+s+s2}{\PYZdq{}}\PY{p}{)}

\PY{n}{knn\PYZus{}opt} \PY{o}{=} \PY{n}{task4\PYZus{}comparison}\PY{o}{.}\PY{n}{loc}\PY{p}{[}\PY{n}{task4\PYZus{}comparison}\PY{p}{[}\PY{l+s+s2}{\PYZdq{}}\PY{l+s+s2}{slug}\PY{l+s+s2}{\PYZdq{}}\PY{p}{]} \PY{o}{==} \PY{l+s+s2}{\PYZdq{}}\PY{l+s+s2}{knn}\PY{l+s+s2}{\PYZdq{}}\PY{p}{]}\PY{o}{.}\PY{n}{iloc}\PY{p}{[}\PY{l+m+mi}{0}\PY{p}{]}
\PY{n}{audit}\PY{p}{(}\PY{l+s+s2}{\PYZdq{}}\PY{l+s+s2}{OK}\PY{l+s+s2}{\PYZdq{}} \PY{k}{if} \PY{n}{knn\PYZus{}opt}\PY{p}{[}\PY{l+s+s2}{\PYZdq{}}\PY{l+s+s2}{delta\PYZus{}test\PYZus{}medical\PYZus{}cost}\PY{l+s+s2}{\PYZdq{}}\PY{p}{]} \PY{o}{\PYZlt{}} \PY{l+m+mi}{0} \PY{o+ow}{and} \PY{n}{knn\PYZus{}opt}\PY{p}{[}\PY{l+s+s2}{\PYZdq{}}\PY{l+s+s2}{test\PYZus{}false\PYZus{}positives}\PY{l+s+s2}{\PYZdq{}}\PY{p}{]} \PY{o}{\PYZlt{}} \PY{l+m+mi}{6} \PY{k}{else} \PY{l+s+s2}{\PYZdq{}}\PY{l+s+s2}{REVISAR}\PY{l+s+s2}{\PYZdq{}}\PY{p}{,} \PY{l+s+s2}{\PYZdq{}}\PY{l+s+s2}{Lectura \PYZhy{} KNN}\PY{l+s+s2}{\PYZdq{}}\PY{p}{,} \PY{l+s+sa}{f}\PY{l+s+s2}{\PYZdq{}}\PY{l+s+s2}{FP test=}\PY{l+s+si}{\PYZob{}}\PY{n}{knn\PYZus{}opt}\PY{p}{[}\PY{l+s+s1}{\PYZsq{}}\PY{l+s+s1}{test\PYZus{}false\PYZus{}positives}\PY{l+s+s1}{\PYZsq{}}\PY{p}{]}\PY{l+s+si}{:}\PY{l+s+s2}{.0f}\PY{l+s+si}{\PYZcb{}}\PY{l+s+s2}{, delta coste test=}\PY{l+s+si}{\PYZob{}}\PY{n}{knn\PYZus{}opt}\PY{p}{[}\PY{l+s+s1}{\PYZsq{}}\PY{l+s+s1}{delta\PYZus{}test\PYZus{}medical\PYZus{}cost}\PY{l+s+s1}{\PYZsq{}}\PY{p}{]}\PY{l+s+si}{:}\PY{l+s+s2}{.4f}\PY{l+s+si}{\PYZcb{}}\PY{l+s+s2}{\PYZdq{}}\PY{p}{)}

\PY{n}{audit}\PY{p}{(}\PY{l+s+s2}{\PYZdq{}}\PY{l+s+s2}{OK}\PY{l+s+s2}{\PYZdq{}} \PY{k}{if} \PY{n}{mlp\PYZus{}row}\PY{p}{[}\PY{l+s+s2}{\PYZdq{}}\PY{l+s+s2}{test\PYZus{}false\PYZus{}negatives}\PY{l+s+s2}{\PYZdq{}}\PY{p}{]} \PY{o}{==} \PY{l+m+mi}{21} \PY{k}{else} \PY{l+s+s2}{\PYZdq{}}\PY{l+s+s2}{REVISAR}\PY{l+s+s2}{\PYZdq{}}\PY{p}{,} \PY{l+s+s2}{\PYZdq{}}\PY{l+s+s2}{Lectura \PYZhy{} MLP débil}\PY{l+s+s2}{\PYZdq{}}\PY{p}{,} \PY{l+s+sa}{f}\PY{l+s+s2}{\PYZdq{}}\PY{l+s+s2}{FN test=}\PY{l+s+si}{\PYZob{}}\PY{n}{mlp\PYZus{}row}\PY{p}{[}\PY{l+s+s1}{\PYZsq{}}\PY{l+s+s1}{test\PYZus{}false\PYZus{}negatives}\PY{l+s+s1}{\PYZsq{}}\PY{p}{]}\PY{l+s+si}{:}\PY{l+s+s2}{.0f}\PY{l+s+si}{\PYZcb{}}\PY{l+s+s2}{, recall test=}\PY{l+s+si}{\PYZob{}}\PY{n}{mlp\PYZus{}row}\PY{p}{[}\PY{l+s+s1}{\PYZsq{}}\PY{l+s+s1}{test\PYZus{}recall\PYZus{}condition\PYZus{}1}\PY{l+s+s1}{\PYZsq{}}\PY{p}{]}\PY{l+s+si}{:}\PY{l+s+s2}{.4f}\PY{l+s+si}{\PYZcb{}}\PY{l+s+s2}{\PYZdq{}}\PY{p}{)}

\PY{n}{ranking\PYZus{}costs} \PY{o}{=} \PY{n}{task4\PYZus{}ranking}\PY{p}{[}\PY{l+s+s2}{\PYZdq{}}\PY{l+s+s2}{cv\PYZus{}medical\PYZus{}cost\PYZus{}mean}\PY{l+s+s2}{\PYZdq{}}\PY{p}{]}\PY{o}{.}\PY{n}{tolist}\PY{p}{(}\PY{p}{)}
\PY{n}{ranking\PYZus{}ok} \PY{o}{=} \PY{n}{ranking\PYZus{}costs} \PY{o}{==} \PY{n+nb}{sorted}\PY{p}{(}\PY{n}{ranking\PYZus{}costs}\PY{p}{)}
\PY{n}{audit}\PY{p}{(}\PY{l+s+s2}{\PYZdq{}}\PY{l+s+s2}{OK}\PY{l+s+s2}{\PYZdq{}} \PY{k}{if} \PY{n}{ranking\PYZus{}ok} \PY{k}{else} \PY{l+s+s2}{\PYZdq{}}\PY{l+s+s2}{REVISAR}\PY{l+s+s2}{\PYZdq{}}\PY{p}{,} \PY{l+s+s2}{\PYZdq{}}\PY{l+s+s2}{Task 4 \PYZhy{} ranking}\PY{l+s+s2}{\PYZdq{}}\PY{p}{,} \PY{l+s+sa}{f}\PY{l+s+s2}{\PYZdq{}}\PY{l+s+s2}{ordenado\PYZus{}por\PYZus{}coste\PYZus{}cv=}\PY{l+s+si}{\PYZob{}}\PY{n}{ranking\PYZus{}ok}\PY{l+s+si}{\PYZcb{}}\PY{l+s+s2}{, costes=}\PY{l+s+si}{\PYZob{}}\PY{p}{[}\PY{n+nb}{round}\PY{p}{(}\PY{n}{x}\PY{p}{,}\PY{+w}{ }\PY{l+m+mi}{4}\PY{p}{)}\PY{+w}{ }\PY{k}{for}\PY{+w}{ }\PY{n}{x}\PY{+w}{ }\PY{o+ow}{in}\PY{+w}{ }\PY{n}{ranking\PYZus{}costs}\PY{p}{]}\PY{l+s+si}{\PYZcb{}}\PY{l+s+s2}{\PYZdq{}}\PY{p}{)}

\PY{n}{audit\PYZus{}df} \PY{o}{=} \PY{n}{pd}\PY{o}{.}\PY{n}{DataFrame}\PY{p}{(}\PY{n}{audit\PYZus{}rows}\PY{p}{)}
\PY{n}{display}\PY{p}{(}\PY{n}{audit\PYZus{}df}\PY{p}{)}

\PY{k}{if} \PY{p}{(}\PY{n}{audit\PYZus{}df}\PY{p}{[}\PY{l+s+s2}{\PYZdq{}}\PY{l+s+s2}{estado}\PY{l+s+s2}{\PYZdq{}}\PY{p}{]} \PY{o}{==} \PY{l+s+s2}{\PYZdq{}}\PY{l+s+s2}{REVISAR}\PY{l+s+s2}{\PYZdq{}}\PY{p}{)}\PY{o}{.}\PY{n}{any}\PY{p}{(}\PY{p}{)}\PY{p}{:}
    \PY{n}{display}\PY{p}{(}\PY{n}{audit\PYZus{}df}\PY{p}{[}\PY{n}{audit\PYZus{}df}\PY{p}{[}\PY{l+s+s2}{\PYZdq{}}\PY{l+s+s2}{estado}\PY{l+s+s2}{\PYZdq{}}\PY{p}{]} \PY{o}{==} \PY{l+s+s2}{\PYZdq{}}\PY{l+s+s2}{REVISAR}\PY{l+s+s2}{\PYZdq{}}\PY{p}{]}\PY{p}{)}
\PY{k}{else}\PY{p}{:}
    \PY{n+nb}{print}\PY{p}{(}\PY{l+s+s2}{\PYZdq{}}\PY{l+s+s2}{Auditoría completada: no se detectaron inconsistencias críticas en los artefactos revisados.}\PY{l+s+s2}{\PYZdq{}}\PY{p}{)}
\end{Verbatim}
\end{tcolorbox}

    
    \begin{Verbatim}[commandchars=\\\{\}]
   estado                      experimento                                            detalle
0      OK                   Task 1 - filas                                          filas=297
1      OK                Task 1 - features                                        features=13
2      OK        Task 1 - objetivo binario                                     valores=[0, 1]
3      OK                   Task 1 - nulos                                            nulos=0
4      OK              Task 1 - duplicados                                       duplicados=0
5      OK         Task 1 - estandarización     max\_abs\_media=4.96e-16, max\_delta\_std=1.11e-16
6      OK                     Task 2 - PCA                   varianza total aproximada=1.0000
7      OK                 Task 2 - K-Means  codo extendido, perfil de clusters y cruce pos{\ldots}
8      OK           Task 2 - Agglomerative     disparidad maxima=7.8584; dendrograma generado
9      OK                  Task 2 - DBSCAN  configuracion con dos clusters conserva ruido={\ldots}
10     OK     Task 3 - Regresión Logística  CV recall=0.8135, test recall=0.8214, FN test={\ldots}
11     OK                 Task 3 - SVM RBF  CV recall=0.8066, test recall=0.8571, FN test={\ldots}
12     OK                 Task 3 - XGBoost  CV recall=0.8064, test recall=0.8214, FN test={\ldots}
13     OK                     Task 3 - KNN  CV recall=0.7853, test recall=0.8214, FN test={\ldots}
14     OK             Task 3 - Naive Bayes  CV recall=0.7282, test recall=0.8214, FN test={\ldots}
15     OK  Task 3 - Árbol de Decisión CART  CV recall=0.7193, test recall=0.7857, FN test={\ldots}
16     OK        Task 3 - Red Neuronal MLP  CV recall=0.5720, test recall=0.2500, FN test={\ldots}
17     OK     Task 4 - Regresion Logistica  params\_coinciden=True, CV cost=0.4217, CV reca{\ldots}
18     OK                 Task 4 - SVM RBF  params\_coinciden=True, CV cost=0.4601, CV reca{\ldots}
19     OK                     Task 4 - KNN  params\_coinciden=True, CV cost=0.5372, CV reca{\ldots}
20     OK    Lectura - SVM base/optimizado                  recall=0.8571, FN=4, coste=0.4667
21     OK    Lectura - Regresión Logística   mejora CV cost=-0.0718, delta recall test=0.0000
22     OK                    Lectura - KNN                FP test=5, delta coste test=-0.0167
23     OK              Lectura - MLP débil                     FN test=21, recall test=0.2500
24     OK                 Task 4 - ranking  ordenado\_por\_coste\_cv=True, costes=[0.4217, 0{\ldots}
    \end{Verbatim}

    
    \begin{Verbatim}[commandchars=\\\{\}]
Auditoría completada: no se detectaron inconsistencias críticas en los
artefactos revisados.
    \end{Verbatim}

    \section{7. Conclusiones y recomendación
final}\label{conclusiones-y-recomendaciuxf3n-final}

El análisis muestra que el dataset tiene tamaño moderado, variables
numéricas y balance razonable entre clases. Los métodos no supervisados
sugieren una estructura latente parcialmente relacionada con
\texttt{condition}, especialmente con K-Means en dos grupos. Esta
lectura se apoya en el método del codo, la disparidad del dendrograma
jerárquico y el cruce posterior con \texttt{condition}. Por eso los
clusters se interpretan como exploración y no como diagnóstico.

En aprendizaje supervisado, Regresión Logística fue el modelo base más
sólido en validación cruzada por recall, ROC-AUC, coste médico e
interpretabilidad. SVM RBF quedó como alternativa prometedora por su
menor coste en el test reservado. XGBoost se añadió como contraste
potente de árboles y mejora claramente al CART simple, aunque no
desplaza a Regresión Logística cuando la decisión principal se apoya en
validación cruzada e interpretabilidad. El MLP queda documentado como
intento débil porque deja demasiados falsos negativos.

Tras corregir la regla de selección de hiperparámetros, Regresión
Logística optimizada queda como el modelo más fuerte por coste médico
medio en validación cruzada. SVM RBF optimizado obtiene el menor coste
de test entre los optimizados y detecta 24 de 28 positivos reservados,
aunque debe interpretarse junto con su menor interpretabilidad. KNN
optimizado queda como alternativa razonable: reduce falsos positivos y
coste de test frente a su versión base, pero no mejora el recall de
test.

La recomendación final es priorizar Regresión Logística optimizada como
opción principal por estabilidad, coste médico en validación e
interpretabilidad. SVM RBF optimizado queda como alternativa competitiva
si se prioriza el menor número de falsos negativos en el test reservado.


    % Add a bibliography block to the postdoc
    
    
    
\end{document}
